\documentclass[12pt,a4paper,oneside]{report}

% ===================== Encoding, language, fonts =====================
\usepackage[utf8]{inputenc}
\usepackage[T1]{fontenc}
\usepackage{lmodern}
\usepackage[english]{babel}
\usepackage{microtype}

% ===================== Layout =====================
\usepackage[margin=2.5cm]{geometry}
\usepackage{setspace}
\onehalfspacing
\usepackage{titlesec}
\usepackage{fancyhdr}
\usepackage{tabularx}
\usepackage{booktabs}
\usepackage{enumitem}

% ===================== Mathematics =====================
\usepackage{amsmath,amssymb,amsthm,mathtools}
\usepackage{bm}
\newtheorem{theorem}{Theorem}[chapter]
\newtheorem{definition}{Definition}[chapter]
\newtheorem{proposition}{Proposition}[chapter]
\newtheorem{remark}{Remark}[chapter]

% ===================== Graphics =====================
\usepackage{graphicx}
\usepackage{xcolor}
\usepackage{tikz}
\usetikzlibrary{arrows.meta,positioning,shapes.geometric,fit,calc,backgrounds,decorations.pathreplacing}
\usepackage{pgfplots}
\pgfplotsset{compat=1.18}
\usepackage{subcaption}
\usepackage{float}

% ===================== Algorithms and code =====================
\usepackage[ruled,vlined,linesnumbered]{algorithm2e}
\usepackage{listings}

% ===================== Bibliography =====================
\usepackage[backend=biber,style=ieee,sorting=nyt]{biblatex}
\addbibresource{references.bib}

% ===================== Links and cross-references (load last) =====================
\usepackage[colorlinks=true,linkcolor=blue!50!black,citecolor=blue!50!black,urlcolor=blue!50!black]{hyperref}
\usepackage[capitalise,noabbrev]{cleveref}

% ===================== Metadata =====================
\hypersetup{
  pdftitle={Design of a Telemonitoring System for Patients with Chronic Diseases (Hypertension, Diabetes) Using Consultation Data, Home Measurements and Possibly Mobile Sensors},
  pdfauthor={Rodrick Ayuni Nsame}
}

% ===================== Sectioning: 1.1.4a-style lettered subsubsections =====================
\setcounter{secnumdepth}{3}
\renewcommand{\thesubsubsection}{\thesubsection\alph{subsubsection}}

% ===================== Plot styling =====================
\pgfplotsset{every axis/.append style={
  tick label style={font=\small}, label style={font=\small},
  legend style={font=\footnotesize}}}
\definecolor{cAfrica}{RGB}{192,57,43}
\definecolor{cRef}{RGB}{41,128,185}
\definecolor{cGrey}{RGB}{127,140,141}

% ===================== Drafting helpers =====================
% Searchable "to verify" footnotes
\newcommand{\verify}[1]{\footnote{\textbf{[To verify]} #1}}
% Back-translated quotation marker: dagger shows quotes still to be replaced by the verbatim English.
% Once all quotes are verified, redefine as: \renewcommand{\bq}[1]{``#1''}
\newcommand{\bq}[1]{``#1''\textsuperscript{\dag}}

\begin{document}

%==========================================================================
% TITLE PAGE
%==========================================================================
\begin{titlepage}
\centering
\setstretch{1.15}

{\normalsize University of Douala}\\[0.25cm]
{\normalsize National Advanced School of Engineering of Douala (ENSPD)}\\[0.25cm]
{\small Laboratory of Computer Engineering, Data Science and Artificial Intelligence}

\vspace{1.2cm}
\rule{\textwidth}{0.6pt}\\[0.6cm]

{\LARGE\bfseries Design of a Telemonitoring System\\[0.2cm]
for Patients with Chronic Diseases\\[0.2cm]
(Hypertension, Diabetes)}\\[0.4cm]
{\large\bfseries Using Consultation Data, Home Measurements\\[0.1cm]
and Possibly Mobile Sensors}

\vspace{0.6cm}
\rule{\textwidth}{0.6pt}

\vspace{1.2cm}
{\large Master of Science Thesis}\\[0.2cm]
{\normalsize Data Science and Artificial Intelligence}

\vspace{1.5cm}
{\normalsize Presented by}\\[0.3cm]
{\Large\bfseries NSAME AYUNI Rodrick}

\vspace{1.5cm}
{\normalsize Thesis supervised by}\\[0.4cm]

\begin{tabularx}{0.9\textwidth}{@{}l X@{}}
\textbf{Dr. Malong}   & Main supervisor, ENSPD \\[0.2cm]
\textbf{Dr. Noulapeu} & Co-supervisor, ENSPD \\
\end{tabularx}

\vfill
{\normalsize Douala, Cameroon \quad $\bullet$ \quad October 14, 2026}

\end{titlepage}

\tableofcontents
\clearpage

%==========================================================================
% CHAPTER 1
%==========================================================================
\chapter{Epidemiological Background and Rationale}\label{ch:background}

%##########################################################################
% SECTION 1.1
%##########################################################################
\section{Burden of Hypertension--Diabetes Comorbidity in Sub-Saharan Africa}\label{sec:burden}

\subsection*{Purpose and methodological framing}

This section establishes the epidemiological basis of the present work. It answers four questions in turn: how large is the problem posed by hypertension and diabetes in sub-Saharan Africa; why existing care systems fail to control it; what scientifically justifies measuring blood pressure and blood glucose jointly rather than monitoring the two conditions separately; and to what extent the current state of the evidence supports the proposal of a telemonitoring system.

A preliminary methodological framing is necessary, because it conditions how every figure below should be read. Available estimates for the region come from four types of sources that are not equivalent and must not be compared without care.

\emph{Model-based estimates}, produced by the World Health Organization (WHO) and the NCD Risk Factor Collaboration (NCD-RisC), yield age-standardised figures for all countries, including those without any recent survey; they offer complete coverage at the price of non-negligible uncertainty for poorly documented countries. \emph{Meta-analyses of population surveys}, such as that of \textcite{ataklte2015}, aggregate real measurements but inherit the heterogeneity of their sources. \emph{National or regional surveys} vary strongly in protocol: WHO STEPS surveys rely on probability sampling, May Measurement Month (MMM) campaigns on self-selected opportunistic screening, and Demographic and Health Surveys (DHS) on household samples; these three designs do not produce comparable prevalences. Finally, \emph{cross-sectional hospital studies}, often small and conducted on patients already engaged in care, are by construction unrepresentative of the general population.

Two further sources of variation are specific to the measurement itself. The first is the diagnostic threshold, 140/90 or 130/80~mmHg. The second is the number of measurement visits: nearly all surveys in the region rely on a single measurement, whereas the clinical diagnosis of hypertension requires confirmation over several visits. \textcite{gafane2025} identify this limitation in most of the surveys they review and report that it leads to overestimating prevalence, since the white-coat effect and spontaneous blood pressure variability artificially inflate the number of subjects classified as hypertensive on an isolated measurement.

The nature of each source is therefore stated systematically in what follows, and the comparability of the main estimates is summarised in \cref{sec:comparability}.

%==========================================================================
\subsection{The epidemiological transition: a shift in the burden of disease}\label{sec:transition}

Sub-Saharan Africa is undergoing an epidemiological transition, that is, a progressive shift of the burden of disease from infectious and nutritional conditions towards noncommunicable diseases (NCDs): cardiovascular disease, diabetes, cancers and chronic respiratory disease. What is distinctive in the region is that this transition does not replace the infectious burden but is superimposed on it, creating what \textcite{hinneh2023} describe as a double burden on health systems that are already under strain.

\subsubsection{Measuring the phenomenon}

The indicator that best captures this shift is the disability-adjusted life year (DALY), which represents one year of healthy life lost, whether through premature death or through living with a disability:
\begin{equation}
  \mathrm{DALY} = \mathrm{YLL} + \mathrm{YLD},
\end{equation}
where YLL (years of life lost) denotes the years lost to premature mortality, computed as the gap between age at death and a reference life expectancy, and YLD (years lived with disability) the years lived with a disability, weighted by a severity coefficient between 0 and 1. For our purpose, the value of this measure is that it captures what mortality alone ignores: a chronic condition such as hypertension or diabetes kills late but disables for a long time and mobilises care resources over several decades. It is precisely this duration that makes longitudinal monitoring relevant and episodic monitoring insufficient.

\subsubsection{Regional orders of magnitude}

The Global Burden of Disease 2017 study, whose results we report as relayed by \textcite{nagavci2022}---and not from the original publication of \textcite{gouda2019}, which should be consulted directly for the final version---provides the following values for sub-Saharan Africa:
\begin{itemize}[nosep]
  \item the share of NCDs in the total burden of disease rose from 18.6\% in 1990 to 29.8\% in 2017;
  \item the total number of DALYs attributable to them increased by 67\% over the same period;
  \item within these diseases, 15\% of DALYs are attributable to cardiovascular disease, of which hypertension is the single most important risk factor;
  \item DALYs attributable to diabetes increased by 126.4\% over three decades.
\end{itemize}

\subsubsection{Two interpretive precautions}

These values call for two precautions, the omission of which is the most common error in the introductory literature on this topic.

First, the rise from 18.6\% to 29.8\% is a \emph{relative share}: it reflects the decline of infectious diseases, a result of HIV, tuberculosis and malaria control programmes, as much as the absolute growth of chronic diseases. It does not directly measure the worsening of the latter.

Second, a substantial fraction of the apparent increase in chronic disease is mechanically explained by population ageing and growth, not by a rise in individual risk. \textcite{who2023hypertension} establishes this explicitly for hypertension: the global age-standardised prevalence barely changed between 1990 and 2019, from 32\% to 33\%, while the absolute number of people affected doubled, from 650 million to 1.3 billion. The report specifies that the total number of adults is growing because of population growth and that the age structure is shifting towards older groups, where prevalence is mechanically higher; these two trends explain most of the change in absolute terms. It also notes that age-standardised prevalence fell in high-income countries, from 38\% to 32\%, and in the European Region, from 45\% to 37\%, these decreases being offset by moderate increases in the Western Pacific (24\% to 28\%) and South-East Asia (29\% to 32\%) Regions.

We therefore state systematically below whether a prevalence is age-standardised or crude. The distinction is not a formality: it determines whether one is describing a worsening of risk or a demographic structure effect, and hence the nature of the response required.

\subsubsection{Diabetes}

For diabetes, \textcite{nagavci2022} report that one adult in twenty-two lives with the disease in the region, that the region records the third-highest diabetes mortality rate in the world, at 111.3 per 100\,000, after the Eastern Mediterranean and South-East Asia Regions, and that International Diabetes Federation projections put the number of people affected at 55 million by 2045. Globally, \textcite{hinneh2023} recall that the burden of diabetes rose from 211.2 million to 476 million people between 1990 and 2017.

Type~2 diabetes poses an even sharper detection problem than hypertension in sub-Saharan Africa, since its diagnosis requires a biological measurement---fasting glucose or glycated haemoglobin (HbA1c)---and not merely a physical measuring device. This technical asymmetry has a direct consequence for the architecture of the proposed system: blood pressure can be acquired autonomously and frequently at home, whereas glycaemic measurements depend on consumables and a more demanding procedure, which requires the two data streams to be handled differently.

%==========================================================================
\subsection{The magnitude of the hypertension burden}\label{sec:hta-burden}

\subsubsection{Operational definition}

For population surveillance, the WHO defines hypertension as a systolic blood pressure (SBP) $\geq 140$~mmHg, or a diastolic blood pressure (DBP) $\geq 90$~mmHg, or use of antihypertensive medication, in adults aged 30--79 years. The report further distinguishes five indicators whose denominators differ and which must not be confused (\cref{tab:who-indicators}).

\begin{table}[htbp]
\centering\small
\caption{WHO surveillance indicators for hypertension (adapted from \cite{who2023hypertension}, Table~1).}\label{tab:who-indicators}
\begin{tabularx}{\textwidth}{@{}>{\raggedright\arraybackslash}p{4.2cm} X >{\raggedright\arraybackslash}p{4.2cm}@{}}
\toprule
\textbf{Indicator} & \textbf{Numerator} & \textbf{Denominator} \\
\midrule
Hypertension prevalence & SBP $\geq$ 140 or DBP $\geq$ 90 or on treatment & all adults 30--79 \\
Uncontrolled hypertension prevalence & SBP $\geq$ 140 or DBP $\geq$ 90 & all adults 30--79 \\
Diagnosis coverage & hypertensive with a prior diagnosis & hypertensive adults 30--79 \\
Treatment coverage & hypertensive on treatment & hypertensive adults 30--79 \\
Effective treatment coverage & on treatment and SBP $<$ 140 and DBP $<$ 90 & hypertensive adults 30--79 \\
\bottomrule
\end{tabularx}
\end{table}

Two points deserve attention. First, the prevalence of uncontrolled hypertension is not the complement of effective treatment coverage: its denominator is all adults, and it includes the undiagnosed, the diagnosed but untreated, and the treated but uncontrolled. Second, a patient whose blood pressure is normal without treatment does not fall into the ``effective treatment'' category, which requires medication. This definition is the one used throughout this thesis; it differs from the 130/80~mmHg threshold adopted by some professional societies, which explains part of the discrepancies between publications.

The WHO also introduces the notion of \emph{elevated systolic blood pressure}, defined as any value above the theoretical minimum-risk level, estimated at 110--115~mmHg. Distinct from hypertension, this notion serves comparative risk assessment and explains why attributable-mortality figures concern a wider population than the hypertensive.

\subsubsection{The global level as a reference}

The WHO global report \parencite{who2023hypertension} establishes that the number of hypertensive adults aged 30--79 rose from 650 million in 1990 to 1.3 billion in 2019, that the age-standardised prevalence stands at 33\% (uncertainty interval 32--35), and that about 78\% of those affected live in low- or middle-income countries (the trend is detailed in \cref{sec:transition}).

The distribution by sex and age is worth reporting, as it informs the design of alert thresholds. Global prevalence is slightly higher in men (34\%) than in women (32\%), but this gap is entirely driven by younger age groups: between 30 and 49 years, prevalence is 19\% in women versus 24\% in men, whereas between 50 and 79 years it reaches 49\% in both sexes. In other words, the female advantage disappears around age fifty, which coincides with the age range where \textcite{hinneh2023} place the maximal risk of comorbidity in diabetic patients.

The report above all establishes that elevated systolic blood pressure is responsible for 10.8 million preventable deaths per year and a burden of 235 million DALYs. In a comparison of eighty-seven behavioural, environmental, occupational and metabolic risk factors, it is the leading risk factor for premature death worldwide, ahead of smoking and hyperglycaemia. Hypertension is thus not one risk factor among others: it is the first of them.

The economic cost is also documented. The WHO notes that affected individuals bear direct medical expenses and lose wages, often during their most productive years, which can impoverish entire families; that hospital and outpatient care following myocardial infarction and stroke weighs on health systems; and that, according to one cited estimate, the economic benefits of improved hypertension treatment programmes exceed their costs by a ratio of about eighteen to one. Raising the global proportion of controlled patients to 50\% would avert 76 million deaths between 2023 and 2050.

\subsubsection{The public policy framework}

These data fit within a normative framework that provides this work with quantified targets. In 2013 the Sixty-sixth World Health Assembly adopted a monitoring framework with nine voluntary global targets, including target~6, which aims for a 25\% relative reduction in hypertension prevalence, and target~8, which provides that at least 50\% of eligible people---aged 40 years and over with a ten-year cardiovascular risk $\geq 20\%$---receive drug therapy and counselling, including glycaemic control, to prevent heart attacks and strokes. Target~8 is the first normative formulation that explicitly associates blood pressure control and glycaemic control within a single objective, and it thus provides direct support for the integrated approach defended here. The whole is further aligned with Sustainable Development Goal target 3.4, which aims at a one-third reduction in premature mortality from NCDs.

\subsubsection{The African Region}

The regional contrast is marked. In the WHO African Region, the number of hypertensive adults aged 30--79 rose from 45 million in 1990 to 106 million in 2019, a factor of 2.36, higher than the global doubling. Age-standardised prevalence reaches 36\% (33--38), three points above the global average.

\textcite{gafane2025}, in a narrative review based on searches in MEDLINE, ERIC, PubMed, EBSCOhost, CINAHL, Health Source, Academic Search Complete and African Journals Online, supplemented by grey literature and covering publications from January 2017 to June 2023 (forty-four studies retained: forty on the hypertension profile and four on community-based strategies), put forward prevalences for sub-Saharan Africa of 48\% in women (95\% CI: 42--54) and 34\% in men (29--39) in 2019.

This estimate cannot however be retained as it stands. The article presents these same two values, 48\% and 34\%, in two different places and for two different quantities, citing the same source. In the introduction they denote prevalence. In the body of the article they denote the diagnosis rate: it states that, according to NCD-RisC, in women diagnosis, treatment and control stood at 48\%, 29\% and 13\%, and that the same trend was observed in men, with 34\% diagnosed hypertension and even lower treatment (22\%) and control (9\%) rates.%
\verify{Quotation back-translated from the French draft; replace with the verbatim English wording of \textcite{gafane2025}.}
That two values coincide exactly, for both sexes and for two distinct indicators, is statistically implausible: one of the two passages contains a transcription error. The second use is moreover consistent with the 43\% diagnosis coverage that the WHO reports for the African Region, with a gap favouring women. We therefore use this source for the care cascade and not for prevalence.

The most defensible estimate drawn from real surveys remains that of the meta-analysis of \textcite{ataklte2015}, which aggregated thirty-three surveys conducted between 2000 and 2013 on more than 110\,000 adults with a mean age of forty years. It establishes a prevalence of 30\% (95\% CI: 27--34), with considerable dispersion between studies, from 15\% to 70\%. The authors show that this dispersion is largely explained by the age structure of the samples and provide a gradient that informs any subsequent comparison (\cref{fig:age-gradient}).

\begin{figure}[htbp]
\centering
\begin{tikzpicture}
\begin{axis}[width=0.65\textwidth,height=6cm,xlabel={Age (years)},ylabel={Predicted prevalence (\%)},
  xmin=25,xmax=65,ymin=0,ymax=50,xtick={30,40,50,60},grid=major,grid style={gray!25},
  nodes near coords,nodes near coords align={above},every node near coord/.append style={font=\footnotesize}]
\addplot[thick,cAfrica,mark=*] coordinates {(30,16)(40,26)(50,35)(60,44)};
\end{axis}
\end{tikzpicture}
\caption{Age gradient of predicted hypertension prevalence in sub-Saharan Africa (values from \textcite{ataklte2015}).}\label{fig:age-gradient}
\end{figure}

This gradient grounds a simple methodological requirement: the mean age of a sample must be known before any comparison of prevalences, otherwise what is due to demography is attributed to geography or time. It also grounds, for the rest of this work, the need to stratify by age the alert thresholds of the designed system: a detection model calibrated on a population of mean age forty will produce a markedly different false-positive rate on a cohort of sixty-year-olds.

\subsubsection{The specific situation of Central Africa}

This point is generally not addressed in the secondary literature, yet it directly conditions the relevance of the intended deployment in Cameroon. Table~1 of \textcite{gafane2025} lists the most recent national and regional surveys, published between 2017 and 2023, for each country of the sub-region (\cref{tab:central-africa}).

\begin{table}[htbp]
\centering\footnotesize
\caption{Most recent hypertension surveys in Central Africa (adapted from \cite{gafane2025}, Table~1). Awareness, treatment and control are expressed as in the source.}\label{tab:central-africa}
\begin{tabularx}{\textwidth}{@{}l >{\raggedright\arraybackslash}X c r r r r r@{}}
\toprule
\textbf{Country} & \textbf{Instrument} & \textbf{Year} & \textbf{$n$} & \textbf{Prev.} & \textbf{Aware.} & \textbf{Treat.} & \textbf{Control} \\
\midrule
Burundi & hospital screening & 2018 & 353 & 16.7\% & -- & -- & -- \\
Cameroon & May Measurement Month & 2019 & 30\,167 & 20.8\% & 29.9\% & 24.0\% & 11.2\% \\
Central African Rep. & WHO STEPS & 2017 & 3\,265 & 42.1\% & 30.2\% & -- & -- \\
Comoros & national survey & 2019 & 2\,602 & 38.4\% & 47.9\% & 45.4\% & 30.2\% \\
DR Congo & May Measurement Month & 2019 & 29\,857 & 25.5\% & 33.1\% & 23.2\% & 11.9\% \\
Congo (Rep.) & May Measurement Month & 2019 & 3\,157 & 33.5\% & 42.6\% & 37.3\% & 23.3\% \\
Gabon & retrospective study & 2013--15 & 2\,104 & 13.5\% & -- & -- & -- \\
Chad, Eq.\ Guinea, S\~ao Tom\'e & no data & -- & -- & -- & -- & -- & -- \\
\bottomrule
\end{tabularx}
\end{table}

Three lessons emerge. First, dispersion is very large, from 13.5\% in Gabon to 42.1\% in the Central African Republic, which reflects the diversity of instruments---opportunistic screening, probabilistic STEPS survey, retrospective study---as much as real differences. Second, three countries of the sub-region have no published data over the period, which illustrates the weakness of documentary coverage. Third, and most important for what follows, hypertension prevalence in the general population of Central Africa does not appear higher than in other sub-Saharan sub-regions: the Cameroonian and Congolese values lie in the lower range. This finding will be decisive when we examine, in \cref{sec:central-estimates}, the claim that Central Africa has the highest prevalence of hypertension among diabetics on the continent.

\subsubsection{Determinants}

\textcite{gafane2025} link the rise of hypertension in the region to the shift from a traditional to a westernised lifestyle: diets rich in salt, sugar and fat, low fruit and vegetable intake, sedentary behaviour, smoking, alcohol consumption and obesity, to which are added socio-demographic factors such as education and income, as well as psychosocial factors. The authors also mention more specific pathophysiological leads, still under investigation: sodium handling and volume overload, endothelial function, circadian blood pressure variation profiles---the so-called nocturnal dipping phenomenon---and the role of human immunodeficiency virus (HIV) infection.

This last point deserves emphasis for the Cameroonian context. The WHO \parencite{who2023hypertension} indicates that hypertension is the most common cardiovascular risk factor among people living with HIV, and that worldwide about 27\% of people on antiretroviral therapy are hypertensive. The ageing of this population, combined with the epidemiological transition, makes it a doubly exposed group whose care follows the same logic of prolonged chronic follow-up.

Geographically, the authors observe that the historical rural--urban gap is narrowing, rural populations progressively adopting the same exposures, while retaining a higher proportion of undiagnosed hypertension owing to lower access to health services. The two large recent multinational studies, SevenCEWA and H3Africa AWI-Gen, differ precisely on this point: some sites include exclusively urban populations, others exclusively rural ones, which precludes direct comparison of their results.

The nature of these statements must finally be qualified: they are observational associations and mechanistic hypotheses formulated in a narrative review, not established causal relationships. The authors themselves acknowledge that important gaps remain in understanding the mechanisms underlying blood pressure elevation over the life course in these populations, and conclude that further research is needed to inform preventive and therapeutic interventions.

%==========================================================================
\subsection{Failure of the care cascade}\label{sec:cascade}

Here lies the strongest justification for the present work. The problem in sub-Saharan Africa is not primarily the frequency of hypertension, barely above the global average in standardised terms, but the fact that the chain of care loses patients at every stage.

\subsubsection{The cascade model}

The WHO structures this indicator in four nested levels, all reported to the same denominator, the set of hypertensive adults aged 30--79:
\begin{equation}
  \text{Hypertensive} \supseteq \text{Diagnosed} \supseteq \text{Treated} \supseteq \text{Controlled}.
\end{equation}

\begin{table}[htbp]
\centering\small
\caption{Hypertension care cascade by WHO region, adults aged 30--79, 2019 (adapted from \cite{who2023hypertension}, Table~2).}\label{tab:cascade}
\begin{tabularx}{\textwidth}{@{}X c c c c@{}}
\toprule
\textbf{Indicator} & \textbf{African Region} & \textbf{World} & \textbf{Americas} & \textbf{Europe} \\
\midrule
Age-standardised prevalence & 36\% & 33\% & 35\% & 37\% \\
Diagnosis coverage & 43\% & 54\% & 70\% & 66\% \\
Treatment coverage & 27\% & 42\% & 60\% & 53\% \\
Effective treatment (control) & 12\% & 21\% & 36\% & 26\% \\
\bottomrule
\end{tabularx}
\end{table}

\begin{figure}[htbp]
\centering
\begin{tikzpicture}[x=0.08cm,y=1cm,every node/.style={font=\footnotesize}]
  % bar widths are proportional to coverage (100 = all hypertensive adults)
  \foreach \w/\lab/\col/\yy in {100/{Hypertensive (100\%)}/cGrey/3,
                                43/{Diagnosed (43\%)}/cRef/2,
                                27/{Treated (27\%)}/cRef!70!cAfrica/1,
                                12/{Controlled (12\%)}/cAfrica/0}{
    \fill[\col!80] (-\w/2,\yy-0.35) rectangle (\w/2,\yy+0.35);
    \node[anchor=east] at (-53,\yy) {\lab};
  }
  \node[anchor=west] at (53,2.5) {$-57\%$};
  \node[anchor=west] at (53,1.5) {$-37\%$ of diagnosed};
  \node[anchor=west] at (53,0.5) {$-56\%$ of treated};
\end{tikzpicture}
\caption{Care cascade of hypertension in the WHO African Region, 2019 (data from \cref{tab:cascade}). Bar widths are proportional to coverage; losses are conditional on the previous stage.}\label{fig:funnel}
\end{figure}

The comparison with the European Region is instructive: its prevalence (37\%) is higher than that of the African Region (36\%), but its control rate is more than twice as high. The African deficit is therefore not epidemiological, it is organisational. This observation underlies the present work: if the problem were mainly overexposure to risk, the response would fall under primary prevention; since it is a failure of the chain of care, the response falls under information and monitoring systems.

The income gradient confirms this diagnosis. Treatment coverage is 26\% in low-income countries versus 58\% in high-income countries, and control 11\% versus 36\%. The WHO notes however that some middle-income countries achieve treatment coverage comparable to the best-performing high-income ones, which shows that wealth is not an absolute constraint.

\subsubsection{Decomposition of conditional losses}

Since the four levels share the same denominator and are nested, the conditional probabilities of loss at each stage can be derived. For the African Region:
\begin{align}
  P(\text{not diagnosed}) &= 1 - 0.43 = \mathbf{57\,\%},\\
  P(\text{not treated}\mid\text{diagnosed}) &= 1 - \frac{0.27}{0.43} = \mathbf{37\,\%},\\
  P(\text{not controlled}\mid\text{treated}) &= 1 - \frac{0.12}{0.27} = \mathbf{56\,\%}.
\end{align}

These three ratios are computed from point estimates with uncertainty intervals---39 to 46\% for diagnosis, 24 to 30\% for treatment, 9 to 14\% for control---and must therefore be read as orders of magnitude. They nevertheless deliver the most important result of this section: attrition is not concentrated at screening, it is spread along the entire chain, and the weakest link in conditional terms is the last one, since more than one treated patient in two does not reach the blood pressure target.

By comparison, the same decomposition applied to the Americas gives 30\% undiagnosed, 14\% of the diagnosed untreated and 40\% of the treated uncontrolled (\cref{fig:losses}). The largest gap between the two regions thus concerns the first two links, but the third remains deficient everywhere, making it a structural problem of hypertension care and not an African specificity.

\begin{figure}[htbp]
\centering
\begin{tikzpicture}
\begin{axis}[ybar,width=0.85\textwidth,height=6.5cm,bar width=16pt,ymin=0,ymax=70,
  ylabel={Conditional loss (\%)},xtick={1,2,3},
  xticklabels={Not diagnosed,{Not treated $\mid$ diagnosed},{Not controlled $\mid$ treated}},
  xticklabel style={font=\footnotesize,text width=3cm,align=center},
  enlarge x limits=0.25,legend pos=north east,
  nodes near coords,every node near coord/.append style={font=\footnotesize}]
\addplot[fill=cAfrica!80,draw=cAfrica] coordinates {(1,57)(2,37)(3,56)};
\addplot[fill=cRef!80,draw=cRef] coordinates {(1,30)(2,14)(3,40)};
\legend{African Region,Americas}
\end{axis}
\end{tikzpicture}
\caption{Conditional losses along the hypertension care cascade, derived from \cref{tab:cascade}.}\label{fig:losses}
\end{figure}

This observation has a direct consequence for the design of the proposed system. A device limited to screening, that is, measuring and alerting, would act only on the first link and leave the next two intact. A longitudinal telemonitoring system, which accompanies the already-diagnosed patient between consultations, acts instead on the second and third links, which together concentrate the majority of avoidable losses in conditional terms.

\subsubsection{The gender dimension}

The WHO \parencite{who2023hypertension} documents a systematic gap between men and women. Globally, the cascade breaks down as shown in \cref{tab:gender}.

\begin{table}[htbp]
\centering\small
\caption{Global care cascade by sex (\cite{who2023hypertension}).}\label{tab:gender}
\begin{tabular}{@{}lcc@{}}
\toprule
\textbf{Stage} & \textbf{Men} & \textbf{Women} \\
\midrule
Diagnosed & 49\% & 59\% \\
Treated & 38\% & 47\% \\
Controlled & 18\% & 23\% \\
\bottomrule
\end{tabular}
\end{table}

In all regions and in all countries but four, hypertensive women are more likely to be treated. The same pattern holds for the probability of being diagnosed and of reaching blood pressure control. This asymmetry, confirmed by \textcite{gafane2025} for sub-Saharan Africa, has a practical implication for the designed system: men are the population in which the marginal gain from automated reminders and remote follow-up is highest, and the recruitment of an evaluation cohort must take care not to reproduce the female participation bias observed in most screening surveys.

\subsubsection{Severity of undetected cases}

A rarely cited result of the WHO report reinforces the argument for early detection. Among people with uncontrolled hypertension who have never been diagnosed, nearly 30\% worldwide have an SBP $\geq 160$~mmHg or a DBP $\geq 100$~mmHg. These values far exceed the 140/90~mmHg diagnostic threshold and, in the words of the report, indicate an urgent need for treatment.

In other words, the undiagnosed population does not consist of borderline or doubtful cases: nearly a third of it warrants immediate treatment initiation. For effective primary-care management, the WHO recommends screening every adult patient, immediate treatment initiation when confirmed pressure exceeds 160/100~mmHg, and a rapid recheck within one to four weeks for all other patients with elevated pressure. This short-term recheck protocol is precisely the kind of task a telemonitoring system can automate, whereas it relies today on the patient's memory and on his or her ability to travel a second time.

\subsubsection{The trajectory towards the 2025 target}

The global prevalence of uncontrolled hypertension fell from 29\% in 2010 to 26\% in 2019. The global voluntary target, a 25\% relative reduction from the 2010 baseline, implies reaching 21\% in 2025.%
\verify{29\% $\times$ 0.75 $\approx$ 21.8\%, not 21\%. Check the exact 2010 baseline and 2025 target value in the WHO report.}
The WHO concludes that this rate of decline is insufficient and that no WHO region or income group is on track to meet the target, only a handful of countries being on the right trajectory.

The report specifies, however, that the gap would be closed if about half of already-diagnosed adults were effectively treated. This is decisive for the positioning of the present work: the global target does not primarily require widening screening, but retaining and effectively controlling the patients already identified. This is exactly the segment of the cascade that a longitudinal telemonitoring system addresses.

\subsubsection{Confirmation from field surveys}

Real surveys corroborate this picture, at levels that vary with the method used (\cref{tab:field}).

\begin{table}[htbp]
\centering\footnotesize
\caption{Awareness, treatment and control of hypertension in sub-Saharan Africa according to source.}\label{tab:field}
\begin{tabularx}{\textwidth}{@{}>{\raggedright\arraybackslash}p{3.3cm} >{\raggedright\arraybackslash}X r c c c@{}}
\toprule
\textbf{Source} & \textbf{Design} & \textbf{$n$} & \textbf{Awareness} & \textbf{Treatment} & \textbf{Control} \\
\midrule
\textcite{ataklte2015} & meta-analysis, 33 surveys 2000--2013 & $>110\,000$ & 27\% (23--31) & 18\% (14--22) & 7\% (5--8) \\
NCD-RisC, women & modelling & -- & 48\% & 29\% & 13\% \\
NCD-RisC, men & modelling & -- & 34\% & 22\% & 9\% \\
MMM 2018--2019 & opportunistic screening & -- & 40.5\% & 32.1\% & 15.4\% \\
SevenCEWA & cross-sectional, 7 communities & 3\,549 & 57.2\% & -- & 47.3\% of treated \\
H3Africa AWI-Gen & cross-sectional, 4 countries & 10\,696 & 47.7\% & -- & -- \\
\bottomrule
\end{tabularx}
\end{table}

The two multinational studies deserve a description, as they provide the most recent and methodologically most homogeneous data of the region. SevenCEWA covered 3\,549 participants aged 18 and over, recruited in seven communities in Kenya, Nigeria, Tanzania and Uganda, with mean age 39.7 years (SD 15.4); a quarter of them (25.4\%) were hypertensive, awareness reached 57.2\%, slightly over half of identified individuals received treatment, and 47.3\% of treated individuals reached the blood pressure target. H3Africa AWI-Gen covered 10\,696 participants aged 40 to 60 years, mean age 49.9 years, recruited in Burkina Faso, Ghana, Kenya and South Africa; prevalence was higher there, consistent with the age gradient of \cref{fig:age-gradient}, and overall awareness was 47.7\%, with strong variation between sites---from 41.6 to 54.1\% in South Africa versus 33\% in Burkina Faso.

Two observations are in order. First, awareness varies twofold, from 27\% to 57\%, depending on whether one considers a meta-analysis of older surveys or a recent community study conducted on selected sites: this gap measures method and period as much as epidemiological reality. Second, the 47.3\% from SevenCEWA is not comparable to the WHO's 12\%, because its denominator differs: it concerns treated patients only, not all hypertensive individuals. \textcite{gafane2025} also stress that the Kenyan sites included in the two studies do not agree with each other, which they attribute to convenience sampling and collection methods, and they conclude that standardised methods and nationally representative surveys are needed.

\subsubsection{Structural causes}

\textcite{gafane2025} distinguish two levels of obstacles. At the \emph{individual} level: low health literacy, particularly regarding hypertension and cardiovascular risk; socio-economic factors affecting the affordability of health insurance and medication; out-of-pocket costs; and distance to care facilities. At the \emph{system} level: availability of medicines, limited access to health facilities, insufficient numbers, training and knowledge of health personnel, and uneven political will in implementing interventions proven elsewhere. The authors cite a study of the Soweto population showing that poor primary-care service quality, notably unpleasant interactions with providers, is in itself a factor of consultation avoidance.

To these obstacles is added \emph{loss to follow-up}, which the WHO \parencite{who2023hypertension} identifies as a major barrier to blood pressure control and which, it indicates, often affects half of patients or more in resource-limited settings. This is precisely the phenomenon that telemonitoring aims to correct, and it is on this indicator that the evaluation of the proposed system should primarily focus.

\subsubsection{Task shifting as an organisational response}

Faced with these constraints, the best-documented response in sub-Saharan Africa is \emph{task shifting}, that is, the delegation of part of care responsibilities to community health workers or peer patients. \textcite{gafane2025} review several evaluations and conclude that, depending on design and implementation, most forms of task shifting improve linkage to health services and medication adherence but yield inconsistent results on long-term blood pressure control. The HealthRise programme, implemented in Brazil, India, South Africa and the United States, illustrates this ambivalence: the South African sites, two of them rural, were the most difficult, owing to a short implementation period, unfavourable socio-demographic factors, low patient literacy and high unemployment. The authors also note that task-shifting studies suffer from high attrition, uneven methodological quality, parallel rather than integrated models, and frequent absence of comparison groups.

This finding is directly relevant to our subject. It establishes that the central question is not to add human resources to the chain, but to maintain continuity of information between actors and over time---a question of information systems before being a question of staffing.

\subsubsection{The normative vacuum}

The scoping review of \textcite{nagavci2022}, conducted according to the PRISMA-ScR extension and based on searches in Ovid MEDLINE, the GIN International Guideline Library, Guideline Central and the TRIP database, supplemented by direct e-mail consultation of the ministries of health of the forty-nine countries of the region and of forty-three public health experts, identified sixty-six clinical practice guidelines on hypertension or diabetes, developed in twenty-three of the forty-nine sub-Saharan countries.

Among them, only two met the two minimal criteria set by the Institute of Medicine for a trustworthy guideline: being based on a systematic review of evidence and including a formal assessment of the certainty of that evidence. Both concern insulin therapy in diabetes and come from the East African Diabetes Study Group. Their quality, assessed with the AGREE-II instrument (twenty-three items in six domains), is shown in \cref{tab:agree}.

\begin{table}[htbp]
\centering\small
\caption{AGREE-II domain scores of the two guidelines meeting the minimal criteria \parencite{nagavci2022}.}\label{tab:agree}
\begin{tabular}{@{}lcc@{}}
\toprule
\textbf{AGREE-II domain} & \textbf{Bahendeka 2018} & \textbf{Bahendeka 2019} \\
\midrule
Scope and purpose & 78\% & 67\% \\
Clarity of presentation & 75\% & 97\% \\
Applicability & 81\% & 83\% \\
Stakeholder involvement & 61\% & 42\% \\
Rigour of development & 56\% & 45\% \\
Editorial independence & 54\% & 42\% \\
\midrule
Overall score & 58.33\% & 50\% \\
\bottomrule
\end{tabular}
\end{table}

No guideline on hypertension met these criteria. This finding must be stated precisely so as not to be over-interpreted: it does not mean that countries of the region lack hypertension care protocols---thirteen documents on this condition and forty on multiple conditions were identified before being excluded---but that none met the methodological criteria of an evidence-based guideline. The sixty-four excluded documents were developed in about twenty countries and published between 1995 and 2021, all in English except one in Amharic, one in French and one in Portuguese.

These results agree with those of an earlier study, cited by the authors, according to which only twenty-six of sixty-two African countries reported using guidelines for the diagnosis and management of hypertension, and with those of another review that identified thirty guidelines developed in thirteen countries. The authors attribute this scarcity to the relatively low research output of the region and to strong capacity disparities between countries, South Africa having the highest capacity, and the Central African Republic, The Gambia, Lesotho and Zambia the lowest. They finally note that developing a trustworthy guideline requires up to three years and considerable financial and human resources, and therefore recommend adopting or adapting existing international guidelines, using GRADE-ADOLOPMENT or ADAPTE methods, rather than full development.

This normative vacuum has a direct consequence for the architecture of the proposed system: alert thresholds and rules of conduct cannot be derived from a methodologically validated Cameroonian national guideline, which does not exist. They will therefore be explicitly anchored to the WHO 2021 guideline on pharmacological treatment of hypertension in adults and to the HEARTS technical package, a choice to be justified in the design chapter.

\subsubsection{Cameroonian data}\label{sec:cameroon-data}

Cameroon, the deployment context of the proposed system, presents the same features, with strong heterogeneity between regions and methods (\cref{tab:cameroon}).

\begin{table}[htbp]
\centering\footnotesize
\caption{Hypertension studies in Cameroon.}\label{tab:cameroon}
\begin{tabularx}{\textwidth}{@{}>{\raggedright\arraybackslash}p{2.6cm} >{\raggedright\arraybackslash}X r c c >{\raggedright\arraybackslash}p{2cm} >{\raggedright\arraybackslash}p{2cm}@{}}
\toprule
\textbf{Study} & \textbf{Design} & \textbf{$n$} & \textbf{Prev.} & \textbf{Aware.} & \textbf{Treatment} & \textbf{Control} \\
\midrule
\textcite{dzudie2021} & MMM 2019, 5 cities, $\geq 18$ y & 30\,167 & 20.8\% & 29.9\% & 24.0\% & 11.2\% \\
\textcite{lemogoum2018} & cross-sectional, multilevel clusters, Far North & 889 & 37.8\% (std.) & 18.6\% & 29.3\% of aware & 16.3\% of treated \\
\bottomrule
\end{tabularx}
\end{table}

The prevalence gap between these two studies---20.8\% versus 37.8\%---illustrates exactly the comparability problem set out at the outset: the first relies on opportunistic screening in an urban setting, with an explicitly flagged selection bias and a single-visit measurement, the second on probabilistic cluster sampling including a rural site, with age standardisation.

The study of \textcite{lemogoum2018}, conducted in Maroua (urban) and Tokombéré (rural) between November 2014 and May 2015 among 889 adults aged 18 and over, 41.5\% of them women, allows a particularly informative calculation. Chaining the successive conditional proportions:
\begin{equation}
  P(\text{controlled}) = 0.186 \times 0.293 \times 0.163 \approx 0.0089,
\end{equation}
that is, about 1\% of all hypertensive individuals of this region have effectively controlled blood pressure.

The contrast between settings is moreover \emph{inverted} depending on the stage considered (\cref{fig:rural-urban}), which is the most instructive result of this study for our purpose.

\begin{figure}[htbp]
\centering
\begin{tikzpicture}
\begin{axis}[ybar,width=0.9\textwidth,height=6.5cm,bar width=14pt,ymin=0,ymax=50,
  ylabel={Percentage (\%)},xtick={1,2,3,4},
  xticklabels={Prevalence,Awareness,{Treatment (of aware)},{Control (of treated)}},
  xticklabel style={font=\footnotesize,text width=2.4cm,align=center},
  enlarge x limits=0.18,legend pos=north east,
  nodes near coords,every node near coord/.append style={font=\footnotesize}]
\addplot[fill=cGrey!70,draw=cGrey] coordinates {(1,34.0)(2,17.9)(3,17.3)(4,22.6)};
\addplot[fill=cAfrica!80,draw=cAfrica] coordinates {(1,41.2)(2,19.5)(3,36.3)(4,4.2)};
\legend{Rural (Tokombéré),Urban (Maroua)}
\end{axis}
\end{tikzpicture}
\caption{Hypertension cascade in the Far North of Cameroon by setting (data from \textcite{lemogoum2018}).}\label{fig:rural-urban}
\end{figure}

Treatment is twice as frequent in the city, but control among treated patients there is five times lower. Such an inversion is explained neither by access to treatment, better in the city, nor by initial severity. It points to a problem of adherence, quality of longitudinal follow-up or dosing adequacy---exactly the class of failure that a telemonitoring system is likely to correct, and that a screening system would leave intact.

Finally, a meta-analysis on the burden of hypertension in Cameroon, aggregating population studies, was published in 2019 in the \emph{Journal of Hypertension}. Its full text could not be consulted, so its results are not used here; consulting it is a priority task for the final version of this chapter, this source being the most appropriate for a consolidated national estimate.

%==========================================================================
\subsection{The hypertension--diabetes comorbidity}\label{sec:comorbidity}

\subsubsection{Pathophysiological basis}

The association between hypertension and type~2 diabetes is not fortuitous. It rests on two complementary mechanisms (\cref{fig:loop}).

The first is shared risk factors: physical inactivity, unbalanced diet and obesity are common to both conditions.

The second is a direct pathophysiological link. Diabetes induces microvascular alterations, in particular a thickening and stiffening of small vessels, a phenomenon termed arteriosclerosis. Now, mean arterial pressure approximately obeys
\begin{equation}
  \mathrm{MAP} \approx Q \times R,
\end{equation}
where $Q$ denotes cardiac output and $R$ total peripheral vascular resistance. Increased arterial stiffness raises $R$ and, consequently, pressure at constant output. Conversely, hypertension accelerates micro- and macrovascular damage in the diabetic patient.

\begin{figure}[htbp]
\centering
\begin{tikzpicture}[node distance=1.0cm and 1.6cm,
  box/.style={draw,rounded corners,align=center,minimum height=0.9cm,font=\small,fill=cRef!10},
  arr/.style={-{Stealth[length=2.5mm]},thick}]
  \node[box,fill=cGrey!15] (risk) {Shared risk factors\\(inactivity, diet, obesity)};
  \node[box,below left=of risk] (dm) {Type 2 diabetes};
  \node[box,below right=of risk] (htn) {Hypertension\\(elevated MAP)};
  \node[box,below=2.2cm of risk] (vasc) {Microvascular damage,\\arterial stiffening $\Rightarrow R\uparrow$};
  \draw[arr] (risk) -- (dm);
  \draw[arr] (risk) -- (htn);
  \draw[arr] (dm) -- (vasc);
  \draw[arr] (vasc) -- node[right,font=\footnotesize]{$\mathrm{MAP}\approx Q\times R$} (htn);
  \draw[arr,cAfrica] (htn.south east) to[bend left=35] node[right,font=\footnotesize,text=cAfrica]{accelerates lesions} (vasc.east);
\end{tikzpicture}
\caption{Mechanisms linking type~2 diabetes and hypertension.}\label{fig:loop}
\end{figure}

\textcite{hinneh2023} note that, beyond neuropathic and nephropathic complications and cardiovascular disease, hypertension is the most frequent comorbidity and the leading cause of death in people with diabetes. They add that the comorbidity degrades blood pressure control and quality of life, increases management costs and complicates care needs, and that affected patients paradoxically have lower chances of receiving appropriate care before complications occur---the accumulation of needs leading to delayed care rather than intensification. This last mechanism is central to our purpose: it establishes that the comorbidity is not merely the sum of two risks, but an aggravating factor of the organisational failure described in \cref{sec:cascade}.

\subsubsection{Quantitative estimate}

The meta-analysis of \textcite{hinneh2023} is the first systematic quantification of the prevalence of hypertension among people with diabetes in Africa. Its protocol deserves detailing, as it conditions the scope of the results.

It relies on a predefined protocol conforming to PRISMA recommendations and registered with PROSPERO under reference CRD42021256221. The search strategy, built on the Joanna Briggs Institute ``population, concept, context'' approach, was run in PubMed, EMBASE and HINARI from inception to March 2023, without publication-year restriction. Inclusion criteria were adults with type~1 or type~2 diabetes, cross-sectional designs only, studies conducted in Africa and published in English, excluding gestational hypertension. The diagnostic threshold was 140/90~mmHg or 130/80~mmHg depending on the study.%
\verify{The French draft attributes both thresholds to the WHO. As stated in \cref{sec:hta-burden}, the WHO uses 140/90~mmHg; 130/80~mmHg comes from other professional societies. Wording adjusted; confirm against \textcite{hinneh2023}.}

The selection flow is as follows: 3\,807 records identified, 564 duplicates removed, 3\,243 titles and abstracts screened, 3\,126 excluded, 117 full texts assessed, to which were added three references from bibliography screening and five additional studies from an updated search, for forty-one studies finally included, totalling 144\,433 participants of mean age 58 years (SD 11), 56\% of them women, with study sample sizes ranging from 80 to 116\,726.

Risk of bias was assessed by two independent reviewers using the MMAT tool; thirty-six of forty-one studies scored 4 or 5, corresponding to high quality and low risk of bias. Pooling used a random-effects model following the DerSimonian--Laird procedure, with Clopper--Pearson intervals, and subgroup analyses by regional block, diabetes type and risk-of-bias score.

The pooled prevalence of hypertension among people with diabetes in Africa is 58.1\% (95\% CI: 52.0--63.2). Among type~2 diabetic patients alone, over twenty-four studies and 15\,476 participants, it rises to 60.8\% (53.4--67.7). Individual prevalences range from 21\% in Ghana to 92\% in South Africa.

The factors associated with hypertension in the included studies are increasing age, overweight and obesity, being employed, duration of diabetes, urban residence and male sex. One study reports that age multiplies the odds of hypertension more than sevenfold (OR 7.26; 95\% CI: 4.20--12.54), the risk being maximal in the 50--60 age range---the same range where the female advantage disappears in the general population (\cref{sec:hta-burden}).

The authors situate their result against two references: an earlier systematic review reporting a prevalence of hypertension in diabetics between 38.5\% and 80\% worldwide, and an estimate of 55.2\% for the general African population. The prevalence in diabetics therefore appears slightly higher than in the general population, the gap being larger in the northern, central and southern sub-regions.

\subsubsection{Sub-regional estimates and their fragility}\label{sec:central-estimates}

\begin{table}[htbp]
\centering\small
\caption{Sub-regional pooled prevalence of hypertension among diabetics in Africa, as published by \textcite{hinneh2023}.}\label{tab:subregional}
\begin{tabular}{@{}lccc@{}}
\toprule
\textbf{Sub-region} & \textbf{Pooled prevalence} & \textbf{95\% CI} & \textbf{Studies announced} \\
\midrule
Central Africa & 77.6\% & 53.0--91.4 & 2 \\
Southern Africa & 69.1\% & 59.8--77.1 & 10 \\
Eastern Africa & 53.0\% & 45.8--59.1 & 15 \\
Western Africa & 51.5\% & 41.8--61.1 & 10 \\
\bottomrule
\end{tabular}
\end{table}

\begin{figure}[htbp]
\centering
\begin{tikzpicture}
\begin{axis}[width=0.85\textwidth,height=6cm,xlabel={Pooled prevalence of hypertension among diabetics (\%)},
  xmin=30,xmax=100,ymin=0.5,ymax=5.5,ytick={1,2,3,4,5},
  yticklabels={Western Africa,Eastern Africa,Southern Africa,Central Africa,All Africa},
  grid=major,grid style={gray!20}]
\addplot[dashed,cGrey,no markers] coordinates {(58.1,0.5)(58.1,5.5)};
% point estimate with exact (asymmetric) 95% CI bounds: += upper offset, -= lower offset
\addplot[only marks,mark=*,mark size=2.5pt,cRef,error bars/.cd,x dir=both,x explicit,error bar style={thick}]
  coordinates {(51.5,1) += (9.6,0) -= (9.7,0)
               (53.0,2) += (6.1,0) -= (7.2,0)
               (69.1,3) += (8.0,0) -= (9.3,0)};
\addplot[only marks,mark=*,mark size=2.5pt,cAfrica,error bars/.cd,x dir=both,x explicit,error bar style={thick}]
  coordinates {(77.6,4) += (13.8,0) -= (24.6,0)};
\addplot[only marks,mark=square*,mark size=2.5pt,black,error bars/.cd,x dir=both,x explicit,error bar style={thick}]
  coordinates {(58.1,5) += (5.1,0) -= (6.1,0)};
\end{axis}
\end{tikzpicture}
\caption{Point estimates and 95\% confidence intervals of \cref{tab:subregional} (exact bounds); the dashed line marks the pooled African estimate. The Central African interval spans 38 points.}\label{fig:forest}
\end{figure}

North Africa (four studies) is omitted from the table for the reasons set out below. These sub-regional estimates call for six reservations, whose weight leads us not to ground the justification of this thesis on the value of 77.6\%, yet the most frequently reproduced in the secondary literature.

\paragraph{First reservation: heterogeneity is extreme.} The $I^2$ index measures the share of observed between-study variability attributable to real heterogeneity rather than sampling chance:
\begin{equation}
  I^2 = \frac{Q-(k-1)}{Q}\times 100,
\end{equation}
where $Q$ is Cochran's statistic and $k$ the number of studies. This index reaches 98\% in the text of the article and 99\% in its Figures~2 and~4, with between-study variances $\tau^2$ of 0.476 and 0.548. A value of this order means that almost all of the observed dispersion is not random---in other words, that the included studies do not measure the same quantity. The authors acknowledge that this heterogeneity persists after accounting for regional variation, that it may be explained by the diversity of diagnostic criteria used for both hypertension and diabetes, and that it limits the generalisability of the pooled estimate. Under these conditions, the point value of 58.1\% must be read as an order of magnitude and not as a precise estimate.

\paragraph{Second reservation: publication bias is probable.} The linear regression test for funnel-plot asymmetry is significant, with $t=-2.22$, 39 degrees of freedom and $p=0.032$. Note that the Methods section of the article announced a Peters test while the Results section reports a linear regression test, an internal inconsistency that calls for caution in interpreting this result.

\paragraph{Third reservation, the most important: the Central African estimate rests on two Cameroonian studies and likely on a classification error.} The authors announce a distribution of their forty-one studies as fifteen for East, ten for West, four for North, two for Central and ten for South Africa. Yet their own Table~1 lists five studies conducted in Central African countries (\cref{tab:central-studies}).

\begin{table}[htbp]
\centering\small
\caption{Central African studies listed in Table~1 of \textcite{hinneh2023}.}\label{tab:central-studies}
\begin{tabular}{@{}llrr@{}}
\toprule
\textbf{Study} & \textbf{Country} & \textbf{$n$} & \textbf{Prevalence} \\
\midrule
\textcite{kemche2020} & Cameroon & 109 & 86.2\% \\
\textcite{choukem2007} & Cameroon & 210 & 66.7\% \\
\textcite{mfeukeu2022} & Cameroon & 112 & 50.4\% \\
\textcite{dzudie2012} & Cameroon & 420 & 50.2\% \\
\textcite{ovono2011} & Gabon & 1\,300 & 35.0\% \\
\bottomrule
\end{tabular}
\end{table}

Pooling only the two high-prevalence studies on the logit scale:
\begin{align}
  \operatorname{logit}(0.862) &= \ln\frac{0.862}{0.138} = 1.832, \qquad
  \operatorname{logit}(0.667) = \ln\frac{0.667}{0.333} = 0.695,\\
  \bar\theta &= \frac{1.832+0.695}{2} = 1.264 \;\Rightarrow\;
  \operatorname{logit}^{-1}(1.264) = \frac{e^{1.264}}{1+e^{1.264}} \approx 0.78,
\end{align}
which closely reproduces the published value of 77.6\%. The weights assigned to these two studies in the global model, 2.2\% and 2.4\% respectively, being almost equal, the unweighted mean is a good approximation of the random-effects result.

The three low-prevalence studies, which total 1\,832 patients and whose values range from 35.0\% to 50.4\%, were therefore classified outside the ``Central Africa'' subgroup. The 77.6\% estimate thus appears to rest on 319 Cameroonian hospital patients, to the exclusion of 1\,832 other patients from the same sub-region whose prevalence is about half. This reconstruction, consistent with the published values, remains to be confirmed on supplementary Figure~S5 of the article, but it suffices to strip this figure of its status as an established fact. The published confidence interval, thirty-eight points wide, bears witness to this (\cref{fig:forest}): it is compatible with a true prevalence of 53\% as well as of 91\%.

This diagnosis is corroborated by independent data (\cref{tab:central-africa}): general-population hypertension prevalence in Central African countries is not higher than in other sub-Saharan sub-regions, recent surveys giving 20.8\% in Cameroon, 25.5\% in the Democratic Republic of the Congo and 33.5\% in the Republic of the Congo. Nothing in the general epidemiology of the sub-region supports the idea that Central Africa is an extreme point of the continent for hypertension among diabetics.

The authors themselves acknowledge this fragility twice. They note that Central African countries were probably underrepresented, which could affect prevalence estimates, notably owing to the restriction to English-language publications while several countries of this sub-region publish in French. They also indicate that the small number of studies identified for the North and Central African regions highlights the limited research on this subject in these sub-regions.

\paragraph{Fourth reservation: the article contains internal inconsistencies on North Africa and its mapping.} Three different values are given for this sub-region: 63.4\% in the abstract and Results section, 53.4\% in the caption of Figure~3, and 65.5\% on the map label of that same figure. More broadly, the labels on the map contradict the text for four of five sub-regions (\cref{tab:map}).

\begin{table}[htbp]
\centering\small
\caption{Text versus map labels in \textcite{hinneh2023}.}\label{tab:map}
\begin{tabular}{@{}lcc@{}}
\toprule
\textbf{Sub-region} & \textbf{Text} & \textbf{Map label} \\
\midrule
Eastern Africa & 53.0\% & 50.7\% \\
Western Africa & 51.5\% & 55.7\% \\
Northern Africa & 63.4\% & 65.5\% \\
Southern Africa & 69.1\% & 74.5\% \\
Central Africa & 77.6\% & 77.6\% \\
\bottomrule
\end{tabular}
\end{table}

The confidence interval for Eastern Africa also differs between the text (45.8 to 59.1) and the caption (46.8 to 59.1). Only Central Africa remains consistent across media. Note moreover that the published interval for North Africa, 37.1 to 69.1 for a point estimate of 63.4\%, is suspiciously asymmetric, the point estimate lying close to the upper bound.

\paragraph{Fifth reservation: the nature of the studies limits the scope of conclusions.} All are cross-sectional, which precludes causal inference: it cannot be established from these data whether hypertension precedes or follows diabetes, nor in what proportion. They are also mostly hospital-based, hence unrepresentative of the general diabetic population. Diagnostic thresholds vary: 80\% of studies use 140/90~mmHg, 10\% use 130/80~mmHg and 10\% do not report it. Diabetes criteria vary similarly, some studies using an HbA1c threshold of 5.7\%, others 6.5\% or 7\%, others a fasting glucose of 5 to 7~mmol/L. Finally, twenty-four studies (59\%) concern type~2 diabetes alone, seventeen (41\%) types 1 and 2 combined, and none type~1 alone, which precludes any conclusion specific to that form.

\paragraph{Sixth reservation: the direction of the measure must be recalled.} The 58.1\% value quantifies the frequency of hypertension \emph{among diabetics}. It says nothing about the inverse frequency, the proportion of diabetics among the hypertensive, which none of the sources used here provides. This asymmetry has a direct practical consequence: it justifies systematically measuring blood pressure in every diabetic patient enrolled in the system, but does not symmetrically justify measuring blood glucose in every hypertensive patient. The system architecture should reflect this asymmetry rather than postulate a symmetry that the data do not establish.

\subsubsection{Cameroonian data on the comorbidity}

\begin{table}[htbp]
\centering\footnotesize
\caption{Studies of hypertension among diabetics in Cameroon.}\label{tab:cameroon-comorbidity}
\begin{tabularx}{\textwidth}{@{}>{\raggedright\arraybackslash}p{2.6cm} >{\raggedright\arraybackslash}X r c c@{}}
\toprule
\textbf{Study} & \textbf{Design} & \textbf{$n$} & \textbf{Mean age} & \textbf{Result} \\
\midrule
\textcite{kemche2020} & hospital, Yaoundé (Central Hospital and Etoug-Ebé Baptist Centre), types 1 and 2 & 109 & 55 y & 86.2\% \\
\textcite{choukem2007} & hospital, types 1 and 2 & 210 & 56.6 y & 66.7\% \\
\textcite{mfeukeu2022} & hospital, type 2 & 112 & 58 y (median) & 50.4\% \\
\textcite{dzudie2012} & hospital, type 2 & 420 & 55.9 y & 50.2\% \\
\textcite{katte2014} & community screening, Bafoussam & 1\,702 & -- & 3.9\% (M) / 5.0\% (F) \\
\bottomrule
\end{tabularx}
\end{table}

The study of \textcite{katte2014} deserves to be presented accurately, as it is frequently miscited. Conducted in Bafoussam, capital of the West Region, in November 2012, it recruited 1\,702 adults---967 women (56.8\%) and 735 men---self-selected following radio announcements, across five screening sites. Its age-standardised results are shown in \cref{tab:katte}.

\begin{table}[htbp]
\centering\small
\caption{Age-standardised results of \textcite{katte2014}.}\label{tab:katte}
\begin{tabular}{@{}lcc@{}}
\toprule
\textbf{Condition} & \textbf{Men} & \textbf{Women} \\
\midrule
Isolated hypertension & 40.4\% (34.7--46.1) & 23.8\% (20.4--27.2) \\
Isolated diabetes & 3.3\% (1.5--5.1) & 5.6\% (3.5--7.7) \\
Hypertension and diabetes combined & 3.9\% (2.6--5.2) & 5.0\% (3.5--6.5) \\
\bottomrule
\end{tabular}
\end{table}

The authors' message is that co-occurrence of the two conditions is as frequent as isolated diabetes, which amounts to saying that about one diabetic in two is also hypertensive in this population ($3.9/(3.3+3.9)\approx 54\%$ in men and $5.0/(5.6+5.0)\approx 47\%$ in women)---a result remarkably convergent with the 58.1\% estimate of \textcite{hinneh2023}, although obtained by an entirely different design. The authors themselves caution against any extrapolation, self-selection making the sample unrepresentative of the city.

The gap between the Cameroonian hospital series, ranging from 50\% to 86\%, and community screening data is explained above all by recruitment mode: patients followed in referral hospitals are older, with longer disease duration, more severe and more often already treated. Note also that the two highest hospital prevalences come from the smallest series (109 and 210 patients), whereas the largest series (420) has the lowest value, a configuration that is precisely what publication bias produces. None of these sources is representative of the Cameroonian diabetic population, which is why they are presented together rather than aggregated.

\subsubsection{What grounds combined monitoring}

If the 77.6\% value cannot be retained, on what then rests the decision to measure blood pressure and blood glucose jointly? On five convergent elements, all solidly established.

\textbf{First}, the pooled estimate of 58.1\%, based on forty-one studies and 144\,433 participants, whose order of magnitude remains robust despite the heterogeneity noted: more than one African diabetic in two is hypertensive. This value is corroborated, by an independent design and on a Cameroonian population, by the results of \textcite{katte2014}.

\textbf{Second}, the pathophysiological mechanism set out in \cref{sec:comorbidity}, which makes this association expected rather than contingent, and the organisational mechanism identified by \textcite{hinneh2023}, whereby the accumulation of needs delays care instead of intensifying it.

\textbf{Third}, the normative position of the WHO. The 2023 report devotes a section to the diabetes--hypertension comorbidity and makes a decisive argument: uncontrolled hypertension and hypercholesterolaemia---not hyperglycaemia---are the risk factors most strongly associated with premature death and disability in people living with diabetes. The WHO adds that treating hypertension is among the most cost-effective interventions in these patients, since it is inexpensive and reduces both cardiovascular risk, the leading cause of death in this population, and the risk of renal failure, a frequent and serious complication of diabetes. It finally recalls that Member States have adopted a global target whereby 80\% of people diagnosed with diabetes should have good blood pressure control. This last point is crucial: it gives the designed system a quantified and institutionally recognised objective, not merely a technical ambition.

\textbf{Fourth}, an organisational convergence. The WHO observes that the components of the HEARTS technical package, developed for hypertension, ``support and can be applied to other health conditions requiring consistent, long-term care, including diabetes, chronic kidney disease and HIV''.%
\verify{Back-translated from the French draft; replace with the verbatim original wording from the WHO report.}
A single system covering several chronic conditions is therefore not a technical convenience, but corresponds to the mode of organisation recommended by the reference normative institution. \textcite{hinneh2023} conclude in the same direction by recommending an integrated care model that allows both conditions to be managed cost-effectively without additional burden on resources and health professionals, as well as team-based care.

\textbf{Fifth}, the bidirectional relationship with chronic kidney disease, which the WHO describes as a vicious circle in which each condition aggravates the other, hypertension being both a major risk factor for progression of renal failure and, often, the first sign of incipient kidney damage. The WHO consequently recommends screening for kidney disease in every newly diagnosed hypertensive person, and specifies that blood pressure targets are stricter in patients with renal failure. This third dimension, although beyond the scope of this thesis, justifies that the retained data architecture be extensible to other biological parameters.

%==========================================================================
\subsection{Comparability of estimates}\label{sec:comparability}

The values used in this section are not directly comparable with one another. \Cref{tab:comparability} makes the nature of each explicit.

\begin{table}[htbp]
\centering\footnotesize
\caption{Nature of the main estimates mobilised in \cref{sec:burden}.}\label{tab:comparability}
\begin{tabularx}{\textwidth}{@{}>{\raggedright\arraybackslash}p{2.6cm} >{\raggedright\arraybackslash}X >{\raggedright\arraybackslash}p{2.8cm} >{\raggedright\arraybackslash}p{2.4cm} r@{}}
\toprule
\textbf{Source} & \textbf{Population and method} & \textbf{Design} & \textbf{Indicator} & \textbf{Value} \\
\midrule
\textcite{who2023hypertension} & African Region, 30--79 y & modelling, age-standardised & hypertension prevalence & 36\% \\
\textcite{ataklte2015} & SSA, 33 surveys 2000--2013, mean age 40 & meta-analysis of population surveys & crude prevalence & 30\% (27--34) \\
\textcite{dzudie2021} & Cameroon, 5 cities, $\geq 18$ y & opportunistic screening, single visit & prevalence & 20.8\% \\
\textcite{lemogoum2018} & Far North Cameroon, $\geq 18$ y & cross-sectional, multilevel clusters & standardised prevalence & 37.8\% \\
\textcite{hinneh2023} & African diabetics, mean age 58 & meta-analysis, mostly hospital studies & hypertension among diabetics & 58.1\% \\
\textcite{katte2014} & Bafoussam, self-selected & community screening & hypertension + diabetes, standardised & 3.9\% (M) / 5.0\% (F) \\
\bottomrule
\end{tabularx}
\end{table}

Five sources of discrepancy must be kept in mind in any comparative reading: the reference population, whether general population, known diabetics or hospital attendees; the diagnostic threshold, 140/90 or 130/80~mmHg; the age structure of the sample, which \textcite{ataklte2015} showed alone explains a large part of the dispersion; the number of measurement visits, a single measurement overestimating prevalence; and the recruitment mode, self-selection and opportunistic screening producing samples systematically different from probability sampling.

None of the sources used is nationally representative for Cameroon. This limitation must be explicitly assumed rather than masked by abusive aggregation: it is in itself an argument for a system that, by collecting standardised longitudinal data on an identified cohort, would produce precisely the type of data the region lacks.

%==========================================================================
\subsection{From the epidemiological finding to the telemonitoring hypothesis}\label{sec:hypothesis}

\subsubsection{Hierarchised synthesis of findings}

The elements gathered in this section converge on five findings, presented in decreasing order of robustness, with the nature of the evidence stated for each.

\textbf{First finding, solidly established by modelling.} The burden is high and growing in absolute terms: 106 million hypertensive adults in the African Region in 2019 versus 45 million in 1990, for an age-standardised prevalence of 36\% versus 33\% globally. The growth in absolute terms is mostly of demographic origin (\cref{sec:transition}).

\textbf{Second finding, solidly established, derived from WHO modelled data.} The care cascade fails at each of its three stages, not at screening alone: 57\% undiagnosed, 37\% of the diagnosed untreated, 56\% of the treated uncontrolled (\cref{fig:losses}). The weakest link in conditional terms is the last. The deficit is organisational, not epidemiological (\cref{sec:cascade}).

\textbf{Third finding, asserted by the WHO and referenced.} Loss to follow-up is a central mechanism of this failure, often affecting more than half of patients in resource-limited settings. The WHO also establishes that the 2025 global target would be met if about half of already-diagnosed adults were effectively treated---that is, without widening screening.

\textbf{Fourth finding, established by meta-analysis with substantial uncertainty on regional detail.} Hypertension and diabetes coexist in a majority of diabetic patients, about 58\% according to the best available estimate, the sub-regional estimate of 77.6\% for Central Africa not being retained for the reasons developed in \cref{sec:central-estimates}.

\textbf{Fifth finding, established by scoping review.} The national normative framework is lacking: no hypertension guideline developed in sub-Saharan Africa meets the minimal methodological criteria of an evidence-based guideline, which requires explicitly anchoring the system to WHO guidelines.

\subsubsection{What these findings do not establish}

These findings establish a \emph{need}. They do not by themselves establish that telemonitoring meets it. Two further elements are required, which the literature provides in a contrasting manner.

\paragraph{Normative support: favourable.} The WHO \parencite{who2023hypertension} explicitly mentions that medical supervision can be provided through innovative methods such as patient telemonitoring and coordination with care team members, so that access to treatment is not delayed. It specifies that, when team-based care is delivered at community level, stable patients with controlled pressure can space out their facility visits, which presupposes multi-month medication dispensing, ideally ninety days, and community or home delivery.

The WHO's detailed requirements for a hypertension information system (secure longitudinal record, HEARTS core indicators, maturity factors for scale-up) are examined in \cref{sec:records}, where they are used to specify what the proposed system must output.

The WHO finally indicates that digital tools can generate lists of patients overdue for follow-up, send automated reminders by message or telephone, and trigger community action, thereby reducing loss to follow-up, and that available data show that digital interventions help improve medication adherence and patient engagement in treatment, thereby increasing blood pressure control.

\paragraph{Empirical support: uncertain.} The systematic review of \textcite{stokes2022} concludes that there is evidence in favour of mobile technology for hypertension management in sub-Saharan Africa, the documented applications concerning blood pressure reduction, improved health knowledge and treatment adherence. It qualifies this finding, however, with three decisive reservations: the heterogeneity of interventions was such that no meta-analysis could be conducted, results being presented narratively; few studies showed a statistically significant benefit over usual care, only three randomised trials reporting moderate blood pressure reductions; and no study using artificial intelligence was identified, the authors suggesting that it could prove beneficial, for example as a diagnostic aid, but that this hypothesis remains untested in the reviewed literature.

The meta-analysis of \textcite{aovare2023}, covering seven randomised trials and 2\,249 African participants---two studies measuring HbA1c, five blood pressure---finds no significant effect of mobile health interventions (\cref{tab:aovare}).

\begin{table}[htbp]
\centering\small
\caption{Pooled effects of mHealth interventions in African randomised trials \parencite{aovare2023}.}\label{tab:aovare}
\begin{tabular}{@{}lccc@{}}
\toprule
\textbf{Outcome} & \textbf{Weighted mean difference} & \textbf{95\% CI} & \textbf{$p$} \\
\midrule
Glycated haemoglobin & $+0.20$ & $-0.40$ to $0.80$ & 0.51 \\
Systolic pressure & $-1.39$~mmHg & $-4.46$ to $1.68$ & 0.37 \\
Diastolic pressure & $+0.36$ & $-1.37$ to $2.05$ & 0.69 \\
\bottomrule
\end{tabular}
\end{table}

A sensitivity analysis excluding one study makes the systolic effect significant, at $-2.22$~mmHg ($-3.94$ to $-0.60$; $p=0.01$), the other two outcomes remaining non-significant. This is however a post hoc analysis depending on a single study, which the authors decline to regard as conclusive. Their conclusion is explicit: there is no conclusive evidence of the efficacy of mobile health interventions on blood pressure reduction and improved glycaemic control in Africa, and larger randomised trials are required.

\subsubsection{Problem statement and working hypothesis}

It is precisely this configuration that defines the space of contribution of the present work. The need is established, the normative framework is explicitly favourable and provides operational indicators, but the efficacy of existing systems is not demonstrated and none of them exploits machine learning methods.

Two explanations of this absence of evidence must be distinguished, and the choice between them determines the orientation of the work. According to the first, mobile health interventions would be genuinely ineffective in this context, in which case designing a new system would be futile. According to the second, the systems evaluated so far---essentially consisting of text-message reminders and therapeutic education, as described by \textcite{stokes2022}---were too weakly integrated into the care chain and made too little use of the longitudinal data collected to produce a measurable effect.

Several elements argue for the second explanation. The recorded interventions act mainly on adherence, that is, on the second link of the cascade, whereas our analysis in \cref{sec:cascade} showed that the third link---effective control in treated patients---is the largest conditional loss, and that this loss depends on therapeutic adjustment as much as on adherence. A system limited to reminders cannot, by construction, act on adjustment. Moreover, the total absence of studies using machine learning methods means that the informational potential of the collected time series---trend detection, early identification of patients heading towards decompensation, personalisation of alert thresholds---remains entirely unexploited in the region.

This thesis therefore hypothesises that the added value lies in the \emph{longitudinal integration} of consultation data, home measurements and mobile sensor data within a single detection model, and not in adding a further communication channel. This hypothesis is made explicit and operationalised in the next chapter, and evaluated in the experimental part.

%##########################################################################
% SECTION 1.2
%##########################################################################
\section{Diagnosis of Conventional Hospital Systems and Record Keeping}\label{sec:diagnosis}

\subsection*{Purpose of the section and status of the sources}

The preceding section established that the chain of care loses patients at each of its three stages, and that the largest conditional loss lies at the last link, effective control among patients who are already treated (\cref{sec:cascade}). This section identifies the concrete mechanisms behind this attrition.

Four main sources are used. Their evidential status differs markedly and must be stated at the outset, because it determines the weight given to each statement. Two primary studies, TELEMED-CAM and Chatrati et al., are discussed where they bear on the argument (\cref{sec:telemedcam,sec:ai-gap}).%
\verify{Quotations marked with $\dagger$ are back-translated from the French draft. Replace each with the verbatim English wording of its source, then redefine \texttt{\textbackslash bq} in the preamble so that the dagger disappears.}

\paragraph{WHO (2023).} \textcite{who2023hypertension} is a normative document produced through an institutional process and supported by 171 references. It carries the main evidential burden for everything concerning the clinical protocol, programme indicators and recommended positions.

\paragraph{Dodoo et al. (2022).} \textcite{dodoo2022}, published in \emph{Health and Technology} 12(1):33--46, is a systematic review conducted under adapted PRISMA guidelines, with a documented search protocol, double assessment and a measure of inter-rater agreement. It is the most robust source on barriers to telemedicine in sub-Saharan Africa.

\paragraph{Liu and Wang (2025).} \textcite{liu2025review}, published in \emph{Frontiers in Public Health} 13:1510456, is titled ``a systematic review'' but is classified by the journal itself under the editorial type ``Review''. The authors state that they used \bq{a rigorous review method including systematic database searching, literature screening against specific inclusion and exclusion criteria, and the adoption of quality assessment tools}, but none of these elements is reported: no protocol registration, no flow diagram, no number of included studies, no list of criteria, no named assessment instrument. The only quantitative information on the corpus is a figure showing the annual number of relevant articles, peaking at about 720 in 2023. We therefore treat this source as a narrative review, useful for structuring the problem and for locating earlier work, and not as aggregated evidence. Its inventory of public chronic-disease datasets (its Section~2.2 and Table~1) does not belong to this diagnosis and is reserved for the experimental chapter.

\paragraph{Navarro Masuet (2023).} \textcite{navarro2023} is a Bachelor thesis in biomedical engineering, defended at CEU San Pablo University in June 2023 under the supervision of Carlos Oscar Sorzano. It is not a peer-reviewed publication. Its value here is twofold: it proposes a clear taxonomy of blood pressure measurement modalities, and it reports an interview with a practitioner. Its factual statements must be cross-checked before being reused; its original contributions, the taxonomy and the interview, are used as such, that is, as an organisational proposal and as a testimony, not as results.

A methodological caveat completes this framing. The claim that consultations shorter than five minutes would be \bq{insufficient to assess physiological variability} is an inference, not a WHO statement: the WHO mentions this duration in another context and draws a different conclusion from it. The argument for the metrological insufficiency of the consultation must be built from the explicit requirements of the diagnostic protocol, which is what the next subsection does.

%==========================================================================
\subsection{The gap between protocol requirements and the capacity of the episodic model}\label{sec:protocol-gap}

\subsubsection{What diagnosis requires}\label{sec:dx-requires}

The diagnostic protocol recommended by the WHO is demanding, and its detail is worth reporting, because the insufficiency of the episodic model follows from it.

Measurement must be performed by trained staff, in a standardised way, with a cuff of appropriate size, the patient seated comfortably, back supported, bladder empty and legs uncrossed. Two readings must be taken at the first encounter, and the result of the second guides the decision whether to schedule a confirmation visit. What follows depends on the level observed (\cref{tab:protocol}).

\begin{table}[htbp]
\centering\small
\caption{Recommended action by confirmed pressure level at the first encounter \parencite{who2023hypertension}.}\label{tab:protocol}
\begin{tabularx}{\textwidth}{@{}>{\raggedright\arraybackslash}p{6.2cm} X@{}}
\toprule
\textbf{Confirmed pressure on both readings of the day} & \textbf{Recommended action} \\
\midrule
SBP $\geq$ 160 or DBP $\geq$ 100~mmHg & start treatment the same day \\
SBP 140--159 or DBP 90--99~mmHg & confirm at an additional visit, usually one to four weeks later \\
\bottomrule
\end{tabularx}
\end{table}

Diagnosis is made only if, at both visits, systolic pressure is $\geq 140$~mmHg and/or diastolic pressure is $\geq 90$~mmHg.

Diagnosis is therefore, by construction, a longitudinal act: it requires two measurement points separated by one to four weeks (the recheck window already noted in \cref{sec:cascade}). A care model that does not guarantee the patient's return for the second visit does not produce imperfect diagnoses; it produces none.

One definitional divergence between our sources must be noted. \textcite{navarro2023} defines hypertension at 130~mmHg systolic or 80~mmHg diastolic, while acknowledging in the same paragraph a consensus to treat from persistent values of 140/90~mmHg, with 130/80~mmHg then as the therapeutic target. This reflects the gap between the European Society of Cardiology classification, which she reproduces, and the WHO surveillance threshold that we retain (\cref{sec:hta-burden}). We keep the WHO definition (140/90~mmHg or ongoing treatment) throughout this thesis, and we flag any figure drawn from another convention.

\subsubsection{What variability imposes}\label{sec:variability}

Blood pressure is, in the words of \textcite{navarro2023}, \bq{a highly dynamic vital sign, easily affected by external factors, producing fluctuations throughout the day and significant variability}. The WHO establishes the same through an empirical observation: in population screening, many people with elevated pressure do not meet the diagnostic criteria when assessed more carefully at follow-up. (The white-coat effect was already listed among the sources of overestimation in \cref{sec:burden}.)

This variability can be formalised. For patient $i$ measured at time $t$,
\begin{equation}
  b_{i,t} = \mu_i + \varepsilon_{i,t},
\end{equation}
where $\mu_i$ is the patient's usual pressure, the only clinically relevant quantity, and $\varepsilon_{i,t}$ gathers circadian variability, variability due to activity and stress, the effect of the measurement environment and instrument error. A single measurement estimates $\mu_i$ with a variance equal to that of $\varepsilon$; $n$ sufficiently spaced measurements, whose errors are independent, reduce it by a factor $n$:
\begin{equation}
  \operatorname{Var}(\hat\mu_i) = \frac{\sigma_\varepsilon^2}{n}.
\end{equation}
Independence is an assumption: two readings taken in the same sitting share their context, so the gain from the second is smaller than the formula suggests, which is one more reason why the protocol also asks for a second visit.

\begin{figure}[htbp]
\centering
\begin{tikzpicture}
\begin{axis}[width=0.8\textwidth,height=6cm,
  xlabel={Number of independent measurements $n$},
  ylabel={$\operatorname{Var}(\hat\mu_i)/\sigma_\varepsilon^2$},
  xmin=0,xmax=34,ymin=0,ymax=1.08,grid=major,grid style={gray!20}]
\addplot[thick,cRef,domain=1:34,samples=100]{1/x};
\addplot[only marks,mark=*,mark size=2.5pt,cAfrica,nodes near coords,point meta=explicit symbolic,
  every node near coord/.append style={font=\footnotesize,anchor=south west}]
  coordinates {(2,0.5) [two visits] (4,0.25) [four visits] (30,0.033) [30 daily readings]};
\end{axis}
\end{tikzpicture}
\caption{Relative variance of the estimate of the usual pressure $\mu_i$ as a function of the number of independent measurements. Errors are assumed independent.}\label{fig:variance}
\end{figure}

This is the logic of the two-reading, two-visit protocol (\cref{fig:variance}). It is also, and this is the central argument of this thesis, the logic that favours home measurement: where a consultation provides at best two to four points per year, home follow-up produces several dozen, which reduces uncertainty about $\mu_i$ and additionally allows $\sigma_\varepsilon$ itself, the patient's blood pressure variability, to be estimated, an independent risk factor.

\textcite{navarro2023} documents two direct clinical consequences of this variability, which are the two misclassification errors of the episodic model:
\begin{itemize}[nosep]
  \item \textbf{White-coat hypertension}, or isolated office hypertension, is elevated pressure in the doctor's office with normal ambulatory or home measurement. The gap between the higher office value and the lower out-of-office value is attributed mainly to an alerting reaction provoked by the measurement itself. Its prevalence in the general population is estimated at about 13\%.
  \item \textbf{Masked hypertension}, or isolated ambulatory hypertension, is the opposite situation in untreated individuals: normal office pressure but elevated home or ambulatory pressure. It is more frequent in young people, men, smokers, overweight people and people with diabetes, as well as in those experiencing anxiety or stress.
\end{itemize}

This last point bears directly on our subject: the target population of the proposed system, patients with diabetes, is precisely the one in which the office measurement is most likely to underestimate true pressure. The misclassification of the episodic model is therefore not symmetric in this population; it leans towards the false negative, that is, the patient who will not be treated although he or she should be.

\textcite{navarro2023} also proposes a taxonomy of measurement modalities along three axes (\cref{fig:taxonomy}), which we adopt because it usefully structures the design choice.

\begin{figure}[htbp]
\centering
\begin{tikzpicture}[node distance=0.35cm and 0.5cm,
  ax/.style={draw,fill=cRef!15,rounded corners,minimum width=2.4cm,minimum height=0.8cm,font=\small\bfseries},
  md/.style={draw,rounded corners,minimum height=0.8cm,font=\small,fill=white}]
  \node[ax] (a1) {Location};
  \node[md,right=of a1] (m11) {Office BP};
  \node[md,right=of m11] (m12) {Out-of-office BP};
  \node[ax,below=0.5cm of a1] (a2) {Frequency};
  \node[md,right=of a2] (m21) {Home (HBPM)};
  \node[md,right=of m21] (m22) {Ambulatory (ABPM)};
  \node[ax,below=0.5cm of a2] (a3) {Technique};
  \node[md,right=of a3] (m31) {Oscillometry};
  \node[md,right=of m31] (m32) {Applanation tonometry};
  \node[md,right=of m32] (m33) {Photoplethysmography};
\end{tikzpicture}
\caption{Taxonomy of blood pressure measurement modalities (after \textcite{navarro2023}). Home and ambulatory measurement are both out-of-office modalities, so the three axes are not strictly orthogonal.}\label{fig:taxonomy}
\end{figure}

Home self-measurement relies on automatic oscillometric devices applied to the brachial artery and can, according to the author, confirm the diagnosis after an elevated office value and identify patients with white-coat or masked hypertension. Ambulatory measurement requires wearing a portable monitor for twelve to twenty-four hours, sometimes up to forty-eight, with recordings every twenty to thirty minutes, including night-time values and short-term variability.

\subsubsection{The clinical testimony}\label{sec:testimony}

The thesis of \textcite{navarro2023} reports an interview with Dr.\ José María Masuet Iglesias, director of the Verge del Toro basic health zone in Mahón, Menorca. This testimony gives a practitioner's wording of the mechanisms described above, and we cite it as such, stating at once its limited scope: it is a single informant, practising in a Spanish health system whose constraints differ markedly from those of a resource-limited system, and whose surname coincides with the author's, which suggests recruitment by proximity rather than purposive sampling. It is therefore not data, but a clinical illustration of mechanisms established elsewhere.

Asked about the main challenge of blood pressure measurement, the practitioner answers that it \bq{lies in the significant variability observed between different types of measurement during the day}. He explains that when doctors ask patients to keep a log of their home pressure, results are mostly satisfactory because measurements are taken in a calm, stress-free environment, but that on re-evaluation in the consultation the values are often considerably higher. He concludes that \bq{this notable discrepancy underlines the need for a continuous monitoring mechanism or device}.

Two other elements of the interview deserve to be retained.

The first is an explicit comparison with diabetes, which highlights an asymmetry that the system design will have to integrate. The practitioner observes that the diagnosis of diabetes is relatively simple, since anyone with durably elevated glucose outside recent food intake is a candidate for diagnosis, whereas assessing blood pressure is a complex process: \bq{measuring blood pressure at isolated moments that do not fully reflect the usual lifestyle---sitting calmly on the sofa after a stressful day---or taking sporadic measurements in situations more representative of daily life, such as during working hours, may not provide sufficient information}. He infers substantial underdiagnosis, which he finds worrying.

The second, the most directly useful, concerns therapeutic evaluation: \bq{since blood pressure is typically measured at specific moments, it is difficult to determine which people are consistently hypertensive; it becomes consequently difficult to verify the effectiveness of a given treatment. Suffering from hypertension is just as harmful as receiving an ineffective treatment that worsens the problem.} This observation describes in clinical terms exactly the mechanism that we identify below as therapeutic inertia (\cref{sec:inertia}): the absence of longitudinal information makes it impossible to assess efficacy, and therefore makes the decision to intensify treatment unjustifiable.

\subsubsection{What follow-up requires}\label{sec:followup-requires}

Recommendation~7 of the WHO 2021 guideline on the pharmacological treatment of hypertension \parencite{who2021guideline} sets the expected follow-up frequency:
\begin{itemize}[nosep]
  \item monthly follow-up after initiating or changing antihypertensive treatment, until the patient reaches the blood pressure target;
  \item follow-up every three to six months for patients whose pressure is controlled.
\end{itemize}
The therapeutic protocol module of the HEARTS technical package \parencite{whoHEARTS} further specifies that a simple protocol must set the return-visit interval, \bq{usually four weeks or less}, to allow rapid titration of doses and drugs until control is reached.

Let us confront this requirement with the data of \cref{sec:burden}. In the African Region, treatment coverage is 27\% and control 12\%, from which we derived that 56\% of treated patients do not reach the target (\cref{sec:cascade}). According to recommendation~7, these patients should be seen monthly until they do.

The gap can be quantified (\cref{tab:contacts}). An uncontrolled patient should receive about twelve contacts per year, a controlled patient two to four. For an active list in which a share $\pi_u = 0.56$ of patients is uncontrolled, the theoretical need is
\begin{equation}
  \bar n = \pi_u \times 12 + (1-\pi_u)\times n_c \approx 7.6 \text{ to } 8.5 \text{ contacts per patient per year}, \qquad n_c \in [2,4].
\end{equation}

\begin{table}[htbp]
\centering\small
\caption{Follow-up contacts per patient per year required by WHO recommendation~7.}\label{tab:contacts}
\begin{tabularx}{\textwidth}{@{}X c@{}}
\toprule
\textbf{Patient profile} & \textbf{Contacts per year} \\
\midrule
Uncontrolled: monthly follow-up until the target is reached & 12 \\
Controlled: follow-up every three to six months & 2 to 4 \\
Active list with 56\% uncontrolled (\cref{sec:cascade}) & 7.6 to 8.5 \\
\bottomrule
\end{tabularx}
\end{table}

If the episodic model delivers one to two contacts per year, it delivers between about an eighth and about a quarter of the recommended frequency.%
\verify{The figure of one to two delivered contacts per year has no source in the draft, where it was introduced as ``available data for the region''. Cite a regional or Cameroonian source, or present it explicitly as an illustrative scenario. The draft's ``one fifth to one quarter'' assumed half of the list uncontrolled; with the 56\% derived in \cref{sec:cascade} the range is about 12--26\%.}
It is this cumulative gap, and not the duration of each consultation taken in isolation, that constitutes the structural deficit.

\subsubsection{The per-consultation time constraint}\label{sec:time}

The WHO indicates that, in many primary-care settings of low- and middle-income countries, patient encounters last less than five minutes, and draws an operational conclusion: digital information systems must therefore be simple and fast to use by front-line providers.

This constraint must be correctly interpreted. It does not mean that five minutes would be insufficient to measure a blood pressure. It means three things.

First, it rules out any solution that would add a documentation burden for the provider: a tool requiring two minutes of data entry would consume 40\% of the consultation time. The WHO states the converse requirement as well: an optimal system replaces paper, so that providers record information only once.

Second, it makes the effective performance of the full protocol within the available time unlikely, which degrades the quality of the estimate of $\mu_i$ beyond the problem of the number of points.

Third, it leaves little room for counselling and therapeutic education, whereas the WHO recalls that these must be provided without delaying the initiation of drug treatment.

This time constraint combines with a workload constraint documented by \textcite{dodoo2022} and revisited in \cref{sec:digital-barriers}: health professionals in the region, facing a considerable workload, perceive telemedicine systems as an additional task generating overload and consuming time. Any design that ignored this perception would meet rejection, whatever its technical merit.

\subsubsection{Therapeutic inertia}\label{sec:inertia}

The WHO introduces a notion that explains why the third link of the cascade is the weakest: therapeutic inertia, defined as \bq{the tendency not to act to intensify treatment despite blood pressure remaining uncontrolled}. It is one of the four critical domains that a quality-improvement strategy must address, alongside medicine supply, retention in care and decentralisation. The WHO indicates that a simple, standardised protocol, by clarifying treatment steps for both patient and clinician, helps reduce it.

Three conditions are necessary for treatment intensification to take place: the patient returns; the provider has the pressure history showing that the target is persistently not reached; and the protocol states unambiguously the next step. The episodic model with a paper record fails on the first two, and the clinical testimony above (\cref{sec:testimony}) states exactly this difficulty.

This is why a device that merely alerted on an isolated value would have only a marginal effect: the decision to intensify rests on a trend, that is, on longitudinal information, and not on a point.

One further element reinforces the need for the drug lever. The WHO reports that the effect of lifestyle modification, assessed in primary-care trials, is modest: about 1.3~mmHg of systolic and 0.5~mmHg of diastolic reduction. It concludes that almost all people diagnosed with hypertension will keep an uncontrolled pressure despite changes in diet and exercise and will need medication. Lifestyle modification is therefore an adjuvant therapy. A telemonitoring system whose function was limited to lifestyle counselling would thus target a lever whose effect size is an order of magnitude smaller than that of drug therapy.%
\verify{The draft states a ratio of ``about one tenth'' without giving the drug effect it is compared with. Cite the expected systolic reduction under drug therapy from the WHO guideline, then state the ratio from those two numbers.}

%==========================================================================
\subsection{Limits of record keeping}\label{sec:records}

\subsubsection{Ineffectiveness of paper supports}

The WHO devotes a section of its chapter on management to this question. It establishes that effective hypertension control programmes require a robust information system allowing blood pressure control to be tracked for the individual patient and progress to be assessed at facility and health-system level. It observes that, in many primary-care facilities of low- and middle-income countries, staff record and transmit data on ineffective paper-based systems.

Three failures are attributed to these supports.

\textbf{Duplication.} The WHO makes the creation of a \bq{secure, accurate and non-duplicated} longitudinal patient record essential, which implies that duplication is the current regime: the same patient exists under several identities across registers and facilities, making it impossible to reconstruct his or her pathway.

\textbf{Loss and inaccuracy.} The report indicates that an optimal system replaces paper so that providers only have to record information once, which reduces the risk of lost and inaccurate data. Repeated re-entry is thus identified as the source of the problem.

\textbf{Unavailability when needed.} The record must be \bq{readily available during clinical visits}. A paper register kept in one facility cannot be consulted in another facility, at home, or remotely.

\subsubsection{Absence of a usable longitudinal history}

This is the most important consequence for our subject, and it follows from the three failures above. The WHO stresses that a longitudinal record is essential for patients with hypertension, diabetes \bq{and other conditions requiring long-term care}.

The wording is significant: what matters is not the recording of the act but the continuity of the series. A paper register is optimised for the act, it attests that a consultation took place, and not for the series. Clinical information therefore exists as isolated, unlinked, non-aggregable points.

This distinction between act data and series data is central to the system design (\cref{fig:act-series}). A pressure of 152/94~mmHg does not have the same meaning depending on whether it follows a stable series around 135/85 or is part of a progressive drift observed over three months. The first case calls for a check, the second for therapeutic intensification. The paper register cannot distinguish the two situations; from an informational viewpoint, this is the reason why therapeutic inertia persists.

\begin{figure}[htbp]
\centering
\begin{tikzpicture}
\begin{axis}[width=0.9\textwidth,height=6.2cm,xlabel={Week},ylabel={Systolic pressure (mmHg)},
  xmin=0,xmax=11,ymin=125,ymax=158,xtick={0,2,...,10},grid=major,grid style={gray!20},
  legend pos=north west,legend style={font=\footnotesize}]
\addplot[thick,cRef,mark=*] coordinates {(0,136)(1,134)(2,137)(3,135)(4,133)(5,136)(6,135)(7,134)(8,137)(9,135)(10,136)(11,152)};
\addplot[thick,cAfrica,mark=square*] coordinates {(0,131)(1,133)(2,132)(3,136)(4,137)(5,140)(6,141)(7,144)(8,146)(9,148)(10,150)(11,152)};
\addplot[dashed,cGrey,no markers,domain=0:11]{140};
\legend{Stable series with isolated excursion,Progressive drift,Threshold (140~mmHg)};
\end{axis}
\end{tikzpicture}
\caption{Same last value, different meaning. Schematic illustration with synthetic values, not patient data: a single reading of 152~mmHg is a check-up case in one series and an intensification case in the other.}\label{fig:act-series}
\end{figure}

\subsubsection{Information silos and obstacles to interoperability}

\textcite{liu2025review} formulate a convergent diagnosis from a more general perspective. They state that \bq{traditional chronic disease monitoring methods rely mainly on manual processes, constrained by a limited scope of population monitoring, data collection limitations, accuracy concerns and a lack of universality of prevention and management strategies}. They add that these methods are \bq{compromised by poor self-management capacity and adherence problems}, and that further factors intervene: \bq{the inconvenience of access to medical services, the phenomenon of information silos, the lack of real-time monitoring and feedback, the uneven distribution of medical resources, the obsolescence of technologies and methods, and the lack of effective prediction and early warning mechanisms}.

In their list of public-health challenges posed by chronic disease, they also mention \bq{disconnections between projects and routine operations} and \bq{inadequate coordination between institutional levels}, two formulations that match precisely the sustainability problem documented by \textcite{dodoo2022} for sub-Saharan Africa.

\textcite{navarro2023} approaches the same difficulty from its technical side: \bq{integrating information from different devices, which may sometimes be of different types, is complex because of the diversity of existing formats, which hinders the unification of all these data}. She observes more broadly that \bq{each health institution seems to incorporate the measures it considers appropriate for telemedicine services, but there is no consensus on the steps to follow or the objectives to achieve}, and that \bq{each device is designed independently, without a guiding principle defining the ultimate goal of the whole process}.

The information silo is therefore not only a consequence of paper: it re-forms in the digital environment as soon as each device produces its own format. This observation sets a design requirement that the present work must explicitly satisfy: define a single, documented data schema, rather than aggregating heterogeneous flows after the fact.

\subsubsection{Quality of digitised data}

\textcite{liu2025review} note that electronic medical records themselves require prior cleaning, the typical artefacts being duplicates, unit errors and values carried over identically from one visit to the next. This must be kept for the experimental part: digitisation does not remove the data-quality problem, it displaces it. Any model trained on real consultation data must include an explicit step to detect these three artefacts, and the identically carried-over value, which falsely simulates clinical stability, is the most dangerous of the three for a trend-detection model.

\subsubsection{What the WHO expects from an information system}\label{sec:hearts-spec}

The report does not stop at the finding: it defines the expected outputs of a hypertension information system, as a short list of core HEARTS indicators. This list is a directly usable functional specification (\cref{tab:hearts}).

\begin{table}[htbp]
\centering\small
\caption{Core HEARTS programme indicators \parencite{who2023hypertension}.}\label{tab:hearts}
\begin{tabularx}{\textwidth}{@{}>{\raggedright\arraybackslash}p{3.6cm} X c@{}}
\toprule
\textbf{Indicator} & \textbf{Definition} & \textbf{Temporal} \\
\midrule
Programme coverage & number of hypertensive patients registered in the catchment area & no \\
Blood pressure control & proportion of registered patients with BP $<$ 140/90~mmHg over the last three months & yes \\
Retention in care & proportion of registered patients with no visit in the last three months, and in the last twelve months & yes \\
Availability & availability of essential antihypertensive medicines and BP measurement devices & no \\
\bottomrule
\end{tabularx}
\end{table}

Two observations are in order. First, two of the four core indicators are temporal: they can only be computed from a dated history, and a paper register cannot produce them except at the cost of an impracticable manual review. Second, the control indicator is itself defined over a three-month sliding window, which confirms that the quantity of interest is a property of the series and not of a single measurement.

The WHO specifies that when programme managers receive these indicators monthly, they can identify lagging facilities and implement targeted improvements. It cites the DHIS2 Hypertension Tracker as an example of a user-friendly mobile application. Among the six maturity domains conditioning the national scale-up of a programme, it also includes the presence of a digital information system for hypertension management, with, as a clinical quality factor, a missed-visit rate below 25\%, defined as failure to return to a scheduled visit within the preceding three or six months. The other five domains are clinical capacity, service delivery, clinical quality, medicine supply and political commitment.

%==========================================================================
\subsection{Loss to follow-up}\label{sec:ltfu}

\subsubsection{Magnitude}

Loss to follow-up, introduced in \cref{sec:cascade}, is a major barrier to hypertension control, often affecting 50\% of patients or more in resource-limited settings according to the WHO. The WHO adds that improving control requires retaining patients in care and keeping them on effective treatment, and that digital information systems can be mobilised as a powerful mechanism to support this retention.

This figure must be located relative to those of \cref{sec:burden}. The 50\% of patients lost to follow-up are not a fourth loss added to the three stages of the cascade: they are its underlying mechanism. A patient lost to follow-up stops being treated, or continues a treatment that has become inadequate without reassessment. He or she therefore feeds simultaneously the loss between diagnosis and treatment and the loss between treatment and control. It is in this sense that retention is the most effective point of application for an intervention.

\subsubsection{Screening without follow-up is inoperative}

The WHO is explicit on this point, and its argument challenges a widespread conception of digital health intervention. It writes that hypertension screening is useful only if closely linked to follow-up in a facility able to measure pressure accurately and to treat effectively, and that in population screening programmes conducted worldwide, effective follow-up of patients with abnormal pressure is the exception rather than the norm. For this reason, detection in health facilities is the priority.

It adds a capacity argument: mass screening programmes increase the workload of the system, not only through the detection process but above all through the need to manage identified cases for life. With a prevalence close to a quarter of the world's adult population, no screening system should be established without first verifying that local resources are sufficient. Adding cases without increasing treatment capacity produces poor-quality care and does not durably improve blood pressure control. The WHO therefore recommends opportunistic measurement in every adult attending a facility, and does not recommend population screening as the primary method of detection.

This position has a direct consequence for the positioning of the present work: it disqualifies a conception of telemonitoring as an instrument of mass screening and validates a conception of telemonitoring as an instrument of retention and follow-up of patients already identified. This is the orientation retained here, and it rests on an explicit normative recommendation.

\subsubsection{The converse demonstration: retention and control are linked}

The report provides, through four national programmes, elements that establish this link empirically (\cref{tab:programmes}).

\begin{table}[htbp]
\centering\footnotesize
\caption{Follow-up mechanisms and reported figures of four national programmes \parencite{who2023hypertension}.}\label{tab:programmes}
\begin{tabularx}{\textwidth}{@{}>{\raggedright\arraybackslash}p{2.1cm} >{\raggedright\arraybackslash}X >{\raggedright\arraybackslash}X@{}}
\toprule
\textbf{Programme} & \textbf{Follow-up mechanism} & \textbf{Figures reported} \\
\midrule
Sri Lanka & HEARTS implemented from 2021 with a simple two-step treatment protocol adapted from the WHO 2021 guideline & Kalutara district, 2023: protocol covering 96\,900 adults, 8\,600 patients enrolled; 86\% of enrolled adults return every three months; control 57\% (facility level) \\
Philippines & barangay community health workers (at least twenty households each, supervised by a nurse or midwife): BP measurement at the station or at home with reporting to the supervisor, education, information on services, dispensing to remote households, tracing of patients who missed the quarterly visit & none in the text used here \\
India & decentralisation of follow-up to Ayushman Bharat health and wellness centres near villages; medicine refills & from 5 states (2018) to 155 districts in 27 states (June 2023); 5.8 million patients enrolled; refills helped more than 38\% of patients in the first COVID-19 wave \\
Cuba & home visits by the care team & early 2023: more than 2 million of about 3 million hypertensive people registered \\
\bottomrule
\end{tabularx}
\end{table}

\textbf{Sri Lanka} is the most direct element at our disposal: 86\% retention coexists with 57\% control. The report specifies that the high retention rate contributes in part to the high control rate. This comparison must however be read with the denominators of \cref{sec:cascade} in mind: the 57\% is control among \emph{enrolled patients} in facilities, whereas the 12\% of \cref{tab:cascade} is effective coverage among \emph{all hypertensive adults} of the region. The nearest African analogue of a programme-level rate is control among treated patients, $12/27 \approx 44\%$, which is closer to the Sri Lankan figure, although still not equivalent.%
\verify{The draft contrasted 57\% with 12\%. Those two figures have different denominators; the 44\% analogue narrows the apparent gap. Check that the WHO report does not give a comparable facility-level African figure.}

\textbf{The Philippines} programme performs five functions along the cascade: BP measurement and recording at the station or at home during home visits, with reporting to the supervisor for collation; education on the prevention and control of noncommunicable diseases; information on the availability of services and medicines; dispensing of medicines to remote rural areas and to households whose patients cannot travel; and tracing of patients who missed their quarterly clinic visit. This arrangement is the manual equivalent, performed by human agents, of the function that the proposed system aims to automate.

\textbf{India} decentralised follow-up to the health and wellness centres, located close to villages. The report also indicates that medicine availability emerged as the main obstacle soon after launch, a situation that \bq{initially discouraged patients from returning to facilities for their follow-up visits} and exposed the programme to a loss of credibility. Retention is therefore itself a dependent variable.

\textbf{Cuba} counts among the distinctive features of its programme home visits by the care team to ensure follow-up.

Of these four examples, only the first reports both retention and control. All four, however, converge on a common feature: they moved part of the follow-up activity out of the hospital facility, towards the home or its immediate vicinity.

%==========================================================================
\subsection{Barriers to digital transition in sub-Saharan Africa}\label{sec:digital-barriers}

The preceding finding does not imply that substituting a digital system for the paper register is easy. \textcite{dodoo2022} conducted on this point the most complete review available for the region.

\subsubsection{Method and caveat on the flow diagram}

The study addresses two research questions: what is the state of progress of telemedicine adoption in sub-Saharan African countries, and what are the obstacles. It adapts the PRISMA guidelines. The search covered eight databases---PubMed, IEEE Xplore Digital Library, Cochrane, ProQuest, Web of Science, Scopus, Global Health and Google Scholar---with an equation combining the terms progress, adoption, telemedicine and digital health on the one hand and regional designations on the other. Only empirical studies and government reports were retained, without restriction of language or publication date. Two experts assessed the relevance of the articles, with a mean inter-rater agreement of 86\% computed by the item-by-item method.

The flow diagram reports 7\,840 records from databases and 50 additional records; 230 records after duplicate removal; 230 records screened, 38 excluded; 192 full texts assessed, 126 excluded; 14 studies retained in qualitative synthesis and 52 in quantitative and mixed synthesis, that is, the 66 announced studies (\cref{fig:prisma-dodoo}).

\begin{figure}[htbp]
\centering
\begin{tikzpicture}[node distance=0.7cm,
  b/.style={draw,rounded corners,align=center,minimum width=5cm,font=\footnotesize,fill=cRef!8},
  s/.style={draw,rounded corners,align=center,minimum width=4.4cm,font=\footnotesize,fill=cGrey!10},
  ar/.style={-{Stealth[length=2.5mm]},thick}]
  \node[b] (id) {Records identified\\ $7\,840 + 50 = 7\,890$};
  \node[b,below=of id] (dup) {Records after ``duplicates removed''\\ 230};
  \node[b,below=of dup] (scr) {Records screened\\ 230};
  \node[b,below=of scr] (ft) {Full texts assessed\\ 192};
  \node[b,below=of ft,fill=cAfrica!12] (inc) {Studies included: 66\\ (14 qualitative; 52 quantitative or mixed)};
  \node[s,dashed,fill=cAfrica!10,right=1.2cm of dup] (gap) {$7\,890-230 = 7\,660$ records\\ unexplained (title and abstract\\ screening not shown)};
  \node[s,right=1.2cm of scr] (ex1) {Excluded: 38};
  \node[s,right=1.2cm of ft] (ex2) {Excluded: 126};
  \draw[ar] (id) -- (dup);
  \draw[ar] (dup) -- (scr);
  \draw[ar] (scr) -- (ft);
  \draw[ar] (ft) -- (inc);
  \draw[ar,dashed] (dup) -- (gap);
  \draw[ar] (scr) -- (ex1);
  \draw[ar] (ft) -- (ex2);
\end{tikzpicture}
\caption{Selection flow of \textcite{dodoo2022} as reported. The drop from 7\,890 records to 230 cannot be duplicate removal alone.}\label{fig:prisma-dodoo}
\end{figure}

This diagram contains an inconsistency that must be signalled: going from 7\,890 identified records to 230 records ``after duplicate removal'' would imply 7\,660 duplicates, which is implausible. The screening step on titles and abstracts, which the text explicitly mentions---\bq{we examined all titles and abstracts related to the search keywords}---is not represented in the diagram, and it is probably what produces this reduction. The rest of the flow is consistent: 230 minus 38 equals 192, and 192 minus 126 equals 66. This reporting imprecision does not invalidate the results but prevents an exact reproduction of the selection.

\subsubsection{The six dimensions of barriers}

The authors find uneven progress between countries and classify obstacles in six dimensions, defined by item-centred coding (\cref{tab:barriers}).

\begin{table}[htbp]
\centering\small
\caption{Six dimensions of barriers to telemedicine in sub-Saharan Africa \parencite{dodoo2022}.}\label{tab:barriers}
\begin{tabularx}{\textwidth}{@{}>{\raggedright\arraybackslash}p{3.4cm} X@{}}
\toprule
\textbf{Dimension} & \textbf{Definition given by the authors} \\
\midrule
Technological & barriers related to information systems and connectivity software \\
Organisational & stakeholder engagement, sustainability of systems \\
Individual & users' experience, knowledge and awareness \\
Financial & financial resources and operating cost of telemedicine \\
Cultural & social constraints and traditional beliefs within the community \\
Legal and regulatory & policies, legislation, norms governing confidentiality and privacy protection \\
\bottomrule
\end{tabularx}
\end{table}

Three of these findings call for comment with regard to the design of the proposed system.

\textbf{The organisational barrier is sustainability.} The authors report that most telemedicine projects met sustainability difficulties, particularly after the end of the project period, and that a number of initiatives were funded mainly by external support and lacked continuity strategies and backing from key stakeholders---government, the private sector, non-governmental organisations and the wider community. They identify the absence of relations with stakeholders outside the health sector as a major obstacle, particularly for system maintenance policy. This is a recurrent failure of digital-health pilots: they demonstrate technical feasibility, then disappear with their funding.

\textbf{The barrier most directly opposable to the present work is workload.} The authors write: \bq{We identified that most health professionals in the region face a considerable workload. This has led health professionals to perceive the telemedicine system as an additional task entailing increased workload and time consumption. These factors constitute a vital challenge and can potentially constrain the progress of telemedicine projects.} Combined with the under-five-minute consultation constraint noted by the WHO (\cref{sec:time}), this observation sets a non-negotiable requirement: the system must reduce the provider's documentation burden, not increase it. A tool requiring additional data entry would hit both constraints at once and would not be adopted.

\textbf{The individual barrier} includes the lack of telemedical expertise among clinicians, which the authors mention alongside ethical and confidentiality concerns. They also note that the store-and-forward method, in which the patient's medical information is sent for expert opinion by e-mail or messaging, proved deficient in user satisfaction for lack of real-time experience.

\subsubsection{The documented Cameroonian context}

The appendix of \textcite{dodoo2022} lists, for each country, the telemedicine collaborations, services and systems identified. Cameroon, classified in the Economic Community of Central African States, appears with the following elements: a collaboration involving the telemedicine task force of the Movement Disorder Society, Laquintinie Hospital in Douala, Laboltel Hospital, the Cameroon telemedicine association and the Ministry of Public Health; the digitalMedLab system, developed in Zurich by a Cameroonian physician and an Australian developer; and the Telemed-cam project, in partnership with the telemedicine centre of Yaoundé General Hospital and rural health centres. The services covered are tele-education, maternal health, wound care and telecardiology; the systems used are the WoundDesk mobile application and the Telemedcam mobile application, in synchronous and asynchronous interaction. The timelines of the article also place in 2005 the use of telemedicine for maternal and child health in Cameroon through a mobile clinic network, and in 2012 the launch of digitalMedLab for wound care.

The authors finally note, as part of the favourable context, that 41 sub-Saharan countries have national digital health strategies and architectures, and that mobile phone connections, estimated at 816 million in 2019, should reach 1.05 billion in 2025, which facilitates wider use of telemedicine applications.

\subsubsection{The TELEMED-CAM trial}\label{sec:telemedcam}

Reference 47 of \textcite{dodoo2022} is the TELEMED-CAM study \parencite{kingue2013}. It is the only Cameroonian trial of hypertension telemanagement identified in the sources at hand, so any claim of originality must be positioned against it. (A later tele-ECG intervention in two district hospitals \parencite{bediang2021} addressed cardiovascular care more broadly and is treated in the state of the art.)%
\verify{Check the exact scope and results of \textcite{bediang2021} against the literature-review entry (source 17) before keeping this sentence.}

\paragraph{Design.} TELEMED-CAM was a prospective interventional pilot conducted in 2008 around Yaoundé (\cref{tab:telemedcam}). A telemedicine centre at Yaoundé General Hospital, with five cardiologists and three nephrologists, supported thirty rural health centres located 50 to 250~km away, each with a general practitioner. Six centres already running a Ministry of Health pilot unit for hypertension and diabetes were assigned to the control group; ten of the remaining twenty-four were drawn at random as intervention centres, which gives twenty control centres. Adults with hypertension above target (SBP $\geq 140$ or DBP $\geq 90$~mmHg, or $\geq 130/80$~mmHg for diabetes or nephropathy) were enrolled consecutively and followed for 24 weeks with twelve visits, one every two weeks.

\begin{table}[htbp]
\centering\footnotesize
\caption{The TELEMED-CAM pilot, as reported by \textcite{kingue2013}.}\label{tab:telemedcam}
\begin{tabularx}{\textwidth}{@{}>{\raggedright\arraybackslash}p{3.3cm} X@{}}
\toprule
\textbf{Item} & \textbf{Content} \\
\midrule
Design and period & prospective interventional pilot, 2008; 24~weeks, 12 post-baseline visits (every two weeks) \\
Setting & telemedicine centre at Yaoundé General Hospital (5 cardiologists, 3 nephrologists); 30 rural health centres 50--250~km away, each with a general practitioner \\
Allocation & 6 centres with an existing Ministry of Health pilot unit assigned to control; 10 of the other 24 drawn at random as intervention centres \\
Participants & 268 adults with hypertension above target: 165 in 10 intervention centres, 103 in 20 control centres \\
Intervention & standardised treatment protocol (generic drugs only), non-drug advice, training of GPs and nurses, GSM telecare: daily SMS or voice-mail reports from the centre, real-time phone feedback \\
BP measurement & at the facility, standard mercury manometer, three readings, mean of the last two \\
Baseline imbalance & intervention arm with higher SBP (169 vs 161~mmHg) and fewer diabetics (9.1\% vs 43.6\%) \\
Main results & more visits attended in the intervention arm; SBP falling from the second visit; at the final visit, stage-3 patients at target 50\% vs 39.1\%, stage I--II patients 65.2\% vs 70\% ($p=0.20$) \\
Limitations acknowledged & non-strictly random allocation; unblinded endpoint assessment; few practitioners and sites; clustering ignored in sample size and analysis; telehealth only one component of the package \\
\bottomrule
\end{tabularx}
\end{table}

\paragraph{Intervention.} The package combined a standardised treatment protocol restricted to generic drugs, non-pharmacological advice, training of GPs and nurses in treatment initiation and titration, and a telecare channel: each intervention centre received a mobile phone, staff reported patient data to the telemedicine centre daily by SMS or voice message, and the centre answered by phone in real time. Over the study, the telemedicine centre received 1\,187 calls from the intervention centres, about six calls per provider per day of about five minutes each, which we estimate at around thirty minutes of telephone time per provider per day.

\paragraph{Results.} Visit attendance was significantly higher in the intervention arm whatever the severity of hypertension (all $p \leq 0.04$). Systolic pressure fell from the second visit in the intervention arm and stayed constant in the usual-care arm. The benefit on control was confined to the most severe stratum: among stage-3 patients at baseline, 50\% of the intervention arm and 39.1\% of the usual-care arm were at target at the final visit, 33.3\% against 27.5\% had improved without reaching it, and 16.7\% against 33.3\% had not improved ($p=0.04$). Among stage I--II patients the proportions at target were similar, 65.2\% and 70\% ($p=0.20$).%
\verify{The source is internally inconsistent on baseline DBP (reported higher in the intervention arm, 100 vs 95~mmHg, then described as significantly lower) and labels the 16.7\% group ``control groups''. Reading 16.7\% as the intervention arm is the only assignment for which both columns sum to about 100\%.}

\paragraph{What it establishes, and what it does not.} TELEMED-CAM shows that a remote-support package for rural primary-care centres was feasible with basic GSM infrastructure in Cameroon, and that it was associated with better visit attendance and, in severe cases, better control. This is consistent with the argument of this section, that retention and decision support are the levers (\cref{sec:ltfu,sec:inertia}). Four limits keep it from settling the question addressed here.
\begin{itemize}[nosep]
  \item The effect cannot be attributed to telehealth: the authors themselves state that the standardised protocol, education materials and training of providers were applied only in the intervention arm.
  \item The baseline imbalance (higher starting SBP, far fewer diabetics) leaves more room for within-arm decline, which regression to the mean can mimic; this is our reading, not the authors'.
  \item Support was clinician-facing and facility-based. Blood pressure was measured at the facility with mercury manometers, which the WHO now recommends replacing (\cref{sec:metrology}), and no home measurement or patient-level series was collected.
  \item The channel loaded providers with daily reporting and phone calls, the kind of added task that \textcite{dodoo2022} identify as a barrier (\cref{sec:digital-barriers}).
\end{itemize}
The present thesis differs on three points: the unit of support is the patient's longitudinal series, not a clinician's phone query; the data combine consultation records, repeated home measurements and possibly mobile sensors within a single schema; and a learned risk model flags decompensation early instead of relying on protocol-based phone advice. TELEMED-CAM used no machine learning, risk score or automated alert.

%==========================================================================
\subsection{The metrological constraint on measuring devices}\label{sec:metrology}

This point was not in the initial plan and must be included, because it conditions the feasibility of the ``home measurements and mobile sensors'' component of the envisaged system. The two available sources diverge on this point, and this divergence is itself instructive.

\subsubsection{The WHO position}

The WHO first establishes that major barriers to obtaining quality measurements are reported in both low- and high-resource settings, and concern device type, accuracy and unawareness of purchasing criteria. It then states three normative positions.

\textbf{Manual devices must be replaced} by independently validated automated electronic devices, because of environmental concerns about mercury, calibration problems of aneroid devices and defects of manual measurement technique. Automated upper-arm cuff devices are strongly recommended, because they eliminate human error---notably that induced by hearing deficit or digit preference---and require less training.

\textbf{Cuffless devices are not recommended.} The WHO writes that many types of cuffless devices have entered the global market, including popular connected wristbands, and that despite their considerable potential to take many measurements day and night without the user having to pay attention, their use is currently not recommended for the management of hypertension, owing to unresolved questions about their accuracy, performance and ability to track pressure variations. It concludes that in 2023 there is no alternative to the validated upper-arm cuff with sufficient consistency and accuracy for diagnosis and management in health facilities.

\textbf{Validation is rare.} Only one in five of the available automated upper-arm or wrist cuff devices has evidence of accuracy validation testing, and validated devices are generally more expensive. The experience of the Region of the Americas is documented: at the launch of HEARTS, not only did countries lack regulations ensuring that only validated devices were marketed, but ministries of health had no register of the types, brands and models used in their primary-care facilities, the Latin American market being dominated by non-validated devices.

\subsubsection{The position of Navarro Masuet and the arbitration}

\textcite{navarro2023} presents a more permissive view. She describes wrist-worn devices as \bq{non-invasive, cuffless blood pressure monitors measuring systolic and diastolic pressures without an inflatable cuff}, qualified as \bq{easy to use, safe and relatively accurate for measurement at rest}. She herself, however, signals the technical limits of each method: oscillometry \bq{has accuracy limitations and is not reliable during physical activity because of motion noise} and \bq{requires validation for specific populations and subjects}; applanation tonometry \bq{requires individual calibration and is not suited to routine clinical settings}. She finally concludes that only wearable devices are available for blood pressure telemonitoring, the attachable, implantable and ingestible categories being empty for this parameter.

The arbitration is unambiguous. The WHO is an institutional normative source, Navarro Masuet a Bachelor thesis; the technical limits that the latter reports herself---lack of accuracy, sensitivity to motion, need for individual calibration---moreover converge with the WHO position rather than contradict it. We retain the WHO position.

\subsubsection{Consequences for the design}

Three consequences must be explicitly assumed:
\begin{enumerate}[nosep]
  \item The ``mobile sensors'' component cannot rely on cuffless devices for blood pressure measurement. These can provide context variables---physical activity, heart rate, sleep---useful as model covariates, but blood pressure must come from a validated automated upper-arm monitor. This restriction does not amputate the project: it moves it from the ambition of unconstrained continuous measurement to that of reliable repeated measurement, which is what the clinical protocol actually requires.
  \item The list of admissible devices must be restricted to devices on a recognised validation register, and this criterion must be part of the system specification.
  \item The cost of the device, higher for validated devices, must be integrated into the economic feasibility analysis alongside the cost of treatment, which the WHO puts at 5 to 10 US dollars per patient per year for quality-assured medicines, amlodipine and hydrochlorothiazide generally being available at 0.01 dollar per tablet or less, but with prices varying more than tenfold between countries and mark-ups recorded in Kenya exceeding 340\% for amlodipine and 280\% for hydrochlorothiazide.
\end{enumerate}

%==========================================================================
\subsection{From the hospital-centred model to the patient-centred model}\label{sec:patient-centred}

\subsubsection{The normative mandate}

The WHO formulates this shift explicitly. In its module on team-based care, after recalling that the world faces a lasting shortage of trained health personnel---the WHO projecting a global deficit of 10 million health workers by 2030, mainly in low- and middle-income countries---and that team-based care is an effective and cost-effective response, the report states that \bq{telemonitoring, community or home blood pressure measurement, and support by health workers are encouraged to improve hypertension management}, within an integrated chronic disease management system, when the medical team judges community care appropriate and patients find it feasible and affordable.

The report's statement on the function of telemonitoring, and its conditions for spacing visits (multi-month dispensing, ideally ninety days, and community or home delivery), were given in \cref{sec:hypothesis} and are not repeated here.

The proposed shift is therefore not a substitution of remote follow-up for consultation, but a redistribution (\cref{fig:redistribution}): the health facility keeps diagnosis, initiation and therapeutic intensification; the home and the community take on repeated measurement, retention and routine medicine dispensing.

\begin{figure}[htbp]
\centering
\begin{tikzpicture}[
  box/.style={draw,rounded corners,align=center,minimum width=5.2cm,minimum height=1.6cm,font=\small},
  arr/.style={-{Stealth[length=2.5mm]},thick}]
  \node[box,fill=cRef!12] (fac) {\textbf{Health facility}\\ diagnosis, treatment initiation,\\ therapeutic intensification};
  \node[box,fill=cAfrica!10,right=3.4cm of fac] (home) {\textbf{Home and community}\\ repeated measurement, retention,\\ routine dispensing};
  \draw[arr] ([yshift=0.3cm]fac.east) -- node[above,font=\footnotesize]{recalls, decisions} ([yshift=0.3cm]home.west);
  \draw[arr] ([yshift=-0.3cm]home.west) -- node[below,font=\footnotesize]{series, alerts} ([yshift=-0.3cm]fac.east);
\end{tikzpicture}
\caption{Redistribution of tasks between the health facility and the home or community, as proposed in this thesis.}\label{fig:redistribution}
\end{figure}

\subsubsection{The formulation of Navarro Masuet}

The thesis of \textcite{navarro2023}, written from the Spanish context, formulates the same shift from a different pressure: not the shortage of health workers and loss to follow-up, but demographic ageing and growing demand. She reports that the proportion of people aged 65 and over in Spain, currently 20.1\% of the population, should peak at 30.4\% around 2050, about 18.2 million people; that health expenditure attributed to the over-65s is about three times higher than for the rest of the population; and that chronic conditions already account for more than 80\% of primary-care consultations.

She deduces that \bq{the use of teleconsultation or telemonitoring allows the development of new forms of care, such as home hospitalisation units}, and that these tools enable \bq{remote monitoring of patients, facilitating the follow-up of their health status and reducing the need for unnecessary travel}.

This convergence between two very dissimilar contexts, a Spanish system under demographic pressure and a sub-Saharan system under resource pressure, deserves to be noted, but also qualified. The two systems reach the same organisational conclusion for different reasons, and implementation constraints are in no way comparable: what is, in Spain, an optimisation choice requires, in Cameroon, first resolving the connectivity, funding, workload and regulatory barriers documented by \textcite{dodoo2022}. Transposition of the model is therefore not automatic, and it is precisely the object of the present work to examine its conditions.

\subsubsection{The absence of artificial intelligence}\label{sec:ai-gap}

Two of our sources converge on a point that defines the space of contribution of this thesis.

\textcite{navarro2023}, at the end of her review, writes: \bq{it is surprising to note how little consideration is given to the technological advances that are revolutionising all sectors with the use of artificial intelligence. Although modern devices use algorithms to provide an interpretation of the data obtained, it is regrettable that more room is not given to this aspect.} She specifies the form that this contribution would take: \bq{there is greater potential for using the capabilities offered by telemonitoring in conjunction with AI. This would involve obtaining, together with the direct or indirect measurements of the diagnostic variable under study, other data on the patient's general condition, environment and habits. With the integration of this tool into telemedicine devices, it would be possible not only to obtain information about the user's current state, but also a more complete, and therefore more useful, view of his or her overall state and foreseeable evolution.}

This formulation describes exactly the architecture proposed in this thesis: the aggregation of clinical measurement, contextual data and behavioural data in a predictive rather than descriptive model.

\textcite{liu2025review}, for their part, report a direct precedent. They cite the work of Chatrati and colleagues, who proposed \bq{a home health monitoring system allowing remote analysis of patients' blood pressure and glucose indices at home. This system predicts the status of hypertensive and diabetic patients by combining conditional decision-making and machine learning, and notifies healthcare providers when an anomaly is detected.}

\textcite{chatrati2022}, published in the \emph{Journal of King Saud University -- Computer and Information Sciences} 34(3):862--870, describe in their abstract a smart home health monitoring system that analyses a patient's blood pressure and glucose readings at home and notifies the healthcare provider when an abnormality is detected. A combination of conditional decision-making and machine-learning approaches is used to predict hypertension and diabetes status, respectively, from the readings; a graphical interface sends categorised alerts and real-time notifications to the registered physician or clinic. Read literally, the abstract assigns rules to blood pressure and a learned classifier to glucose.%
\verify{Only the abstract and the bibliographic record were available; the full text was not consulted. Check the dataset(s) and whether any real patient data or home deployment was used, the thresholds of the rule-based part, the reported performance, whether windows of readings (trends) or single readings are modelled, and how labels are defined, then complete this paragraph.}
This is the closest technical precedent to the system envisaged here. What the abstract does not settle, and what will decide how close it really is, is whether the model sees trends or single readings, and whether it was evaluated on real home data. It will be examined in detail in the technical state of the art.

These findings, together with the absence of studies using artificial intelligence noted for sub-Saharan Africa (\cref{sec:hypothesis}), outline a precise situation: the combination is reported as technically feasible elsewhere, it is called for by the specialised literature, and it has not been evaluated in the context at hand.

\subsubsection{Summary diagnosis}

\begin{table}[htbp]
\centering\footnotesize
\caption{Failures of the current model and the functions required of the proposed system.}\label{tab:synthesis}
\begin{tabularx}{\textwidth}{@{}>{\raggedright\arraybackslash}p{4.4cm} >{\raggedright\arraybackslash}X >{\raggedright\arraybackslash}X@{}}
\toprule
\textbf{Failure of the current model} & \textbf{Mechanism} & \textbf{Required function} \\
\midrule
Diagnosis requires two visits one to four weeks apart & non-return $\Rightarrow$ diagnosis unconfirmed & automated recall and recording of the second measurement \\
Single measurement not representative of $\mu_i$ & white-coat hypertension ($\sim$13\%) and masked hypertension, the latter more frequent in diabetics & aggregation of repeated measurements, estimation of $\mu_i$ and $\sigma_\varepsilon$ \\
Monthly follow-up required for the uncontrolled, about 1 to 2 contacts per year assumed delivered & therapeutic under-titration & continuity of contact without travel \\
Consultation $<$ 5~min and perceived workload & rejection of any tool adding a task & single data entry, interface with no additional burden \\
Therapeutic inertia & impossibility of verifying treatment efficacy & time series and trend detection, not isolated events \\
Duplicated, non-longitudinal, unavailable paper record & loss of the time dimension & single accessible longitudinal record \\
Information silos and heterogeneous formats & fragmentation re-formed in the digital environment & single documented data schema \\
HEARTS indicators not computable & steering blind & automatic production of the four core indicators \\
Loss to follow-up $\geq$ 50\% & simultaneous rupture of links 2 and 3 & identification of overdue patients, recall \\
Shortage of 10 million health workers by 2030 & impossible to densify in-person follow-up & shift part of the follow-up to the home \\
Dependence on external funding & disappearance of the system at the end of the project & documentation of cost and workload in steady state \\
\bottomrule
\end{tabularx}
\end{table}

%==========================================================================
\subsection{Scope, limitations and problem statement}\label{sec:scope-problem}

\subsubsection{Four reservations}

\textbf{The nature of the sources.} The WHO report is a normative and synthesis document, not a primary study. Its statements on the ineffectiveness of paper supports and on the extent of loss to follow-up refer to references that we have not consulted directly; they have the value of an institutional position grounded in a review, not of an experimental result. The figure of 50\% loss to follow-up is moreover introduced by the adverb ``often'' and has no interval: it is an order of magnitude, not an estimate.

\textbf{Heterogeneity of the level of evidence.} Of the four main sources of this section, only one, \textcite{dodoo2022}, is a systematic review with a documented protocol. \textcite{liu2025review} present themselves as such without reporting the elements. \textcite{navarro2023} is a Bachelor thesis, and the clinical testimony it reports comes from a single informant recruited by proximity. Statements drawn from these last two sources are used as formulations of the problem or as illustrations, never as evidence, and are systematically cross-checked against the WHO when they concern a factual point.

\textbf{Generalisation.} The WHO findings concern low- and middle-income countries as a whole, and those of \textcite{dodoo2022} a sub-Saharan Africa whose uneven character they precisely stress. Those of \textcite{navarro2023} concern Spain and those of \textcite{liu2025review} an unspecified international corpus. None of the sources used documents specifically the organisation of care and record keeping in Cameroonian facilities managing chronic diseases; TELEMED-CAM (\cref{sec:telemedcam}) documents rural primary-care centres and a telecare package, not their record keeping. This gap will have to be filled by the field-survey phase of the present work, which is in itself a contribution: characterising the real information flows of such a facility yields data that do not exist in the literature consulted.

\textbf{Causality.} The examples of Sri Lanka, the Philippines, India and Cuba establish an association between high retention and high control, not a causal relationship. These programmes modified several components simultaneously---simplified protocol, medicine supply, team-based care, information system, community follow-up---so that it is impossible to attribute the result to any one of them. TELEMED-CAM shows the same limit at trial scale: the protocol, education and training were bundled with the telecare channel. The Indian example illustrates it negatively: the supply disruption discouraged patients from returning, which shows that retention is itself a dependent variable. A telemonitoring system deployed in a context of medicine stock-outs would probably not improve control, whatever its technical performance. This limit will have to be integrated into the evaluation protocol as a control variable documenting medicine availability over the observation period.

\subsubsection{Problem statement}

At the end of this diagnosis, the problem that this thesis intends to address can be formulated as follows.

The current care model produces clinical information that is point-in-time, discontinuous and non-aggregable, whereas the two conditions considered require longitudinal, repeated information that can be used as a trend. This gap simultaneously explains three failures that the sources allow us to establish separately: the insufficiency of diagnosis, which requires two visits and out-of-office measurements to rule out white-coat and masked hypertension; therapeutic inertia, which results from the impossibility of verifying the persistence of an uncontrolled pressure and hence the efficacy of a treatment; and the impossibility of steering programmes with the indicators that the WHO defines as core, two of four of which are temporal.

The research question is therefore not whether the register should be digitised---the WHO has settled this and specified the expected outputs---but to determine which longitudinal information to collect, at what frequency, with which admissible devices, and what processing of this information makes it possible to detect early enough the patients heading towards decompensation to trigger a useful clinical action, under four constraints established in this section: the metrological validation of measuring devices, the consultation time under five minutes, the workload perceived by providers who reject any additional task, and the requirement of sustainability beyond project funding.

This question is addressed in the next section, devoted to the state of the art of telemonitoring systems and detection methods.

%##########################################################################
% SECTION 1.3
%##########################################################################
\section{Modalities of Out-of-Hospital Follow-up: Mobile Data Entry (mHealth) and Self-Measurement}\label{sec:outpatient}

\subsection*{Purpose of the section and status of the sources}

The preceding section established that the episodic consultation model produces point-in-time, discontinuous information, whereas hypertension and diabetes require longitudinal information usable as a trend. This section examines the concrete modalities by which that longitudinal information can be produced outside the care facility: home self-measurement, data entry through a mobile application, and the mechanisms intended to support patient adherence to these practices.

Four sources are mobilised, whose evidential status is again very uneven.

\paragraph{Zhou et al. (2024).} \textcite{zhou2024}, published in the \emph{Journal of Medical Internet Research} 26:e54978, is a systematic review with meta-analysis of randomised controlled trials, conforming to PRISMA recommendations and registered with PROSPERO under reference CRD42023483746. It provides the pooled effect estimate and the analysis of active ingredients.

\paragraph{Mihevc et al. (2025).} \textcite{mihevc2025}, published in \emph{JMIR mHealth and uHealth} 13:e59733, is a multicentre randomised controlled trial, registered in the ISRCTN registry under number 31471852, with a separately published protocol and reporting conforming to the CONSORT recommendations. It is the highest level of evidence in this section and the only source addressing the hypertension--type~2 diabetes comorbidity.

\paragraph{Alessa et al. (2019).} \textcite{alessa2019}, published in \emph{JMIR mHealth and uHealth} 7(5):e13645, is an exhaustive content analysis of the commercial applications available. It measures no efficacy whatsoever: it describes a state of the market. This distinction is essential and we maintain it throughout the section.

\paragraph{Navarro Masuet (2023).} \textcite{navarro2023} remains a Bachelor thesis in biomedical engineering. We retain its taxonomy of telemedicine modalities, which usefully structures the exposition, and not its factual statements, which we cross-check systematically.

%==========================================================================
\subsection{Conceptual framework: where self-measurement is situated}\label{sec:framework}

\subsubsection{The absence of a consensual taxonomy}

\textcite{navarro2023} notes at the outset a difficulty that any work on this subject encounters: \bq{there is currently no universally accepted classification of the types and applications of telemedicine}. She identifies three competing proposals in the literature---a division into synchronous telemedicine, asynchronous telemedicine and telemonitoring; a division into four categories, teleconsultation, telediagnosis, telemonitoring and telesurgery; and a division adding teleradiology---and concludes that no consensus exists.

She consequently proposes a simplified framework in five branches, which we adopt for clarity of exposition while recalling that it is an author's proposal and not an established norm (\cref{tab:branches}).

\begin{table}[htbp]
\centering\small
\caption{Branches of telemedicine as proposed by \textcite{navarro2023}.}\label{tab:branches}
\begin{tabularx}{\textwidth}{@{}>{\raggedright\arraybackslash}p{3.1cm} X@{}}
\toprule
\textbf{Branch} & \textbf{Definition retained} \\
\midrule
Teleconsultation & direct contact between patient and physician through a platform, without an intermediary \\
Telediagnosis & remote transmission of records and examination results to a specialist \\
Telemonitoring & digital transmission of physiological data from the home to care centres, with clinical feedback \\
Teletreatment & performance of therapeutic acts at a distance \\
Tele-education & continuing education of patients and professionals \\
\bottomrule
\end{tabularx}
\end{table}

\subsubsection{Synchronous and asynchronous}

The distinction that is most operative for our object is however not that of the branches but that of the modes of interaction, which \textcite{navarro2023} defines as follows.

The \emph{asynchronous} mode, commonly designated by the expression \emph{store-and-forward}, consists in recording medical information in one place and then transmitting it to another place for later evaluation by a specialist. It allows the transfer of pre-recorded information without simultaneous presence of the parties. It is widely employed in dermatology, neurology, ophthalmology, cardiology and otorhinolaryngology.

The \emph{synchronous} mode, also designated as real-time, interactive or live, presupposes an immediate interaction between participants located in two or more places, hence a simultaneous presence. It is particularly useful for rapid decisions, notably in emergency situations, and allows a second opinion to be obtained immediately.

This distinction has a direct consequence for the design of the proposed system: home self-measurement belongs by construction to the asynchronous mode. The patient records, the system transmits, the provider evaluates later. The performance of the device therefore does not depend on simultaneity, but on the quality of the stored information and on the relevance of the trigger that will bring the provider back to the record. This observation fixes the central technical problem of the following section: in an asynchronous flow, it is the triggering rule that carries the clinical value, since it determines what rises to human attention and what does not.

It must be noted here that \textcite{dodoo2022}, whose results we reported in \cref{sec:digital-barriers}, signal that the store-and-forward method proved deficient in user satisfaction in sub-Saharan Africa for want of a real-time experience, and recommend recourse to videoconferencing. This reservation does not disqualify the asynchronous mode for the follow-up of vital parameters---which does not lend itself to real time---but it indicates that a purely asynchronous device will have to provide a return channel towards a human interaction, failing which the patient will lose the sense of being followed. This is precisely what the architecture of \textcite{mihevc2025} achieves, examined below, through the introduction of triggered teleconsultations.

\subsubsection{The definition of telemonitoring}

\textcite{navarro2023} defines telemonitoring as \bq{an innovative strategy that digitally transmits disease-related and physiological data from patients' homes to care centres, which provide clinical feedback}. She attributes four effects to it: early detection of disease decompensation, reduction of mortality and hospitalisations, empowerment of patients in their own management, and reduction of physical contacts, waiting times and costs, with a lightening of the medical staff's workload.

These attributions are taken from the literature she cites and do not constitute her own results; we treat them as hypotheses to be verified. \Cref{sec:efficacy} confronts precisely the first of them---early detection---with the results of the trials.

%==========================================================================
\subsection{Self-measurement protocols}\label{sec:protocols}

\subsubsection{Blood pressure self-measurement}

The WHO, whose protocol we reported in \cref{sec:dx-requires}, sets the conditions of measurement: trained staff or patient, cuff of appropriate size, seated position with back supported, empty bladder, legs uncrossed, two readings of which the second is retained. It further establishes that only independently validated automated electronic upper-arm cuff devices are recommended, and that no cuffless alternative was recommended in 2023 (\cref{sec:metrology}).

\textcite{mihevc2025} operationalise this protocol in a real telemonitoring device, and their description constitutes the most precise specification at our disposal: patients took two readings in the morning and in the evening, after five minutes of rest in a seated position, with back and arm supported. A one-minute interval was programmed by the device between the two readings, and the second was retained as valid.

Two elements of this description deserve to be retained for the design. First, the interval between the two readings is imposed by the device and not left to the patient's initiative, which removes a major source of non-compliance with the protocol. Second, the validity rule---retain the second reading---is applied automatically, and not by a computation demanded of the patient.

For establishing the reference value at the assessment time points, the authors apply a stricter rule still: the second morning reading and the second evening reading are averaged over seven days. This is the direct application of the variance-reduction logic set out in \cref{sec:variability}: with fourteen retained measurements, the standard error of the estimator of the usual pressure is divided by $\sqrt{14} \approx 3.7$ relative to an isolated measurement.

\subsubsection{Capillary blood glucose self-monitoring}

\textcite{mihevc2025} describe a glycaemic protocol articulated in two regimes.

In the \emph{routine} regime, the patient performs two measurements: one in the morning while fasting and one ninety minutes after lunch.

In the \emph{verification} regime, triggered by an abnormal value, the patient performs a one-day glycaemic profile comprising six measurements: fasting, 90 minutes after breakfast, before lunch, 90 minutes after lunch, before dinner, and 90 minutes after dinner.

The therapeutic targets retained, adjusted for an elderly population, are a fasting glucose below 7~mmol/L and a postprandial glucose below 10~mmol/L. For Cameroonian practice, which usually expresses blood glucose in grams per litre, these thresholds correspond to about 1.26~g/L and 1.80~g/L, the conversion factor being $1\ \text{mmol/L} \approx 0.18\ \text{g/L}$.

The blood pressure target retained is a home pressure below 135/85~mmHg---and not 140/90~mmHg---in accordance with the convention whereby self-measurement thresholds are five points below consultation thresholds, the absence of a white-coat effect lowering the values measured at home.

\subsubsection{The decision architecture of the joint protocol}\label{sec:triage}

The most directly exploitable contribution of \textcite{mihevc2025} is their three-level classification system, applied symmetrically to both conditions (\cref{tab:triage-levels}).

\begin{table}[htbp]
\centering\small
\caption{Three-level classification of \textcite{mihevc2025}.}\label{tab:triage-levels}
\begin{tabularx}{\textwidth}{@{}>{\raggedright\arraybackslash}p{1.7cm} >{\raggedright\arraybackslash}p{3.6cm} X@{}}
\toprule
\textbf{Code} & \textbf{Interpretation} & \textbf{Action triggered} \\
\midrule
Green & normal values & no change to the follow-up protocol \\
Yellow & abnormal values & intensified surveillance: seven-day blood pressure profile, or one-day six-measurement glycaemic profile \\
Red & critically abnormal values & immediate teleconsultation with the general practitioner, further assessment or treatment modification \\
\bottomrule
\end{tabularx}
\end{table}

The operational thresholds, as presented in Figure~1 of the article, are reproduced in \cref{tab:thresholds}. We have reorganised them by decision level and by meal timing, because the original presentation leaves the reading of the glycaemic intervals ambiguous.

\begin{table}[htbp]
\centering\footnotesize
\caption{Operational thresholds of the two protocols \parencite{mihevc2025}, reorganised by decision level. BP in mmHg, glucose in mmol/L.}\label{tab:thresholds}
\begin{tabularx}{\textwidth}{@{}>{\raggedright\arraybackslash}p{2.4cm} >{\raggedright\arraybackslash}X >{\raggedright\arraybackslash}X@{}}
\toprule
\textbf{Level} & \textbf{Blood pressure (current reading)} & \textbf{Blood glucose (current reading)} \\
\midrule
Red (low) & SBP $<$ 110 or DBP $<$ 60 & $<$ 4 (any timing) \\
Green & 110--134 / 60--84 & 4--7 before a meal; 7--10 after a meal \\
Yellow & 135--174 / 85--104 & 7.1--10 before a meal; 10.1--14 after a meal \\
Red (high) & SBP $\geq$ 175 or DBP $\geq$ 105 & $>$ 10 before a meal; $>$ 14 after a meal \\
\midrule
On the profile mean & $<$ 110/60 or $\geq$ 135/85 $\Rightarrow$ teleconsultation & $>$ 7 before a meal or $>$ 10 after a meal $\Rightarrow$ teleconsultation \\
\bottomrule
\end{tabularx}
\end{table}

This architecture possesses three remarkable properties, which it is appropriate to state explicitly because they directly orient the design of the proposed system.

\textbf{It is bidirectional.} Low thresholds trigger as much as high ones. For diabetes this is an obvious clinical necessity---hypoglycaemia is an acute risk---but the device applies the same logic to blood pressure, hypotension in the treated elderly patient being a risk of falls. A system detecting only elevated values would be clinically incomplete.

\textbf{It is two-stage.} An isolated abnormal value does not trigger a clinical action but an additional collection of information. Only in the light of the profile mean is the teleconsultation triggered. This is exactly the operational answer to the variance problem set out in \cref{sec:variability}: the device distinguishes noise from trend by increasing $n$ before deciding, rather than by lowering the threshold (\cref{fig:triage}).

\textbf{It is economical in human attention.} Only the red class mobilises a physician. The yellow class mobilises the patient. In a context where, as \textcite{dodoo2022} establish, professionals perceive telemonitoring as an additional workload (\cref{sec:digital-barriers}), this distribution is a condition of adoption and not a refinement.

\begin{figure}[htbp]
\centering
\begin{tikzpicture}[node distance=0.55cm and 0.9cm,
  st/.style={draw,rounded corners,align=center,font=\footnotesize,minimum height=0.85cm,minimum width=2.5cm},
  dec/.style={draw,diamond,aspect=2,align=center,font=\footnotesize,inner sep=1pt,fill=cGrey!12},
  ar/.style={-{Stealth[length=2.5mm]},thick}]
  \node[st,fill=cGrey!12] (meas) {Routine reading\\ (low frequency)};
  \node[dec,right=of meas] (cl) {Class?};
  \node[st,fill=green!14,above right=0.5cm and 1.0cm of cl] (gr) {\textbf{Green}\\ no change};
  \node[st,fill=yellow!28,right=1.0cm of cl] (ye) {\textbf{Yellow}\\ verification profile\\ (7 days / 6 readings)};
  \node[st,fill=cAfrica!18,below right=0.5cm and 1.0cm of cl] (re) {\textbf{Red}\\ immediate\\ teleconsultation};
  \node[st,fill=cRef!12,right=1.0cm of ye] (mn) {Profile mean\\ vs target};
  \draw[ar] (meas) -- (cl);
  \draw[ar] (cl) -- (gr);
  \draw[ar] (cl) -- (ye);
  \draw[ar] (cl) -- (re);
  \draw[ar] (ye) -- (mn);
  \draw[ar] (mn.south) |- node[below right,font=\scriptsize,pos=0.75]{off target} (re.east);
  \draw[ar] (mn.north) |- node[above right,font=\scriptsize,pos=0.75]{on target} (gr.east);
\end{tikzpicture}
\caption{The two-stage triggering architecture of \textcite{mihevc2025}. An isolated abnormal value increases the number of measurements instead of mobilising a clinician; only the mean of the verification profile reaches a clinical decision.}\label{fig:triage}
\end{figure}

\subsubsection{The question of frequency}

The sources diverge on the frequency of the routine regime, and this divergence calls for clarification before any reuse.

The abstract and the body of the article of \textcite{mihevc2025} indicate that participants measured their blood pressure twice a week and their blood glucose once a week. Figure~1 of the same article, by contrast, carries the labels ``weekly blood pressure measurements'' and ``monthly blood glucose measurements''. These two statements are incompatible. We retain the textual description, confirmed by the abstract, and we signal the figure as erroneous, without being able to settle the matter definitively without recourse to the separately published protocol.%
\verify{Resolve this frequency divergence against the published protocol of \textcite{mihevc2025} (ISRCTN31471852) before reusing either value as a design parameter.}

The frequency trade-off must in any event be posed explicitly, because it sets two contrary requirements against each other. Variance reduction argues for maximal frequency, since $\operatorname{Var}(\hat\mu) = \sigma_\varepsilon^2 / n$. Adherence argues for minimal frequency, every measurement demanded being an occasion for abandonment. The device of \textcite{mihevc2025} resolves this tension by modulation: a low routine frequency, sufficient for surveillance, and a high frequency triggered only when a signal justifies it. This is a solution the proposed system should adopt, because it optimises both criteria simultaneously instead of sacrificing one of them.

One residual ambiguity of the threshold table must be noted, because any implementation will have to resolve it: the blood pressure rules are stated on two variables, and the source does not say what happens when the two disagree---a systolic value in the yellow band together with a diastolic value in the low red band, for instance. The safe convention, which we adopt, is that the most severe of the two classifications prevails, red before yellow and yellow before green. This is our decision, not a statement of the source.

%==========================================================================
\subsection{Mobile data entry: what applications actually do}\label{sec:apps}

\subsubsection{A state of the market, not a measure of efficacy}

\textcite{alessa2019} surveyed the entirety of the applications supporting hypertension self-management available in January 2018 on the two main stores, in the English language, in the United Kingdom. Their protocol is documented: identification of 775 applications---495 on Google Play, 280 on the App Store---exclusion of 564 for duplication, language, irrelevance, intended use by professionals or attachment to a specific study, then exclusion of 25 for failure to function, yielding 186 included applications, of which 106 on Android and 80 on iPhone. Each application was installed and run for two days in order to observe the reminders and notifications actually triggered. Thirty-three presented so-called comprehensive functionality, that is, at least three functions, and thirty were subjected to in-depth analysis.

It must be stressed that this study measures no clinical efficacy. It describes what the market offers. Confusing this inventory with an evaluation would be a reading error, and the authors themselves insist on this point, concluding that \bq{the technical usability and effectiveness of hypertension self-management support applications need to be assessed as a matter of urgency}.

\subsubsection{The distribution of functionality}

\begin{table}[htbp]
\centering\small
\caption{Functions present in the 186 applications surveyed \parencite{alessa2019}, Table~1. Percentages recomputed from the counts.}\label{tab:app-features}
\begin{tabularx}{\textwidth}{@{}X r r r@{}}
\toprule
\textbf{Function} & \textbf{iPhone ($n=80$)} & \textbf{Android ($n=106$)} & \textbf{Total ($n=186$)} \\
\midrule
Educational information & 34 (42.5\%) & 76 (71.7\%) & 110 (59.1\%) \\
Self-monitoring & 56 (70.0\%) & 43 (40.6\%) & 99 (53.2\%) \\
Feedback & 36 (45.0\%) & 16 (15.1\%) & 52 (28.0\%) \\
Data export & 29 (36.3\%) & 22 (20.8\%) & 51 (27.4\%) \\
Reminder & 23 (28.8\%) & 16 (15.1\%) & 39 (21.0\%) \\
Goal setting & 8 (10.0\%) & 2 (1.9\%) & 10 (5.4\%) \\
Stress management & 3 (3.8\%) & 2 (1.9\%) & 5 (2.7\%) \\
Communication with others & 4 (5.0\%) & 1 (0.9\%) & 5 (2.7\%) \\
\bottomrule
\end{tabularx}
\end{table}

The percentages of \cref{tab:app-features} have been recomputed from the counts, because three of the values published in the source do not follow from its own numerators: educational information on iPhone is given as 43.8\% where $34/80 = 42.5\%$, on Android as 70.8\% where $76/106 = 71.7\%$, and feedback on iPhone as 43.8\% where $36/80 = 45.0\%$. The counts themselves are internally consistent, since each platform pair sums to the stated total, so we retain the counts and present the recomputed proportions (\cref{fig:app-features}).

\begin{figure}[htbp]
\centering
\begin{tikzpicture}
\begin{axis}[ybar,width=\textwidth,height=6.4cm,bar width=7pt,ymin=0,ymax=80,
  ylabel={Applications offering the function (\%)},
  symbolic x coords={Education,Self-monitoring,Feedback,Export,Reminder,Goals,Stress,Communication},
  xtick=data,xticklabel style={rotate=35,anchor=east,font=\scriptsize},
  enlarge x limits=0.08,legend pos=north east,legend style={font=\scriptsize}]
\addplot[fill=cRef!75,draw=cRef] coordinates
  {(Education,42.5)(Self-monitoring,70.0)(Feedback,45.0)(Export,36.3)
   (Reminder,28.8)(Goals,10.0)(Stress,3.8)(Communication,5.0)};
\addplot[fill=cAfrica!75,draw=cAfrica] coordinates
  {(Education,71.7)(Self-monitoring,40.6)(Feedback,15.1)(Export,20.8)
   (Reminder,15.1)(Goals,1.9)(Stress,1.9)(Communication,0.9)};
\legend{iPhone ($n=80$),Android ($n=106$)}
\end{axis}
\end{tikzpicture}
\caption{Distribution of functions by platform, recomputed from the counts of \textcite{alessa2019}. The two platforms differ in profile: Android favours educational content, iPhone self-monitoring and feedback.}\label{fig:app-features}
\end{figure}

The most important result of this table is not the frequency of each function but their combination. One hundred and eight applications out of 186, that is 58.1\%, offer only one function. Forty-five, that is 24.2\%, offer two. Only 33, that is 17.7\%, offer three or more, and none offers all eight. The most frequent combination is self-monitoring together with automatic feedback, present in 24 applications out of 186, that is 12.9\%.

This finding matters because the authors recall, on the strength of an earlier review, that applications with comprehensive functionality are more likely to lower blood pressure than those offering only one or two. The market therefore predominantly offers tools whose configuration is, a priori, the least effective.

\subsubsection{The data-entry problem}

The most directly operational detail for the present work concerns the mode of data entry. Among the 99 applications offering self-monitoring, 94 allow blood pressure tracking. Seven received the values automatically from the monitor while retaining a manual entry option; of the remainder, 73 required manual entry and 11 allowed automatic transfer.%
\verify{The counts do not close: $7 + 73 + 11 = 91$, not 94, and the source describes the 73 and the 11 as drawn from ``the remaining 84''. The three unaccounted applications are plausibly the three that claim to turn the smartphone itself into a measuring device (\cref{sec:cuffless}), which would give $7+73+11+3 = 94$. Confirm against Table~1 and the Results section of \textcite{alessa2019} before citing these figures.}

In other words, in the great majority of commercial applications the patient must retype what the device has just displayed. This redundancy has three consequences. It creates a burden, hence a motive for abandonment. It introduces a source of error---transposition of digits, inversion of systolic and diastolic---that nothing subsequently allows to be detected. And it makes the patient the sole guarantor of the time stamp, hence of the temporal coherence of the series.

The architecture retained by \textcite{mihevc2025} resolves this problem entirely: an A\&D~651 BLE blood pressure monitor and a Contour Next One glucose meter transmit by Bluetooth to a mobile application serving as a concentrator, which retransmits to a cloud platform over a 4G or 5G network, the whole being encrypted and compliant with the General Data Protection Regulation. No data entry is demanded of the patient. This choice must be adopted in the proposed system, and it is moreover a condition of the validity of the data collected: a manually entered series does not satisfy the quality requirements identified in \cref{sec:records}, where we noted that the identically carried-over value is the most dangerous artefact for a trend-detection model.

\subsubsection{Cuffless measurement: convergence of the sources}\label{sec:cuffless}

\textcite{alessa2019} note that three applications out of 186, that is 1.6\%, claim to turn the smartphone itself into a blood pressure measuring device, presumably by means of a cuffless technique consisting in placing a finger on the camera. The authors observe that, despite this claim, none of these applications reports any element of reliability or validity.

They develop this point in their discussion in terms that match exactly the WHO position reported in \cref{sec:metrology}: cuffless devices based on smartphone applications \bq{have recently been shown to be highly inaccurate and unfeasible, and may negatively affect patients' health and safety}. They add, citing an earlier study, that users have a strong propensity to download these applications and to rate them favourably, which is precisely the mechanism by which a dangerous tool spreads. They conclude that physicians should be aware of the use their patients make of these applications and should promote only validated devices.

Three independent sources---the WHO, \textcite{alessa2019}, and the technical reservations that \textcite{navarro2023} herself formulates on oscillometry and applanation tonometry---therefore converge on this point. The restriction to validated upper-arm monitors is established and is not reopened in the present work.

\subsubsection{Absence of professional endorsement and of evaluation}

Two further findings of \textcite{alessa2019} complete this picture of the market. Six applications out of 186, that is 3.2\%, claim a contribution from health professionals or medical organisations to their development. None reports any involvement of end users, that is, of hypertensive patients. None appears to have been the subject of a scientific evaluation, and none of the thirty examined in depth is approved by the Food and Drug Administration or carries the CE marking as a medical device.

Eighteen of the thirty, that is 60\%, have a user rating, and thirteen of those, that is 72.2\%, reach four stars or more. The authors note that the user rating is a questionable indicator of usability, an earlier study of diabetes management applications having shown a poor correlation between rating and actual usability.

\subsubsection{Data protection}

Of the thirty applications with comprehensive functionality, only twenty, that is 66.7\%, have a privacy policy in English; of the other ten, only one provides a link to such a policy, and that link is non-functional (\cref{tab:privacy}).

\begin{table}[htbp]
\centering\small
\caption{Content of the privacy policies of the twenty applications that have one \parencite{alessa2019}, Table~3.}\label{tab:privacy}
\begin{tabularx}{\textwidth}{@{}X c@{}}
\toprule
\textbf{Question} & \textbf{Applications ($n=20$)} \\
\midrule
Collects personally identifiable information & 16 (80\%) \\
Shares data with third parties & 14 (70\%) \\
States the data-security measures applied & 9 (45\%) \\
Policy available in another language & 1 (5\%) \\
\bottomrule
\end{tabularx}
\end{table}

The authors note that three applications state that they do not share data \bq{except in exceptional circumstances}, a vague formulation that they interpret as undeclared sharing. They situate this finding in continuity with an earlier study which established that one fifth of the applications in the library of the British National Health Service had no privacy policy and that the majority breached data-protection rules.

This point is not incidental for the present work. It joins the legal and regulatory barrier identified by \textcite{dodoo2022} in sub-Saharan Africa (\cref{sec:digital-barriers}), and it requires that the questions of consent, of data localisation and of data protection be treated explicitly in the design chapter, and not deferred to a statement of principle.

%==========================================================================
\subsection{Behaviour change techniques}\label{sec:bct}

\subsubsection{Definition and taxonomy}

A behaviour change technique, or BCT, is defined by \textcite{zhou2024} as an observable, replicable and irreducible intervention component, designed to alter or redirect the causal processes that regulate behaviour. The value of such a decomposition is that it allows the active ingredients of an intervention to be identified, results to be synthesised across studies, and what works to be reproduced.

The reference taxonomy is that of Michie and colleagues, version~1 \parencite{michie2013}, which organises 93 techniques into 16 categories. Both sources mobilised here employ it, which makes their results directly comparable.

\textcite{alessa2019} complement this coding with a second instrument, the Theoretical Domains Framework \parencite{cane2012}, which identifies the psychological mechanisms of action by which the techniques operate.

\subsubsection{What commercial applications contain}

\textcite{alessa2019} coded the thirty applications with comprehensive functionality and identified 24 distinct techniques in them, with inter-rater agreement measured by a Cohen's kappa of 0.85.

Three techniques are present in all thirty applications (\cref{tab:bct-apps}).

\begin{table}[htbp]
\centering\small
\caption{The three behaviour change techniques present in all thirty applications analysed in depth \parencite{alessa2019}.}\label{tab:bct-apps}
\begin{tabularx}{\textwidth}{@{}>{\raggedright\arraybackslash}p{4.0cm} X@{}}
\toprule
\textbf{Technique} & \textbf{Most common implementation in the applications} \\
\midrule
Self-monitoring of behaviour & frequent recording by the user of his or her own health behaviour in a log: pressure, weight, mood, medication intake \\
Prompts and cues & triggering of an alarm at the moment of performing a task, in order to prompt the behaviour \\
Action planning & scheduling of a reminder to perform a task---take a medicine, measure pressure, exercise---at a specified time \\
\bottomrule
\end{tabularx}
\end{table}

Next come feedback on behaviour, present in 25 applications, and monitoring of behaviour by others without feedback, present in 24. At the other extreme, social comparison and demonstration of the behaviour each appear in only one application.

An arithmetic inconsistency must be signalled in this article. The authors write that the total number of techniques per application \bq{ranges between 6 and 17, with a mean of 18.4 (standard deviation 2.6) and a median of 9}. A mean of 18.4 is arithmetically impossible when the maximum observed is 17. The median of 9 suggests an actual mean in the neighbourhood of 9 or 10. The value of 18.4 is reproduced as such in the abstract of the article. We do not retain it and we recommend citing the median of 9, the only value consistent with the stated range.

On the side of psychological mechanisms, ten domains are mobilised, from five to nine per application, with a median of six. Five domains are present in all thirty applications: behavioural regulation, knowledge, goals, memory, attention and decision processes, and beliefs about consequences. Four domains are absent from all applications: intentions, optimism, professional role and identity, and environmental context and resources. The authors note moreover that none of the thirty applications claims an explicit theoretical foundation, whereas self-management programmes grounded in a theory are more likely to be effective.

\subsubsection{What evaluated interventions contain}

\textcite{zhou2024} coded the twenty randomised trials included in their review and identified 21 distinct techniques in them, belonging to 12 categories, with a mean of 8.7 techniques per intervention (standard deviation 3.8) and a range of 2 to 17, that is between 2\% and 18\% of the 93 techniques of the taxonomy (\cref{tab:bct-trials}).

\begin{table}[htbp]
\centering\small
\caption{Most frequently employed behaviour change techniques in the twenty trials \parencite{zhou2024}.}\label{tab:bct-trials}
\begin{tabularx}{\textwidth}{@{}X c@{}}
\toprule
\textbf{Technique} & \textbf{Frequency} \\
\midrule
Self-monitoring of outcomes of behaviour & 17/20 (85\%) \\
Feedback on outcomes of behaviour & 15/20 (75\%) \\
Instruction on how to perform the behaviour & 15/20 (75\%) \\
Pharmacological support & 15/20 (75\%) \\
Biofeedback & 14/20 (70\%) \\
Prompts and cues & 13/20 (65\%) \\
Credible source & 13/20 (65\%) \\
Action planning & 11/20 (55\%) \\
\bottomrule
\end{tabularx}
\end{table}

The groups of techniques employed in at least ten studies are those of goals and planning, feedback and monitoring, shaping knowledge, associations, comparison of outcomes, and self-regulation. Four techniques appear in only two studies: problem solving, information about health consequences, behaviour substitution, and conserving mental resources.

Only six studies out of twenty, that is 30\%, declare a theoretical foundation, mobilising seven theory instances: self-determination, planned behaviour, self-regulation, self-efficacy, cognitive behavioural therapy and the behaviour change wheel. One study combines two, the other five rest on a single one.%
\verify{Six theories are named for seven declared instances; one name is missing from the French draft. Complete the list from \textcite{zhou2024}.}
The authors describe this absence of theoretical foundation as a \bq{common phenomenon in current behaviour change interventions} calling for urgent attention.

\subsubsection{Comparison of the two inventories}

The confrontation of the two sources is instructive. The commercial applications of \textcite{alessa2019} contain a median of 9 techniques; the evaluated interventions of \textcite{zhou2024} contain 8.7 on average. The two sets are therefore comparable in technical density. What separates them is not the behavioural content but everything else: professional involvement, scientific evaluation, clinical supervision and integration into a care pathway. A commercial tool contains about as many behavioural levers as a tested intervention; what it lacks is the arrangement that turns those levers into a clinical result.

This observation shifts the design question. What is at stake is not stacking up behaviour change techniques---the market already does that---but building the rest.

%==========================================================================
\subsection{Measured efficacy}\label{sec:efficacy}

\subsubsection{The meta-analysis of Zhou et al.}

Twenty studies were included in the systematic review and sixteen in the meta-analysis, originating from nine countries: China and the United States for six studies each, Iran for two, then the United Kingdom, Spain, Germany, Japan, Palestine and Jordan for one each. All were published after 2017, mainly between 2020 and 2023. Sample sizes per study range from 30 to 636 participants and intervention durations from 1.5 to 12 months. In total 4\,168 participants were enrolled, with a retention rate of 85.2\%, that is 3\,551 of 4\,168.

The pooled results and their appraisal by the GRADE method are given in \cref{tab:zhou-results}.
\begin{table}[htbp]
\centering\footnotesize
\caption{Pooled effects of mobile self-management interventions on blood pressure \parencite{zhou2024}.}\label{tab:zhou-results}
\begin{tabularx}{\textwidth}{@{}X c c >{\raggedright\arraybackslash}p{2.9cm} c c@{}}
\toprule
\textbf{Outcome} & \textbf{Studies} & \textbf{Interv./control} & \textbf{Mean difference (95\% CI)} & \textbf{$I^2$} & \textbf{GRADE} \\
\midrule
Systolic pressure & 16 & 1\,568 / 1\,516 & $-5.78$~mmHg ($-7.97$ to $-3.59$) & 82\% & Low \\
Diastolic pressure & 15 & 1\,359 / 1\,316 & $-3.28$~mmHg ($-4.39$ to $-2.17$) & 58\% & Moderate \\
\bottomrule
\end{tabularx}
\end{table}

Egger's test indicates no significant publication bias ($P = 0.16$).

The certainty of the evidence must be stressed. Fifteen studies out of twenty, that is 75\%, are judged at high risk of bias, one presents some concerns and only four are at low risk. The risk bears mainly on the randomisation process---ten studies out of twenty, that is 50\%, poorly rated for failure to perform or to describe randomisation and allocation concealment---and on outcome measurement---eight studies out of twenty, that is 40\%, at high risk, mainly because of subjective self-reported outcomes and the absence of blinding of assessors. The authors further mention high dropout rates, the absence of intention-to-treat analyses, and selective reporting of predefined outcomes.

Consequently, the reduction of 5.78~mmHg in systolic pressure rests on low-certainty evidence. This qualifier must accompany any citation of that figure.

\subsubsection{The determinants of the effect}

Four subgroup analyses were conducted on systolic pressure. They constitute the most useful contribution of the article for the design of a device (\cref{tab:subgroups,fig:subgroups}).

\begin{table}[htbp]
\centering\small
\caption{Subgroup analyses on systolic pressure \parencite{zhou2024}. BCT: behaviour change technique.}\label{tab:subgroups}
\begin{tabularx}{\textwidth}{@{}>{\raggedright\arraybackslash}p{3.3cm} >{\raggedright\arraybackslash}p{2.5cm} X c@{}}
\toprule
\textbf{Factor} & \textbf{Modality} & \textbf{Effect on SBP (95\% CI)} & \textbf{$P$} \\
\midrule
Risk factors addressed & multiple & $-6.50$~mmHg ($-9.00$ to $-3.99$) & \multirow{2}{*}{.007} \\
                       & single   & $-1.54$~mmHg ($-4.15$ to $1.06$) & \\
\addlinespace
Number of techniques & $\geq$ 11 BCT & $-9.68$~mmHg ($-13.49$ to $-5.87$) & \multirow{2}{*}{$<$.001} \\
                     & $<$ 11 BCT    & $-2.88$~mmHg ($-3.90$ to $-1.86$) & \\
\addlinespace
Theoretical foundation & with theory & $-10.06$~mmHg ($-16.42$ to $-3.70$) & \multirow{2}{*}{.07} \\
                       & without     & $-4.13$~mmHg ($-5.50$ to $-2.75$) & \\
\addlinespace
Duration & 3 months & $-8.87$~mmHg ($-10.90$ to $-6.83$) & \multirow{2}{*}{.09} \\
         & $\geq$ 6 months & $-5.76$~mmHg ($-8.74$ to $-2.77$) & \\
\bottomrule
\end{tabularx}
\end{table}

\begin{figure}[htbp]
\centering
\begin{tikzpicture}
\begin{axis}[ybar,width=\textwidth,height=6.2cm,bar width=13pt,ymin=0,ymax=14,
  ylabel={Magnitude of SBP reduction (mmHg)},
  symbolic x coords={Risk factors,Techniques,Theory,Duration},
  xtick=data,xticklabel style={font=\footnotesize},
  enlarge x limits=0.22,legend pos=north west,legend style={font=\scriptsize},
  nodes near coords,every node near coord/.append style={font=\scriptsize}]
\addplot[fill=cAfrica!80,draw=cAfrica] coordinates
  {(Risk factors,6.50)(Techniques,9.68)(Theory,10.06)(Duration,8.87)};
\addplot[fill=cGrey!60,draw=cGrey] coordinates
  {(Risk factors,1.54)(Techniques,2.88)(Theory,4.13)(Duration,5.76)};
\legend{Multiple / $\geq$11 BCT / with theory / 3 months,
        Single / $<$11 BCT / without theory / $\geq$6 months}
\end{axis}
\end{tikzpicture}
\caption{Subgroup analyses of \textcite{zhou2024}, plotted as magnitudes of reduction (all effects are decreases). Confidence intervals are given in \cref{tab:subgroups}; residual heterogeneity remains high in every subgroup, so these comparisons are exploratory.}\label{fig:subgroups}
\end{figure}

Two of these comparisons reach statistical significance and two approach it without reaching it.

The number of techniques produces the sharpest gap: interventions mobilising at least eleven techniques obtain a reduction more than three times greater than those mobilising fewer. The breadth of risk factors addressed produces a gap of the same order, single-factor interventions failing to reach significance.

The result concerning duration is counter-intuitive and deserves comment, since it contradicts the intuition that longer follow-up would produce a more durable effect. Interventions of three months are more effective than those of six months or more. The authors explain this by habit-formation theory: habit strength reaches a peak of automaticity around the twelfth week and then weakens progressively. They infer that maintaining motivation over time constitutes an unresolved design problem.

One limit to this interpretation must however be signalled. Heterogeneity remains considerable in every subgroup---86\% for interventions longer than six months, 92\% for those grounded in a theory---which the authors explicitly acknowledge. These analyses are therefore exploratory and do not permit the gap to be attributed causally to the factor examined.

\subsubsection{The trial of Mihevc et al.}\label{sec:mihevc}

This trial is the most directly relevant to the present work, since it concerns the very population targeted by the proposed system.

\paragraph{Design.} A multicentre, prospective, open-label, parallel-group trial, randomised 1:1 by an independent third party, conducted in three Slovenian primary-care centres between March 2021 and June 2023 within the SCUBY project. One urban centre, Ljubljana, serving about 300\,000 inhabitants with a gross domestic product of 108\% of the European average; two rural centres, Trebnje and Slovenj Gradec, serving about 50\,000 inhabitants with a gross domestic product of 74\% of the European average.

\paragraph{Population.} Adults aged 65 and over with both confirmed hypertension---seven-day home mean $\geq 135/85$~mmHg---and confirmed type~2 diabetes, both present for at least one year, able to use the equipment. Exclusions: diabetes requiring insulin, gestational or type~1 diabetes, dementia, inability to use the equipment.

\paragraph{Intervention.} Telemonitoring in addition to usual care, and not in substitution for it. Each patient received one hour of training delivered by a nurse on the measurement protocol and the use of the equipment. A telemedicine centre coordinated by a general practitioner reviewed the transmitted values and served as the link between patients and their treating physicians.

\paragraph{Flow of participants.} One hundred and sixty patients assessed, 32 excluded---6 for unmet criteria, 26 for refusal to participate---and 128 randomised into two groups of 64. In the telemonitoring group, seven patients discontinued the intervention within the first month through inability to use the equipment correctly, one at three months for psychological burden, and one at nine months after the death of a spouse, leaving 55 patients analysed. In the control group, one death at eight months and one emigration at twelve months left 62 patients analysed. Twelve-month follow-up was completed by 117 participants of 128, that is 91.4\%.

\paragraph{Population analysed.} Mean age 71.3 years (standard deviation 4.7; range 65--88), 60.7\% men, 53.8\% rural residents. Mean duration of diabetes 9.5 years, of hypertension 14.5 years. Baseline systolic pressure 136.7~mmHg (14.1), baseline glycated haemoglobin 7.2\% (1.0), body mass index 30.4, obesity in 47.9\%, chronic kidney disease in 20.5\%. No significant difference between groups at inclusion.

\paragraph{Main results.} They are given in \cref{tab:mihevc-results} and plotted in \cref{fig:mihevc}.

\begin{table}[htbp]
\centering\footnotesize
\caption{Main results of \textcite{mihevc2025}. Means with standard deviations in parentheses.}\label{tab:mihevc-results}
\begin{tabularx}{\textwidth}{@{}>{\raggedright\arraybackslash}p{2.1cm} >{\raggedright\arraybackslash}p{2.1cm} c c c >{\raggedright\arraybackslash}p{2.6cm} c@{}}
\toprule
\textbf{Outcome} & \textbf{Group} & \textbf{Baseline} & \textbf{6 months} & \textbf{12 months} & \textbf{Between-group difference at 12 months} & \textbf{$P$} \\
\midrule
Systolic pressure & Telemonitoring ($n=55$) & 135.8 (14.9) & 124.7 (11.2) & 126.2 (10.4) & $-6.9$~mmHg ($-11$ to $-2.7$) & .001 \\
                  & Control ($n=62$) & 137.7 (13.1) & 137.1 (9.4) & 134.9 (9.6) & & \\
\addlinespace
Glycated haemoglobin & Telemonitoring & 7.2\% (1.2) & --- & 6.7\% (0.8) & $-0.5\%$ ($-0.8$ to $-0.2$) & .002 \\
                     & Control & 7.1\% (0.8) & --- & 6.7\% (0.8) & & \\
\bottomrule
\end{tabularx}
\end{table}

\begin{figure}[htbp]
\centering
\begin{tikzpicture}
\begin{axis}[width=0.85\textwidth,height=6.4cm,
  xlabel={Months since inclusion},ylabel={Systolic pressure (mmHg)},
  xmin=-0.6,xmax=12.6,ymin=120,ymax=142,xtick={0,6,12},
  grid=major,grid style={gray!20},legend pos=south west,legend style={font=\footnotesize}]
\addplot[thick,cAfrica,mark=*] coordinates {(0,135.8)(6,124.7)(12,126.2)};
\addplot[thick,cRef,mark=square*] coordinates {(0,137.7)(6,137.1)(12,134.9)};
\addplot[dashed,cGrey,no markers,domain=-0.6:12.6]{135};
\legend{Telemonitoring ($n=55$),Usual care ($n=62$),Home target (135~mmHg)}
\node[font=\scriptsize,anchor=west,text=cAfrica] at (axis cs:6.2,127.5) {$-10.5$ at 6 months};
\node[font=\scriptsize,anchor=east,text=cAfrica] at (axis cs:11.8,130.5) {$-6.9$ at 12 months};
\end{axis}
\end{tikzpicture}
\caption{Course of systolic pressure in \textcite{mihevc2025}. The annotated between-group differences are differences in change from baseline, not differences of the plotted means; about a third of the effect is lost between six and twelve months.}\label{fig:mihevc}
\end{figure}

One point of reading must be made explicit, because it is not stated in the source and a reader comparing the figures would otherwise find them inconsistent. The reported between-group differences are not the differences of the reported means, which are $124.7 - 137.1 = -12.4$~mmHg at six months and $126.2 - 134.9 = -8.7$~mmHg at twelve months. They are differences in \emph{change from baseline}, and on that reading they reproduce the published values exactly:
\begin{align}
  \text{6 months:} \quad & (124.7-135.8) - (137.1-137.7) = -11.1 + 0.6 = -10.5~\text{mmHg},\\
  \text{12 months:} \quad & (126.2-135.8) - (134.9-137.7) = -9.6 + 2.8 = -6.8 \approx -6.9~\text{mmHg}.
\end{align}
This derivation is ours; it confirms that the effect reported is adjusted for the baseline imbalance of 1.9~mmHg in favour of the intervention arm, and it is the reading we retain.

The same reading, however, does not reconcile the glycated haemoglobin figures. The telemonitoring arm falls by 0.5 points (7.2 to 6.7) and the control arm by 0.4 points (7.1 to 6.7), which gives a difference in change of $-0.1$ point, not the $-0.5$ point reported with a confidence interval excluding zero.%
\verify{One of the four HbA1c values is inconsistent with the reported between-group difference of $-0.5\%$ (95\% CI $-0.8$ to $-0.2$; $P=.002$). The most likely error is the control-arm endpoint, since the authors later compare their $-0.5\%$ favourably with the literature. Check the four values and the model used against Table~3 of \textcite{mihevc2025} before citing them.}

The mixed-design analysis of variance confirms a significant change over time in systolic pressure in both groups ($F_{2,114} = 26.93$; $P < .001$; $\eta^2 = 0.19$) and, above all, a significant time $\times$ group interaction ($F_{2,114} = 15.08$; $P < .01$; $\eta^2 = 0.12$), that is, the two groups did not evolve in the same way.

\paragraph{Secondary results.} None survives adjustment for multiple comparisons by Holm's method. Fasting glucose showed a between-group difference of $-0.8$~mmol/L ($-1.4$ to $-0.2$; $P = .006$) becoming non-significant after adjustment ($P = .06$). Diastolic pressure showed a difference of $-3.4$~mmHg at six months ($P = .007$) falling back to $-1.7$~mmHg at twelve months ($P = .15$, adjusted $P = .81$); the authors note that this result does not reach the threshold of 5~mmHg associated with a 20\% reduction in the risk of major cardiovascular events, and therefore judge it clinically irrelevant. Lipid profile, body mass index, glomerular filtration rate and diabetes-related quality of life: no significant difference.

\paragraph{Clinical significance.} The authors report that a reduction in systolic pressure at six months of the magnitude observed corresponds to a risk of major cardiovascular events lower by about 20\%, and that at twelve months it corresponds to a 10 to 15\% reduction in cardiovascular disease risk. A reduction of 0.5\% in glycated haemoglobin is clinically significant and associated with an estimated 17\% reduction in cardiovascular disease risk.

They finally situate their glycaemic result relative to the literature: previously reported reductions in glycated haemoglobin ranged from $-0.03\%$ to $-0.32\%$ in patients with type~2 diabetes alone and from $+0.14\%$ to $-0.33\%$ in those with both conditions, and a recent meta-analysis places the effect of telemonitoring in primary care on this population at twelve months at $-0.21\%$ ($-0.37$ to $-0.06$). Their result of $-0.5\%$ is therefore above the state of the art.

\subsubsection{Synthesis of the available effect estimates}

\begin{table}[htbp]
\centering\footnotesize
\caption{Available effect estimates for telemonitoring and mobile health on blood pressure and glycaemic control.}\label{tab:effect-synthesis}
\begin{tabularx}{\textwidth}{@{}>{\raggedright\arraybackslash}p{2.3cm} >{\raggedright\arraybackslash}p{2.6cm} >{\raggedright\arraybackslash}p{2.4cm} >{\raggedright\arraybackslash}X >{\raggedright\arraybackslash}p{2.4cm}@{}}
\toprule
\textbf{Source} & \textbf{Design} & \textbf{Population} & \textbf{Effect on SBP} & \textbf{Glycaemic effect} \\
\midrule
\textcite{zhou2024} & meta-analysis, 16 RCTs, 3\,084 patients & hypertensive, 9 countries & $-5.78$~mmHg ($-7.97$ to $-3.59$), low certainty & not assessed \\
\textcite{mihevc2025} & multicentre RCT, 117 patients & $\geq$ 65 y, HTN + T2D, Slovenia & $-6.9$~mmHg ($-11$ to $-2.7$) & HbA1c $-0.5\%$ ($-0.8$ to $-0.2$) \\
\textcite{aovare2023}, cf.\ \cref{sec:hypothesis} & meta-analysis, 7 RCTs, 2\,249 patients & Africa & $-1.39$~mmHg ($-4.46$ to $1.68$), not significant & HbA1c $+0.20\%$, not significant \\
Meta-analysis cited by \textcite{mihevc2025} & primary care & HTN + T2D & --- & HbA1c $-0.21\%$ ($-0.37$ to $-0.06$) \\
\bottomrule
\end{tabularx}
\end{table}

\begin{figure}[htbp]
\centering
\begin{tikzpicture}
\begin{axis}[width=0.85\textwidth,height=5.2cm,
  xlabel={Effect on systolic pressure (mmHg)},
  xmin=-13,xmax=4,ymin=0.5,ymax=3.5,ytick={1,2,3},
  yticklabels={{\footnotesize Aovare 2023 (Africa)},{\footnotesize Mihevc 2025 (Slovenia)},{\footnotesize Zhou 2024 (9 countries)}},
  grid=major,grid style={gray!20}]
\addplot[dashed,cGrey,no markers] coordinates {(0,0.5)(0,3.5)};
\addplot[only marks,mark=*,mark size=2.5pt,cAfrica,error bars/.cd,x dir=both,x explicit,error bar style={thick}]
  coordinates {(-1.39,1) += (3.07,0) -= (3.07,0)};
\addplot[only marks,mark=square*,mark size=2.5pt,cRef,error bars/.cd,x dir=both,x explicit,error bar style={thick}]
  coordinates {(-6.90,2) += (4.20,0) -= (4.10,0)
               (-5.78,3) += (2.19,0) -= (2.19,0)};
\end{axis}
\end{tikzpicture}
\caption{The three pooled or trial estimates of \cref{tab:effect-synthesis} on systolic pressure. The dashed line marks the null effect; the only estimate obtained in Africa is the only one whose interval includes it.}\label{fig:effect-synthesis}
\end{figure}

This table constitutes the most important result of the section, and it must be read as a whole (\cref{fig:effect-synthesis}). Telemonitoring devices produce, in high- and middle-income countries, a modest but real and clinically relevant effect on systolic pressure and glycaemic control. The only body of data bearing specifically on Africa finds no significant effect.

This gap is not interpretable as it stands. Three explanations remain open. The interventions evaluated in Africa may have been of poorer design---\textcite{stokes2022} described them as consisting essentially of text-message reminders and therapeutic education, which corresponds, in Zhou's grid, to single-factor interventions poor in techniques, that is, precisely to the subgroup without significant effect. The conditions of implementation may differ---medicine availability, connectivity, provider workload. Or the effect may be genuinely absent in this context. The present thesis cannot settle this question, but it must state it, because it constitutes the central uncertainty of its object.

It is worth noting that \textcite{zhou2024} cite \textcite{aovare2023} among their references, without however discussing this geographical discordance.

%==========================================================================
\subsection{Is the mechanism of action behaviour change?}\label{sec:mechanism}

This subsection does not appear in the initial plan and must appear in it, because it calls into question one of that plan's implicit assumptions.

\subsubsection{The result of Mihevc et al.}

The authors measured, in an exploratory analysis, the course of six behavioural risk factors over twelve months: regular physical activity, harmful alcohol consumption habits, adding salt at the table, number of daily meals, daily vegetable consumption and smoking.

None of these six factors shows a significant change, either in the telemonitoring group or in the control group. Regular physical activity even decreases by 10.8\% in the intervention group and by 13.5\% in the control group. Smoking increases in both groups.

The authors draw from this the following conclusion, which must be cited precisely: \bq{telemonitoring alone did not significantly improve diabetes-related quality of life or lead to positive lifestyle changes. This finding is important, because patients achieved significant control of blood pressure and blood glucose without making substantial lifestyle changes. The positive results may be attributable to other factors such as improved communication, medication adherence, early detection of exacerbations, and treatment adjustments based on accurate telemonitoring data.}

\subsubsection{Scope and limits of this result}

This finding is at once the most important and the most fragile of the section, and it is appropriate to be precise about both aspects.

Its fragility derives first from the fact that lifestyle changes were self-reported, as the authors acknowledge, which is an insensitive measure. It derives next from the fact that their intervention included no structured therapeutic education component---the article mentions this explicitly---so that the absence of behavioural change could reflect the absence of a behavioural intervention rather than its ineffectiveness. It derives finally from the fact that an exploratory analysis on six outcomes without appropriate correction does not permit the conclusion that there is no effect, only that no effect was detected.

Its scope, by contrast, remains considerable, because the inference it authorises is logically robust independently of these limits: a significant clinical effect was obtained, and it is not explicable by behaviour change, which did not occur. Whatever the reason for that absence, the clinical result must have another cause.

\subsubsection{Convergence with the diagnosis of Section 1.2}

The four alternative causes that the authors propose---improved communication, medication adherence, early detection of exacerbations, treatment adjustments based on accurate data---match exactly the diagnosis established in \cref{sec:diagnosis}.

We showed there that the largest conditional loss of the cascade lies at the last link, 56\% of treated patients not reaching the target, and that this loss is a matter of therapeutic inertia, defined by the WHO as the tendency not to act to intensify treatment despite pressure remaining uncontrolled (\cref{sec:inertia}). We showed that three conditions are necessary for intensification: that the patient returns, that the provider has the history establishing persistence, and that the protocol indicates the next step. We finally reported the testimony of a practitioner according to whom, pressure being measured only at isolated moments, it becomes difficult to verify the efficacy of a given treatment (\cref{sec:testimony}).

The result of \textcite{mihevc2025} is the empirical confirmation of that analysis: telemonitoring acts by removing the informational obstacle to therapeutic adjustment, and not by modifying the patient's behaviour.

\subsubsection{Consequences for the design of the system}

If this interpretation is correct, it substantially reorients the design.

\textbf{The priority function of the system is not to motivate the patient but to inform the provider.} Behaviour change techniques remain useful---\textcite{zhou2024} establish that a high number of techniques is associated with a greater effect---but they do not constitute the principal mechanism. A device that invested the bulk of its complexity in gamification, motivational messages or lifestyle counselling would be targeting a secondary lever, all the more so since the WHO establishes, as we reported in \cref{sec:inertia}, that the effect of lifestyle modification is of the order of 1.3~mmHg systolic, that is about one fifth of the effect observed in the trial of \textcite{mihevc2025}.

\textbf{The output of the system must be a therapeutic decision aid.} What must be produced is not a patient dashboard but information allowing the provider to establish, with sufficient reliability, that the target is persistently not being reached. That is the very definition of a detection problem on a time series, and it is the object of the learning model that this thesis proposes to build.

\textbf{Behaviour change techniques must nevertheless be present and documented.} \Textcite{zhou2024} establish a gap of 2.88 to 9.68~mmHg between interventions poor and rich in techniques. The three techniques present in all applications and in the majority of trials---self-monitoring of behaviour, prompts and cues, action planning---constitute the minimal base, and they are inexpensive to implement since they reduce to recording, reminding and scheduling.

%==========================================================================
\subsection{Conditions of feasibility}\label{sec:feasibility}

\subsubsection{The usage barrier}

The most instructive result of the trial of \textcite{mihevc2025} for a deployment in a resource-limited context is also the least commented upon. Seven patients of the 64 in the telemonitoring group, that is 10.9\% of the randomised arm, abandoned within the first month through inability to use the equipment correctly, and this despite one hour of individual training delivered by a nurse. These seven patients are seven of the nine withdrawals recorded in that arm.%
\verify{The French draft also described these seven patients as ``12.7\% of study withdrawals''. That figure corresponds to $7/55$, the analysed sample, not to the withdrawals, of which there were nine in that arm and eleven in the trial. Use 10.9\% of the randomised arm and drop the 12.7\%, or recover the denominator the source intends.}

This barrier is the main cause of attrition in the trial, far ahead of medical causes. It occurs in an elderly population---mean age 71.3 years---but in a high-income country, with equipment supplied, individual training and ongoing support. It would be imprudent to anticipate a lower rate in a context of lower digital literacy.

Three consequences follow for the present work. Training must be treated as a component of the device and not as a preliminary formality; one hour is manifestly insufficient for a fraction of the population. The non-adoption rate must be anticipated in the sizing of any evaluation cohort. And patient selection must be treated explicitly---the authors moreover conclude that \bq{complex interventions may pose difficulties for older people, which suggests the need for careful patient selection in order to ensure that the benefits outweigh the potential burdens}.

\subsubsection{Retention}

The two sources converge on high retention rates: 85.2\% in the twenty trials pooled by \textcite{zhou2024}, and 91.4\% at twelve months in the trial of \textcite{mihevc2025}. These values contrast strongly with the 50\% or more loss to follow-up that the WHO reports in routine care in resource-limited settings (\cref{sec:ltfu}).

This comparison must however be handled with caution, because it sets non-comparable quantities against each other. Trial participants are volunteers, selected, followed under a protocol and solicited by a research team. The retention rate of a trial does not estimate the retention rate of a routine programme. It establishes only that abandonment is not an insurmountable obstacle when the device is properly supported, which is useful but limited information.

\subsubsection{Duration and erosion of the effect}

Two convergent observations signal a phenomenon of erosion. In \textcite{zhou2024}, interventions of three months produce a greater effect than those of six months or more, a phenomenon explained by the peak of habit automaticity around the twelfth week. In \textcite{mihevc2025}, the between-group difference in systolic pressure falls from $-10.5$~mmHg at six months to $-6.9$~mmHg at twelve months, a loss of about one third of the effect (\cref{fig:mihevc}).

These two observations, arising from different devices, suggest that maintaining the effect over time constitutes an unresolved design problem. For a system intended for chronic diseases followed over decades, this is a major limitation, and it will have to be discussed explicitly in the evaluation of the proposed device rather than masked by a short observation horizon.

\subsubsection{The context of technological access}

\textcite{navarro2023} documents, for Spain, a rapid progression of access: 96.1\% of households had internet access in 2022 against 82.7\% in 2017; 99.9\% had at least one type of telephone; teleconsultations rose from 5.4\% of all consultations in 2019, that is 18.1 million, to 34.6\% in 2020, that is 126.9 million, then 37.9\% in 2021, that is 167 million.

These figures are not transposable to the Cameroonian context and we do not mobilise them as such. The relevant datum is the one reported by \textcite{dodoo2022} for sub-Saharan Africa: 816 million mobile telephone connections in 2019, with a projection of 1.05 billion in 2025, and 41 countries endowed with national digital health strategies and architectures (\cref{sec:digital-barriers}). The access constraint is therefore not the absence of terminals but, as those same authors establish, the conjunction of the technological, organisational, financial, individual, cultural and regulatory barriers examined in \cref{sec:digital-barriers}.

%==========================================================================
\subsection{Scope, limits and specification}\label{sec:spec13}

\subsubsection{Limits of this section}

\textbf{Geographical generalisation.} The twenty trials of \textcite{zhou2024} come from nine countries, none of them in sub-Saharan Africa. The trial of \textcite{mihevc2025} was conducted in Slovenia. The survey of \textcite{alessa2019} covers the English-language British market. None of the sources in this section documents self-measurement or mobile data entry in an African context, and the only source available on that context, reported in \cref{sec:hypothesis}, finds no significant effect.

\textbf{Certainty of the evidence.} The estimate of $-5.78$~mmHg carries low GRADE certainty, 75\% of the included studies being at high risk of bias. The estimate of $-3.28$~mmHg benefits from moderate certainty.

\textbf{Scope of the reference trial.} The trial of \textcite{mihevc2025} concerns 117 analysed patients, exclusively aged 65 and over, exclusively not on insulin, and it is open-label, that is, unblinded. Its power was calculated to detect a difference of at least 6.2~mmHg, and the observed effect of 6.9~mmHg lies just above that threshold.

\textbf{Coding of techniques.} Both sources acknowledge that the coding of behaviour change techniques depends on intervention descriptions that are often under-reported or coarse, and remains subject to coder judgement despite good agreement indices. \Textcite{zhou2024} add that they were unable to examine quantitatively the effect of an isolated technique or of a particular combination.

\textbf{Inclusion bias of Zhou et al.} Their criteria excluded interventions based principally on telephone calls, text messages and web programmes, as well as applications offering only notification, self-monitoring and advice. The authors acknowledge that these excluded interventions also contained behaviour change techniques, which may have biased the assessment. This exclusion is particularly significant for our object, since the devices evaluated in Africa are precisely of that type---which means that the contrast of \cref{tab:effect-synthesis} is partly an artefact of the inclusion criteria, the two bodies of evidence not having selected the same kind of intervention.

\subsubsection{Specification arising from this section}

\begin{table}[htbp]
\centering\footnotesize
\caption{Design requirements established in \cref{sec:outpatient} and their basis.}\label{tab:spec13}
\begin{tabularx}{\textwidth}{@{}X >{\raggedright\arraybackslash}p{5.6cm}@{}}
\toprule
\textbf{Requirement} & \textbf{Basis} \\
\midrule
Validated automated upper-arm monitor, never cuffless measurement & \textcite{who2023hypertension}; \textcite{alessa2019} \\
Automatic transmission, no manual entry & \textcite{alessa2019}: 73 of 84 applications require entry; \textcite{mihevc2025}: Bluetooth to a concentrator \\
Measurement protocol imposed by the device: 2 readings, 1-minute interval, second retained & \textcite{mihevc2025} \\
Reference value = seven-day mean of the second morning and evening readings & \textcite{mihevc2025} \\
Three-level classification, bidirectional (high and low thresholds) & \textcite{mihevc2025} \\
Two-stage architecture: anomaly $\rightarrow$ verification profile $\rightarrow$ clinical decision & \textcite{mihevc2025} \\
Most severe of the systolic and diastolic classifications prevails & our decision; the source is silent (\cref{sec:triage}) \\
Low routine frequency, triggered intensification & \textcite{mihevc2025} \\
At least 11 behaviour change techniques, including the 3 universal ones & \textcite{zhou2024}: $-9.68$ vs $-2.88$~mmHg \\
Multifactorial, not single-factor intervention & \textcite{zhou2024}: $-6.50$ vs $-1.54$~mmHg \\
Explicit and declared theoretical foundation & \textcite{zhou2024}: only 30\% of trials declare one \\
Priority output = therapeutic decision aid for the provider & \textcite{mihevc2025}: effect obtained without behavioural change \\
Privacy policy, data localisation and protection made explicit & \textcite{alessa2019}: 70\% of applications share with third parties \\
Training treated as a component of the device, not a precondition & \textcite{mihevc2025}: 10.9\% abandonment through inability to use \\
Assessment of effect maintenance beyond 12 months & \textcite{zhou2024} and \textcite{mihevc2025}: documented erosion \\
\bottomrule
\end{tabularx}
\end{table}

\subsubsection{Position of the problem for the following section}

This section establishes that the technical modalities of out-of-hospital follow-up are documented, that an operational joint pressure--glucose protocol exists and has been evaluated, and that it produces a significant clinical effect in a population comparable to the target population.

It also establishes three things that the proposed system will have to address and that none of the sources resolves.

\textbf{First}, the decision architecture of \textcite{mihevc2025} rests on fixed thresholds applied uniformly to all patients. Yet \cref{sec:burden} showed that prevalence and risk vary strongly with age, and \cref{sec:variability} that the usual pressure $\mu_i$ and the variability $\sigma_{\varepsilon,i}$ are specific to each individual. A fixed threshold ignores that information. This is precisely the space in which a learning model brings value: personalisation of the triggering threshold from the patient's own history.

\textbf{Second}, the device evaluated triggers on the current value or on a mean, that is, on a level. It does not trigger on a trend. A patient whose pressure drifts steadily without yet having crossed the threshold triggers nothing, whereas that is precisely the patient in whom early intervention would be most useful---the configuration illustrated in \cref{fig:act-series}.

\textbf{Third}, none of the interventions surveyed, neither by \textcite{zhou2024} nor by \textcite{stokes2022}, mobilises any machine learning method.

It is to these three questions that the following section responds, devoted to the state of the art of detection methods applied to physiological time series.

%##########################################################################
% SECTION 1.4
%##########################################################################
\section{Clinical Effectiveness Evidence and Field Experience with Telemonitoring}\label{sec:evidence}

\subsection*{Purpose and hierarchy of sources}

The preceding sections established the need, characterised the failures of the episodic care model, and specified the technical modalities of out-of-hospital follow-up. This section addresses the question on which the legitimacy of the entire work depends: does telemonitoring produce a measurable clinical effect, and under what conditions?

Six sources are used, two of them already analysed earlier in the chapter.

\paragraph{Mao et al. (2020).} \textcite{mao2020}, published in \emph{BMJ Open Diabetes Research \& Care} 8(1):e001225, is a systematic review with meta-analysis covering 51 randomised controlled trials and 13\,054 participants. It is the broadest available evidence base on this question, and the only one that explicitly stratifies by level of economic development.

\paragraph{Mihevc et al. (2025).} \textcite{mihevc2025}, already analysed in \cref{sec:mihevc}, is the randomised controlled trial conducted on precisely the population this thesis targets.

\paragraph{Bediang et al. (2021).} \textcite{bediang2021}, published in \emph{Studies in Health Technology and Informatics} 281:615--619, is a controlled multicentre study conducted in Cameroon. Its level of evidence is markedly lower than the two preceding sources---it is not randomised---but it is one of only two published evaluations of a cardiovascular telemedicine intervention in Cameroonian primary care available to us.

\paragraph{Kingue et al. (2013).} \textcite{kingue2013}, published in the \emph{Pan African Medical Journal} 15:153, is the TELEMED-CAM study, a prospective interventional study of a hypertension management package including telemanagement in rural Cameroon. Its evidential level is intermediate: allocation was partly random at the centre level but not strictly so, and clustering was not accounted for. It is the second of the two Cameroonian evaluations.

\paragraph{Segun-Omosehin et al. (2025).} \textcite{segun2025}, published in \emph{npj Cardiovascular Health} 2(1):20, is a Perspective article on the digital health technology landscape of sub-Saharan Africa and its cardiovascular implications. It is not a systematic review: it declares no search strategy, no inclusion criteria and no quality appraisal, and a portion of its citations are grey literature, press items and organisational websites. We use it for the contextual data it assembles---connectivity, specialist density, documented programmes---and not as a source of effect estimates.

\paragraph{Dodoo et al. (2022).} \textcite{dodoo2022}, already analysed in \cref{sec:digital-barriers}, documents the state of deployment across sub-Saharan Africa.

%==========================================================================
\subsection{The aggregate clinical effect}\label{sec:aggregate-effect}

\subsubsection{The design of Mao et al.}

The search covered PubMed, ResearchGate, Embase and Cochrane for articles published between January 2008 and June 2019, restricted to randomised controlled trials, human studies and English-language publications. Of 1\,747 articles identified, 51 were retained, totalling 13\,054 participants. Sample sizes per study ranged from 34 to 1\,665 and intervention durations from one month to five years---eleven studies of three months or less, twenty-three of three to six months, seventeen of more than six months.

Quality was assessed using the Jadad scale: 37 of 51 studies (72.55\%) described in detail the generation of the randomisation sequence and the handling of missing data; 44 of 51 (86.27\%) reported the follow-up process and explained reasons for withdrawal. The authors note that blinding was not feasible given the nature of the interventions.

The authors distinguish five mobile health modalities: text messaging, telephone calls, wearable or portable monitoring devices, mobile applications, and telemedicine. A fixed-effects model was used when $I^2 < 50\%$ and a random-effects model above that threshold.

\subsubsection{Primary results}

\begin{table}[htbp]
\centering\small
\caption{Pooled effects reported by \textcite{mao2020}.}\label{tab:mao-results}
\begin{tabularx}{\textwidth}{@{}X c c >{\raggedright\arraybackslash}p{4.0cm} c@{}}
\toprule
\textbf{Outcome} & \textbf{Studies} & \textbf{Participants} & \textbf{Weighted mean difference (95\% CI)} & \textbf{$I^2$} \\
\midrule
Glycated haemoglobin & 40 & 8\,006 & $-0.39\%$ ($-0.50$ to $-0.29$) & 62.7\% \\
Fasting blood glucose & 15 & --- & $-0.52$~mmol/L ($-0.93$ to $-0.12$) & 57.6\% \\
Systolic blood pressure & 30 & 9\,476 & $-2.99$~mmHg ($-4.19$ to $-1.80$) & 67.3\% \\
Diastolic blood pressure & 28 & 8\,506 & $-1.14$~mmHg ($-1.86$ to $-0.42$) & 57.1\% \\
\bottomrule
\end{tabularx}
\end{table}

Begg's and Egger's tests detected no significant publication bias for fasting glucose ($p = 0.16$), systolic pressure ($p = 0.439$) or diastolic pressure ($p = 0.857$). For glycated haemoglobin, Egger's test was significant ($p = 0.036$), but trim-and-fill analysis produced no correction, as no trimming proved necessary.

\subsubsection{Consolidated synthesis of available estimates}\label{sec:consolidated}

\Cref{tab:sbp-synthesis,tab:glyc-synthesis} assemble every effect estimate encountered in this work, extending the partial synthesis of \cref{tab:effect-synthesis} with the subgroups of \textcite{mao2020}. They constitute the central result of this section (\cref{fig:sbp-forest}).

\begin{table}[htbp]
\centering\footnotesize
\caption{Consolidated estimates of the effect on systolic blood pressure.}\label{tab:sbp-synthesis}
\begin{tabularx}{\textwidth}{@{}>{\raggedright\arraybackslash}p{3.2cm} >{\raggedright\arraybackslash}p{3.0cm} >{\raggedright\arraybackslash}X >{\raggedright\arraybackslash}p{3.3cm} c@{}}
\toprule
\textbf{Source} & \textbf{Design} & \textbf{Population} & \textbf{Mean difference (95\% CI)} & \textbf{Sig.} \\
\midrule
\textcite{mihevc2025} & RCT, 117 patients & $\geq$ 65 y, HTN + T2D, Slovenia & $-6.9$~mmHg ($-11$ to $-2.7$) & yes \\
\textcite{zhou2024} & meta-analysis, 16 RCTs & hypertensive, 9 countries & $-5.78$~mmHg ($-7.97$ to $-3.59$) & yes, low certainty \\
\textcite{mao2020} & meta-analysis, 30 RCTs, 9\,476 & diabetes and/or HTN & $-2.99$~mmHg ($-4.19$ to $-1.80$) & yes \\
Mao, developed countries & subgroup & --- & $-5.72$~mmHg ($-7.46$ to $-3.99$) & yes \\
Mao, developing countries & subgroup & --- & $+0.25$~mmHg ($-3.10$ to $3.59$) & no \\
Mao, combined intervention & subgroup & --- & $-6.17$~mmHg ($-8.83$ to $-3.50$) & yes \\
Mao, mHealth alone & subgroup & --- & $-2.16$~mmHg ($-5.07$ to $0.75$) & no \\
\textcite{aovare2023} & meta-analysis, 7 RCTs, 2\,249 & Africa & $-1.39$~mmHg ($-4.46$ to $1.68$) & no \\
\bottomrule
\end{tabularx}
\end{table}

\begin{figure}[htbp]
\centering
\begin{tikzpicture}
\begin{axis}[width=0.86\textwidth,height=8.2cm,
  xlabel={Effect on systolic blood pressure (mmHg)},
  xmin=-13,xmax=5,ymin=0.4,ymax=8.6,ytick={1,2,3,4,5,6,7,8},
  yticklabels={{\scriptsize Aovare 2023, Africa},
               {\scriptsize Mao, mHealth alone},
               {\scriptsize Mao, combined},
               {\scriptsize Mao, developing},
               {\scriptsize Mao, developed},
               {\scriptsize Mao 2020, overall},
               {\scriptsize Zhou 2024},
               {\scriptsize Mihevc 2025}},
  grid=major,grid style={gray!20}]
\addplot[dashed,cGrey,no markers] coordinates {(0,0.4)(0,8.6)};
% significant estimates
\addplot[only marks,mark=square*,mark size=2.4pt,cRef,error bars/.cd,x dir=both,x explicit,error bar style={thick}]
  coordinates {(-6.90,8) += (4.20,0) -= (4.10,0)
               (-5.78,7) += (2.19,0) -= (2.19,0)
               (-2.99,6) += (1.19,0) -= (1.20,0)
               (-5.72,5) += (1.73,0) -= (1.74,0)
               (-6.17,3) += (2.67,0) -= (2.66,0)};
% non-significant estimates
\addplot[only marks,mark=*,mark size=2.4pt,cAfrica,error bars/.cd,x dir=both,x explicit,error bar style={thick}]
  coordinates {(0.25,4) += (3.34,0) -= (3.35,0)
               (-2.16,2) += (2.91,0) -= (2.91,0)
               (-1.39,1) += (3.07,0) -= (3.07,0)};
\end{axis}
\end{tikzpicture}
\caption{Consolidated estimates of \cref{tab:sbp-synthesis}. Squares mark intervals excluding the null, circles those including it. The three non-significant estimates are the three that describe interventions without professional oversight or conducted in resource-limited settings.}\label{fig:sbp-forest}
\end{figure}

\begin{table}[htbp]
\centering\footnotesize
\caption{Consolidated estimates of the effect on glycaemic control.}\label{tab:glyc-synthesis}
\begin{tabularx}{\textwidth}{@{}>{\raggedright\arraybackslash}p{3.6cm} >{\raggedright\arraybackslash}X >{\raggedright\arraybackslash}p{3.4cm} c@{}}
\toprule
\textbf{Source} & \textbf{Outcome} & \textbf{Mean difference (95\% CI)} & \textbf{Sig.} \\
\midrule
\textcite{mihevc2025} & HbA1c & $-0.5\%$ ($-0.8$ to $-0.2$) & yes \\
\textcite{mao2020} & HbA1c, overall & $-0.39\%$ ($-0.50$ to $-0.29$) & yes \\
Mao, developed countries & HbA1c & $-0.35\%$ ($-0.46$ to $-0.24$) & yes \\
Mao, developing countries & HbA1c & $-0.52\%$ ($-0.78$ to $-0.26$) & yes \\
Mao, type 2 diabetes & HbA1c & $-0.40\%$ ($-0.52$ to $-0.28$) & yes \\
Mao, type 1 diabetes & HbA1c & $-0.30\%$ ($-0.47$ to $-0.12$) & yes \\
Meta-analysis cited by \textcite{mihevc2025} & HbA1c, HTN + T2D, primary care & $-0.21\%$ ($-0.37$ to $-0.06$) & yes \\
\textcite{aovare2023} & HbA1c, Africa & $+0.20\%$ ($-0.40$ to $0.80$) & no \\
\bottomrule
\end{tabularx}
\end{table}

\subsubsection{Interpreting the magnitude}

These values appear modest and must be referred to their clinical significance, which \textcite{mihevc2025} make explicit and which we reported in \cref{sec:mihevc}: a systolic reduction of about 10~mmHg corresponds to roughly 20\% lower risk of major cardiovascular events, a 7~mmHg reduction to a 10--15\% reduction in cardiovascular disease risk, and a 0.5\% reduction in glycated haemoglobin to an estimated 17\% reduction in cardiovascular risk.

It must be stated, however, that these correspondences come from the epidemiological literature and not from the trials themselves. None of the sources used here measures a morbidity or mortality endpoint directly. The trials reviewed measure intermediate outcomes---pressure, glycated haemoglobin---and the inference to clinical benefit rests on a relationship established elsewhere. This qualification must be carried into the conclusion of this thesis: telemonitoring has been shown to lower blood pressure, not to reduce mortality.

A second observation emerges from the tables. The four general estimates converge within a band of $-3$ to $-7$~mmHg for systolic pressure and $-0.2$ to $-0.5\%$ for glycated haemoglobin. Convergence across different designs, periods and populations is an argument for robustness. But it conceals a heterogeneity that the subgroup analyses reveal, and that is the subject of the following subsection.

%==========================================================================
\subsection{The decisive determinant: professional oversight}\label{sec:oversight}

\subsubsection{The result from Mao et al.}

The authors compared interventions combining mobile health with professional management---delivered by pharmacists, nutritionists, full-time nurses or sports coaches---against interventions relying exclusively on mobile health:
\begin{align}
  \text{combined intervention:} \quad & -6.17~\text{mmHg} \quad (-8.83 \text{ to } -3.50),\\
  \text{mHealth alone:} \quad & -2.16~\text{mmHg} \quad (-5.07 \text{ to } +0.75).
\end{align}

The confidence interval for the second modality contains zero: mobile health deployed on its own produces no statistically significant effect on systolic blood pressure. The authors state the conclusion unambiguously, that \bq{despite the immense convenience provided by the continuous development of high technology, greater benefits can be produced just by combining human intelligence with artificial intelligence}, and, in their general conclusion, that \bq{the combination of mHealth intervention with professional management is crucial in order to achieve optimum clinical effectiveness}.

This result is corroborated by their sensitivity analysis. Heterogeneity in the blood pressure outcomes falls appreciably---$I^2$ dropping from 67.3\% to 60.1\% and then 49.6\%---when two particular trials are excluded. On inspection, the authors find that in both of these studies disease management was delivered by professional pharmacists, and conclude that \bq{professional intervention could strengthen disease management, which may also be a cause of heterogeneity}. In other words, the variable explaining a substantial share of between-trial dispersion is the degree of human oversight.

\subsubsection{Convergence of sources}

This result does not stand alone. Four sources used in this work converge on the same conclusion by entirely independent routes.

\Textcite{mihevc2025} found that blood pressure and glucose control were achieved without any measurable behaviour change, and attributed the effect to improved communication, medication adherence, early detection of exacerbations and therapy adjustments based on accurate telemonitoring data (\cref{sec:mechanism}). Their intervention included a telemedicine centre coordinated by a general practitioner who reviewed the transmitted values.

\Textcite{who2023hypertension} presents telemonitoring not as a self-management tool but as a mode of medical oversight, a position reported in \cref{sec:hypothesis}.

\Textcite{zhou2024} establish that single-risk-factor interventions fail to reach statistical significance ($-1.54$~mmHg, CI $-4.15$ to $1.06$) while multifactorial interventions do reach it ($-6.50$~mmHg, CI $-9.00$ to $-3.99$), as reported in \cref{tab:subgroups}. Multifactorial design presupposes a coordination the patient alone cannot organise.

\Textcite{bediang2021}, examined below, observed that remote cardiology expertise modified the local provider's diagnosis in 46.2\% of cases, against 17.9\% in control centres.

\subsubsection{Reframing the African puzzle}\label{sec:african-puzzle}

\Cref{sec:outpatient} closed on an open question: why do interventions evaluated in Africa produce no significant effect when the effect is established elsewhere? Three explanations were considered---poorer intervention design, different implementation conditions, or a genuine absence of effect.

\Textcite{mao2020} partly resolve this. They observe that the three trials conducted in developing countries and included in their blood pressure analysis used mobile health exclusively, without associated professional management, whereas in developed countries mobile health is usually combined with professional healthcare management providing disease-related care during the intervention. They conclude explicitly that \bq{this explains why mHealth intervention performed in developed countries demonstrates a higher efficacy with respect to SBP control}.

The geographical gap is therefore, at least in part, a design gap rather than a context effect. What the interventions evaluated in resource-limited settings lack is not the technology but the professional at the end of the chain. This also supplies the missing half of the argument begun in \cref{sec:spec13}, where we noted that the contrast between the African and non-African estimates is partly an artefact of what each review chose to include.

This reframing is the most important result of this section for the positioning of the present thesis. It turns a discouraging finding---telemonitoring supposedly does not work in Africa---into a design specification: telemonitoring must be built as a clinical decision-support system articulated to a clinician, not as a self-management tool handed to the patient.

%==========================================================================
\subsection{The economic gradient and the African evidence deficit}\label{sec:evidence-deficit}

\subsubsection{The scale of the deficit}

Of the 51 studies in \textcite{mao2020}, 36 (70.59\%) were conducted in developed countries and 15 (29.41\%) in developing countries. The fifteen break down as follows: China 5, Iran 3, Egypt 1, India 1, Honduras and Mexico 1, Malaysia 1, Turkey 1, Poland 1, South Africa 1.

Two of the fifty-one studies, that is 3.9\%, were conducted in Africa, and only one of them, that is 2.0\%, in sub-Saharan Africa---the South African study, Egypt belonging to the WHO Eastern Mediterranean Region and to North Africa. This figure must appear explicitly in the thesis: it establishes quantitatively that the global evidence base on telemonitoring effectiveness for chronic disease is, for the region in question, almost non-existent.

This echoes the observation of \textcite{bediang2021} themselves, that \bq{although this evaluation culture is gradually being implemented, there is little robust scientific work on evaluating the effectiveness of telemedicine projects implemented in Africa}.

\subsubsection{An asymmetric gradient by outcome}

The result of \textcite{mao2020} by development level contains an asymmetry that must be rendered accurately, on pain of misstatement.

For systolic blood pressure, the effect is present in developed countries ($-5.72$~mmHg) and absent in developing countries ($+0.25$~mmHg, interval containing zero).

For glycated haemoglobin, the relationship inverts: the effect is larger in developing countries ($-0.52\%$, CI $-0.78$ to $-0.26$) than in developed ones ($-0.35\%$, CI $-0.46$ to $-0.24$), both significant.

This asymmetry suggests that the obstacle encountered in developing countries lies neither in the system's capacity to collect data nor in patient adherence---otherwise both outcomes would be equally affected---but in something specific to blood pressure control. The hypothesis most consistent with the full body of sources is that glycaemic control benefits more from self-management---diet, physical activity, adherence---while blood pressure control depends more on medication adjustment, which requires a prescriber. \Textcite{mao2020} support this reading when they explain the larger effect in type~2 diabetes by noting that mobile health interventions \bq{chiefly achieve blood sugar control by promoting favorable lifestyle and eating habits}.

This hypothesis remains to be verified and we do not present it as established. It is, however, directly testable, and it supplies a relevant secondary research question for the present work. It also sits in tension with one element of \cref{sec:inertia}: the WHO reports that lifestyle modification lowers systolic pressure by only about 1.3~mmHg, so a self-management mechanism cannot account for a large blood pressure effect anywhere---which is consistent with the asymmetry, since it is precisely the blood pressure effect that disappears when the professional is absent.

\subsubsection{What can and cannot be concluded}

It can be concluded that the effect of telemonitoring on systolic blood pressure is not demonstrated in resource-limited settings, and that this absence is plausibly linked to the lack of professional oversight in the interventions evaluated.

It cannot be concluded that telemonitoring is ineffective in these settings. The three trials underpinning the ``developing countries'' subgroup of \textcite{mao2020} are few, as are the seven African trials of \textcite{aovare2023}. Absence of evidence is not evidence of absence, particularly since the two subgroups concerned are precisely those in which the interventions were most stripped-down.

\subsubsection{The regional landscape: capacity, specialists, sustainability}\label{sec:regional-landscape}

The quantitative deficit established above is situated by \textcite{segun2025} within a broader structural account. Given the nature of this source---a Perspective, not a systematic review---the figures it reports are presented here as orders of magnitude drawn from heterogeneous references, and not as pooled estimates.

\paragraph{Connectivity is no longer the principal obstacle.} More than 80\% of the sub-Saharan African population holds a mobile telephone subscription. Mobile internet penetration stood at 25\% in 2023, with internet use of any kind at 37\%, and smartphone adoption is projected to reach 88\% by 2030. These figures converge with those of \textcite{dodoo2022} reported in \cref{sec:digital-barriers}.

\paragraph{Specialist scarcity is quantified, and it is extreme.} The authors report that up to 18\% of sub-Saharan African countries had no registered cardiologist as of 2020, and that Nigeria has approximately one cardiologist per 475\,000 inhabitants. Up to 80\% of global cardiovascular deaths occur in developing countries, with more than one million cardiovascular deaths in sub-Saharan Africa.

These two figures must be read together, because their conjunction states the problem of this thesis with unusual clarity. The population is reachable; the specialist is not. A system that merely transmits data to a cardiologist reproduces the scarcity it was meant to relieve. A system that sorts---that determines which cases require that scarce attention and which do not---addresses it. This is the argument for an automated triage function, and it is a resource argument before it is a technical one.

\paragraph{Infrastructural constraints are precisely located.} Approximately 550 million people lack reliable access to electricity and the internet, and roughly 100\,000 health facilities lack reliable access. Health expenditure represents 5 to 6\% of gross domestic product in sub-Saharan Africa, against 10 to 16\% in wealthier countries. A meta-analysis of five studies found that about half of Ethiopian health workers possessed sufficient digital literacy.

\paragraph{Documented experience in Cameroon.} The authors cite remote interpretation of electrocardiograms and echocardiograms in Cameroon and Uganda, and describe the Cardio Pad, a Cameroonian device allowing rural health workers to record three-, six- or twelve-lead electrocardiograms and transmit them to urban cardiologists. They also cite TELEMED-CAM as having been associated with improved blood pressure control in Cameroon, a study examined in detail in \cref{sec:telemedcam}.

\paragraph{Sustainability, illustrated by a failure at scale.} \Cref{sec:digital-barriers} reported the finding of \textcite{dodoo2022} that most telemedicine projects in the region encountered difficulties once their funding period ended. \Textcite{segun2025} supply a concrete case: Babyl, a Rwandan telehealth platform, reached approximately 2.5 million registered users before ceasing operations in August 2023 for financial reasons. They add that MomConnect costs approximately one million dollars annually to operate.

The Babyl case deserves to be retained, because it refutes an implicit assumption. A system may demonstrate technical feasibility, achieve national scale and acquire millions of users, and still disappear. Adoption is not sustainability. Any proposal issuing from the present work must therefore document its operating cost per patient and the workload it imposes on clinical staff, and not solely its technical performance.

\paragraph{The state of artificial intelligence in the region.} The authors state that very few studies or implementations of artificial intelligence exist in sub-Saharan Africa, that few have been tested in randomised trials, and that there are \bq{very sparse to no data} on the effectiveness of several widely promoted tools, with no published validation study in an African population. They identify two causes: the limited availability of large labelled datasets for training, and the risk that externally developed algorithms perform poorly in African populations. They cite as exceptions artificial-intelligence-guided echocardiography enabling novice operators to screen for rheumatic mitral valve disease in Uganda, and the SPEC-AI Nigeria randomised trial, in which guided screening roughly doubled the detection of cardiomyopathy relative to usual obstetric care.

This last observation is doubly significant for the present work. It confirms the gap identified at the close of \cref{sec:outpatient}---machine learning remains essentially untested in this context. And it identifies the condition under which that gap can be closed honestly: a model trained on data external to the target population cannot be assumed valid within it. This bears directly on the design of the experimental phase of this thesis, and it condemns the approach taken by \textcite{chatrati2022}, who train a diabetes classifier on a North American dataset without discussing its transferability (\cref{sec:chatrati}).

%==========================================================================
\subsection{Field experience in sub-Saharan Africa}\label{sec:field-ssa}

\subsubsection{State of deployment}

The review of \textcite{dodoo2022}, whose method, six barrier dimensions and Cameroonian inventory were reported in \cref{sec:digital-barriers}, finds uneven progress across countries. Two of its findings bear directly on the evidence examined in this section, and are recalled here rather than restated: sustainability, most projects having encountered difficulties once external funding ended; and perceived workload, health professionals viewing telemedicine as an additional task. The first governs how the present proposal must be evaluated, the second how it must be designed.

\subsubsection{The Cameroonian study of Bediang et al.}\label{sec:bediang}

\paragraph{Design.} Controlled multicentre study conducted in 2016 and 2017 in four Cameroonian health facilities: two intervention centres, the district hospitals of Mbouda and Akonolinga, where tele-ECG was implemented, and two control centres, the district hospitals of Foumbot and Sa'a, where it was not. Control centres were matched to intervention centres on level of care, location and distance to the nearest large town with a higher-level facility.

\paragraph{Population.} Patients with at least one cardiovascular risk factor---age 45 or above, or blood pressure of at least 140/90~mmHg---or one or more cardiovascular complaints. Hospitalised patients and those in seriously altered general condition were excluded.

\paragraph{Intervention.} In intervention centres, in addition to usual primary care, patients were offered an on-site resting electrocardiogram using a portable twelve-channel device with dedicated software producing PDF reports. This examination was systematically paired with a request for remote expertise from a cardiologist through an internet tele-expertise platform, to which the local provider transmitted clinical data, diagnostic hypothesis, the electrocardiogram and the current management plan. The cardiologist was required to provide adapted expertise within a maximum of 72 hours. In control centres, patients received usual care only, with a recommendation to consult a cardiologist in Bafoussam or Yaoundé.

\paragraph{Sample size and primary outcome.} The authors anticipated a 25-point difference favouring intervention centres on the primary outcome, with 1:1 allocation, 80\% power, 5\% alpha and a two-sided test, yielding a minimum of 140 patients. The primary outcome was the rate of patient access to an electrocardiogram with cardiologist expertise, analysed by intention to treat: any patient for whom it could not be determined at three months whether they had undergone the examination and received expertise was counted as not having had access.

\subsubsection{Reported results}

One hundred and seventy-one participants were recruited, 93 (54.4\%) in intervention centres and 78 (45.6\%) in control centres. For the primary outcome, 92 of 93 participants (98.9\%) had access to an electrocardiogram with expertise in intervention centres, against 26 of 78 (33.3\%) in control centres, a relative risk of 2.97 (95\% CI 2.17--4.06), $p < 0.01$.

\begin{table}[htbp]
\centering\footnotesize
\caption{Secondary outcomes reported by \textcite{bediang2021}.}\label{tab:bediang-secondary}
\begin{tabularx}{\textwidth}{@{}X >{\raggedleft\arraybackslash}p{3.2cm} >{\raggedleft\arraybackslash}p{3.2cm} c@{}}
\toprule
\textbf{Outcome} & \textbf{Intervention ($n=93$)} & \textbf{Control ($n=78$)} & \textbf{$p$} \\
\midrule
Attendance at 3-month appointment & 92 (98.9\%) & 26 (33.3\%) & $<$ 0.01 \\
Diagnosis modified by the expert & 43 (46.2\%) & 14 (17.9\%) & $<$ 0.01 \\
Inappropriate drug prescriptions & 32 (34.4\%) & 24 (30.8\%) & 0.75 \\
Referrals deemed unnecessary by expert & 22 (23.7\%) & 9 (11.5\%) & 0.01 \\
Favourable BP evolution at 3 months & 38 (40.9\%) & 14 (17.9\%) & $<$ 0.01 \\
Time to perform ECG (days) & $1.3 \pm 4.1$ (0--20) & $7.0 \pm 14.5$ (0--69) & $<$ 0.01 \\
Time to receive expertise (days) & $10.3 \pm 11.6$ (0--50) & $11.2 \pm 27.9$ (1--121) & 0.82 \\
Treatment compliance at 3 months (median, IQR) & 4 (3--5) & 4 (3--5) & 0.08 \\
Participant satisfaction at 3 months (median, IQR) & 5 (5--5) & 4 (3--5) & $<$ 0.01 \\
\bottomrule
\end{tabularx}
\end{table}

Note that the primary outcome and the first secondary outcome carry identical figures, 92 of 93 against 26 of 78. The two are the same measurement: access to the examination required attending the appointment, so attendance and access cannot diverge. This identity is not an error, but it means the study reports one quantity twice, and it anticipates the tautology objection raised below.

The authors conclude that \bq{telemedicine is effective for the management of patients with cardiovascular diseases in primary health care} and that it \bq{could improve healthcare providers' clinical processes, clinical outcomes of patients and their satisfaction}.

\subsubsection{Critical appraisal}\label{sec:bediang-appraisal}

This study must be cited, as it is one of only two Cameroonian evaluations available. It must be cited with five reservations, the fourth of which is decisive.

\textbf{First, this is not a randomised trial.} Centres were matched, not randomly allocated, and patients were assigned by their place of consultation. The risk of residual confounding is therefore high, and the authors do not discuss it.

\textbf{Second, the groups differ at baseline.} Exertional dyspnoea affects 18.3\% of intervention patients against 37.2\% of controls ($p = 0.01$), and the age risk factor 84.9\% against 73.1\% ($p = 0.05$). The groups are therefore not comparable on two clinically relevant characteristics, and no adjustment is reported.

\textbf{Third, the primary outcome is tautological.} Intervention centres are defined as those where the electrocardiograph and tele-expertise platform were installed. Finding that 98.9\% of patients there had access to an electrocardiogram, against 33.3\% where no device exists, measures service availability, not effectiveness. The relative risk of 2.97 is a deployment indicator, not a clinical result.

\textbf{Fourth, and most importantly, the blood pressure result reverses once differential attrition is corrected.} The table reports 38 of 93 patients (40.9\%) with favourable blood pressure evolution in the intervention group, against 14 of 78 (17.9\%) in controls. Yet the same table reports three-month attendance of 92 of 93 against 26 of 78. Since the denominators of the blood pressure outcome are the total group sizes---the percentages confirm this, $38/93 = 40.9\%$ and $14/78 = 17.9\%$---control patients absent from the appointment are counted as not having improved. Restricting the calculation to patients actually reviewed:
\begin{equation}
  \text{intervention:} \; \frac{38}{92} = 41.3\%, \qquad
  \text{control:} \; \frac{14}{26} = 53.8\%.
\end{equation}

\begin{figure}[htbp]
\centering
\begin{tikzpicture}
\begin{axis}[ybar,width=0.78\textwidth,height=6cm,bar width=18pt,ymin=0,ymax=65,
  ylabel={Favourable BP evolution at 3 months (\%)},
  symbolic x coords={As reported (ITT),Among patients reviewed},
  xtick=data,xticklabel style={font=\footnotesize},
  enlarge x limits=0.45,legend pos=north west,legend style={font=\footnotesize},
  nodes near coords,every node near coord/.append style={font=\footnotesize}]
\addplot[fill=cRef!75,draw=cRef] coordinates {(As reported (ITT),40.9)(Among patients reviewed,41.3)};
\addplot[fill=cAfrica!75,draw=cAfrica] coordinates {(As reported (ITT),17.9)(Among patients reviewed,53.8)};
\legend{Intervention,Control}
\end{axis}
\end{tikzpicture}
\caption{The blood pressure result of \textcite{bediang2021} before and after restriction to patients actually reviewed at three months. The intervention proportion is unchanged; the control proportion triples, because 52 of its 78 patients never attended and were counted as not improved.}\label{fig:bediang-attrition}
\end{figure}

The apparent effect disappears and reverses (\cref{fig:bediang-attrition}). Among patients actually assessed at three months, the proportion with favourable blood pressure evolution is higher in the control group. Two readings are possible: either the displayed denominator is indeed the full sample, in which case the reported blood pressure effect is an artefact of the difference in attendance; or the authors made a denominator error. In either case, the claim that this study demonstrates improved blood pressure control is not defensible as it stands, and this thesis does not repeat it. The authors are themselves cautious: they write that the study \bq{highlighted telemedicine's potential to improve clinical outcomes} and that \bq{this requires more robust studies}.

A caution on the reversed figure is nonetheless in order. The 53.8\% rests on 26 control patients, and the subset of controls who attended is self-selected---plausibly the more motivated or the more symptomatic. The corrected comparison therefore does not establish that usual care performs better; it establishes only that the reported comparison does not establish the opposite.

\textbf{Fifth, the timing protocol was not met and the gain on expertise is nil.} The cardiologist was required to respond within 72 hours; the observed mean delay is 10.3 days, with a maximum of 50 days. More importantly, this delay is not statistically different from that of control centres (11.2 days; $p = 0.82$). The system therefore accelerated the performance of the examination---1.3 against 7.0 days, $p < 0.01$---but not the obtaining of specialist advice. This distinction matters greatly for system design: the bottleneck was not transmission but expert availability.

It should also be noted that treatment adherence does not differ significantly between groups ($p = 0.08$), as the authors themselves acknowledge, and that inappropriate prescriptions do not differ either ($p = 0.75$).

\subsubsection{What the study nonetheless establishes}

Two results survive these reservations, and they matter.

\textbf{Tele-expertise substantially alters clinical decisions.} The local provider's diagnosis was modified by the expert in 46.2\% of cases, against 17.9\% in control centres. Nearly one primary-care diagnosis in two was therefore revised by the cardiologist. This result does not depend on attrition, since it concerns patients who received expertise, and it directly documents the added value of remote specialist advice---converging with the finding of \textcite{mao2020} on the importance of professional oversight (\cref{sec:oversight}).

\textbf{Patient satisfaction is markedly higher:} median 5 of 5 with an interquartile range of 5 to 5, against 4 with a range of 3 to 5 ($p < 0.01$). The authors offer an explanation worth retaining: the effect \bq{may be due to the high patients' confidence in the local healthcare provider whom they know to be supported by the expert}. Tele-expertise would thus also act by reinforcing the credibility of the front-line provider in the patient's eyes.

\subsubsection{The earlier Cameroonian trial: TELEMED-CAM}\label{sec:telemedcam}

\Textcite{bediang2021} are not the only Cameroonian evaluation available. Eight years earlier, \textcite{kingue2013} conducted a study bearing directly on hypertension telemanagement, and its design differs in a way that makes the two complementary rather than redundant: Bediang evaluates remote expertise for diagnosis, Kingue evaluates remote support for treatment.

\paragraph{Design and setting.} Prospective interventional study conducted from January to December 2008. A telemedicine centre at Yaoundé General Hospital, staffed by five cardiologists and three nephrologists, served thirty rural health centres located between 50 and 250~km from Yaoundé. Six centres already hosting a Ministry of Health hypertension pilot unit were allocated to the control group; of the twenty-four remaining, ten were randomly selected for the intervention.%
\verify{Six control centres and ten intervention centres leave fourteen of the thirty centres unaccounted for. The participant counts are consistent with six control centres (103 participants, about 17 per centre, against 165 over ten intervention centres), so the fourteen were presumably not enrolled. Confirm against the Methods section of \textcite{kingue2013}.}

\paragraph{Participants.} 268 in total, 165 in the intervention group and 103 in usual care. Adults with uncontrolled hypertension resident in the area for at least twelve months were eligible; pregnant women and patients with severe target-organ damage were excluded. The power calculation targeted a 10~mmHg difference in systolic pressure, assuming a standard deviation of 18~mmHg, a two-sided alpha of 0.01 and 90\% power, requiring at least 96 participants per arm. The computation reproduces: $2(z_{0.005}+z_{0.10})^2\sigma^2/\Delta^2 = 2(2.576+1.282)^2(18^2)/10^2 \approx 96$. The control arm, at 103, therefore clears that requirement only narrowly.

\paragraph{The intervention was a package, and this is essential to its interpretation.} It comprised four distinct components: restriction of prescribing to generic antihypertensives drawn from six classes, with a thiazide diuretic as default first agent unless contraindicated and an angiotensin-converting-enzyme inhibitor preferred in diabetes or renal disease; lifestyle counselling covering weight, sodium, smoking, alcohol and physical activity; training of general practitioners and nurses at intervention centres in treatment protocols and in telecommunication; and telemanagement proper. Control-centre staff received neither the protocol training nor the telecommunication training.

\paragraph{Telemanagement in operation.} Each intervention centre received a GSM handset. Providers transmitted patient data daily by text message or voicemail to the telemedicine centre, which returned prompt telephone feedback. Visits occurred every two weeks, twelve visits over twenty-four weeks. The volume of activity is documented: 1\,187 calls from intervention clinics, 791 unique provider calls, approximately six calls per provider per day lasting about five minutes each.

This last figure deserves emphasis, because it quantifies something no other source in this chapter supplies: the workload that a telemanagement system imposes on the remote expert. Thirty minutes per provider per day, multiplied across ten centres, represents a substantial and recurrent commitment of specialist time. Read alongside the specialist density reported by \textcite{segun2025}---up to 18\% of countries with no cardiologist at all---this establishes that a model requiring systematic human review of every transmission does not scale. It is a further argument for automated triage.

\paragraph{Results.} Baseline imbalance is substantial and was not adjusted for. Intervention participants had higher systolic pressure (169 against 161~mmHg, $p=0.01$), higher diastolic pressure (100 against 95~mmHg, $p=0.01$), and far fewer diabetics (9.1\% against 43.6\%, $p=0.02$).

Adherence to scheduled visits was significantly higher in the intervention group at every grade of hypertension ($p \leq 0.04$). Systolic pressure fell significantly in the intervention group from the second visit onward while remaining flat in the control group; diastolic pressure was lower in the control group at the final two visits.

\begin{table}[htbp]
\centering\small
\caption{Outcomes at the final visit \parencite{kingue2013}.}\label{tab:telemedcam-results}
\begin{tabularx}{\textwidth}{@{}X c c c@{}}
\toprule
\textbf{Outcome at final visit} & \textbf{Intervention} & \textbf{Usual care} & \textbf{$p$} \\
\midrule
Stage III hypertension: at target & 50\% & 39.1\% & --- \\
Stage III: improved but not at target & 33.3\% & 27.5\% & --- \\
Stage III: no improvement & 16.7\% & 33.3\% & 0.04 \\
Stage I--II hypertension: at target & 65.2\% & 70\% & 0.20 \\
\bottomrule
\end{tabularx}
\end{table}

The effect is confined to severe hypertension. For stages I and II, the difference is not significant, and the point estimate favours usual care. This restriction is not emphasised by the authors but it bears directly on the present work: it suggests that the benefit of telemanagement concentrates where the margin for therapeutic adjustment is widest.%
\verify{The source labels the 16.7\% group ``control groups''. Reading 16.7\% as the intervention arm is the only assignment under which both columns sum to about 100\%. Confirm against Table~4 of \textcite{kingue2013}.}

\paragraph{Reservations.} The authors list most of these themselves, with commendable candour. Allocation was not strictly random and assessment was not blinded. Clustering was not accounted for in either the sample size calculation or the analysis: with ten intervention clusters and six control clusters, this inflates the effective sample size and understates the standard errors, so the reported $p$-values are optimistic. Changes in treatment prompted by telecare calls were not recorded. Telehealth was one component of a package, so its independent effect cannot be isolated---the intervention simultaneously changed the formulary, trained the providers and added the telephone link, and attributing the result to telemonitoring is not supportable. Few practitioners at few sites participated. The study was funded by a Sanofi-Aventis grant.

To these must be added three observations the authors do not make. The baseline systolic imbalance of 8~mmHg favours the intervention group under any change-from-baseline analysis through regression to the mean, and no adjustment is reported. The diabetes imbalance---9.1\% against 43.6\%---is extreme, unexplained, and clinically material, diabetes being a determinant of blood pressure response. And the Methods describe a twelve-month study period while the intervention lasted twenty-four weeks; the source does not reconcile this.

\paragraph{A connection to be verified.} \Cref{sec:hypothesis} reported that the sensitivity analysis of \textcite{aovare2023}, excluding one trial, rendered the pooled systolic effect for Africa statistically significant at $-2.22$~mmHg. The excluded trial is attributed to Kingue et al. If it is the present study---which is probable given the topic, the period and the setting---then TELEMED-CAM would be the Cameroonian trial that drew the African meta-analysis toward the null.%
\verify{This identification is not asserted. The full text of \textcite{aovare2023} was not available to us; confirm from its sensitivity-analysis table which trial was excluded before using this connection anywhere in the thesis.}

\subsubsection{What the two Cameroonian studies establish jointly}\label{sec:cameroon-joint}

\begin{table}[htbp]
\centering\footnotesize
\caption{The two Cameroonian evaluations compared.}\label{tab:cameroon-pair}
\begin{tabularx}{\textwidth}{@{}>{\raggedright\arraybackslash}p{2.6cm} >{\raggedright\arraybackslash}X >{\raggedright\arraybackslash}X@{}}
\toprule
 & \textbf{\textcite{kingue2013}} & \textbf{\textcite{bediang2021}} \\
\midrule
Period & 2008 & 2016--2017 \\
Function evaluated & telemanagement of treatment & tele-expertise for diagnosis \\
Design & prospective, partly randomised by centre & controlled, matched centres, not randomised \\
Sites & 30 rural health centres, hub at Yaoundé & 4 district hospitals \\
Participants & 268 & 171 \\
Duration & 24 weeks, 12 visits & 3 months, 2 visits \\
Technology & GSM handset, SMS and voicemail & portable 12-lead ECG, web platform \\
Principal result & BP control improved in stage III only & access to ECG with expertise, 98.9\% vs 33.3\% \\
Retention result & adherence to visits higher at every grade & attendance at 3 months, 98.9\% vs 33.3\% \\
Confirmed weakness & clustering ignored, package effect, industry funding & not randomised, groups not comparable, differential attrition \\
\bottomrule
\end{tabularx}
\end{table}

Three conclusions may be drawn from the pair, and only three.

\textbf{First, retention is the most robust finding.} Both studies report a substantial improvement in patient attendance, by different mechanisms and eight years apart. This is the one result that neither the absence of randomisation in Bediang nor the package design in Kingue undermines, since the follow-up contact is itself the intervention.

\textbf{Second, the effect on blood pressure control is not established in Cameroon.} The apparent result of \textcite{bediang2021} dissolves on correction for differential attrition (\cref{sec:bediang-appraisal}). The result of \textcite{kingue2013} is confined to stage III, is non-significant for stages I and II, and rests on an analysis that ignores clustering and leaves an 8~mmHg baseline imbalance unadjusted. The honest statement is that no methodologically sound Cameroonian evidence of a telemonitoring effect on blood pressure control currently exists.

\textbf{Third, the value added is located in the clinical decision, not in the patient.} Bediang's expert modified the local provider's diagnosis in 46.2\% of cases. Kingue's intervention operated through provider training, protocol standardisation and telephone support to the clinician---not through any patient-facing device. Neither study placed a tool in the patient's hands. Both converge with the finding of \textcite{mao2020} on the necessity of professional oversight, and with that of \textcite{mihevc2025} that the effect arises from therapeutic adjustment rather than behaviour change (\cref{sec:mechanism}).

%==========================================================================
\subsection{Retention and access to specialists}\label{sec:retention-access}

\subsubsection{Retention}

The most striking result of \textcite{bediang2021} is not the primary outcome but the three-month attendance rate: 98.9\% against 33.3\%. It must be read with the same caution as the rest---the absence of randomisation precludes causal attribution, and the follow-up appointment was an integral part of the intervention protocol---but its magnitude is considerable and it echoes independent observations (\cref{tab:retention}).

\begin{table}[htbp]
\centering\small
\caption{Retention reported across settings.}\label{tab:retention}
\begin{tabularx}{\textwidth}{@{}>{\raggedright\arraybackslash}p{3.6cm} >{\raggedright\arraybackslash}p{3.4cm} X@{}}
\toprule
\textbf{Source} & \textbf{Setting} & \textbf{Retention} \\
\midrule
\textcite{who2023hypertension} & routine care, limited resources & loss to follow-up in 50\% of patients or more \\
\textcite{bediang2021} & controls, Cameroon & 33.3\% attendance at 3 months \\
\textcite{bediang2021} & intervention, Cameroon & 98.9\% attendance at 3 months \\
\textcite{kingue2013} & intervention, rural Cameroon & visit adherence significantly higher at every grade ($p \leq 0.04$) \\
\textcite{zhou2024} & 20 trials, 4\,168 participants & 85.2\% retention \\
\textcite{mihevc2025} & trial, Slovenia & 91.4\% at 12 months \\
HEARTS programme, Sri Lanka & routine & 86\% returning every 3 months \\
\bottomrule
\end{tabularx}
\end{table}

The 33.3\% observed in Cameroonian control centres is lower than what the WHO describes as the ordinary situation in resource-limited settings, and it supplies a local measurement of a phenomenon documented until now only through general statements (\cref{sec:ltfu}).

One methodological caveat applies, already stated in \cref{sec:feasibility}: the rates drawn from trials---85.2\% and 91.4\%---concern volunteer participants, selected and prompted by a research team, and do not estimate retention in a routine programme. Only the 86\% from the Sri Lankan programme comes from routine activity.

\subsubsection{Access to specialists}\label{sec:specialist-access}

The distinctive contribution of the Cameroonian studies is to document this dimension, absent from the European and Asian trials. The control centres of \textcite{bediang2021} referred patients to a cardiologist in Bafoussam or Yaoundé; two-thirds never got there.

\Textcite{segun2025} supply the structural measurement that explains why. Up to 18\% of sub-Saharan African countries had no registered cardiologist as of 2020, and Nigeria, one of the better-resourced countries in the region, has approximately one cardiologist per 475\,000 inhabitants. The referral pathway on which control-centre patients depended is not merely long; in much of the region it terminates nowhere.

This connects the present section to \cref{sec:cascade}, where we established that 57\% of hypertensive adults in the WHO African Region are undiagnosed and that 37\% of those diagnosed are untreated. Distance to the specialist is one concrete mechanism of these losses, and \textcite{gafane2025} identified it among individual barriers alongside cost and health literacy.

Two limitations revealed by the Cameroonian studies themselves must however be set against this.

\textbf{Remote access does not create expert availability.} The time to obtain expertise in \textcite{bediang2021} was 10.3 days on average, against a protocol maximum of 72 hours, and not statistically different from the control centres (11.2 days; $p=0.82$). The system accelerated the examination---1.3 against 7.0 days, $p<0.01$---but not the specialist opinion. Physical access was replaced by remote access, but the latter remains constrained by a scarce resource that technology does not manufacture.

\textbf{Systematic review does not scale.} \Textcite{kingue2013} quantify the cost of the alternative: approximately six calls per provider per day, of about five minutes each, with 1\,187 calls logged over the study. A model in which every transmission is reviewed by a specialist consumes expert time in proportion to the number of patients enrolled, and is therefore self-limiting at exactly the point where it would be most useful.

These two observations converge on a single design requirement. A telemonitoring system deployed at scale in this region must triage: it must determine which transmissions warrant specialist attention and which do not, and it must do so reliably enough that the clinician can trust the ones it filters out. This is not an optional refinement. It is the condition under which the system is deployable at all, and it is the function that the detection model developed in this thesis is intended to serve.

%==========================================================================
\subsection{Scope, limitations and positioning}\label{sec:positioning}

\subsubsection{Hierarchised synthesis}

\textbf{First finding, established by meta-analysis of 51 randomised trials.} Telemonitoring and mobile health interventions produce a statistically significant effect on all four intermediate outcomes of hypertension and diabetes control: $-2.99$~mmHg systolic, $-1.14$~mmHg diastolic, $-0.39\%$ glycated haemoglobin, $-0.52$~mmol/L fasting glucose.

\textbf{Second finding, established by subgroup analysis and corroborated by three independent sources.} This effect depends on professional oversight: $-6.17$~mmHg for combined interventions against a non-significant $-2.16$~mmHg for mobile health alone.

\textbf{Third finding, established by a randomised trial on the target population.} A joint blood pressure and glucose home monitoring system, backed by a telemedicine centre coordinated by a physician, produces over twelve months a 6.9~mmHg reduction in systolic pressure and a 0.5\% reduction in glycated haemoglobin in older patients with both hypertension and type~2 diabetes.

\textbf{Fourth finding, established by enumeration.} The evidence base is almost empty for sub-Saharan Africa: one study out of 51 in \textcite{mao2020}, none outside South Africa; and the only two syntheses covering the region find no significant effect on blood pressure.

\textbf{Fifth finding, of low evidential strength.} The two available Cameroonian studies establish that tele-expertise massively improves access to specialist examination and modifies the local provider's diagnosis in nearly one case in two, and that both interventions improved attendance; neither establishes, contrary to the claim of one of them, an improvement in blood pressure control.

\subsubsection{Limitations of this section}

\textbf{Intermediate outcomes.} No source measures mortality or morbidity. Inference to clinical benefit rests on relationships established in other literatures.

\textbf{Heterogeneity.} All meta-analyses report moderate to high heterogeneity: $I^2$ of 57\% to 67\% in \textcite{mao2020}, 82\% in \textcite{zhou2024} for systolic pressure. Aggregate estimates are averages over very different interventions.

\textbf{Publication bias.} Egger's test is significant for glycated haemoglobin in \textcite{mao2020} ($p = 0.036$), even though trim-and-fill correction changed nothing.

\textbf{Observation period.} \Textcite{mao2020} themselves conclude that \bq{the long-term effects of such interventions remain unknown and will need to be tested in a more representative population over a longer time period}. This compounds the effect erosion documented in \cref{sec:feasibility}.

\textbf{Evidential strength of the first Cameroonian source.} \Textcite{bediang2021} is a non-randomised study, with groups not comparable at baseline, a tautological primary outcome, differential attrition untreated and the timing protocol not met.

\textbf{Evidential strength of the second Cameroonian source.} \Textcite{kingue2013} is partly randomised at the centre level, does not account for clustering in either sample size or analysis, leaves a substantial baseline imbalance unadjusted, evaluates a package of which telemanagement is one component, and was funded by a pharmaceutical grant. Its blood pressure result is confined to stage III hypertension.

\textbf{Evidential strength of the regional landscape source.} \Textcite{segun2025} is a Perspective article, not a systematic review. It declares no search strategy, no inclusion criteria and no quality appraisal, and a portion of its references are grey literature, press items or organisational materials. Its figures are used here as orders of magnitude and contextual indicators, never as effect estimates.

\subsubsection{Positioning of the thesis}

The elements assembled across the four sections of this chapter now permit a precise statement of the intended contribution.

\textbf{What is established.} The need is substantial and documented. The failure of the episodic model is characterised and its mechanism identified. The technical modalities of out-of-hospital follow-up are specified down to measurement protocol detail. The effectiveness of telemonitoring is demonstrated in high-income countries, including on precisely this thesis's population. The condition of that effectiveness---articulation to a clinical decision-maker---is identified.

\textbf{What is not.} Effectiveness in sub-Saharan Africa is not demonstrated, and the evidence base there is almost non-existent: one study out of 51 in \textcite{mao2020}, none outside South Africa, and the only two Cameroonian evaluations both methodologically compromised. No clinically evaluated telemonitoring system applies machine learning to the longitudinal series of an individual patient; the sole published system combining home blood pressure and glucose readings with a learned model, that of \textcite{chatrati2022}, classifies a single reading on a non-transferable dataset, leaves isolated systolic hypertension unclassified, and has never been evaluated on patients (\cref{sec:chatrati}). No system personalises its alert thresholds to the individual. None detects drift rather than threshold crossing. And none triages the scarce resource---specialist time---which \textcite{kingue2013} show to be consumed at roughly thirty minutes per provider per day under systematic review, and which \textcite{segun2025} show to be absent altogether in up to 18\% of the region's countries.

One further constraint must be stated, because it is the condition of honesty for the experimental phase. \Textcite{segun2025} observe that externally developed algorithms risk performing poorly in African populations, and that the shortage of large labelled datasets is itself a barrier. A model whose training data come from elsewhere cannot be presumed valid here. This thesis must therefore either train on data collected in context, or document explicitly and quantitatively the transfer assumption it makes---and it must not repeat the silence of \textcite{chatrati2022} on that question.

\textbf{The contribution.} These gaps jointly define the space for this thesis. A telemonitoring system designed for the Cameroonian context, articulated to a clinician rather than handed to the patient, able to learn each patient's own parameters from their history, to detect drift before a fixed threshold is crossed, and to prioritise the demands placed on clinical staff, would address all of them simultaneously. The following part examines the methods capable of performing this detection function.


%##########################################################################
% CHAPTER 2
%##########################################################################
\chapter{Technological State of the Art: IoT Sensors, Cloud Architectures and Predictive AI}\label{ch:techsota}

\section{Multi-Source Acquisition: IoT/IoMT, Biomedical Sensors, Mobile Entry}\label{sec:acquisition}

\subsection*{Purpose and status of the sources}

\Cref{ch:background} established what the system must do: produce a longitudinal, repeated, trend-exploitable record of blood pressure and blood glucose, articulated to a clinician rather than handed to the patient, using validated measurement devices, without adding documentary burden, and able to triage so that scarce specialist attention is spent only where needed. This chapter examines how that is technically achieved. The present section addresses the first link in the chain: acquisition.

\paragraph{Efendi et al. (2025).} \textcite{efendi2025}, published in \emph{Health Science Reports} 8:e70498, is an original research article describing and validating a three-layer IoT architecture for elderly monitoring. It is the only source in this corpus that reports a complete implementation with instrumentation, cloud layer, device-level validation and a learned classifier. It is therefore examined in detail, including its weaknesses, which are instructive.

\paragraph{Liu and Wang (2025).} \textcite{liu2025review} survey IoT mobile sensing devices for chronic disease management. As established in \cref{sec:diagnosis}, the article presents itself as a systematic review while reporting no protocol, no flow diagram, no inclusion criteria and no quality instrument; we therefore treat it as a narrative survey, useful for mapping the space of device categories, not as pooled evidence.

\paragraph{The research framing document.} \textcite{cadrage2026} is not a bibliographic source. It is the framing document prepared in response to the supervisory committee, and it states the intended contributions, hypotheses, limitations and dataset design of the present work. It is used here to confront the technical state of the art with the design actually proposed.

\paragraph{A source announced and not obtained.} Tai et al.\ (2024), \emph{Transactions on Emerging Telecommunications Technologies}, was announced in the plan for this section and is not available to us. A one-row summary of it exists in our own working literature table, which is not the article. Its three-tier architecture proposal is therefore not cited in this thesis.%
\verify{Obtain the full text of Tai et al.\ (2024) and either cite it properly or leave this sentence as the record of its absence. A claim recorded in one's own summary table cannot be cited as a source.}

%==========================================================================
\subsection{Three acquisition pipelines, three different kinds of data}\label{sec:pipelines}

The system draws on three sources that differ not only in their technology but, more importantly for what follows, in their statistical properties. Treating them as a single undifferentiated input stream would be the first modelling error, and the distinction must be established before any architecture is described (\cref{fig:pipelines}).

\begin{figure}[htbp]
\centering
\begin{tikzpicture}[node distance=0.5cm and 1.1cm,
  pl/.style={draw,rounded corners,align=center,font=\footnotesize,
             minimum width=3.5cm,minimum height=1.5cm},
  tgt/.style={draw,rounded corners,align=center,font=\footnotesize,fill=cGrey!12,
             minimum width=3.2cm,minimum height=1.5cm},
  ar/.style={-{Stealth[length=2.5mm]},thick}]
  \node[pl,fill=cRef!14] (p1) {\textbf{Pipeline 1}\\ Consultation\\ 1--2 per year};
  \node[pl,fill=cAfrica!12,below=of p1] (p2) {\textbf{Pipeline 2}\\ Self-report\\ several per week};
  \node[pl,fill=cGrey!14,below=of p2] (p3) {\textbf{Pipeline 3}\\ Sensors\\ up to 1~Hz};
  \node[tgt,right=2.4cm of p1] (gt) {ground truth\\ and labelling};
  \node[tgt,right=2.4cm of p2] (ser) {principal\\ longitudinal series};
  \node[tgt,right=2.4cm of p3] (cov) {contextual\\ covariates};
  \draw[ar] (p1) -- node[above,font=\scriptsize]{MNAR} (gt);
  \draw[ar] (p2) -- node[above,font=\scriptsize]{MNAR} (ser);
  \draw[ar] (p3) -- node[above,font=\scriptsize]{MCAR} (cov);
\end{tikzpicture}
\caption{The three acquisition pipelines, their sampling rates and their roles. Two of the three carry missingness that depends on the patient's state (MNAR); only the instrumented pipeline fails at random (MCAR), and it is the one that cannot measure the parameters of interest.}\label{fig:pipelines}
\end{figure}

\subsubsection{Pipeline 1: clinical consultation data}

\textbf{Content.} Hospital history, laboratory results, vital signs recorded at the visit, background variables and clinical context. The framing document specifies the record structure: systolic and diastolic blood pressure, blood glucose, heart rate and measurement timestamp, together with background variables---age, sex, body mass index---and clinical context---history, treatments.

\textbf{Properties.} This pipeline is low-frequency, high-reliability, and irregularly sampled by clinical rather than temporal logic. A measurement exists because a consultation occurred, and a consultation occurred because the patient attended---which \cref{sec:followup-requires} established happens for perhaps one to two contacts per year against the seven to eight the protocol requires. The sampling process is therefore not independent of the patient's state: a patient who deteriorates may consult more often, or may disappear entirely. In Rubin's terminology this is data missing not at random, and it has direct consequences for any model trained on the series: the missingness pattern is itself informative, and a model that ignores it will be fitted on a sample systematically unrepresentative of the states it is meant to detect.

\textbf{Reference status.} Within the system, this pipeline supplies the ground truth. The framing document states that risk level is derived from recognised clinical thresholds---ESH 2023 for blood pressure, ADA standards for blood glucose---and validated by the partner physician.%
\verify{Add formal bibliographic entries for the 2023 European Society of Hypertension guideline and the American Diabetes Association Standards of Care edition actually used, and state the edition year. Both are cited as label sources and must be referenced as such.}

\subsubsection{Pipeline 2: mobile self-report}

\textbf{Content.} Manual patient entry of home measurements.

\textbf{Properties.} This pipeline is medium-frequency, variable-reliability, and subject to behavioural artefacts. \Cref{sec:apps} established its failure modes precisely: in the commercial applications reviewed by \textcite{alessa2019}, 73 of 84 required manual entry of a value the device had just displayed; in \textcite{chatrati2022}, the entire acquisition is desktop keyboard entry (\cref{sec:chatrati}). Three consequences follow. Transcription error---digit transposition, systolic--diastolic inversion---is undetectable downstream. The timestamp records the moment of entry, not the moment of measurement, so a patient entering three days of readings on a Sunday evening produces a series whose temporal structure is fictitious. And selective reporting is possible: a patient may omit a reading they find discouraging, which biases the series in precisely the direction that matters clinically.

The last of these deserves emphasis, because it is not a noise problem but a bias problem. Noise can be averaged away---this is the whole content of $\operatorname{Var}(\hat\mu_i) = \sigma_\varepsilon^2/n$ established in \cref{sec:variability}---whereas selective omission cannot, and it systematically removes the extreme values a detection model exists to find.

\subsubsection{Pipeline 3: connected biomedical sensors}

\textbf{Content.} Instrumented acquisition transmitting without human intervention. In the architecture described by \textcite{efendi2025}, a MAX30100 optical sensor for heart rate and oxygen saturation and a MAX30205 digital temperature sensor, both read by an ESP32 microcontroller.

\textbf{Properties.} This pipeline is high-frequency, instrument-limited, and free of transcription and reporting bias. \Textcite{efendi2025} collect one record per second over ten minutes, yielding 600 observations per participant. Reliability is no longer a function of patient behaviour but of sensor accuracy, signal conditioning and transmission integrity, and its failures---a flat battery, a lost connection---are plausibly independent of the patient's clinical state.

Its limitation is different in kind: it measures what the sensor can measure, which, as the next subsection shows, does not include blood pressure or blood glucose at an acceptable level of validation.

\subsubsection{The properties that matter for modelling}

\begin{table}[htbp]
\centering\footnotesize
\caption{Statistical properties of the three acquisition pipelines.}\label{tab:pipelines}
\begin{tabularx}{\textwidth}{@{}>{\raggedright\arraybackslash}p{3.1cm} >{\raggedright\arraybackslash}X >{\raggedright\arraybackslash}X >{\raggedright\arraybackslash}X@{}}
\toprule
\textbf{Property} & \textbf{P1: consultation} & \textbf{P2: self-report} & \textbf{P3: sensors} \\
\midrule
Sampling frequency & 1--2 per year & several per week & up to 1~Hz \\
Sampling regularity & irregular, event-driven & irregular, behaviour-driven & regular \\
Missingness mechanism & MNAR: attendance depends on state & MNAR: omission possible & MCAR: technical failure \\
Dominant error & measurement conditions & transcription, timestamp, selection & instrument accuracy, movement artefact \\
Parameters available & BP, glucose, HR, laboratory & BP, glucose, HR & HR, SpO$_2$, temperature \\
Blood pressure available & yes, reference quality & yes, if validated cuff & \textbf{no} \\
Blood glucose available & yes, laboratory & yes, if glucometer & \textbf{no} \\
Role in the system & ground truth and labelling & principal longitudinal series & contextual covariates \\
\bottomrule
\end{tabularx}
\end{table}

The final row of \cref{tab:pipelines} is the conclusion of this subsection and it should govern the architecture. The two parameters that define the clinical problem---blood pressure and blood glucose---are not obtainable from the optical sensor pipeline. That pipeline supplies heart rate, oxygen saturation and temperature, which are covariates: useful for characterising the context of a reading, for detecting movement, for enriching the feature space, but not substitutes for the measurements on which the clinical decision rests.

The record structure of the framing document confirms this allocation: the fields listed are systolic pressure, diastolic pressure, blood glucose, heart rate and timestamp. Of these, only heart rate can come from an optical sensor.

%==========================================================================
\subsection{Physical principles and hardware constraints}\label{sec:hardware}

\subsubsection{Photoplethysmography}\label{sec:ppg}

Photoplethysmography is a non-invasive optical method for measuring variations in blood volume in peripheral tissue. \Textcite{navarro2023} defines it as the use of a light-emitting diode and a photodetector to measure changes in blood volume: light is emitted into the skin and the reflected or transmitted fraction detected, capturing the pulsatile variation in blood flow driven principally by the heartbeat. She distinguishes transmissive sensors, commonly placed on fingertips where tissue is thin, from reflective sensors, more practical on the wrist, arm, chest or forehead where thickness makes transmission impractical.

\Textcite{efendi2025} describe the same mechanism for the MAX30100: the sensor works \bq{by shining both lights on the finger and detecting the amount of reflected light with a photo-detector}, and they name this method photoplethysmography explicitly.

The underlying physics is Beer--Lambert attenuation. For light of incident intensity $I_0$ traversing an absorbing medium,
\begin{equation}
  I = I_0 \, e^{-\varepsilon c d},
\end{equation}
where $\varepsilon$ is the molar extinction coefficient of the absorber at the wavelength used, $c$ its concentration and $d$ the optical path length. The measured signal decomposes into a large quasi-static component, conventionally written DC, arising from tissue, bone, venous blood and the non-pulsatile arterial fraction, and a small pulsatile component, written AC, arising from the arterial volume change at each systole (\cref{fig:ppg}). Heart rate is recovered from the periodicity of the AC component alone; it does not require absolute calibration, which is why heart rate from photoplethysmography is comparatively robust.

\begin{figure}[htbp]
\centering
\begin{tikzpicture}
\begin{axis}[width=0.86\textwidth,height=5.4cm,
  xlabel={Time (s)},ylabel={Detected intensity (a.u.)},
  xmin=0,xmax=3,ymin=0,ymax=1.15,ytick=\empty,
  axis lines=left,enlarge y limits=false]
\addplot[thick,cRef,domain=0:3,samples=400]
  {0.72 + 0.06*sin(deg(2*pi*1.2*x)) + 0.02*sin(deg(4*pi*1.2*x))};
\addplot[dashed,cGrey,no markers,domain=0:3] {0.72};
\draw[{Stealth[length=2mm]}-{Stealth[length=2mm]},cAfrica,thick]
  (axis cs:2.6,0) -- (axis cs:2.6,0.72) node[midway,right,font=\footnotesize,text=cAfrica]{DC};
\draw[{Stealth[length=2mm]}-{Stealth[length=2mm]},cAfrica,thick]
  (axis cs:2.85,0.64) -- (axis cs:2.85,0.80) node[midway,right,font=\footnotesize,text=cAfrica]{AC};
\end{axis}
\end{tikzpicture}
\caption{Schematic decomposition of a photoplethysmographic signal into its quasi-static (DC) and pulsatile (AC) components. Synthetic illustration, not recorded data. Heart rate follows from the period of the AC component and needs no absolute calibration; oxygen saturation needs the AC/DC ratio at two wavelengths and an empirical curve.}\label{fig:ppg}
\end{figure}

\subsubsection{Dual wavelength and oxygen saturation}

Oxygen saturation requires a second wavelength. \Textcite{efendi2025} state that the MAX30100 light-emitting diodes \bq{have wavelengths of 660 nm and 880 nm, respectively}.

The principle rests on a differential absorption property. Oxyhaemoglobin and deoxyhaemoglobin have different extinction spectra: at 660~nm, in the red, deoxyhaemoglobin absorbs more strongly; in the near infrared, around 880 to 940~nm, oxyhaemoglobin absorbs more strongly. Measuring the pulsatile and static components at both wavelengths yields a ratio of ratios,
\begin{equation}
  R = \frac{AC_{660}/DC_{660}}{AC_{880}/DC_{880}},
\end{equation}
from which saturation is derived by an empirical calibration curve, conventionally of the form $\mathrm{SpO_2} \approx a - bR$, with coefficients established on human volunteers under controlled desaturation.

Two observations follow, and both are design constraints.

\textbf{First, the calibration is empirical, not physical.} The coefficients $a$ and $b$ are fitted, not derived. A sensor used outside the population and conditions of its calibration has no guaranteed accuracy, and skin pigmentation is one documented source of divergence in pulse oximetry. For a system intended for deployment in Cameroon this is not an academic point, and it must be acknowledged even though none of the sources in this corpus addresses it. It is also the same argument, in a different register, as the transferability warning that \textcite{segun2025} apply to externally trained algorithms (\cref{sec:regional-landscape}): a calibration is a model fitted on a population, and it travels no better than any other.

\textbf{Second, the description given by \textcite{efendi2025} is physically imprecise.} They write that \bq{the saturated oxygen in arterial blood has the property of absorbing infrared radiation. The higher the hemoglobin level, the more IR light is absorbed.} This conflates two distinct quantities: the concentration of haemoglobin and the saturation of haemoglobin with oxygen. What the ratio-of-ratios method exploits is the relative difference in absorption between the oxygenated and deoxygenated forms at two wavelengths, not the total haemoglobin concentration---which is in fact largely cancelled by the ratio. The statement as written would imply that an anaemic patient reads as desaturated, which is not the case. We note this because the sentence may otherwise be reproduced uncritically.

\subsubsection{The MAX30100}

\begin{table}[htbp]
\centering\small
\caption{MAX30100 characteristics as reported by \textcite{efendi2025}.}\label{tab:max30100}
\begin{tabularx}{\textwidth}{@{}>{\raggedright\arraybackslash}p{4.0cm} X@{}}
\toprule
\textbf{Characteristic} & \textbf{Value} \\
\midrule
Function & heart rate and oxygen saturation \\
Principle & reflective dual-wavelength photoplethysmography \\
Wavelengths & 660~nm (red) and 880~nm (infrared) \\
Supply voltage & 3.3~V to 5~V, optimal \\
Interface & I\textsuperscript{2}C (SCL clock line, SDA data line) \\
Pin count & 7: VIN, SCL, SDA, INT, IRD, RD, GND \\
INT pin & interrupt, pulse detection \\
IRD pin & infrared LED indicator \\
\bottomrule
\end{tabularx}
\end{table}

\subsubsection{The MAX30205}

\begin{table}[htbp]
\centering\small
\caption{MAX30205 characteristics as reported by \textcite{efendi2025}.}\label{tab:max30205}
\begin{tabularx}{\textwidth}{@{}>{\raggedright\arraybackslash}p{4.0cm} X@{}}
\toprule
\textbf{Characteristic} & \textbf{Value} \\
\midrule
Function & human body temperature \\
Measurement range & up to 50~\textdegree C \\
Accuracy & up to 0.1~\textdegree C \\
Supply voltage & 2.7~V to 3.3~V \\
Operating current & 600~\textmu A, optimal \\
Converter & high-resolution sigma--delta analogue-to-digital converter \\
Interface & I\textsuperscript{2}C-compatible two-wire serial \\
Pin count & 8: VCC, GND, SDA, SCL, OS, A0, A1, A2 \\
A0--A2 & address pins \\
Additional functions & over-temperature alarm, interrupt, shutdown \\
Command set & write byte, read byte, send byte, receive byte \\
\bottomrule
\end{tabularx}
\end{table}

The 0.1~\textdegree C accuracy figure should be read against the clinical requirement. The reference table of \textcite{efendi2025} gives a normal range of 36.4--37.6~\textdegree C for the 50--65 age band and 35.8--36.9~\textdegree C above 66. A resolution of 0.1~\textdegree C is therefore adequate to place a reading within a range spanning 1.1 to 1.2~\textdegree C, but it offers little margin: a reading of 36.9~\textdegree C in an older patient sits exactly on the boundary.%
\verify{The reference ranges of \textcite{efendi2025} are internally inconsistent: their Table~1 and Table~9 disagree on normal heart rate by age, giving 85--119 and 85--145~BPM respectively at age 50. Use this source for instrumentation specifications only, and take physiological reference ranges from a clinical source.}

\subsubsection{The ESP32}

\begin{table}[htbp]
\centering\small
\caption{ESP32 characteristics as reported by \textcite{efendi2025}.}\label{tab:esp32}
\begin{tabularx}{\textwidth}{@{}>{\raggedright\arraybackslash}p{4.0cm} X@{}}
\toprule
\textbf{Characteristic} & \textbf{Value} \\
\midrule
Processor & Tensilica Xtensa \\
ROM & 448~kB \\
SRAM & 520~kB \\
Sleep-mode memory & RTC fast SRAM and RTC slow SRAM, 8~kB each \\
Pin count & 36, including EN, GND, VIN, 3.3~V, GPIO \\
Connectivity & integrated Bluetooth and Wi-Fi \\
Security standard & IEEE 802.11, supporting WAPI, WFA, WPA/WPA2 \\
\bottomrule
\end{tabularx}
\end{table}

Two of these specifications bear directly on deployment in the target context. The sleep-mode memory determines the ability to retain state through low-power intervals, which matters where mains electricity is intermittent---and \cref{sec:regional-landscape} reported that approximately 550 million people in sub-Saharan Africa lack reliable access to electricity and the internet. The integrated Wi-Fi with WPA/WPA2 support determines whether transmission can be encrypted at the link layer, which bears on the legal and regulatory barrier that \textcite{dodoo2022} identified as one of six (\cref{sec:digital-barriers}), and on the data-protection requirement carried forward from \cref{sec:spec13}.

%==========================================================================
\subsection{A documented reference implementation}\label{sec:reference-impl}

\subsubsection{The three-layer architecture}

\Textcite{efendi2025} structure their system into the three layers conventional in IoT construction (\cref{fig:iot-layers}).

\begin{figure}[htbp]
\centering
\begin{tikzpicture}[node distance=0.9cm,
  lay/.style={draw,rounded corners,align=center,font=\footnotesize,
              minimum width=11cm,minimum height=1.25cm},
  ar/.style={-{Stealth[length=2.5mm]},very thick}]
  \node[lay,fill=cGrey!14] (phy) {\textbf{Physical layer} \quad MAX30100, MAX30205 \quad---\quad heart rate, SpO$_2$, temperature};
  \node[lay,fill=cRef!14,above=of phy] (net) {\textbf{Network layer} \quad ESP32 with Wi-Fi, Firebase \quad---\quad transmission, storage};
  \node[lay,fill=cAfrica!12,above=of net] (app) {\textbf{Application layer} \quad mobile application, website \quad---\quad presentation, classification};
  \draw[ar] (phy) -- node[right,font=\footnotesize]{MQTT (publish--subscribe, low bandwidth)} (net);
  \draw[ar] (net) -- node[right,font=\footnotesize]{HTTP (request--response)} (app);
\end{tikzpicture}
\caption{The three-layer architecture of \textcite{efendi2025}, with the two transport protocols used in sequence.}\label{fig:iot-layers}
\end{figure}

Transmission uses two protocols in sequence. The sensor-to-cloud path uses MQTT, a lightweight messaging protocol that the authors select because it \bq{enables resource-constrained devices to send sensor data to a gateway which acts as a translator, then processes and stores the data with Firebase functionalities such as Cloud Functions or Real-time Database}. The cloud-to-client path uses HTTP: \bq{the data which has been stored to Firebase is accessed by mobile applications and websites using HTTP protocol}.

The choice of MQTT upstream is appropriate and worth retaining. It is a publish--subscribe protocol designed for low bandwidth and unreliable links, and it is markedly more economical than HTTP for frequent small payloads---which is the exact profile of sensor telemetry over an intermittent mobile connection.

Firebase is selected for four properties the authors enumerate: real-time database services synchronising data instantaneously across clients and platforms; built-in authentication and authorisation ensuring secure access control; transmission in JSON; and persistent WebSocket connections, so that \bq{whenever current sensor data arrives, Firebase sends it to linked devices}.

The fourth property is the one that matters most for the target context: offline operation. The authors state that data are \bq{stored locally and real-time events continue to occur even when the device is offline}, and that \bq{the real-time data will be synchronised between local data change remotely when the device is offline and auto merging when the device is going to be online}. Given that the framing document identifies intermittent connectivity in the Cameroonian context as a source of irregular and missing measurements, a storage layer with native offline-first behaviour and automatic merge on reconnection addresses a stated constraint directly.

\subsubsection{Device validation}\label{sec:device-validation}

The authors validate the instrumentation against reference instruments before using it. The protocol is as follows.

\textbf{Participants.} Six individuals aged 50 to 65. Two of them---aged 63 and 65---have a chronic disease and are classified as not stable.

\textbf{Procedure.} Data collected every second for ten minutes in a seated resting position, yielding 600 records per participant. Readings from the monitoring device are compared against measurement tools, namely a digital oximeter and a digital thermometer. An interview established each participant's health history, and written consent was obtained.

\textbf{Metric.} Mean absolute percentage error, with a predefined acceptance threshold of 3\%:
\begin{equation}
  \mathrm{MAPE} = \frac{100}{n}\sum_{k=1}^{n}\left|\frac{A_k - P_k}{A_k}\right|.
\end{equation}

\begin{table}[htbp]
\centering\small
\caption{Validation results reported by \textcite{efendi2025}.}\label{tab:efendi-validation}
\begin{tabularx}{\textwidth}{@{}X c c c c@{}}
\toprule
\textbf{Parameter} & \textbf{Device} & \textbf{Reference tools} & \textbf{Gap} & \textbf{Percentage error} \\
\midrule
Heart rate (BPM) & 84.37 & 85.79 & 1.42 & 1.68\% \\
Oxygen saturation (\%) & 96.85 & 96.30 & 0.55 & 0.57\% \\
Body temperature (\textdegree C) & 36.53 & 36.69 & 0.16 & 0.44\% \\
\midrule
Overall MAPE & & & & 0.90\% \\
Accuracy ($100\% - \mathrm{MAPE}$) & & & & 99.10\% \\
\bottomrule
\end{tabularx}
\end{table}

The experimental protocol also defines four real-world scenarios---sleeping, walking, sitting, cardio exercising---though the validation results reported are from the seated condition.

\subsubsection{Critical appraisal}\label{sec:efendi-appraisal}

Five reservations must accompany any citation of these figures.

\textbf{First, the denominator of the error is the device, not the reference.} The reported percentage errors can be reproduced only by dividing the gap by the device reading. For heart rate, $1.42/84.37 = 1.68\%$, which is the published value, whereas $1.42/85.79 = 1.66\%$. The quantity written $A_k$ in the formula---conventionally the actual or reference value---is therefore the device output, and the comparator instrument occupies the position of the prediction. The study thus measures how far the reference deviates from the device. For the other two parameters the two conventions round to the same figure, so heart rate is the only row that discriminates; but the inversion is systematic and it reverses the logic of a validation. The overall figure is the unweighted mean of the three, $(1.68+0.57+0.44)/3 = 0.90\%$, computed on the parameter means rather than on paired observations, which removes dispersion before the statistic is formed.%
\verify{Confirm from the Methods of \textcite{efendi2025} whether MAPE was computed per observation and then averaged, or on the aggregate means as the published figures suggest. The two give different values and only the first is a meaningful agreement statistic.}

\textbf{Second, agreement with an unvalidated reference is not accuracy.} The comparator instruments are a digital oximeter and a digital thermometer, neither identified by make or model, and neither reported as independently validated. \Cref{sec:metrology} established, on WHO authority, that only about one device in five on the market carries evidence of accuracy validation. What the study demonstrates is concordance between two consumer devices, which is a necessary but not a sufficient condition for accuracy.

\textbf{Third, MAPE is not an appropriate agreement statistic for this purpose.} It averages the absolute relative difference and therefore conceals the structure of the disagreement: a device reading systematically 1.4~BPM high and a device oscillating randomly by $\pm 1.4$~BPM produce the same MAPE but have entirely different clinical implications. The standard instrument for method comparison is Bland--Altman analysis, which reports bias and limits of agreement separately; it is not used here. The authors' own figures reveal what the summary statistic hides: they note \bq{a fairly wide gap between the two devices, especially at time 300 to 360} for heart rate, and \bq{a significant gap in the 340 and 380 time range} for oxygen saturation. These transient divergences disappear entirely in a 0.90\% average.

\textbf{Fourth, the validation conditions are the most favourable possible.} Six participants, ten minutes, seated at rest. \Textcite{navarro2023} notes that oscillometric and optical methods are \bq{not reliable during physical activity due to movement noise}. The four scenarios including walking and cardio exercise are described in the methods but no comparative error figures are reported for them.

\textbf{Fifth, and decisively for the present work, neither blood pressure nor blood glucose is measured.} The validated parameters are heart rate, oxygen saturation and temperature. This architecture, however well executed, does not acquire the two variables on which the clinical problem of this thesis rests.

\subsubsection{The machine learning component, and a structural flaw worth studying}\label{sec:leakage}

\Textcite{efendi2025} add a supervised classifier predicting whether the user's condition is stable or not stable.

\textbf{Dataset.} 4\,500 records, combining participant data with data from a public Kaggle repository. Three features---heart rate, oxygen saturation, body temperature---and a binary label. The classes are balanced exactly: 2\,250 records each.

\textbf{Preprocessing.} The authors apply a transformation in which each feature is recoded as 0, 1 or 2 to represent low, normal and high conditions, with reference to their Tables~9 to~11, and the label is recoded as 1 for stable and 0 for not stable.

\textbf{Split and training.} 70\% training (3\,150), 10\% validation (450), 20\% testing (900). Five models compared: AdaBoost, XGBoost, decision tree, random forest and a Bayesian classifier. Configuration: 300 epochs, learning rate 0.025, sigmoid activation, Adam optimiser, on an Intel Core i7-4770 at 3.40~GHz with 8~GB of memory under Python in a Jupyter notebook.

\textbf{Results.} All five models exceed 0.9 accuracy. XGBoost is selected with accuracy 0.973 and $F_1$ 0.970; AdaBoost follows at 0.967 accuracy and 0.946 $F_1$; the decision tree, random forest and Bayesian classifiers reach $F_1$ scores of 0.933, 0.929 and 0.928. Loss values are below 0.1 throughout, with XGBoost lowest at 0.025 and the Bayesian classifier highest at 0.062.

\textbf{The flaw.} The label \emph{stable} is defined as the condition in which all three parameters fall within the normal ranges of Tables~9 to~11. The feature transformation recodes each parameter as low, normal or high using those same ranges. After transformation, therefore, the label is a deterministic function of the three transformed features:
\begin{equation}
  y = 1 \iff (x_1 = 1) \wedge (x_2 = 1) \wedge (x_3 = 1).
\end{equation}

\begin{figure}[htbp]
\centering
\begin{tikzpicture}[node distance=0.6cm and 1.3cm,
  b/.style={draw,rounded corners,align=center,font=\footnotesize,minimum height=0.95cm,minimum width=2.9cm},
  ar/.style={-{Stealth[length=2.5mm]},thick}]
  \node[b,fill=cGrey!14] (raw) {Raw values\\ HR, SpO$_2$, temp.};
  \node[b,fill=cAfrica!16,right=of raw] (thr) {Normal ranges\\ (Tables 9--11)};
  \node[b,fill=cRef!14,above right=0.4cm and 1.3cm of thr] (feat) {Features $x_1,x_2,x_3$\\ recoded 0/1/2};
  \node[b,fill=cRef!14,below right=0.4cm and 1.3cm of thr] (lab) {Label $y$\\ stable / not stable};
  \node[b,fill=cGrey!10,right=4.6cm of raw,yshift=-2.6cm] (mod) {Classifier};
  \draw[ar] (raw) -- (thr);
  \draw[ar] (thr) -- (feat);
  \draw[ar] (thr) -- (lab);
  \draw[ar] (feat.east) -| ([xshift=1.0cm]feat.east) |- (mod.east);
  \draw[ar] (lab) -- (mod);
  \node[font=\footnotesize,text=cAfrica,align=center] at ([yshift=-1.5cm]thr.south)
    {the same rule generates\\ both inputs and target};
\end{tikzpicture}
\caption{The leakage mechanism in \textcite{efendi2025}. The normal ranges that define the label also define the feature encoding, so the classifier recovers a rule already present in its inputs.}\label{fig:leakage}
\end{figure}

This is label leakage (\cref{fig:leakage}). The model is not learning a clinical relationship; it is recovering the rule that generated the labels, from features into which that rule has already been encoded. A decision tree of depth three should achieve 100\% accuracy on this task. That XGBoost reaches 97.3\% rather than 100\% indicates inconsistency somewhere in the labelling or the transformation, not predictive difficulty.

Two further features of the setup confirm the diagnosis. The perfect class balance of 2\,250 against 2\,250 does not arise in clinical data, where, as the framing document notes for the present work, critical states are rare; it arises from construction. And the comparison of five classifiers is uninformative when all five are solving a lookup problem: the 0.045 spread in $F_1$ between best and worst measures nothing of consequence.

This matters for the present thesis well beyond the appraisal of one paper. It is the precise error that must not be repeated. The framing document states that the risk label---stable, moderate, high, critical---is derived from recognised clinical thresholds, ESH 2023 for blood pressure and ADA standards for blood glucose, and validated by the partner physician. If the model's features then include the current values of blood pressure and blood glucose, the same circularity arises: the model will reproduce the threshold rule and report a high score that demonstrates nothing.

Two of the framing document's own commitments already point the way out, and should be stated explicitly as the defence against this failure.

The first is the anticipatory dual-task formulation: forecasting physiological degradation at 24 to 48 hours rather than classifying the current state. A model that predicts a future threshold crossing from a past window cannot leak its label, because the label does not exist in the input window. This is not merely a more ambitious objective; it is the formulation that makes the problem genuine. It is also, in different words, the distinction drawn in \cref{sec:records} between the act and the series: a classifier of the current reading is an instrument of the act, a forecaster over a window is an instrument of the series.

The second is the feature engineering the framing document describes---rolling aggregates, trends, variability, deviations from thresholds. Of these, deviations from thresholds computed at the same instant as the label is precisely the leaking feature, while the same quantity computed on a lagged window is legitimate. The distinction must be made explicit in the methodology and enforced in code.

%==========================================================================
\subsection{Low-cost alternatives and a contradiction to resolve}\label{sec:lowcost}

\subsubsection{What the framing document proposes}

The framing document states the acquisition design as a deliberately spare collection device---manual entry and phone-camera photoplethysmography---adapted to infrastructure constraints, which makes the approach deployable where expensive medical equipment is unavailable. Its technological limitations section adds that this low-cost acquisition device is less precise than dedicated medical equipment and introduces measurement noise.

\subsubsection{What the corpus says about camera-based measurement}

The sources assembled across \cref{ch:background} and this section converge, and they converge against part of this proposal.

\Textcite{who2023hypertension} states that cuffless devices, including popular connected wristbands, are not currently recommended for hypertension management, because of unresolved questions regarding accuracy, performance and ability to track blood pressure change, and that as of 2023 no alternative to the validated upper-arm cuff exists with sufficient consistency and accuracy for diagnosis and management (\cref{sec:metrology}).

\Textcite{alessa2019} found three applications out of 186 claiming to turn the smartphone into a blood pressure measuring device, presumably by a cuffless technique placing a finger over the camera, and report that none provided any evidence of reliability or validity. In discussion they state that such smartphone-based cuffless devices \bq{have recently been shown to be highly inaccurate and unfeasible and may negatively affect patients' health and safety}, adding that users nonetheless show a strong inclination to download and favourably rate them (\cref{sec:cuffless}).

\Textcite{navarro2023} notes that oscillometry \bq{presents limitations in accuracy, and it is not reliable during physical activity due to movement noise} and \bq{requires validation for specific populations}, and that applanation tonometry \bq{requires individual calibration and is not suitable for routine clinical settings}.

No source in this corpus provides any validation evidence for camera-based photoplethysmography.

\subsubsection{Resolving the contradiction}\label{sec:contradiction}

The contradiction is real but narrower than it first appears, and it dissolves once the parameter is specified.

Photoplethysmography measures pulsatile blood volume. From that signal, heart rate is recovered from periodicity alone, requiring no absolute calibration---this is the robust case (\cref{sec:ppg}). Oxygen saturation requires two wavelengths and an empirical calibration curve; a smartphone camera with white-light illumination does not provide the controlled dual-wavelength illumination the method assumes, and any value derived from it is of unknown validity. Blood pressure is not measured by photoplethysmography at all: it is inferred, typically from pulse wave characteristics, through a model requiring individual calibration---which is exactly what \textcite{navarro2023} says of applanation tonometry and what the WHO says is not yet acceptable.

The design is therefore defensible if and only if it is stated with the precision of \cref{tab:param-sources}.

\begin{table}[htbp]
\centering\small
\caption{Admissible source for each parameter of the proposed system.}\label{tab:param-sources}
\begin{tabularx}{\textwidth}{@{}>{\raggedright\arraybackslash}p{3.6cm} >{\raggedright\arraybackslash}X >{\raggedright\arraybackslash}p{4.2cm}@{}}
\toprule
\textbf{Parameter} & \textbf{Acceptable source} & \textbf{Not acceptable} \\
\midrule
Systolic and diastolic pressure & validated automated upper-arm cuff; manual entry or Bluetooth & camera PPG, cuffless wristband, any inferential method \\
Blood glucose & glucometer; manual entry or Bluetooth & any optical method \\
Heart rate & camera PPG, cuff-derived, or optical sensor & --- \\
Oxygen saturation & validated pulse oximeter & camera PPG \\
Temperature & digital thermometer or MAX30205 & --- \\
\bottomrule
\end{tabularx}
\end{table}

Read this way, the low-cost acquisition contribution survives intact. The device collects what low-cost optics can legitimately collect---heart rate, a contextual covariate---while blood pressure and blood glucose enter through manual entry from validated instruments, which is the pipeline the framing document's own record structure implies and which carries no inferential claim.

What must not appear in this thesis is any claim that blood pressure is obtained by camera. It would contradict the WHO, \textcite{alessa2019} and the physics, and it is the single assertion most likely to be challenged at the defence.

%==========================================================================
\subsection{The device landscape, and what it omits}\label{sec:device-landscape}

\Textcite{liu2025review} survey the categories of IoT mobile sensing devices in chronic disease management. Their taxonomy is useful for situating the proposed system, subject to the reservation on their method stated at the outset (\cref{tab:device-categories}).

\begin{table}[htbp]
\centering\footnotesize
\caption{Categories of IoT mobile sensing devices reported by \textcite{liu2025review}.}\label{tab:device-categories}
\begin{tabularx}{\textwidth}{@{}>{\raggedright\arraybackslash}p{3.5cm} X@{}}
\toprule
\textbf{Category} & \textbf{Description and examples reported} \\
\midrule
Portable medical devices & battery-powered, handheld, compact: portable scanners collecting ECG, temperature, oxygen saturation, electromyography, respiratory rate; portable magnetic resonance scanners \\
Wearable devices & smartwatches and health tracking bands recording heart rate, blood pressure and blood glucose in real time, with analysis and health recommendations \\
Intelligent clinical devices & mobile nursing management, infusion monitoring, smart wards, bedside intelligent interaction \\
Remote health devices & wireless monitoring platforms for remote dynamic blood pressure monitoring and remote wireless health check management \\
Implantable devices & pacemakers, implantable cardioverter-defibrillators \\
Ingestible devices & capsule endoscopes, remotely controlled drug delivery \\
\bottomrule
\end{tabularx}
\end{table}

Two entries bear directly on this work.

They report that Vybornova and colleagues developed an IoT-based non-invasive cuffless blood pressure monitor measuring optical photoplethysmographic pulse wave signals on the wrist, calculating systolic and diastolic values using pulse wave analysis. This is cited without any accuracy figures and without reference to validation status. Given the WHO position, it should be cited only alongside that position, and never as evidence that cuffless measurement is ready for clinical use.

They report the smart home system of \textcite{chatrati2022}, analysing home blood pressure and glucose indices, which \cref{sec:chatrati} examined in detail and found to rest on manual entry, conjunctive rules with coverage gaps, a non-transferable training set and an unjustified alert mapping.

They also report, in their statement of the problem, that traditional chronic disease monitoring predominantly relies on manual processes, suffers from the phenomenon of information silos, and lacks real-time monitoring and feedback and effective prediction and early warning mechanisms---the diagnosis carried forward from \cref{sec:records}.

\textbf{What the landscape omits.} Across both this survey and the architecture of \textcite{efendi2025}, no device category supplies clinically validated blood pressure or blood glucose without a cuff and without a finger-stick respectively. The entire instrumented pipeline, however elaborate, measures around the problem rather than into it. This is not a gap in the literature; it is a physical constraint, and the architecture of the proposed system must accommodate it rather than attempt to engineer past it.

%==========================================================================
\subsection{Specification and the decision this section forces}\label{sec:spec21}

\subsubsection{Acquisition specification}

\begin{table}[htbp]
\centering\footnotesize
\caption{Acquisition requirements established in \cref{sec:acquisition} and their basis.}\label{tab:spec21}
\begin{tabularx}{\textwidth}{@{}X >{\raggedright\arraybackslash}p{5.9cm}@{}}
\toprule
\textbf{Requirement} & \textbf{Basis} \\
\midrule
Three pipelines modelled separately, with distinct missingness mechanisms & \cref{sec:pipelines}; MNAR in pipelines 1 and 2 \\
Blood pressure from a validated automated upper-arm cuff only & \textcite{who2023hypertension}; \textcite{alessa2019} \\
Blood glucose from a glucometer only & \cref{sec:device-landscape}; no optical alternative exists \\
Camera PPG restricted to heart rate; never blood pressure or SpO$_2$ & \textcite{who2023hypertension}; \textcite{alessa2019}; \textcite{navarro2023} \\
Timestamp of measurement stored distinctly from timestamp of entry & \cref{sec:pipelines}; fictitious temporal structure otherwise \\
MQTT upstream, HTTP downstream & \textcite{efendi2025} \\
Offline-first storage with automatic merge on reconnection & \textcite{efendi2025}; intermittent connectivity stated in \textcite{cadrage2026} \\
Link-layer encryption (WPA/WPA2) and authenticated cloud access & \textcite{efendi2025}; regulatory barrier in \textcite{dodoo2022} \\
Method comparison by Bland--Altman, not MAPE alone & \cref{sec:efendi-appraisal}; MAPE conceals bias structure \\
Error computed with the reference instrument in the denominator & \cref{sec:efendi-appraisal}; inverted in \textcite{efendi2025} \\
Sensor validation under movement, not seated rest only & \cref{sec:efendi-appraisal}; \textcite{navarro2023} on movement noise \\
Lagged features only; no threshold-derived feature at label time & \cref{sec:leakage}; label leakage in \textcite{efendi2025} \\
Reference instrument for validation must itself be validated & \cref{sec:metrology}; one device in five carries validation evidence \\
\bottomrule
\end{tabularx}
\end{table}

\subsubsection{The acquisition decision}\label{sec:acq-decision}

The plan for this chapter describes an instrumented pipeline built on the MAX30100, MAX30205 and ESP32. The framing document describes a deliberately spare device using manual entry and camera photoplethysmography. These are not variants of one design; they are two different systems with different costs, different failure modes and different claims. Three positions are coherent, and this thesis must adopt one explicitly.

\textbf{Position A, spare acquisition.} No microcontroller hardware. Blood pressure and blood glucose by manual entry from validated instruments; heart rate optionally by camera photoplethysmography. The IoMT hardware pipeline is then presented as state of the art examined and deliberately not adopted, with the reasons given: cost, maintenance, the sustainability failure documented in \cref{sec:regional-landscape}, and the fact that the sensors do not measure the parameters of interest.

\textbf{Position B, instrumented acquisition.} An ESP32-based device is built, measuring heart rate, oxygen saturation and temperature as covariates, alongside manual entry of blood pressure and blood glucose. This adds a hardware development workload and a device-validation obligation---Bland--Altman against a validated reference, under movement---and it contradicts the framing document as currently written, which would require amendment.

\textbf{Position C, hybrid and staged.} The system is specified to accept the sensor pipeline, with the data model and transport layer provided for, but the evaluated prototype uses spare acquisition. The hardware is then documented as the extension path.

\textbf{The position adopted is A.} It is the one consistent with the deployability argument that runs through \cref{ch:background}: the instrumented pipeline acquires none of the parameters on which the clinical decision rests (\cref{tab:pipelines}), it carries a validation obligation that this work is not equipped to discharge, and the sustainability failures documented in \cref{sec:regional-landscape} fall hardest on designs that require maintained devices in the field. The hardware examined in \cref{sec:hardware,sec:reference-impl} is therefore reported as state of the art and deliberately not adopted, and the transport and storage choices retained from \textcite{efendi2025}---MQTT upstream, offline-first storage with automatic merge---are retained independently of the sensors, because they address connectivity rather than instrumentation.

% ---------------------------------------------------------------
% TO SWITCH TO POSITION C, replace the paragraph above with:
% \textbf{The position adopted is C.} The data model and transport layer are
% specified to accept the sensor pipeline described in \cref{sec:reference-impl},
% and the evaluated prototype uses spare acquisition. The hardware is documented
% as the extension path, and the device-validation obligation set out in
% \cref{tab:spec21} applies to any future version that activates it.
% ---------------------------------------------------------------

This decision determines the treatment of the following section: the gateway and transport layers are specified rather than implemented, and \cref{sec:cloud} examines them on that basis.

%##########################################################################
% SECTION 2.2
%##########################################################################
\section{Cloud Infrastructure, Real-Time Transmission and Storage}\label{sec:cloud}

\subsection*{Purpose and status of the sources}

\Cref{sec:acquisition} established what is acquired and by what means. This section addresses what happens to the data between the point of measurement and the point of clinical decision: how it is layered, transported, stored and pushed. The governing constraint, carried from \cref{ch:background}, is that the system must triage---it must determine which readings warrant clinical attention---and this section shows that the architectural literature already provides the place where that function belongs.

\paragraph{Mucchi et al. (2021).} \textcite{mucchi2021}, published in the \emph{International Journal of Environmental Research and Public Health} 18:7459, presents an information and communication technology reference architecture for the telemonitoring of chronic diseases, developed jointly by the Department of Information Engineering of the University of Florence and the Nursing and Midwifery Department of the Local Health Unit Toscana Centro, and based on the high-level architecture defined by the ETSI SmartBAN Working Group, of which the authors are members. It reports a preliminary usability evaluation. It is the most complete architectural specification in this corpus.

\paragraph{Efendi et al. (2025).} \textcite{efendi2025}, examined in \cref{sec:reference-impl}, is the only source reporting a working implementation with measured device performance.

\paragraph{Liu and Wang (2025).} \textcite{liu2025review} supply a seven-layer platform diagram for chronic disease monitoring with machine learning. The reservation stated in \cref{sec:diagnosis} stands: the paper declares a systematic review protocol it does not report, and we treat it as a narrative survey.

\paragraph{Yazdani et al. (2018).} \textcite{yazdani2018}, published in \emph{Acta Informatica Medica} 26(3):195--200, proposes a scalable architecture for telemonitoring chronic diseases supporting multiple clinical decision support systems on a common platform. It is a conceptual design expressed in UML, with no implementation and no evaluation, derived from a literature review of 1\,326 articles reduced to 73 full texts covering 2005 to 2018. Its technical choices---XML for data exchange, SMS reminders, client--server without a cloud layer---are dated relative to the others, and it is cited here for its modular decomposition and its scalability argument, not as a reference for cloud infrastructure.

%==========================================================================
\subsection{The multi-tier model}\label{sec:multitier}

\subsubsection{The four-layer reference architecture}

\Textcite{mucchi2021} define a reference communication network architecture in four layers. We reproduce their definitions precisely, because the terminology of this chapter derives from them (\cref{fig:four-layer}).

\textbf{BODY layer.} A set of body sensors and devices gathering the patient's parameters. The authors distinguish in-body devices---tiny sensors that are implanted or swallowed---from on-body devices worn by the patient for bio-parameter measurement, and they name the examples explicitly: blood pressure, pulse oximeter, ECG sensors.

\textbf{ENVIRONMENT layer.} Indoor and outdoor environmental sensors for information gathering---temperature, air quality, humidity---and user positioning, by satellite positioning or video monitoring, allowing location-based and personalised services based on the user's lifestyle.

\textbf{EDGE layer.} Responsible for short-term analysis of the acquired data and for user--system interactions, typically through a smartphone application. In the authors' words, it \bq{allows the quick processing and identification of anomalous patient's conditions and the consequent warning notification}. Structurally it comprises a gateway, a wide-area network device containing a reception and transmission unit, and a hub containing a data collection unit, a data processing unit, a local storage unit and a transmission unit.

\textbf{CLOUD layer.} In charge of the eHealth service and of long-term data analysis through the adoption of advanced algorithms for data cross-correlation. Its units are a cloud repository unit, a control unit, a cloud data processing unit, eHealth services, data analysis and alert management.

The division of labour between the last two layers is the single most useful idea in this section, and it is developed in \cref{sec:edge-triage}.

\begin{figure}[htbp]
\centering
\begin{tikzpicture}[node distance=0.75cm,
  lay/.style={draw,rounded corners,align=center,font=\footnotesize,
              minimum width=10.5cm,minimum height=1.35cm},
  ar/.style={-{Stealth[length=2.5mm]},very thick},
  tag/.style={font=\scriptsize,text=cGrey!80!black}]
  \node[lay,fill=cGrey!14] (body) {\textbf{BODY} \quad validated cuff, glucometer, optional optical sensor\\ \scriptsize in-body devices not applicable to this system};
  \node[lay,fill=cGrey!8,above=0.45cm of body] (env) {\textbf{ENVIRONMENT} \quad \emph{not instantiated in this system}};
  \node[lay,fill=cRef!14,above=of env] (edge) {\textbf{EDGE} \quad gateway + hub: collection, \emph{processing}, \emph{local storage}, transmission\\ \scriptsize short-term analysis, anomaly detection, local pre-alert};
  \node[lay,fill=cAfrica!12,above=of edge] (cloud) {\textbf{CLOUD} \quad repository, control, processing, eHealth services,\\ data analysis, \emph{alert management} \quad---\quad long-term analysis};
  \node[lay,fill=cGrey!20,above=of cloud] (clin) {\textbf{CLINICIAN} \quad priority-ordered queue};
  \draw[ar] (body) -- (env);
  \draw[ar] (env) -- node[right,tag]{Inter-ComLayer-1 / 2} (edge);
  \draw[ar] (edge) -- node[right,tag]{Inter-ComLayer-3: MQTT over 3G/4G} (cloud);
  \draw[ar] (cloud) -- node[right,tag]{HTTP/REST, WebSocket push} (clin);
\end{tikzpicture}
\caption{The four-layer reference architecture of \textcite{mucchi2021}, instantiated for the present system. The environment layer is retained in the model and not instantiated, since the system has no environmental sensors. Italicised units are the ones that carry the triage function.}\label{fig:four-layer}
\end{figure}

\subsubsection{The three implementation layers of Efendi et al.}

\begin{table}[htbp]
\centering\small
\caption{Implementation layers of \textcite{efendi2025}.}\label{tab:efendi-layers}
\begin{tabularx}{\textwidth}{@{}>{\raggedright\arraybackslash}p{2.6cm} >{\raggedright\arraybackslash}p{5.2cm} X@{}}
\toprule
\textbf{Layer} & \textbf{Components} & \textbf{Function} \\
\midrule
Physical & MAX30100, MAX30205 & sensing \\
Network & ESP32 with Wi-Fi, Firebase & transmission, storage \\
Application & mobile application, website, classifier & presentation, prediction \\
\bottomrule
\end{tabularx}
\end{table}

This is the same stack with the edge layer collapsed into the network layer and the environment layer absent. The compression is consistent with the implementation: the ESP32 performs no local analysis, it transmits.

\subsubsection{The seven-layer platform of Liu and Wang}

\begin{table}[htbp]
\centering\footnotesize
\caption{Seven-layer platform reported by \textcite{liu2025review}.}\label{tab:liu-layers}
\begin{tabularx}{\textwidth}{@{}>{\raggedright\arraybackslash}p{2.8cm} X@{}}
\toprule
\textbf{Layer} & \textbf{Contents reported} \\
\midrule
Access & connectivity to IoT devices, mobile applications and other data sources; blood pressure monitor, blood glucose meter, smartwatches \\
Data transfer & wireless communication: Bluetooth, Wi-Fi, 4G, 5G \\
Data storage & web (Java, Tomcat, SpringCloud); big data (Hadoop, Spark); database (MySQL, Redis cache); DevOps continuous integration and delivery; Linux server side \\
Data processing & data cleansing, data conversion, feature engineering, data extraction, data integration \\
Data analysis & support vector machine, decision tree, $k$-means, gradient boosting, AdaBoost, neural networks \\
Application service & master data services (remote monitoring, intelligent follow-up, online consultation); business applications (medication reminders, user feedback, research support) \\
Decision support & multi-source data fusion, feature learning, personalised medicine, interactive analysis, integration of interdisciplinary knowledge \\
\bottomrule
\end{tabularx}
\end{table}

Two cross-cutting systems span all seven: a standard specification support system and an information security assurance system. Between the data analysis and application service layers, the authors place a multi-point triggering and dynamic response mechanism.

The explicit presence of a blood pressure monitor and a blood glucose meter at the access layer confirms the device allocation established in \cref{tab:param-sources}: these parameters enter the system as instrument readings, not as inferences from optical sensors.

\subsubsection{The three functional layers of Yazdani et al.}

Their architecture has a user base layer, where physicians, the healthcare team, patients and system administrators connect to the server through a client interface installable on a tablet or phone or presented as a web page; a function base layer containing the architectural modules; and a database layer storing physiological data alongside the physician's prescriptions, diagnoses and the recommendations of the decision support system.

The function base comprises five modules, and this decomposition is the paper's contribution (\cref{tab:yazdani-modules}).

\begin{table}[htbp]
\centering\footnotesize
\caption{Modules of the function base layer \parencite{yazdani2018}.}\label{tab:yazdani-modules}
\begin{tabularx}{\textwidth}{@{}>{\raggedright\arraybackslash}p{3.3cm} X@{}}
\toprule
\textbf{Module} & \textbf{Function} \\
\midrule
Membership management & registration of patient, physician and healthcare team; definition of the checkup period, the intervals at which patient data must be collected and sent; rule selection, choosing the decision support system and rules appropriate to the disease and symptoms \\
Communication transmission & ensuring data aggregation tools are attached and data are sent \\
Supervisor & determining which decision support system corresponds to the data received, by parsing metadata attached to the patient record \\
CDSS setting & defining the decision support systems and their rules; adding and updating them \\
Notification & communicating analysis results to medical staff and informing patients of prescriptions \\
\bottomrule
\end{tabularx}
\end{table}

The data flow is also specified: \bq{if no data is sent within the time limits set by the physician, an SMS reminder will be sent to the patient. The SMS reminder is sent repeatedly as long as the patient enters a new data.} The server verifies and stores the data, determines the related decision support system, analyses the patient data, and classifies the result as normal or alert. The physician can view patient status on a priority basis, this prioritisation being based on that status.

Two elements here are directly relevant and should be retained despite the paper's age. The checkup period as a configurable parameter operationalises the WHO follow-up requirement established in \cref{sec:followup-requires}---monthly until target is reached, three-monthly thereafter---as a system setting rather than a fixed constant. And the priority-ordered physician view is the user-facing form of the triage function.

\subsubsection{Comparative mapping}

\begin{table}[htbp]
\centering\footnotesize
\caption{Mapping of the four architectures onto the reference layers of \textcite{mucchi2021}.}\label{tab:layer-mapping}
\begin{tabularx}{\textwidth}{@{}>{\raggedright\arraybackslash}p{3.0cm} >{\raggedright\arraybackslash}X >{\raggedright\arraybackslash}X >{\raggedright\arraybackslash}X@{}}
\toprule
\textbf{\textcite{mucchi2021}} & \textbf{\textcite{efendi2025}} & \textbf{\textcite{liu2025review}} & \textbf{\textcite{yazdani2018}} \\
\midrule
BODY & physical layer & access layer & --- \\
ENVIRONMENT & --- & access layer & --- \\
EDGE & \emph{absent} & data transfer layer (transport only) & \emph{absent} \\
CLOUD & network layer & storage + processing + analysis & database layer \\
eHealth services (within cloud) & application layer & application service + decision support & user base + function base \\
Alert management & classifier in application layer & multi-point triggering & notification module \\
\bottomrule
\end{tabularx}
\end{table}

\subsubsection{What the comparison establishes}

The four architectures agree on sensing, transport, storage and presentation. They disagree on whether an edge layer exists, and that disagreement is the substantive finding of this subsection.

\Textcite{mucchi2021} give the edge layer a data processing unit and a local storage unit and assign it short-term analysis and anomaly warning. \Textcite{liu2025review} reduce their equivalent layer to transport only. \Textcite{efendi2025} and \textcite{yazdani2018} have no edge layer: the ESP32 transmits without processing, and the client interface of Yazdani et al.\ is a presentation surface.

The consequence is architectural, not cosmetic. In an edgeless design, every reading travels to the cloud and every decision is taken there. In an edge design, the local tier can filter, aggregate and raise a first-level warning without a round trip. Given that \cref{sec:regional-landscape} established approximately 550 million people in sub-Saharan Africa lacking reliable electricity and internet, and that \textcite{mucchi2021} themselves parameterise for \bq{very limited connectivity}, the edge design is not an optimisation but a precondition.

%==========================================================================
\subsection{Communication protocols}\label{sec:protocols-cloud}

\subsubsection{The communication layer taxonomy}

\Textcite{mucchi2021} decompose the links into three intra-layer and three inter-layer communications (\cref{tab:comm-links}).

\begin{table}[htbp]
\centering\footnotesize
\caption{Communication links defined by \textcite{mucchi2021}.}\label{tab:comm-links}
\begin{tabularx}{\textwidth}{@{}>{\raggedright\arraybackslash}p{3.0cm} X >{\raggedright\arraybackslash}p{4.1cm}@{}}
\toprule
\textbf{Link} & \textbf{Between} & \textbf{Technology} \\
\midrule
Intra-Com-1 & within the body: nano-devices, implants & molecular communications \\
Intra-Com-2 & in-body $\leftrightarrow$ on-body; sink nodes collect from in-body devices and deliver off-body & not specified \\
Intra-Com-3 & optional: hub $\leftrightarrow$ gateway, where the gateway is not integrated into the data collector & not specified \\
Inter-ComLayer-1 & BODY $\leftrightarrow$ EDGE: on-body sensors to hub or gateway, e.g.\ a smartphone & not specified \\
Inter-ComLayer-2 & ENVIRONMENT $\leftrightarrow$ EDGE & not specified \\
Inter-ComLayer-3 & EDGE $\leftrightarrow$ CLOUD & 3G/4G/5G, Wi-Fi, SigFox, low-power wide-area network \\
\bottomrule
\end{tabularx}
\end{table}

Only the last of these is named with specific technologies, and only the last is relevant to a system whose body layer reduces to a Bluetooth cuff and a glucometer. Intra-Com-1 and Intra-Com-2 concern implanted and ingestible devices and do not apply here; this thesis states so rather than reproducing a taxonomy whose upper half is inert.

\subsubsection{MQTT upstream}

\Textcite{efendi2025} select MQTT for the sensor-to-cloud path, stating that it \bq{enables resource-constrained devices to send sensor data to a gateway which acts as a translator, then processes and stores the data with Firebase functionalities such as Cloud Functions or Real-time Database}, and describing it as \bq{a messaging protocol which is commonly used in IoT devices}.

The rationale is sound. MQTT is a publish--subscribe protocol: a client publishes to a topic on a broker, and subscribers receive without the publisher knowing who they are. Its header overhead is minimal relative to HTTP, it maintains a single long-lived connection rather than opening one per message, and it defines quality-of-service levels permitting at-least-once or exactly-once delivery over an unreliable link.

A reservation must be recorded. No source in this corpus quantifies the advantage. \Textcite{efendi2025} assert that MQTT suits resource-constrained devices; they do not measure bandwidth, latency, power consumption or delivery reliability against HTTP, and they report no comparison. The choice is therefore justified by convention and by protocol design, not by evidence presented here. If this thesis claims an efficiency benefit, it must either cite a benchmark outside this corpus or measure it.

\subsubsection{HTTP and REST downstream}

\Textcite{efendi2025} route the cloud-to-client path over HTTP: \bq{the data which has been stored to Firebase is accessed by mobile applications and websites using HTTP protocol}. Firebase, they note, \bq{can support various protocols including HTTP and MQTT}, and \bq{data transmitted by IoT devices will be stored as JSON which is then accessed by users in real-time}.

The asymmetry is deliberate and should be preserved. Upstream traffic is high-frequency, small-payload, from a constrained device over an unreliable link---the MQTT profile. Downstream traffic is request-driven, from a capable device, and must interoperate with web clients---the HTTP profile.

\subsubsection{WebSockets and real-time push}

The real-time property does not come from HTTP but from a persistent connection beneath it. \Textcite{efendi2025} state that Firebase \bq{provides Realtime Database and designed for real-time updates. Whenever current sensor data arrives, Firebase sends it to linked devices via persistent WebSocket connections. This enables clients to examine the health vitals virtually immediately on mobile application and website.}

The distinction between polling and push matters for this system's economics as much as its responsiveness. A polling client issues requests whether or not data have changed; a pushed client consumes resources in proportion to events. For a cohort transmitting a handful of readings per week, polling would dominate the traffic entirely.

\subsubsection{The XML choice in Yazdani et al., and why it is not adopted}

\Textcite{yazdani2018} use XML for patient data exchange, and give four reasons: that the label elements of XML content are independent of communication protocols, making XML documents suitable for internet and public data transfer networks; that XML is advantageous for editor applications, middleware and application tools, improving efficiency and yielding financial savings in information systems; that users can send and request patient data, view status and receive prescriptions and notifications anywhere with mobile or Wi-Fi coverage; and that metadata added to the XML file allows the supervisor module to determine which decision support system should analyse the patient data, this being performed by an XML parser.

The fourth reason is architecturally interesting---self-describing payloads enabling routing without a central registry---but the format choice is dated. The three later sources all use JSON, which is more compact, natively supported by the storage layer and by every mobile runtime, and does not require a parser stage. The idea to retain from \textcite{yazdani2018} is metadata-driven routing, not XML.

\subsubsection{Protocol selection}

\begin{table}[htbp]
\centering\small
\caption{Protocol selected for each link, and the constraint it answers.}\label{tab:protocol-selection}
\begin{tabularx}{\textwidth}{@{}>{\raggedright\arraybackslash}p{3.5cm} >{\raggedright\arraybackslash}p{3.2cm} X@{}}
\toprule
\textbf{Link} & \textbf{Protocol} & \textbf{Constraint it answers} \\
\midrule
Device $\rightarrow$ gateway & Bluetooth Low Energy & validated cuff and glucometer transmit without manual entry (\cref{sec:acquisition}) \\
Gateway $\rightarrow$ cloud & MQTT over 3G/4G & constrained device, unreliable link, small frequent payloads \\
Cloud $\rightarrow$ client & HTTP/REST & interoperability, request-driven access \\
Cloud $\rightarrow$ client, real time & WebSocket & event-proportional traffic, immediate alert delivery \\
Payload format & JSON & native to storage and runtime; no parser stage \\
\bottomrule
\end{tabularx}
\end{table}

%==========================================================================
\subsection{The cloud backend}\label{sec:backend}

\subsubsection{Why Firebase, according to the source}

\Textcite{efendi2025} enumerate four properties. Real-time database services enabling IoT devices to synchronise data instantaneously across multiple clients and platforms, which they describe as \bq{particularly valuable for applications requiring timely data updates, such as sensor readings and device status notifications}. Built-in authentication and authorisation mechanisms ensuring secure access control for IoT devices and users, which \bq{enhances the integrity and confidentiality of data transmitted and stored in the cloud, mitigating potential security risks associated with IoT deployments}. Scalability, which the authors summarise by saying that Firebase \bq{empowers IoT applications with scalability, real-time data processing, and seamless integration}. And offline operation, treated separately below.

\subsubsection{The NoSQL JSON schema, and a modelling consequence nobody states}\label{sec:schema}

Data are stored as a JSON tree. \Textcite{efendi2025} reproduce the live structure, which in their system is as flat as it could be: a single node holding \texttt{heart-rate}, \texttt{oxygen-saturation} and \texttt{temperature} with their current values.

This is a document store, not a relational database. Three consequences follow for the present work, and none of the four sources addresses them.

\textbf{A JSON tree is not natively a time series.} The structure above holds the current value and nothing else. To hold a series, the key design must carry time---typically a push key or timestamp-keyed child node per reading---and that design decision determines whether a range query over a window is cheap or expensive. Since the framing document specifies feature engineering by rolling aggregates, trends and variability, the data model must be designed around the window query, not around the current reading. The structure of \textcite{efendi2025} would require a full subtree read to compute a seven-day mean---which is precisely the quantity the measurement protocol of \cref{sec:protocols} defines as the reference value.

\textbf{Denormalisation shifts the burden to write time.} Document stores conventionally duplicate data across the shapes each client needs, so that reads are single-node. This is efficient for display and poor for ad hoc analysis, and it means the schema must anticipate the queries. A model-training export is an ad hoc analysis.

\textbf{Schemalessness is not an advantage here.} The record structure is fixed by the framing document---systolic, diastolic, blood glucose, heart rate, timestamp, plus background and context. A validated schema enforced at write time would catch the unit errors and carried-forward values that \textcite{liu2025review} identify as the characteristic defects of electronic medical records (\cref{sec:records}). Firebase supports this through security rules; the opportunity should be taken rather than relying on client-side validation.

\subsubsection{Offline-first synchronisation}\label{sec:offline}

This is the property that matters most for the target context, and the source states it directly: data are \bq{stored locally and real-time events continue to occur even when the device is offline}, and \bq{the real-time data will be synchronised between local data change remotely when the device is offline and auto merging when the device is going to be online}.

Set against the framing document's statement that intermittent connectivity in the Cameroonian context leads to irregular and missing measurements, and against the 550 million figure from \cref{sec:regional-landscape}, a storage layer whose default behaviour is local write with deferred merge removes one cause of missingness entirely: readings taken while offline are not lost, they are delayed.

One caution. Deferred synchronisation means the arrival timestamp differs from the measurement timestamp, sometimes by days. \Cref{sec:pipelines} identified this as a defect of manual entry; it reappears here as a property of the transport. The schema must therefore carry both times explicitly, and every temporal feature must be computed on the measurement time. A model trained on arrival order would learn the connectivity pattern rather than the physiology---and since connectivity correlates with geography and income, that model would encode exactly the inequalities the system exists to reduce.

\subsubsection{Alerting and notification}\label{sec:alerting}

The four sources place the alert function differently, and the differences are instructive.

\Textcite{mucchi2021} place alert management in the cloud layer, while assigning identification of anomalous patient conditions and the consequent warning notification to the edge layer---that is, a fast local warning and a considered cloud-side decision. Their application implements this in the professional profile, which performs \bq{pre-processing of the acquired data for a pre-alert generation in case of anomalous values and transmission to the cloud repository} and \bq{generation and management of alerts/warnings}.

\Textcite{yazdani2018} place it in a dedicated notification module with two distinct functions, patient notification and physician notification, and add the SMS reminder for missing data.

\Textcite{liu2025review} place it in the multi-point triggering and dynamic response mechanism between analysis and application.

\Textcite{efendi2025} do not implement alerting at all: their application layer classifies into stable or not stable and displays the result.

The convergent design is a two-level alert: a local pre-alert raised at the edge on the current reading, and a considered alert raised in the cloud on the series. This is structurally identical to the two-stage clinical protocol of \textcite{mihevc2025} documented in \cref{sec:triage}---an abnormal reading triggers a verification profile, and only the profile triggers a teleconsultation. The architecture and the clinical protocol agree, which is a strong argument for adopting both.

\subsubsection{What none of the sources reports}\label{sec:cloud-gaps}

\textbf{Operating cost.} No source states what any of these architectures costs to run per patient per year. Given that \cref{sec:regional-landscape} documented Babyl ceasing operations in August 2023 for financial reasons after reaching approximately 2.5 million registered users, and that \textcite{dodoo2022} identified sustainability as the characteristic failure of telemedicine projects in the region, this silence is a gap this thesis names. The billing unit of a real-time database is the persistent connection and the bytes synchronised, and a design that holds an open WebSocket per client consumes resources continuously rather than per reading. This must be quantified for the proposed system rather than assumed negligible.

\textbf{Vendor dependence.} Firebase is a proprietary managed service. None of the sources discusses what happens if pricing, terms or availability change, or whether data can be exported and the system rehosted. For a system intended to outlive a thesis, this is a design risk, and it is the technical counterpart of the sustainability barrier.

\textbf{Data residency.} The database endpoint reported by \textcite{efendi2025} is hosted in the \texttt{asia-southeast1} region. Where patient data physically reside is a regulatory question, and \textcite{dodoo2022} identified the legal and regulatory dimension as one of six barriers in sub-Saharan Africa. This thesis must state the chosen region and justify it.

%==========================================================================
\subsection{The edge layer and the triage requirement}\label{sec:edge-triage}

This subsection is where the architectural literature answers the requirement that \cref{ch:background} established.

\subsubsection{The requirement, restated}

\Cref{sec:specialist-access} concluded that a telemonitoring system deployed at scale in this region must triage, and gave three quantitative reasons. \Textcite{bediang2021} found that remote expertise took 10.3 days on average against a 72-hour protocol and was not significantly faster than the control pathway, the bottleneck being specialist availability rather than transmission. \Textcite{kingue2013} quantified the cost of systematic review at approximately six calls per provider per day, five minutes each, across 1\,187 calls. \Textcite{segun2025} reported that up to 18\% of sub-Saharan African countries had no registered cardiologist as of 2020.

A system that forwards every reading to a clinician reproduces this scarcity. A system that sorts relieves it.

\subsubsection{Where the function belongs}

The division of \textcite{mucchi2021} assigns the edge layer short-term analysis, quick processing, identification of anomalous conditions and consequent warning notification, and the cloud layer long-term data analysis through advanced algorithms for data cross-correlation. This is precisely a triage split:
\begin{equation*}
  \underbrace{\text{reading} \rightarrow \text{edge}}_{\text{anomalous? cheap, local}}
  \;\longrightarrow\;
  \underbrace{\text{cloud}}_{\text{a trend? costly, global}}
  \;\longrightarrow\;
  \underbrace{\text{clinician}}_{\text{decision? scarce}}
\end{equation*}

\begin{figure}[htbp]
\centering
\begin{tikzpicture}[
  st/.style={draw,rounded corners,align=center,font=\footnotesize,
             minimum height=1.5cm},
  ar/.style={-{Stealth[length=3mm]},very thick}]
  \node[st,fill=cGrey!12,minimum width=3.6cm] (e) {\textbf{Edge}\\ current reading\\ \scriptsize abundant, immediate};
  \node[st,fill=cRef!14,minimum width=2.9cm,right=1.5cm of e] (c) {\textbf{Cloud}\\ the series\\ \scriptsize costly, considered};
  \node[st,fill=cAfrica!16,minimum width=2.2cm,right=1.5cm of c] (d) {\textbf{Clinician}\\ \scriptsize scarce};
  \draw[ar] (e) -- node[above,font=\scriptsize]{pre-alert} node[below,font=\scriptsize]{most readings stop here} (c);
  \draw[ar] (c) -- node[above,font=\scriptsize]{alert} node[below,font=\scriptsize]{few reach here} (d);
\end{tikzpicture}
\caption{Triage as three filters in series, each stage costlier and scarcer than the last. The design objective is to maximise what each stage resolves, so that the final stage is reached only when necessary.}\label{fig:triage-cascade}
\end{figure}

\Textcite{yazdani2018} supply the complementary element at the clinician end: \bq{the physician can have access to patient data at any time through the patient status and is able to view patients' status on a priority basis. This prioritisation is done based on the patient status.} Triage is not only filtering but ordering: the clinician sees a queue sorted by urgency, not a chronological stream.

\subsubsection{Where the learned model sits}

\Textcite{liu2025review} place their data analysis layer---support vector machine, decision tree, $k$-means, gradient boosting, AdaBoost, neural networks---between data processing and application service, with a multi-point triggering and dynamic response mechanism connecting it upward.

This is the correct placement for the model developed in this thesis, and it differs from that of \textcite{efendi2025}, who place their classifier in the application layer. The difference is not pedantic: a model in the application layer classifies for display, while a model in the analysis layer classifies for routing. The anticipatory formulation of the framing document---forecasting degradation at 24 to 48 hours---is a routing function, since its output determines who is contacted and when.

The rule selection and CDSS setting modules of \textcite{yazdani2018} suggest a further refinement worth considering. Their architecture allows different decision rules to be attached to different diseases, configured by medical staff through an administrator interface, so that the platform supports several decision support systems rather than being rebuilt for each. For a system covering both hypertension and diabetes, with thresholds anchored in the ESH and ADA references respectively (\cref{sec:pipelines}), this separation of the rule base from the platform is a sound design and a defensible extensibility claim.

%==========================================================================
\subsection{Security, privacy and the gap between recommendation and practice}\label{sec:security}

\subsubsection{What is recommended}

\Textcite{mucchi2021} set out the principle clearly: \bq{health data can flow from patient to cloud and from cloud to others (health operators) by using different channels (often wireless), which not always can be assumed secure. A trusted and secure domain for the health data should be assured, so that there is no segment which has a lower security than others.} They add that beyond classical cryptography, \bq{modern security techniques should be added, such as physical-layer security and crypto-key extraction from health signals}, that access policy should be addressed \bq{by making the patient aware of where the data is and who is requesting access to it}, and that blockchain \bq{can also be used to trace the access log to data and for non-repudiation of actions ordered or made on the patient's data}.

\Textcite{liu2025review} devote a section to the same question, surveying differential privacy frameworks using convolutional autoencoders and generative adversarial networks, $k$-anonymity through the Mondrian algorithm, RFID-based secure authentication for remote patient monitoring, lightweight encryption with access tokens, permissioned blockchain using elliptic curve digital signatures, and private set intersection. They observe that regulatory compliance is an imperative for healthcare providers, necessitating a prioritisation of data governance security.

\Textcite{efendi2025} report that the ESP32 supports IEEE 802.11 security including WAPI, WFA and WPA/WPA2, and that Firebase provides built-in authentication and authorisation.

\subsubsection{What was implemented}\label{sec:security-gap}

\Textcite{mucchi2021} state it themselves, and the admission deserves to be quoted: \bq{this paper does not focus on this topic, for the sake of space, but anyway this paragraph wants to summarise the main characteristics of a good policy of security and privacy when a tele-medicine architecture is on. Concerning our experimental activity, we only relied on the application layer security protocols, and on the fact that the data where stored in a private server without public access.}

The gap between the recommendation---physical-layer security, key extraction from biosignals, blockchain audit trails---and the practice---application-layer transport encryption and a server that is not publicly addressable---is the characteristic pattern of this literature. It is not a criticism of the authors, who are candid about it. It is a warning that security sections in this field frequently describe what ought to be done rather than what was done, and that this thesis states its own implementation plainly rather than enumerating techniques it does not use.

\subsubsection{The minimum defensible position}

\begin{table}[htbp]
\centering\small
\caption{Minimum security and privacy requirements, and their basis.}\label{tab:security-spec}
\begin{tabularx}{\textwidth}{@{}X >{\raggedright\arraybackslash}p{5.4cm}@{}}
\toprule
\textbf{Requirement} & \textbf{Basis} \\
\midrule
Transport encryption end to end, with no segment weaker than the others & \textcite{mucchi2021} \\
Authenticated access, role-separated (patient, nurse, physician) & \textcite{efendi2025}; user roles in \textcite{yazdani2018} \\
Direct identifiers removed before any modelling & \textcite{cadrage2026} \\
Declared data residency region & endpoint of \textcite{efendi2025}; regulatory barrier in \textcite{dodoo2022} \\
Access traceability: who read what, when & \textcite{mucchi2021} \\
Documented consent, with purpose limited to research & \textcite{cadrage2026} \\
Written statement of what was implemented, distinct from what is recommended & \cref{sec:security-gap} \\
\bottomrule
\end{tabularx}
\end{table}

\Textcite{mucchi2021} also record the applicable legal instrument in their jurisdiction, Italian Legislative Decree 196/2003. The Cameroonian equivalent must be identified and cited; none of the sources in this corpus provides it.%
\verify{Identify and cite the Cameroonian legal instrument governing personal and health data processing, together with any national health research ethics requirement applicable to the cohort. This cannot be drawn from the present corpus and must come from national sources.}

%==========================================================================
\subsection{What these architectures have actually been validated on}\label{sec:arch-validation}

This subsection is necessary because the four sources are frequently cited as if they were equivalent, and they are not.

\begin{table}[htbp]
\centering\footnotesize
\caption{Implementation and validation status of the four architectural sources.}\label{tab:arch-status}
\begin{tabularx}{\textwidth}{@{}>{\raggedright\arraybackslash}p{3.2cm} >{\raggedright\arraybackslash}X >{\raggedright\arraybackslash}X >{\raggedright\arraybackslash}p{2.6cm}@{}}
\toprule
\textbf{Source} & \textbf{Implementation} & \textbf{Validation} & \textbf{Real patients} \\
\midrule
\textcite{yazdani2018} & none: UML design only & none & no \\
\textcite{liu2025review} & none: survey diagram & none & no \\
\textcite{efendi2025} & full hardware and cloud & 6 participants, 10 min, seated & no: non-patient volunteers, two with a chronic disease \\
\textcite{mucchi2021} & Android application, cloud repository & 10 nurses, 20 days, simulated patients & no \\
\bottomrule
\end{tabularx}
\end{table}

\subsubsection{The validation of Mucchi et al.}

The protocol is documented carefully. Ten family and community nurses of the Florence city district were selected, with inclusion criteria of at least six months' length of service and possession of an Android smartphone; no requirement was placed on information technology skills, the authors noting that one aim was to develop an application as user-friendly as possible. A training session presented the protocol and functionalities, the application was installed on the nurses' personal phones, and simulated vital parameters were entered. Over the following twenty days each nurse was asked to simulate at least one patient check-up per day, entering parameters and checking the historical record including alert or warning messages. A questionnaire followed, assessed on the mean opinion score scale.

No real patients were involved. The authors state that the second phase, \bq{whose Ethical Committee approval process is still ongoing, will involve the introduction of the App on an experimental basis in clinical practice}.

\subsubsection{The result the authors do not discuss}\label{sec:usability-result}

\begin{table}[htbp]
\centering\footnotesize
\caption{Usability questionnaire results \parencite{mucchi2021}, Table~1. The scale runs from insufficient to excellent; on the mean opinion score convention, ``sufficient'' ranks below ``fair''.}\label{tab:mucchi-usability}
\begin{tabularx}{\textwidth}{@{}X c c c c c@{}}
\toprule
\textbf{Investigated aspect} & \textbf{Insuff.} & \textbf{Suff.} & \textbf{Fair} & \textbf{Good} & \textbf{Excel.} \\
\midrule
Nurse ICT skills & --- & --- & 67\% & 33\% & --- \\
App functionalities w.r.t. operating conditions & --- & --- & 33\% & 50\% & 17\% \\
App usability in operative contexts & --- & --- & 33\% & 50\% & 17\% \\
Approval of the service and the app & --- & --- & 67\% & 33\% & --- \\
System efficiency for managing chronic diseases & 16\% & 33\% & 17\% & 17\% & 17\% \\
\bottomrule
\end{tabularx}
\end{table}

\begin{figure}[htbp]
\centering
\begin{tikzpicture}
\begin{axis}[xbar stacked,width=0.70\textwidth,height=6cm,
  xmin=0,xmax=100,xlabel={Respondents (\%)},
  symbolic y coords={Efficiency for chronic disease,Approval of service,Usability in context,Functionalities,Nurse ICT skills},
  ytick=data,yticklabel style={font=\scriptsize},
  legend style={font=\scriptsize,at={(0.5,-0.22)},anchor=north,legend columns=5},
  bar width=15pt]
\addplot[fill=cAfrica,draw=cAfrica] coordinates
  {(0,Nurse ICT skills)(0,Functionalities)(0,Usability in context)(0,Approval of service)(16,Efficiency for chronic disease)};
\addplot[fill=cAfrica!45,draw=cAfrica!60] coordinates
  {(0,Nurse ICT skills)(0,Functionalities)(0,Usability in context)(0,Approval of service)(33,Efficiency for chronic disease)};
\addplot[fill=cGrey!45,draw=cGrey] coordinates
  {(67,Nurse ICT skills)(33,Functionalities)(33,Usability in context)(67,Approval of service)(17,Efficiency for chronic disease)};
\addplot[fill=cRef!55,draw=cRef] coordinates
  {(33,Nurse ICT skills)(50,Functionalities)(50,Usability in context)(33,Approval of service)(17,Efficiency for chronic disease)};
\addplot[fill=cRef,draw=cRef!60!black] coordinates
  {(0,Nurse ICT skills)(17,Functionalities)(17,Usability in context)(0,Approval of service)(17,Efficiency for chronic disease)};
\legend{Insufficient,Sufficient,Fair,Good,Excellent}
\end{axis}
\end{tikzpicture}
\caption{The usability results of \cref{tab:mucchi-usability}. Only the bottom row---efficiency for the clinical purpose---receives any rating below ``fair'', and it is the only aspect spread across all five categories.}\label{fig:usability}
\end{figure}

The final row is the only aspect on which any respondent returned an insufficient rating, and the published percentages place half of the respondents at sufficient or below. It is also the only aspect spread across all five categories. The results discussion of \textcite{mucchi2021} states that the main tested functionalities of the application can effectively support the nursing staff, which is supported by the second and third rows, and does not comment on the fifth.

A caution on the denominator is required before any of these figures is quoted. Ten nurses were recruited, but every percentage in the table is an exact multiple of one sixth: 16.7, 33.3, 50.0, 66.7. The five categories of the final row correspond to one, two, one, one and one respondent out of six. The questionnaire therefore appears to have been answered by six respondents, not ten, and the figures should be reported as counts out of six---three of six rating efficiency at sufficient or below---rather than as the percentages 16\% and 33\%, whose sum of 49\% is a rounding artefact of three sixths.%
\verify{Confirm the number of questionnaire respondents in \textcite{mucchi2021}. Ten nurses were enrolled, but every value in their Table~1 is a multiple of $1/6$, which implies six respondents. If six is correct, report counts rather than percentages throughout.}

This matters because the fifth row is the aspect closest to the clinical question. Usability was rated well; efficiency for the purpose was not. This thesis cites this architecture for its structure, which is sound and well specified, and does not cite it as evidence that such a system is effective---which its own usability data do not support and which it never claims.

The first row is also worth retaining: two-thirds of respondents rated their own information technology skills only fair and none rated them excellent, in a high-income European setting with no skill requirement for inclusion. Read with the 10.9\% equipment-related dropout in \textcite{mihevc2025} (\cref{sec:feasibility}) and with the individual barrier of \textcite{dodoo2022}, this converges on a consistent finding: digital competence of front-line staff is a binding constraint, not a transitional one.

\subsubsection{The feedback the nurses gave}

\Textcite{mucchi2021} report the comments, and three bear directly on the present design: \bq{the alerts notification should be improved especially in case of urgent situations; the integration with the protocols adopted in different real operating conditions should be considered}; \bq{interfacing with the current management information systems in use in the healthcare sector may be useful}; and \bq{the possibility of interfacing with other healthcare professionals and the need for an immediate feedback from the doctor in the event of an alert should be considered}.

The first and third are the same request in different words: the alert must reach a decision-maker, and reach them fast. This is the professional-oversight finding of \cref{sec:oversight} arriving from a fourth direction---through the stated preference of the nurses using the system.

%==========================================================================
\subsection{Specification}\label{sec:spec22}

\begin{table}[htbp]
\centering\footnotesize
\caption{Architectural requirements established in \cref{sec:cloud} and their basis.}\label{tab:spec22}
\begin{tabularx}{\textwidth}{@{}X >{\raggedright\arraybackslash}p{5.9cm}@{}}
\toprule
\textbf{Requirement} & \textbf{Basis} \\
\midrule
Four-layer architecture with an explicit edge tier & \textcite{mucchi2021}; edge absent in \textcite{efendi2025} and \textcite{yazdani2018} \\
Edge performs short-term analysis and local pre-alert & \textcite{mucchi2021} \\
Cloud performs long-term analysis and considered alert & \textcite{mucchi2021} \\
Two-level alerting mirroring the clinical two-stage protocol & \textcite{mucchi2021}; \textcite{mihevc2025}, \cref{sec:triage} \\
Learned model in the analysis layer, not the application layer & \textcite{liu2025review}; routing, not display \\
Clinician view ordered by priority, not chronology & \textcite{yazdani2018} \\
Configurable checkup period per patient & \textcite{yazdani2018}; WHO follow-up intervals, \cref{sec:followup-requires} \\
Rule base separable from platform, to carry both ESH and ADA thresholds & CDSS setting module of \textcite{yazdani2018} \\
MQTT upstream, HTTP/REST downstream, WebSocket for push & \textcite{efendi2025} \\
JSON payloads; XML not adopted & \textcite{efendi2025}; contra \textcite{yazdani2018} \\
Schema keyed for window queries, not current value only & \cref{sec:schema}; rolling-aggregate features \\
Measurement time and arrival time stored separately & \cref{sec:offline}; deferred sync decouples them \\
Server-side schema validation at write time & \cref{sec:schema}; EMR defects in \textcite{liu2025review} \\
Offline-first local storage with automatic merge & \textcite{efendi2025}; connectivity limitation in \textcite{cadrage2026} \\
Local storage unit at the edge & hub composition in \textcite{mucchi2021} \\
End-to-end encryption with no weaker segment & \textcite{mucchi2021} \\
Role-separated authentication: patient, nurse, physician & \textcite{efendi2025}; \textcite{yazdani2018} \\
Declared data residency and applicable national instrument & endpoint of \textcite{efendi2025}; \textcite{dodoo2022} \\
Access traceability & \textcite{mucchi2021} \\
Implemented security stated separately from recommended security & \cref{sec:security-gap} \\
Operating cost per patient quantified & \cref{sec:cloud-gaps}; Babyl collapse, \cref{sec:regional-landscape} \\
Export path and vendor independence assessed & \cref{sec:cloud-gaps} \\
Front-line digital competence treated as a binding constraint & \cref{sec:usability-result}; \textcite{mihevc2025} \\
\bottomrule
\end{tabularx}
\end{table}

%==========================================================================
\subsection{What this section establishes, and what it does not}\label{sec:cloud-conclusion}

\textbf{Established.} The multi-tier model is stable across four independent sources and admits a clean mapping (\cref{tab:layer-mapping}). The protocol choices are coherent and defensible. The cloud backend offers the offline-first behaviour the deployment context requires. And the edge--cloud division of labour in the reference architecture is structurally the triage function that \cref{ch:background} demonstrated to be necessary---which means the architecture required by the clinical analysis and the architecture proposed by the engineering literature coincide. That convergence is the main result of this section.

\textbf{Not established.} No architecture in this corpus has been validated on real patients in routine care. Two have no implementation at all. The one architectural source with a usability evaluation returned half of its respondents rating system efficiency for chronic disease management at sufficient or below, a result its authors do not discuss. No source quantifies the protocol efficiency it claims, nor the operating cost, nor the behaviour of any of these designs under the intermittent connectivity they nonetheless parameterise for. The architecture is well specified and poorly evidenced, and this thesis says so.

%##########################################################################
% SECTION 2.3
%##########################################################################
\section{Machine Learning Algorithms and Predictive AI for Risk Assessment}\label{sec:ml}

\subsection*{Purpose and status of the sources}

\Cref{sec:acquisition,sec:cloud} established what is acquired and how it travels. This section addresses what is computed on it. The governing requirement, carried from \cref{ch:background} and refined in \cref{sec:edge-triage}, is that the model exists to triage: to decide which patients reach a clinician's attention and in what order. That purpose determines which algorithms are candidates, which metrics matter, and which failure modes are unacceptable.

\paragraph{Mendoza-Cano et al. (2025).} \textcite{mendoza2025}, published in \emph{PLOS ONE} 20(9):e0331565, is a cross-sectional study identifying predictors of uncontrolled hypertension in 1\,308 Mexican adults from the nationally representative ENSANUT 2022 survey, comparing LASSO-selected random forest, XGBoost and logistic regression with SHAP interpretability. Its dataset is publicly archived and its code published. It is the only source in this corpus that applies these algorithms to hypertension control as a genuine prediction problem and reports full discrimination metrics with confidence intervals.

\paragraph{Efendi et al. (2025).} \textcite{efendi2025}, examined in \cref{sec:reference-impl,sec:cloud}, compares five classifiers on an IoT-derived dataset.

\paragraph{Liu and Wang (2025).} \textcite{liu2025review} survey machine learning applied to chronic disease monitoring. The reservation stated in \cref{sec:diagnosis} stands.

\paragraph{The research framing document.} \textcite{cadrage2026} is used here to confront the state of the art with the contributions and hypotheses actually proposed.

%==========================================================================
\subsection{The algorithm families and their inductive biases}\label{sec:families}

Before comparing reported performances it is necessary to state what each family assumes, because the comparison in \cref{sec:reported} is otherwise uninterpretable.

\subsubsection{Decision trees}

A decision tree recursively partitions the feature space by axis-aligned splits, choosing at each node the split that most reduces an impurity criterion. For a binary outcome with proportion $p$ of positives in a node, the Gini impurity is
\begin{equation}
  G = 2p(1-p),
\end{equation}
and a split is selected to minimise the weighted impurity of the children. Trees are interpretable and require no feature scaling, and they handle mixed categorical and continuous inputs natively. Their weakness is variance: a small change in the data can produce a different tree.

\subsubsection{Bagging: random forest}

A random forest fits many trees on bootstrap resamples and, at each split, considers only a random subset of features; predictions are the majority vote. The variance of an average of $B$ identically distributed variables with pairwise correlation $\rho$ and variance $\sigma^2$ is
\begin{equation}
  \rho\sigma^2 + \frac{1-\rho}{B}\,\sigma^2,
\end{equation}
so the random feature subsetting, which lowers $\rho$, is what makes the ensemble work rather than the averaging alone: the second term vanishes as $B$ grows, the first does not. \Textcite{mendoza2025} describe the method as \bq{ensemble learning methods which construct decision trees during training and output the mode of the classes}, noting its \bq{robustness against overfitting} and \bq{ability to handle large datasets with high dimensionality}. They configured 500 trees, three variables considered per split, and node size one.

\subsubsection{Boosting: AdaBoost and XGBoost}

Boosting fits trees sequentially, each correcting the residual errors of its predecessors. XGBoost minimises a regularised objective
\begin{equation}
  \mathcal{L} = \sum_i \ell(y_i, \hat{y}_i) + \sum_k \Omega(f_k),
  \qquad \Omega(f) = \gamma T + \tfrac{1}{2}\lambda\lVert w \rVert^2,
\end{equation}
where $T$ is the number of leaves and $w$ the leaf weights, so that complexity is penalised explicitly. \Textcite{mendoza2025} configured maximum depth six, learning rate 0.3, 100 boosting rounds and a fixed random seed of 123 for reproducibility.

Boosting typically outperforms bagging on tabular data of moderate size, but it is more sensitive to noise and to class imbalance, because misclassified examples receive increasing weight---and in an imbalanced problem the misclassified examples are disproportionately the rare class, or disproportionately the noise.

\subsubsection{Support vector machines}

A support vector machine seeks the separating hyperplane maximising the margin, solving
\begin{equation}
  \min_{w,b} \; \tfrac{1}{2}\lVert w \rVert^2 + C\sum_i \xi_i
  \quad\text{subject to}\quad
  y_i(w^{\top} x_i + b) \geq 1 - \xi_i, \;\; \xi_i \geq 0,
\end{equation}
with non-linear boundaries obtained by kernel substitution. \Textcite{liu2025review} report support vector machines widely used in this field, citing an accuracy of 92.19\% for dementia prediction and 92.19\% for respiratory pathology classification among their examples.%
\verify{The same figure, 92.19\%, is reported for two unrelated applications. One of the two is probably a transcription error in \textcite{liu2025review}. Check the two primary references before reproducing either value.}
Support vector machines scale poorly with sample size and require feature scaling, and they do not produce calibrated probabilities without an additional step.

\subsubsection{Recurrent networks: LSTM}

A long short-term memory cell maintains a state vector updated through gates that control what is retained, discarded and emitted:
\begin{align}
  f_t &= \sigma(W_f[h_{t-1}, x_t] + b_f), \qquad
  i_t = \sigma(W_i[h_{t-1}, x_t] + b_i), \qquad
  o_t = \sigma(W_o[h_{t-1}, x_t] + b_o),\\
  c_t &= f_t \odot c_{t-1} + i_t \odot \tanh(W_c[h_{t-1}, x_t] + b_c),
  \qquad h_t = o_t \odot \tanh(c_t).
\end{align}
The gating permits gradient flow across long sequences, which is why such networks are the conventional baseline for clinical time series.

\Textcite{liu2025review} report one directly relevant application: Reddy and Delen used a recurrent LSTM network to predict readmission of lupus patients, with the deep method outperforming traditional classification and achieving an area under the curve of 0.70. This figure deserves emphasis, because it is the only realistic recurrent-network result in the corpus and it sits far below the 95--100\% accuracies that dominate the surrounding literature.

\subsubsection{Comparison of inductive biases}

\begin{table}[htbp]
\centering\footnotesize
\caption{Inductive biases of the five candidate families.}\label{tab:inductive-bias}
\begin{tabularx}{\textwidth}{@{}>{\raggedright\arraybackslash}p{2.3cm} >{\raggedright\arraybackslash}X >{\raggedright\arraybackslash}X >{\raggedright\arraybackslash}X >{\raggedright\arraybackslash}p{2.3cm}@{}}
\toprule
\textbf{Family} & \textbf{Assumes} & \textbf{Handles} & \textbf{Fails on} & \textbf{Interpretability} \\
\midrule
Decision tree & axis-aligned splits suffice & mixed types, missing values & high variance & direct \\
Random forest & ensembling reduces variance & high dimension, noise & smooth boundaries & importance measures \\
XGBoost & residuals are learnable sequentially & tabular structure & noise, imbalance & SHAP \\
SVM & margin separability after kernel mapping & small $n$, high $p$ & large $n$, calibration & low \\
LSTM & temporal dependence carries signal & sequences, irregular lags & small $n$, needs volume & low \\
\bottomrule
\end{tabularx}
\end{table}

The decisive line of \cref{tab:inductive-bias} is the last. Only the recurrent network treats the input as a sequence. The other four receive a feature vector and have no notion of order; any temporal information they use must be engineered into that vector as lags, rolling means, slopes or variability. This is precisely the comparison the framing document poses in its third contribution and first hypothesis.

%==========================================================================
\subsection{What the literature reports, and the contradiction}\label{sec:reported}

\subsubsection{Efendi et al.: XGBoost best}

\begin{table}[htbp]
\centering\small
\caption{Classifier comparison reported by \textcite{efendi2025}, Table~16.}\label{tab:efendi-models}
\begin{tabularx}{\textwidth}{@{}X c c c c@{}}
\toprule
\textbf{Model} & \textbf{Precision} & \textbf{Recall} & \textbf{Accuracy} & \textbf{$F_1$} \\
\midrule
XGBoost & 0.973 & 0.968 & 0.973 & 0.970 \\
AdaBoost & 0.956 & 0.937 & 0.967 & 0.946 \\
Decision tree & 0.941 & 0.925 & 0.931 & 0.933 \\
Random forest & 0.936 & 0.922 & 0.946 & 0.929 \\
Bayesian classifier & 0.925 & 0.931 & 0.931 & 0.928 \\
\bottomrule
\end{tabularx}
\end{table}

All five exceed 0.92 on every metric. XGBoost is selected and embedded in the deployed system.

\subsubsection{Mendoza-Cano et al.: XGBoost worst}

\begin{table}[htbp]
\centering\footnotesize
\caption{Model comparison reported by \textcite{mendoza2025}, Table~3.}\label{tab:mendoza-models}
\begin{tabularx}{\textwidth}{@{}X >{\raggedright\arraybackslash}p{2.3cm} c c c c c@{}}
\toprule
\textbf{Model} & \textbf{AUC (95\% CI)} & \textbf{Prec.} & \textbf{Sens.} & \textbf{Spec.} & \textbf{Acc.} & \textbf{$F_1$} \\
\midrule
Random forest & 0.75 (0.72--0.77) & 0.94 & 0.51 & 0.98 & 0.84 & 0.659 \\
XGBoost & 0.54 (0.49--0.60) & 0.42 & 0.23 & 0.85 & 0.65 & 0.294 \\
Logistic regression & 0.61 (0.59--0.64) & 0.70 & 0.98 & 0.03 & 0.70 & 0.818 \\
\bottomrule
\end{tabularx}
\end{table}

\begin{figure}[htbp]
\centering
\begin{tikzpicture}
\begin{axis}[width=0.78\textwidth,height=4.2cm,
  xlabel={Area under the ROC curve},
  xmin=0.40,xmax=0.85,ymin=0.4,ymax=3.6,ytick={1,2,3},
  yticklabels={{\footnotesize Logistic regression},{\footnotesize XGBoost},{\footnotesize Random forest}},
  grid=major,grid style={gray!20}]
\addplot[dashed,cAfrica,thick,no markers] coordinates {(0.5,0.4)(0.5,3.6)};
\addplot[only marks,mark=square*,mark size=2.6pt,cRef,error bars/.cd,x dir=both,x explicit,error bar style={thick}]
  coordinates {(0.75,3) += (0.02,0) -= (0.03,0)
               (0.61,1) += (0.03,0) -= (0.02,0)};
\addplot[only marks,mark=*,mark size=2.6pt,cAfrica,error bars/.cd,x dir=both,x explicit,error bar style={thick}]
  coordinates {(0.54,2) += (0.06,0) -= (0.05,0)};
\node[font=\scriptsize,text=cAfrica,anchor=south] at (axis cs:0.5,3.15) {chance};
\end{axis}
\end{tikzpicture}
\caption{Discrimination reported by \textcite{mendoza2025} with 95\% confidence intervals. The XGBoost interval spans the chance line: on this dataset the model is not distinguishable from a coin flip.}\label{fig:mendoza-auc}
\end{figure}

The confidence interval for the XGBoost area under the curve spans 0.49 to 0.60 and therefore includes 0.50 (\cref{fig:mendoza-auc}). The authors state it plainly: \bq{XGBoost documented the lowest predictive performance}.

\subsubsection{The ambient literature}

\Textcite{liu2025review} report a long series of accuracies from the chronic-disease literature: gradient boosting at 99\% for chronic kidney disease risk; an artificial neural network for kidney disease detection nearing 100\%; a deep recurrent network with supervised multi-labelling achieving 99.76\% in an intrusion detection system for the internet of medical things; a deep-learning-assisted AdaBoost variant at 97.2\% accuracy for breast cancer; a convolutional network at 84.5\% against $k$-nearest neighbours at 83\% for disease prediction; AdaBoost and random forest exhibiting accuracy over 95\%.

The distribution of these figures is itself informative. A literature in which reported accuracies cluster between 95\% and 100\% is not a literature reporting hard problems. \Textcite{liu2025review} themselves note that the limitations of machine learning methods include their black-box nature and dependence on a large amount of labelled data, and that the efficacy of a model \bq{is contingent upon the dataset's representation rather than its volume}.

\subsubsection{The contradiction stated}

\begin{table}[htbp]
\centering\small
\caption{The same algorithm on two health-monitoring classification tasks, both published in 2025.}\label{tab:contradiction}
\begin{tabularx}{\textwidth}{@{}>{\raggedright\arraybackslash}p{3.6cm} >{\raggedright\arraybackslash}X >{\raggedright\arraybackslash}X@{}}
\toprule
 & \textbf{\textcite{efendi2025}} & \textbf{\textcite{mendoza2025}} \\
\midrule
XGBoost accuracy & 0.973 & 0.65 \\
XGBoost $F_1$ & 0.970 & 0.294 \\
XGBoost AUC & not reported & 0.54 (0.49--0.60) \\
Best model & XGBoost & random forest \\
Outcome predicted & ``stable'' vs ``not stable'' & uncontrolled hypertension \\
Sample & 4\,500 records, balanced 50/50 & 1\,308 adults, 30.3\% positive \\
Features & 3 vital signs, discretised & demographic and administrative \\
\bottomrule
\end{tabularx}
\end{table}

\begin{figure}[htbp]
\centering
\begin{tikzpicture}
\begin{axis}[ybar,width=0.62\textwidth,height=5.4cm,bar width=34pt,
  ymin=0,ymax=1.08,ylabel={$F_1$ reported for XGBoost},
  symbolic x coords={Efendi 2025,Mendoza-Cano 2025},
  xtick=data,xticklabel style={font=\footnotesize},
  enlarge x limits=0.6,
  nodes near coords,every node near coord/.append style={font=\footnotesize}]
\addplot[fill=cGrey!55,draw=cGrey] coordinates {(Efendi 2025,0.970)(Mendoza-Cano 2025,0.294)};
\end{axis}
\end{tikzpicture}
\caption{The same algorithm, the same year, two peer-reviewed health-monitoring tasks, and a difference of 0.676 in $F_1$. The explanation lies in how each task was constructed, not in the algorithm.}\label{fig:xgb-contradiction}
\end{figure}

Both are peer-reviewed, both published in 2025, both use the same algorithm on a health-monitoring classification task, and they disagree by 0.676 of $F_1$ (\cref{fig:xgb-contradiction}). A thesis that cited only the first would conclude that XGBoost is the obvious choice. A thesis that cited only the second would conclude it is unusable. Neither conclusion is correct, and the reason is the subject of the next subsection.

%==========================================================================
\subsection{Resolving the contradiction: the task, not the algorithm}\label{sec:task-not-algo}

\subsubsection{Efendi et al. measure label recovery, not prediction}

\Cref{sec:leakage} established the mechanism. The class label \emph{stable} is defined as the condition in which heart rate, oxygen saturation and temperature all fall within the normal ranges of the paper's reference tables. The feature transformation then recodes each parameter as 0, 1 or 2---low, normal, high---with reference to those same ranges. After transformation, the label is a deterministic function of the three transformed features, $y = 1 \iff (x_1 = 1) \wedge (x_2 = 1) \wedge (x_3 = 1)$.

Any classifier capable of representing a three-way conjunction achieves perfect accuracy on this task. The reported spread of 0.042 in $F_1$ across five algorithms measures nothing about their relative merit; it measures implementation noise on a problem that is solved by construction.

Three further observations confirm that the result is not a measurement of predictive capability.

\textbf{The deployed model contradicts the paper's own thresholds.} In the real-world experiment, a single 50-year-old participant was monitored in four conditions. Sleeping, the system recorded 70.1~BPM and predicted \emph{stable}. One reference table gives the normal range at age 50 as 85--119~BPM, the other 85--145~BPM. A value of 70.1~BPM lies below the lower bound of both, and the model nonetheless returned \emph{stable}. During cardio exercise it recorded 121~BPM and again returned \emph{stable}, which the second table permits but the first does not. The classifier therefore does not implement the labelling rule it was trained on, or the ranges stated in the paper are not those used.

\textbf{The training configuration is not applicable to the models trained.} The authors report 300 epochs, a learning rate of 0.025, a sigmoid activation function and the Adam optimiser, and plot training loss over 300 epochs for all five models including the decision tree, the random forest and the Bayesian classifier. None of these three trains by gradient descent; none has epochs, a sigmoid activation or an Adam optimiser. The paper also describes XGBoost as a deep learning model in its contribution table; XGBoost is an ensemble of gradient-boosted decision trees and is not a deep learning method.

\textbf{The reported metrics are not mutually consistent.} On a test set of 900 records balanced 50/50, precision and recall determine accuracy. For random forest, a recall of 0.922 and a precision of 0.936 imply an accuracy of 0.930, not the 0.946 reported; for AdaBoost, 0.937 and 0.956 imply 0.947, not 0.967. The remaining three rows are consistent within rounding.%
\verify{Recompute the confusion matrix of each model in \textcite{efendi2025}, Table~16, against the stated 50/50 test balance. Two of the five rows do not close. Either the test split is not balanced as stated, or the metrics were computed on different subsets.}

These are not pedantic objections. They indicate that the comparison was not conducted under conditions in which it could be informative, and they mean that \cref{tab:efendi-models} cannot be cited as evidence that XGBoost outperforms random forest.

\subsubsection{Mendoza-Cano et al. measure genuine prediction from weak features}

The opposite situation obtains. The outcome---uncontrolled blood pressure, defined as systolic $\geq 130$~mmHg or diastolic $\geq 80$~mmHg on the average of the second and third of three readings---is measured independently of the predictors. The predictors, selected by LASSO at $\lambda = 0.0056528$ with ten-fold cross-validation, are sex, age, socioeconomic status, years since diagnosis, diabetes history, source of medical care and adherence to a nutritional plan. Physical activity engagement was eliminated with a zero coefficient.

None of these is a physiological measurement. They are demographic and administrative variables. The question being asked is whether one can predict who has uncontrolled blood pressure without measuring anything about their physiology---and the answer, on 1\,308 participants, is: weakly. Random forest reaches an area under the curve of 0.75; logistic regression 0.61; XGBoost 0.54.

This is not a failure of the algorithms. It is a correct measurement of how much signal demographic variables carry about blood pressure control, and the answer is some, but not much. Note also that this outcome is defined at the 130/80~mmHg threshold rather than the 140/90~mmHg convention retained throughout this thesis (\cref{sec:hta-burden}), so its prevalence of 30.3\% is not directly comparable with the control figures of \cref{tab:cascade}.

\subsubsection{The principle, and why it governs this thesis}\label{sec:principle}

\begin{equation}
  \text{reported performance} = f(\text{task difficulty}, \text{label integrity}, \text{feature informativeness}, \text{algorithm}).
\end{equation}

The algorithm is the last and smallest term. A reported accuracy figure is uninterpretable without the other three, and comparing accuracies across studies with different tasks is meaningless.

For the present work this yields a direct and uncomfortable implication. The framing document specifies that the risk label---stable, moderate, high, critical---is derived from recognised clinical thresholds, validated by the partner physician. If the model's features then include the current blood pressure and blood glucose values, the task becomes that of \textcite{efendi2025}: recovering a threshold rule from the variables that define it. The resulting accuracy will be high and will demonstrate nothing.

The anticipatory formulation of the framing document---forecasting physiological degradation at 24 to 48 hours---is the escape, and it must be presented as such rather than as an additional ambition. Formally, the leaking task is
\begin{equation}
  \hat{y}_t = g(x_t) \quad\text{where}\quad y_t = h(x_t),
\end{equation}
and the legitimate task is
\begin{equation}
  \hat{y}_{t+\Delta} = g(x_{t-k}, \dots, x_t), \qquad \Delta \in [24\,\text{h}, 48\,\text{h}],
\end{equation}
in which the label lies outside the input window by construction and cannot be recovered from it (\cref{fig:leak-vs-forecast}). The expected accuracy of the second task is far lower than the first, and that lower figure is the honest one. This thesis anticipates the difference explicitly, because a jury that sees 0.97 on a threshold-classification task will ask the question posed above.

\begin{figure}[htbp]
\centering
\begin{tikzpicture}[font=\footnotesize]
  % leaking task
  \draw[thick,-{Stealth[length=2mm]}] (0,1.6) -- (8.6,1.6);
  \fill[cAfrica!25] (3.6,1.35) rectangle (4.4,1.85);
  \node[anchor=south,font=\scriptsize] at (4.0,1.9) {$x_t$};
  \node[circle,fill=cAfrica,inner sep=2pt] at (4.0,1.6) {};
  \node[anchor=north,font=\scriptsize,text=cAfrica] at (4.0,1.3) {$y_t = h(x_t)$};
  \node[anchor=west,font=\footnotesize\bfseries] at (8.7,1.6) {leaks};
  \node[anchor=east] at (-0.15,1.6) {(a)};
  % legitimate task
  \draw[thick,-{Stealth[length=2mm]}] (0,0) -- (8.6,0);
  \fill[cRef!22] (1.2,-0.25) rectangle (4.4,0.25);
  \node[anchor=south,font=\scriptsize] at (2.8,0.3) {input window $x_{t-k}\dots x_t$};
  \foreach \x in {1.4,2.0,2.6,3.2,3.8,4.3} \node[circle,fill=cRef,inner sep=1.5pt] at (\x,0) {};
  \draw[dashed,cGrey] (4.4,-0.45) -- (4.4,0.45);
  \node[circle,fill=cAfrica,inner sep=2pt] at (6.6,0) {};
  \node[anchor=north,font=\scriptsize,text=cAfrica] at (6.6,-0.15) {$y_{t+\Delta}$};
  \draw[{Stealth[length=2mm]}-{Stealth[length=2mm]},cGrey] (4.4,0.55) -- (6.6,0.55)
    node[midway,above,font=\scriptsize,text=cGrey]{$\Delta$};
  \node[anchor=west,font=\footnotesize\bfseries] at (8.7,0) {does not};
  \node[anchor=east] at (-0.15,0) {(b)};
\end{tikzpicture}
\caption{(a) Classifying the current state from the current reading, when the label is a rule applied to that same reading. (b) Forecasting a state outside the input window. In (b) the label cannot be recovered from the input, so the measured performance is a measurement.}\label{fig:leak-vs-forecast}
\end{figure}

%==========================================================================
\subsection{Metrics: what to report, and three errors to avoid}\label{sec:metrics}

\subsubsection{Definitions}

With TP, FP, TN and FN the entries of the confusion matrix,
\begin{gather}
  \text{Precision} = \frac{TP}{TP+FP}, \qquad
  \text{Recall (sensitivity)} = \frac{TP}{TP+FN}, \qquad
  \text{Specificity} = \frac{TN}{TN+FP},\\
  F_1 = \frac{2\,\text{Precision}\cdot\text{Recall}}{\text{Precision}+\text{Recall}}, \qquad
  \text{Accuracy} = \frac{TP+TN}{TP+TN+FP+FN}.
\end{gather}
The area under the ROC curve is the probability that a randomly chosen positive is ranked above a randomly chosen negative; 0.5 is chance.

\subsubsection{The accuracy trap}\label{sec:accuracy-trap}

With prevalence $\pi$, accuracy decomposes as
\begin{equation}
  \text{Accuracy} = \pi\,\text{Sensitivity} + (1-\pi)\,\text{Specificity}.
\end{equation}
When $\pi$ is small, specificity dominates and a model that never predicts the positive class attains accuracy $1-\pi$. The framing document states that in the target data critical states are rare. Accuracy must therefore not be the headline metric of this thesis.

\Textcite{mendoza2025} illustrate the trap concretely. Their logistic regression achieves an accuracy of 0.70, which is approximately the prevalence of controlled hypertension in their sample ($912/1308 = 0.697$), while its specificity is 0.03. It is a model that has learned to answer the same way almost every time.

\subsubsection{First error: the positive class is not stated, and is not consistent}\label{sec:positive-class}

Applying the decomposition above to \cref{tab:mendoza-models}, with uncontrolled prevalence $\pi = 396/1308 = 0.303$:
\begin{align}
  \text{random forest:}\quad & 0.303 \times 0.51 + 0.697 \times 0.98 = 0.838 \approx 0.84 \;\checkmark\\
  \text{XGBoost:}\quad & 0.303 \times 0.23 + 0.697 \times 0.85 = 0.662 \approx 0.65 \;\checkmark\\
  \text{logistic regression:}\quad & 0.303 \times 0.98 + 0.697 \times 0.03 = 0.318 \neq 0.70.
\end{align}
The $F_1$ values confirm the first two: $2(0.94)(0.51)/1.45 = 0.661$ against a reported 0.659, and $2(0.42)(0.23)/0.65 = 0.297$ against a reported 0.294. For the logistic regression row, taking \emph{controlled} hypertension as the positive class instead reproduces every published value: accuracy $0.697 \times 0.98 + 0.303 \times 0.03 = 0.692 \approx 0.70$, precision equal to the controlled prevalence 0.697 against a reported 0.70, and $F_1 = 2(0.70)(0.98)/1.68 = 0.817$ against a reported 0.818.

The three rows of the same table are therefore not computed with the same positive class. Two treat uncontrolled hypertension as positive; the third treats controlled hypertension as positive. The paper never states which convention is used. The consequence is that the logistic regression row is not comparable with the other two, and its apparently excellent $F_1$ of 0.818 describes performance on the clinically uninteresting class.

\subsubsection{Second error: sensitivity reported as specificity}\label{sec:sens-spec}

The Results section of \textcite{mendoza2025} states correctly that \bq{Random Forest achieved the highest precision (0.94, 95\% CI 0.91--0.97), minimizing false positives, whereas Logistic Regression demonstrated high sensitivity (0.98, 95\% CI 0.98--0.99)}.

The Discussion then states, of random forest, that \bq{its precision (0.94, 95\% CI 0.91--0.97) was notably high, ensuring reliable identification of at-risk individuals, while its sensitivity (0.98, 95\% CI 0.97--0.99) confirmed its effectiveness in detecting cases of uncontrolled HTN}.

\Cref{tab:mendoza-models} gives random forest a sensitivity of 0.51 (0.46--0.55) and a specificity of 0.98 (0.97--0.99). The Discussion has substituted the second for the first.

The correction is material. Random forest detects 51\% of uncontrolled cases. It misses almost half of them. The claim that its sensitivity confirmed its effectiveness in detecting cases of uncontrolled hypertension is the reverse of what the paper's own table shows. This thesis cites the table and not that sentence.

For a triage system this is the central result of the section. The best model in the only genuine prediction study available misses half the patients it exists to find.

\subsubsection{Third error: loss curves for models that do not minimise a loss by iteration}

\Textcite{efendi2025} report a training loss over 300 epochs for the decision tree, random forest and Bayesian classifier alongside AdaBoost and XGBoost. The first three are not fitted by iterative gradient minimisation and have no epoch-indexed loss. The quantity plotted is therefore not defined for those models, and the comparison of loss values across the five---0.025 for XGBoost, 0.062 for the Bayesian classifier---is not reproduced here.

\subsubsection{Confidence intervals}

\Textcite{mendoza2025} report an interval for every metric; \textcite{efendi2025} report none. The difference matters. An area under the curve of 0.54 would be reported as a number slightly above chance; the interval 0.49--0.60 establishes that it is not distinguishable from chance (\cref{fig:mendoza-auc}). With the sample the framing document specifies, intervals in the present work will be wide, and reporting them is what makes a modest result honest rather than a weak one disguised.

%==========================================================================
\subsection{Interpretability}\label{sec:interpretability}

\Textcite{mendoza2025} apply two complementary approaches, and both are directly transferable.

\textbf{Random forest variable importance}, by percentage increase in mean squared error on permutation and by increase in node purity. The two rankings differ, which is itself instructive: by permutation the order is years since diagnosis (11.9), age (9.4), source of medical care (4.6), then diabetes history and sex; by node purity it is age (189.8), source of medical care (42.1), years since diagnosis (32.9), then diabetes history and sex.

\textbf{SHAP values for XGBoost}, giving the marginal contribution of each predictor: age 0.115, source of medical care 0.065, years since diagnosis 0.055, sex 0.035, with diabetes history lowest.

Two points for the present work. First, variable importance and SHAP need not agree, and reporting only one invites the suspicion that it was chosen. Second---and this requires care---SHAP explains a model, not a phenomenon. Here it explains an XGBoost model whose discrimination is not distinguishable from chance. Explaining the decisions of a model that does not discriminate is an exercise of limited value, and the paper does not note this. If the present work reports SHAP, it must do so for a model whose discrimination has first been shown to be adequate.

\Textcite{liu2025review} observe more generally that the black-box nature of machine learning is among its principal limitations, and that regulatory and clinical acceptance depends on interpretability. \Cref{sec:specialist-access} reached the same conclusion from the clinical side: a triage system is adopted only if the clinician can understand why a patient was flagged and, more importantly, trust the ones it did not flag.

%==========================================================================
\subsection{Class imbalance and the rare critical state}\label{sec:imbalance}

The framing document anticipates this correctly: it identifies strong class imbalance, critical states being rare, as a data limitation, and proposes as methodological responses a metric centred on recall of the critical class, and oversampling.

The results of \textcite{mendoza2025} demonstrate why this is not optional. At a prevalence of 30.3\%---which is not extreme---the sensitivity of random forest falls to 0.51 and that of XGBoost to 0.23. At the prevalence the framing document anticipates for critical states, considerably lower than 30\%, the untreated behaviour of these algorithms will be to predict the majority class almost always.

The design response has three components, and the first is the most important.

\textbf{Choose the operating point from the cost of error, not from the default threshold.} A classifier outputs a score; the 0.5 threshold is a convention with no clinical meaning. For triage, a false negative is a patient in deterioration who is not seen; a false positive is a clinician's time spent on a patient who turns out to be stable. \Cref{sec:specialist-access} established that clinician time is the scarce resource, so false positives are not free---but \cref{sec:cascade} established that the dominant failure of the existing system is non-detection. The threshold should be set to a target sensitivity on the critical class and the resulting alert volume reported, rather than the reverse.

\textbf{Weight or resample during training}, as the framing document proposes, while noting that resampling changes the calibration of the predicted probabilities and that any probability reported to a clinician must be recalibrated afterwards.

\textbf{Report per-class metrics, never aggregates alone.} A macro-averaged $F_1$ conceals exactly the failure that matters. The random forest row of \cref{tab:mendoza-models}---precision 0.94, sensitivity 0.51---shows what a single aggregate would have hidden.

%==========================================================================
\subsection{From classification to forecasting}\label{sec:forecasting}

The hypotheses of the framing document can now be assessed against the state of the art.

\subsubsection{The three hypotheses against the evidence}

\textbf{H1: a deep temporal model obtains better recall on high and critical risk states than robust classical models, on locally collected telemonitoring data.} The only comparable datum in the corpus is the recurrent network of Reddy and Delen for lupus readmission, reported by \textcite{liu2025review} as outperforming traditional classification with an area under the curve of 0.70. This is a modest figure, and it is the realistic order of magnitude for a genuine clinical prediction task on a sequence. It supports the plausibility of H1 while setting expectations well below the 95--100\% band that dominates the ambient literature.

\textbf{H2: short-horizon forecasting of physiological parameters allows anticipation of degradation events with a clinically useful lead time.} No source in this corpus attempts this. The architectures of \cref{sec:cloud} provide for it---\textcite{mucchi2021} assign the cloud layer long-term data analysis through advanced algorithms for data cross-correlation---but none implements it. H2 is therefore untested in the available literature, which is the strongest claim to novelty this thesis can make, and also its principal risk.

\textbf{H3: given the limited size and irregularity of the local data, the gain of the deep model over classical models may remain modest.} This is the hypothesis best supported by the state of the art. \Textcite{liu2025review} note that deep learning methods present high computational resource demands, increased model complexity and the need for data privacy protection, and that efficacy is contingent upon the representation of the dataset rather than its volume. With the sample size the framing document targets, a model with the parameter count of a recurrent network is at substantial risk of overfitting, and a negative result would be a genuine finding rather than a failure.

\subsubsection{The horizon must match the sampling interval}\label{sec:horizon}

One constraint on H2 follows directly from the acquisition decision of \cref{sec:acq-decision} and is not stated in the framing document, although it governs whether the hypothesis is testable at all.

Under the spare acquisition adopted, blood pressure and blood glucose arrive through pipeline 2 at the routine frequency established in \cref{sec:protocols}: of the order of two readings per week. The expected number of readings falling inside a forecast horizon of length $\Delta$ is then
\begin{equation}
  \mathbb{E}[N_\Delta] = \lambda\Delta, \qquad \lambda = \frac{2}{168\,\text{h}},
\end{equation}
which gives $\mathbb{E}[N_{24\text{h}}] = 0.29$ and $\mathbb{E}[N_{48\text{h}}] = 0.57$. Fewer than one reading falls inside the horizon on average, so for most prediction points there is no observation to predict (\cref{fig:sampling-horizon}). Pipeline 1 is worse still: at one to two consultations per year it contributes nothing at this horizon. A 24 to 48 hour forecast is therefore not estimable from routine-frequency data.

\begin{figure}[htbp]
\centering
\begin{tikzpicture}
\begin{axis}[width=0.76\textwidth,height=5cm,
  xlabel={Forecast horizon $\Delta$ (hours)},
  ylabel={Expected readings in horizon},
  xmin=0,xmax=180,ymin=0,ymax=5.2,
  grid=major,grid style={gray!20},legend pos=north west,legend style={font=\scriptsize}]
\addplot[thick,cAfrica,domain=0:180,samples=2]{2*x/168};
\addplot[thick,cRef,domain=0:180,samples=2]{14*x/168};
\addplot[dashed,cGrey,domain=0:180,samples=2]{1};
\draw[dotted,thick,cGrey] (axis cs:48,0) -- (axis cs:48,5.2);
\node[font=\scriptsize,text=cGrey,anchor=west,rotate=90] at (axis cs:49,2.6) {48~h};
\legend{routine, 2 readings/week, intensified profile, 2 readings/day, one reading}
\end{axis}
\end{tikzpicture}
\caption{Expected number of observations inside a forecast horizon, by sampling regime. At routine frequency the 48-hour horizon contains 0.57 readings on average; the intensified verification profile of \cref{sec:triage} contains four.}\label{fig:sampling-horizon}
\end{figure}

Three resolutions are available and one must be adopted. The horizon may be lengthened to match the sampling interval, predicting the next scheduled reading or the seven-day mean that \cref{sec:protocols} defines as the reference value---which preserves the anticipatory character while making the target estimable. The routine frequency may be raised to daily, which costs adherence against the trade-off set out in \cref{sec:protocols}. Or the forecast may be defined only within the intensified verification regime, where two readings per day give $\mathbb{E}[N_{48\text{h}}] = 4$---which is coherent with the two-stage architecture of \cref{sec:alerting}, since that is precisely the regime in which a clinical decision is pending.

The third is the most defensible, because it aligns the forecasting task with the moment at which its output would change a decision. This thesis adopts it, and H2 is restated accordingly: short-horizon forecasting applies to patients in the verification regime, and the routine regime is served by trend detection over the longer interval rather than by a 48-hour forecast.%
\verify{Confirm the routine sampling frequency actually achieved in the collected cohort before fixing the horizon. If observed adherence yields fewer than two readings per week, even the seven-day target becomes marginal and the design must fall back on consultation-anchored prediction.}

\subsubsection{A methodological commitment to add}

The framing document specifies group $k$-fold cross-validation by patient, which is correct and prevents the most common leakage in clinical modelling. For a forecasting task it is not sufficient: the split must also be temporal. Training on later data and testing on earlier data leaks the future, even across patients, if any global preprocessing---normalisation, imputation, feature scaling---is fitted on the full dataset. The split must therefore be by patient and forward in time, and every preprocessing step must be fitted on the training fold alone.

%==========================================================================
\subsection{Specification for the experimental phase}\label{sec:spec23}

\begin{table}[htbp]
\centering\footnotesize
\caption{Requirements for the experimental phase established in \cref{sec:ml}.}\label{tab:spec23}
\begin{tabularx}{\textwidth}{@{}X >{\raggedright\arraybackslash}p{5.9cm}@{}}
\toprule
\textbf{Requirement} & \textbf{Basis} \\
\midrule
Target is $\hat{y}_{t+\Delta}$, not $\hat{y}_t$ & label leakage in \textcite{efendi2025}; \cref{sec:principle} \\
Horizon matched to the sampling regime; forecast defined in the verification regime & \cref{sec:horizon} \\
No feature computed from the same instant as the label & \cref{sec:principle} \\
Preprocessing fitted on the training fold only & \cref{sec:forecasting} \\
Split by patient and forward in time & framing document, extended \\
Baselines in order: majority class, threshold rule, logistic regression, random forest, XGBoost, LSTM & \textcite{mendoza2025} include logistic regression as reference \\
The clinical threshold rule is a mandatory baseline & a model that does not beat the ESH and ADA thresholds has no justification \\
Headline metric: recall on the critical class & framing document; \cref{sec:imbalance} \\
Accuracy never reported alone & \cref{sec:accuracy-trap} \\
Positive class stated explicitly for every metric & \cref{sec:positive-class} \\
Per-class metrics reported, not only aggregates & \cref{sec:sens-spec} \\
Confidence intervals on every metric & \textcite{mendoza2025}; absent in \textcite{efendi2025} \\
AUC reported with interval, chance level marked & XGBoost interval 0.49--0.60 includes 0.50 \\
Operating threshold chosen from error cost, then alert volume reported & \cref{sec:imbalance} \\
Probabilities recalibrated after any resampling & \cref{sec:imbalance} \\
Interpretability reported only for a model with adequate discrimination & \cref{sec:interpretability} \\
Both permutation importance and SHAP, with disagreement discussed & \textcite{mendoza2025} \\
Fixed random seed, stated & \textcite{mendoza2025}: seed 123 \\
Code, feature definitions and labelling rule published & \textcite{mendoza2025} archived data and code \\
Training configuration appropriate to each model family & \cref{sec:task-not-algo} \\
\bottomrule
\end{tabularx}
\end{table}

The penultimate line deserves a note. \Textcite{mendoza2025} deposited their dataset and their code in public repositories. The framing document states that the data are real, anonymised and supplied by a partner physician, with direct identifiers removed before any processing, collected with the agreement of the physician and the institution for research purposes only. Full data publication will not be possible. Publishing the code, the feature definitions and the labelling rule is possible, and is committed to here, since it is what allows the result to be checked without exposing the patients.

%==========================================================================
\subsection{What this section establishes}\label{sec:ml-conclusion}

\textbf{Established.} The algorithm families and their inductive biases are well characterised, and only one of the candidates treats the input as a sequence. Reported performance in this literature is governed by task construction far more than by algorithm choice, and the two principal sources demonstrate both poles of that observation: a task solved by construction yields 0.97, and a genuine prediction task from weak features yields 0.54 to 0.75. The metric apparatus required for an honest report is available and is largely absent from the sources.

\textbf{Not established.} No source predicts future physiological state. No source works on a clinical time series of blood pressure and blood glucose. No source operates in a low-resource setting. The best sensitivity achieved on genuine prediction of uncontrolled hypertension anywhere in this corpus is 0.51, and the model that achieved it is a random forest---not the gradient-boosted ensemble that the surrounding literature treats as the default choice.

%##########################################################################
% SECTION 2.4
%##########################################################################
\section{Clinical Decision Support Systems and Alert Generation}\label{sec:cdss}

\subsection*{Purpose and status of the sources}

\Cref{sec:acquisition,sec:cloud,sec:ml} established what is acquired, how it travels, and what is computed on it. This section addresses the last link: how a computed score becomes a clinical action. It is the section in which the triage requirement established in \cref{ch:background} and located architecturally in \cref{sec:edge-triage} is finally specified as a mechanism.

\paragraph{Yazdani et al. (2018).} \textcite{yazdani2018}, examined in \cref{sec:multitier}, is the only source in this corpus whose explicit subject is the coupling of telemonitoring with clinical decision support. Its limitations stand: a UML design with no implementation and no evaluation, using XML and SMS. Its contribution is a complete decision flow and a modular decomposition, and these remain valid independently of the technology in which they were expressed.

\paragraph{Liu and Wang (2025).} \textcite{liu2025review} describe a multi-point trigger mechanism coupling a rule engine to machine learning. The reservation on their method stands; we use the source for the techniques it names, each of which is independently verifiable, and not as pooled evidence.

\paragraph{Zhou et al. (2024).} \textcite{zhou2024}, examined in \cref{sec:efficacy}, supplies the only quantitative evidence in this corpus on the effectiveness of prompts, feedback and reminders---that is, on the patient-facing half of the feedback loop. It says nothing about alert mechanisms directed at clinicians, which is the principal subject of this section, and it is not cited for that purpose.

\paragraph{Mihevc et al. (2025).} \textcite{mihevc2025}, examined in \cref{sec:mihevc}, supplies the only clinically validated alert protocol.

%==========================================================================
\subsection{What a decision support telemonitoring system is}\label{sec:dsts-def}

\subsubsection{Definition}

\Textcite{yazdani2018} give the definition we adopt: \bq{Decision Support Telemedicine Systems (DSTS) is a combination of clinical decision support systems and telemedicine. A hybrid system that transmits information from the patient's home to the health monitoring center and allows them to support their monitoring data is an example of DSTS.}

They situate its necessity in a problem \cref{ch:background} established independently: \bq{producing a large amount of patient data that requires careful evaluations can lead to unacceptable increase in workloads of physicians. Therefore, a combined decision support tools and telemedicine can be used as a possible solution to overcome human resource limitation and analyze generated data.}

This is the same argument reached in \cref{sec:specialist-access} from the clinical side---the six calls per provider per day of \textcite{kingue2013}, the 18\% of countries without a cardiologist reported by \textcite{segun2025}---arriving here from the informatics side. Decision support is not an enhancement of telemonitoring; it is the condition under which telemonitoring does not consume more clinical time than it saves.

\subsubsection{The scalability problem}

The stated contribution of \textcite{yazdani2018} is a diagnosis worth retaining. They observe that existing decision-support telemonitoring architectures are \bq{largely confined to a specific disease and it means these architectures are not scalable}, so that \bq{if that architecture is to be used to monitor the health status of other diseases, it should be re-established with the new knowledge base (the new decision support system). This issue will result in the creation of multiple island systems for monitoring the status of patients.}

Their response is to separate the platform from the knowledge base, so that several decision support systems can be loaded onto a common infrastructure, each with its own rule set, selected per patient according to disease and symptoms.

For a thesis whose object covers two pathologies with two distinct guideline sources (\cref{sec:pipelines}), this separation is not an abstraction. A system in which the hypertension thresholds are hard-coded alongside the glycaemic thresholds is one that must be rewritten to add a third condition. One in which the rule base is a configurable artefact is extensible by configuration.

\subsubsection{Three families of decision support}

\begin{table}[htbp]
\centering\footnotesize
\caption{Families of clinical decision support and their trade-offs.}\label{tab:cdss-families}
\begin{tabularx}{\textwidth}{@{}>{\raggedright\arraybackslash}p{2.0cm} >{\raggedright\arraybackslash}p{3.2cm} >{\raggedright\arraybackslash}X >{\raggedright\arraybackslash}X@{}}
\toprule
\textbf{Family} & \textbf{Basis} & \textbf{Strength} & \textbf{Weakness} \\
\midrule
Rule-based & explicit thresholds from clinical guidelines & transparent, auditable, no training data, defensible to a clinician & no personalisation, no trend detection, coverage gaps \\
Model-based & learned from data & personalisation, trend detection, adaptive & opaque, needs data, may not transfer \\
Hybrid & rules as a floor, model as a refinement & guideline conformity plus personalisation & two components to validate \\
\bottomrule
\end{tabularx}
\end{table}

The sources distribute across these. \Textcite{chatrati2022} are purely rule-based for hypertension and purely model-based for diabetes, with no coordination between the two (\cref{sec:chatrati}). \Textcite{efendi2025} are purely model-based, on a model that reproduces a rule (\cref{sec:leakage}). \Textcite{mihevc2025} are purely rule-based, and are the only ones clinically validated. \Textcite{liu2025review} advocate the hybrid: \bq{by integrating a sophisticated rule engine with machine learning algorithms, the chronic disease management platform is capable of dynamically monitoring and conducting real-time analysis of patients' health status}.

The hybrid is the position this thesis adopts, and \cref{sec:spec23} already established why in a different form: the clinical threshold rule is a mandatory baseline, because a model that does not outperform the guideline thresholds has no justification for existing. Making the rule a component rather than merely a baseline means the system degrades to guideline conformity when the model is uncertain, rather than degrading to nothing.

%==========================================================================
\subsection{The reference decision flow}\label{sec:decision-flow}

\subsubsection{The activity diagram of Yazdani et al.}

Their specification covers three actors---patient as client, supervisor as server, physician as client---and it is the most detailed account of a telemonitoring decision cycle in this corpus. Traced step by step:

\textbf{Acquisition.} \bq{Patient data is collected by sensors or manually by the patient interface. The data are proportionate with the intervals that the physician has determined, and are sent to the server for storage and processing.}

\textbf{Non-transmission handling.} \bq{If no data is sent within the time limits set by the physician, a SMS reminder will be sent to the patient. The SMS reminder is sent repeatedly as long as the patient enters a new data.}

\textbf{Verification and routing.} \bq{The server, after verifying the data and store them in the database, automatically analyzes the data to determine the related CDSS.}

\textbf{Classification.} \bq{After determining the CDSS, patient data will be analyzed. The result of this analysis is to determine the patient's condition in one of the two Normal or Alert status.}

\textbf{Branch on normal.} The patient is notified, the record is updated, no clinician is involved.

\textbf{Branch on alert.} A treatment proposal is generated and the physician notified: \bq{the CDSS in the Alert mode displays the proposed prescription to the physician. After reviewing the patient's data, the physician can approve the proposed CDSS prescription or prescribe his own recommendation.}

\textbf{Physician sub-decision.} Three outcomes are modelled: therapy unchanged; new therapy required, in which case the physician prescribes; or more information needed, in which case \bq{the patient will be informed and coordinated to have face to face meeting with the physician through the patient interface}.

\textbf{Closure.} \bq{All events in this telehealth monitoring architecture, including physician and CDSS prescriptions, are stored as patient record in the database.}

\subsubsection{The three decision points}

\begin{table}[htbp]
\centering\small
\caption{Decision points of the reference flow \parencite{yazdani2018}.}\label{tab:decision-points}
\begin{tabularx}{\textwidth}{@{}>{\raggedright\arraybackslash}p{5.0cm} >{\raggedright\arraybackslash}p{2.2cm} X@{}}
\toprule
\textbf{Point} & \textbf{Decided by} & \textbf{Output} \\
\midrule
Did data arrive on schedule? & system & reminder, repeated until compliance \\
Is the patient normal or alert? & CDSS & notify patient, or escalate with proposal \\
Is the proposal accepted? & physician & approve, override, or request consultation \\
\bottomrule
\end{tabularx}
\end{table}

The third is the one that matters most and is the most frequently omitted in engineering designs. The system proposes; the physician disposes. No prescription reaches the patient without passing a clinician. \Textcite{yazdani2018} do not argue for this; they simply build it in. The argument is supplied by \cref{sec:ml}: the best model of \textcite{mendoza2025} detects 51\% of uncontrolled cases and the second best is indistinguishable from chance. A system that acted autonomously on such a model would be acting on a coin flip half the time.

The first decision point is also significant and connects to \cref{ch:background}. Non-transmission is itself an event that triggers an action. \Cref{sec:ltfu} established loss to follow-up at 50\% or more in low-resource settings, and \cref{sec:bediang} measured 33.3\% three-month attendance in Cameroonian control centres. A system that only reacts to data received is blind to exactly the patients most at risk. The repeated SMS of \textcite{yazdani2018}, and the HEARTS indicator for patients with no visit in the last three and twelve months (\cref{sec:hearts-spec}), are the same mechanism: silence must be alarmed.

\subsubsection{Comparison with the only validated protocol}

\begin{table}[htbp]
\centering\footnotesize
\caption{The reference flow against the only clinically validated protocol.}\label{tab:flow-compare}
\begin{tabularx}{\textwidth}{@{}>{\raggedright\arraybackslash}p{2.6cm} >{\raggedright\arraybackslash}X >{\raggedright\arraybackslash}X@{}}
\toprule
 & \textbf{\textcite{yazdani2018}} & \textbf{\textcite{mihevc2025}} \\
\midrule
Classification & two states: normal, alert & three states: green, yellow, red \\
On normal & notify patient & no change \\
On intermediate & --- & verification profile: 7-day BP profile or 1-day 6-point glucose profile \\
On critical & notify physician with proposal & immediate teleconsultation \\
Who decides & physician approves or overrides & general practitioner evaluates, may modify therapy \\
Validated & no & yes: RCT, $-6.9$~mmHg, $-0.5\%$ HbA1c \\
\bottomrule
\end{tabularx}
\end{table}

The yellow tier is the critical addition. In a two-state system, every reading that is not normal escalates to a clinician. In the three-state system, an abnormal reading escalates to more data, and only the aggregate of that data escalates to a clinician. This is the architectural two-stage split of \cref{sec:edge-triage} expressed as a clinical protocol, and it is the single most transferable design element in the corpus---not least because it is the only one attached to a demonstrated clinical effect.

\subsubsection{The convergent flow}

Combining the two, with the model placed where \cref{sec:ml} established it belongs, gives the flow of \cref{fig:decision-flow}.

\begin{figure}[htbp]
\centering
\begin{tikzpicture}[node distance=0.45cm and 0.95cm,
  b/.style={draw,rounded corners,align=center,font=\scriptsize,minimum height=0.8cm,minimum width=2.3cm},
  d/.style={draw,diamond,aspect=2.2,align=center,font=\scriptsize,inner sep=0pt,fill=cGrey!12},
  ar/.style={-{Stealth[length=2mm]},thick}]
  \node[b,fill=cGrey!14] (r) {reading};
  \node[d,right=of r] (rule) {rule};
  \node[b,fill=green!14,above right=0.35cm and 0.9cm of rule] (g) {green:\\ record};
  \node[b,fill=yellow!28,right=0.9cm of rule] (y) {yellow:\\ verification profile};
  \node[b,fill=cAfrica!18,below right=0.35cm and 0.9cm of rule] (rd) {red:\\ escalate now};
  \node[d,right=of y] (mod) {model\\ on series};
  \node[b,fill=cRef!14,right=0.9cm of mod] (esc) {prioritised\\ queue};
  \node[b,fill=cGrey!20,right=0.9cm of esc] (phy) {physician:\\ approve / override /\\ request consultation};
  \draw[ar] (r) -- (rule);
  \draw[ar] (rule) -- (g);
  \draw[ar] (rule) -- (y);
  \draw[ar] (rule) -- (rd);
  \draw[ar] (y) -- (mod);
  \draw[ar] (mod) -- node[above,font=\tiny]{persists} (esc);
  \draw[ar] (mod.north) |- node[near end,above,font=\tiny]{resolves} (g.east);
  \draw[ar] (rd.east) -| ([xshift=0.45cm]rd.east) |- (esc.south);
  \draw[ar] (esc) -- (phy);
  \draw[ar,dashed,cAfrica] (phy.north) -- ++(0,0.75) -| node[near start,above,font=\tiny,text=cAfrica]{decision stored as adjudicated label} (mod.north east);
\end{tikzpicture}
\caption{The convergent decision flow: rule-based triage at the edge, model-based assessment of the verification profile, prioritised escalation, and clinician disposition. The dashed return path is the closed loop of \cref{sec:closed-loop}.}\label{fig:decision-flow}
\end{figure}

A parallel branch runs alongside: non-transmission within the configured interval raises a reminder, repeated, and escalates as non-adherence if unanswered.

%==========================================================================
\subsection{The rule engine coupled to the model}\label{sec:rule-model}

\Textcite{liu2025review} provide the only technical specification of this coupling in the corpus. Four elements are directly usable.

\subsubsection{Characterise the series before choosing the algorithm}\label{sec:ts-diagnostics}

\bq{The process of feature analysis necessitates the identification of variation patterns in time-series data derived from historical monitoring indicators, such as periodicity, drift, and stationarity, to select the most suitable algorithms. For example, autocorrelation functions can be used to identify periodic changes, while median smoothing and trend component detection can capture data drift. Stationarity changes are assessed using unit root tests to determine if the time series is stationary.}

This prescribes a diagnostic step that precedes modelling. Blood pressure has a known circadian structure---the nocturnal dipping reported in \cref{sec:hta-burden}---so periodicity is expected; therapeutic drift is the clinical quantity of interest (\cref{sec:records}); and stationarity determines whether differencing is required before any trend feature is meaningful. This diagnostic belongs in the experimental chapter before any model is fitted, and it is the kind of step that distinguishes a considered analysis from an undirected sweep of classifiers.

A caution attaches. These tools assume regularly sampled series. The framing document anticipates irregular and missing measurements from intermittent connectivity, and \cref{sec:pipelines} established that pipeline 2 has missingness that is not at random. Autocorrelation and unit-root tests on irregularly spaced data require either resampling onto a regular grid, which fabricates observations, or methods designed for unevenly spaced series. The choice must be stated, because resampling a series of two true readings per week onto a daily grid manufactures five interpolated points per week, and any trend feature computed on that grid is five-sevenths fiction.

\subsubsection{Anomaly detection as a distinct problem}

\bq{When selecting an algorithm, it is essential to choose the appropriate anomaly detection algorithm based on the data's distribution characteristics. Commonly utilized algorithms in this field include box plots, absolute median differences, and extreme value theory, among others.}

These are unsupervised and distribution-based, and they complement rather than compete with the supervised classifier of \cref{sec:ml}. The median absolute deviation in particular is directly useful for the per-patient normalisation the framing document proposes, since it is robust to the outliers a telemonitoring series will contain:
\begin{equation}
  \mathrm{MAD}_i = \operatorname{median}\bigl(\lvert x_{i,t} - \operatorname{median}(x_i)\rvert\bigr),
  \qquad
  z_{i,t} = \frac{x_{i,t} - \operatorname{median}(x_i)}{1.4826\,\mathrm{MAD}_i}.
\end{equation}
This yields the personalised threshold identified at the close of \cref{sec:spec13} as the gap in the fixed-threshold protocol of \textcite{mihevc2025}: a deviation is assessed relative to the individual patient's own distribution rather than against a population constant. It is also the empirical counterpart of the quantity $\sigma_{\varepsilon,i}$ introduced in \cref{sec:variability}.

\Textcite{liu2025review} add that \bq{machine learning algorithms, such as random forests or autoencoders, are employed to promptly capture anomalies in data and trigger early warning systems}.

\subsubsection{The offline and online split}

\bq{Two crucial aspects of system design include model training and real-time detection. Real-time detection typically employs stream processing technologies such as Flink for online anomaly detection, while offline training involves reading training data from a data warehouse, training the model, and saving it for use in real-time detection.}

This is the standard separation between a batch training path and a low-latency serving path, and it maps directly onto the architecture of \cref{sec:cloud}: offline training in the cloud layer, where \textcite{mucchi2021} place long-term data analysis, and online detection at the edge, where they place quick processing and identification of anomalous conditions.

The specific technology named is not appropriate here. A stream processing framework is dimensioned for throughput that a cohort transmitting a few readings per week will never approach. The separation is the transferable element; the tooling is not. For this work, offline training on a periodic schedule and online evaluation of a saved model at the point of reception is the same architecture at the appropriate scale.

\subsubsection{The adaptive trigger}

\bq{An adaptive triggering strategy, based on multi-source data, is proposed. This strategy involves real-time monitoring of patients' health statuses by integrating data from various channels such as electronic health records, wearable device monitoring data, lifestyle habits, and biomarkers. The comprehensive health profile established for each patient is then transmitted to a central processing system, thereby creating an intelligent multi-point triggering disease prediction and assessment system.}

The term \emph{multi-point} carries two meanings in their framework, and the distinction is useful: multiple data sources triggering jointly, and multiple monitoring points---hospital, community, family---contributing. The first is precisely the three-pipeline integration specified in \cref{sec:pipelines}.

\Textcite{liu2025review} also report a directly relevant precedent: Zhang and colleagues developed an ontology-based framework amalgamating type~2 diabetes patient data with medical domain knowledge and assessment criteria, and a clinical decision support system that \bq{autonomously selects and modifies standard assessment protocols to align with individual patient conditions}. This is the rule-selection module of \textcite{yazdani2018} implemented, and it is the closest published instance of the configurable-knowledge-base design.

%==========================================================================
\subsection{Alert prioritisation, and the arithmetic of alert volume}\label{sec:alert-volume}

This subsection contains the most important quantitative result of the chapter, and no source in this corpus states it.

\subsubsection{Prioritisation as specified}

\Textcite{yazdani2018}: \bq{the physician can have access to patient data at any time through the patient status and is able to view patients' status on a priority basis. This prioritisation is done based on the patient status (Normal and Alert).} They report a precedent in which a system identifies patients with gestational diabetes and prioritises them on the basis of their metabolic status, sending automatic recommendations from a rule-based database but routing insulin dosing decisions to the physician first.

The principle is sound but the resolution is coarse. A binary sort produces two buckets, and when the alert bucket contains more patients than the clinician can see in a session, the sort has not helped: it has merely relabelled the problem. Prioritisation requires an ordering, not a partition, and an ordering requires a continuous score---which is what the model of \cref{sec:ml} produces and what the rule of \textcite{mihevc2025} does not.

\subsubsection{The arithmetic}\label{sec:volume-arithmetic}

Consider a cohort of $N$ patients each transmitting $m$ readings per week, with $\pi$ the proportion of readings that genuinely warrant clinical attention, and a classifier of sensitivity $\mathrm{Se}$ and specificity $\mathrm{Sp}$. Alerts per week number
\begin{equation}
  A = N m \bigl[\pi\,\mathrm{Se} + (1-\pi)(1-\mathrm{Sp})\bigr],
\end{equation}
of which the proportion that are genuine is the positive predictive value
\begin{equation}
  \mathrm{PPV} = \frac{\pi\,\mathrm{Se}}{\pi\,\mathrm{Se} + (1-\pi)(1-\mathrm{Sp})}.
\end{equation}

The second term of the denominator is the one that governs, and it does not depend on sensitivity at all. False alerts are generated by the volume of normal readings multiplied by the false-positive rate, and in a monitoring system the volume of normal readings is overwhelming.

Substituting the measured values of \textcite{mendoza2025} at their prevalence $\pi = 0.303$ reproduces their reported precisions, which confirms the model: for random forest, $0.303(0.51) / [0.303(0.51) + 0.697(0.02)] = 0.92$ against a reported 0.94; for XGBoost, $0.303(0.23)/[0.303(0.23) + 0.697(0.15)] = 0.40$ against a reported 0.42.

Now apply it to a telemonitoring setting, where critical events are rare. Taking $\pi = 0.02$---two per cent of readings genuinely critical, which is optimistic given the framing document's statement that critical states are rare---gives \cref{tab:ppv}.

\begin{table}[htbp]
\centering\small
\caption{Positive predictive value and false-alert burden at $\pi = 0.02$.}\label{tab:ppv}
\begin{tabularx}{\textwidth}{@{}c c X c@{}}
\toprule
\textbf{Se} & \textbf{Sp} & \textbf{PPV at $\pi = 0.02$} & \textbf{False alerts per 1\,000 readings} \\
\midrule
0.51 & 0.98 & 0.34 & 19.6 \\
0.90 & 0.98 & 0.48 & 19.6 \\
0.90 & 0.99 & 0.65 & 9.8 \\
0.90 & 0.995 & 0.79 & 4.9 \\
0.90 & 0.999 & 0.95 & 1.0 \\
\bottomrule
\end{tabularx}
\end{table}

\begin{figure}[htbp]
\centering
\begin{tikzpicture}
\begin{axis}[width=0.82\textwidth,height=6cm,
  xlabel={Prevalence of genuinely critical readings, $\pi$},
  ylabel={Positive predictive value},
  xmin=0,xmax=0.35,ymin=0,ymax=1.02,
  grid=major,grid style={gray!20},
  legend pos=south east,legend style={font=\scriptsize}]
\addplot[thick,cAfrica,domain=0.002:0.35,samples=200]{0.9*x/(0.9*x + (1-x)*0.02)};
\addplot[thick,cRef,domain=0.002:0.35,samples=200]{0.9*x/(0.9*x + (1-x)*0.01)};
\addplot[thick,cGrey,domain=0.002:0.35,samples=200]{0.9*x/(0.9*x + (1-x)*0.005)};
\addplot[dashed,black,domain=0:0.35,samples=2]{0.8};
\draw[dotted,thick,cGrey] (axis cs:0.02,0) -- (axis cs:0.02,1.02);
\node[font=\scriptsize,anchor=west,rotate=90,text=cGrey] at (axis cs:0.022,0.45) {$\pi=0.02$};
\legend{$\mathrm{Sp}=0.98$,$\mathrm{Sp}=0.99$,$\mathrm{Sp}=0.995$}
\end{axis}
\end{tikzpicture}
\caption{Positive predictive value against prevalence at $\mathrm{Se}=0.90$, for three specificities. The dashed line marks an 80\% precision target. At $\pi=0.02$ it is met only beyond $\mathrm{Sp}=0.995$.}\label{fig:ppv-prevalence}
\end{figure}

At two per cent prevalence, reaching 80\% precision requires 99.5\% specificity. Solving for the requirement directly,
\begin{equation}
  \mathrm{Sp} \;\geq\; 1 - \frac{\pi\,\mathrm{Se}}{1-\pi}\cdot\frac{1-\mathrm{PPV}}{\mathrm{PPV}}.
\end{equation}
This is the quantitative form of the triage requirement, and it is considerably more demanding than anything the literature reports. The best model of \textcite{mendoza2025} reaches a specificity of 0.98 at a prevalence fifteen times higher.

\subsubsection{Why sampling frequency decides the problem}\label{sec:frequency-decides}

The term $Nm$ is the multiplier, and it varies across the architectures of this chapter by four orders of magnitude (\cref{tab:alert-designs,fig:alert-volume}).

\begin{table}[htbp]
\centering\small
\caption{False-alert burden by acquisition design, at $\mathrm{Sp} = 0.98$ and $\pi = 0.02$.}\label{tab:alert-designs}
\begin{tabularx}{\textwidth}{@{}X c c@{}}
\toprule
\textbf{Design} & \textbf{Readings per patient per week} & \textbf{False alerts per patient per week} \\
\midrule
\textcite{mihevc2025}: 2 BP + 1 glucose & 3 & 0.06 \\
Manual entry, daily & 7 & 0.14 \\
\textcite{efendi2025}: 1~Hz continuous & 604\,800 & 11\,854 \\
\bottomrule
\end{tabularx}
\end{table}

\begin{figure}[htbp]
\centering
\begin{tikzpicture}
\begin{loglogaxis}[width=0.78\textwidth,height=5.6cm,
  xlabel={Readings per patient per week},
  ylabel={False alerts per patient per week},
  xmin=1,xmax=2e6,ymin=1e-2,ymax=1e5,
  grid=major,grid style={gray!20}]
\addplot[thick,cRef,domain=1:2e6,samples=2,no marks]{0.0196*x};
\addplot[only marks,mark=*,mark size=2.6pt,cAfrica]
  coordinates {(3,0.0588)(7,0.1372)(604800,11854)};
\node[font=\scriptsize,anchor=west] at (axis cs:3.6,0.0588) {Mihevc protocol};
\node[font=\scriptsize,anchor=west] at (axis cs:8.5,0.1372) {daily entry};
\node[font=\scriptsize,anchor=east] at (axis cs:400000,11854) {1~Hz sensor};
\end{loglogaxis}
\end{tikzpicture}
\caption{False-alert burden scales linearly with sampling frequency at fixed specificity. The clinical protocols and the sensor pipeline differ by more than five orders of magnitude on this axis.}\label{fig:alert-volume}
\end{figure}

For a cohort of thirty patients on the protocol of \textcite{mihevc2025}, the system generates roughly two false alerts per week---entirely manageable. The same specificity applied to continuous sensor streaming generates over eleven thousand false alerts per patient per week, which is not a tuning problem but a categorical impossibility.

This resolves a tension that has run through the chapter. \Cref{sec:acquisition} found that the optical sensor pipeline cannot measure the parameters of clinical interest. The present subsection establishes a second and independent reason to restrict its role: per-reading alerting at sensor frequency is arithmetically unusable regardless of model quality. High-frequency data can be aggregated into features---a daily mean, a variability measure, a trend slope---and alerted on at the aggregate level, but it cannot be alerted on reading by reading. The data processing unit at the edge, in the architecture of \textcite{mucchi2021}, is exactly where that aggregation belongs.

\subsubsection{What the sources say about alert burden, and what they do not}

No source in this corpus uses the term \emph{alert fatigue} or measures it. This thesis does not cite these sources for a phenomenon they do not address. What they do establish, which is adjacent and usable:

\Textcite{dodoo2022}, reported in \cref{sec:digital-barriers}, found that health professionals in the region view telemedicine as an additional task entailing increased workload. An alerting system whose precision is low converts this perception into fact.

\Textcite{mucchi2021}, reported in \cref{sec:usability-result}, record nurses requesting that alert notification be improved especially in urgent situations, and that immediate feedback from the doctor in the event of an alert be considered. They were asking for alerts that are fewer and more consequential, not more numerous.

\Textcite{bediang2021}, reported in \cref{sec:bediang}, supply an indirect measurement that deserves notice. In the intervention centres, 23.7\% of local providers' referrals were judged unnecessary by the remote expert, against 11.5\% in control centres ($p = 0.01$). The system that improved access to expertise also doubled the proportion of inappropriate escalations recorded. The authors do not comment on this. It is the closest thing in this corpus to a measured instance of over-escalation accompanying a decision support system, and it is cited here as such---with two reservations. The study is not randomised and the groups differ at baseline. More importantly, the ascertainment mechanism is not symmetric: in the intervention arm a cardiologist reviewed every case and could judge a referral unnecessary, whereas no such reviewer existed in the control arm, so the difference may measure how closely referrals were examined rather than how often they were unwarranted.%
\verify{Establish from \textcite{bediang2021} how referrals in the control centres were adjudicated as unnecessary, given that those centres had no remote expert. If the two arms were assessed by different procedures, the 23.7\% against 11.5\% comparison measures ascertainment and must not be presented as over-escalation.}

\Textcite{who2023hypertension} place an upper bound on tolerable workload from the other direction: the readiness-to-scale criterion requires a missed-visit rate below 25\% (\cref{sec:hearts-spec}), which presupposes that follow-up is actionable at the scale of the cohort.

\subsubsection{The design response}\label{sec:volume-response}

Four measures follow, and they are testable rather than aspirational.

\textbf{Alert on aggregates, not on readings.} Every alert is raised on a statistic over a window---a seven-day mean, as in the verification profile of \textcite{mihevc2025}---never on a single value, except at the red tier where a single extreme value justifies immediate escalation.

\textbf{Use the yellow tier as the volume regulator.} An intermediate reading costs the patient one verification profile and costs the clinician nothing. This is the mechanism by which the arithmetic above is made tractable. Averaging $n$ readings divides the variance of the estimate by $n$ (\cref{sec:variability}), so the separation between the normal and abnormal distributions, measured in standard deviations, grows as $\sqrt{n}$. The whole receiver operating characteristic curve moves outward, and a higher specificity becomes attainable at the same sensitivity. With the fourteen readings of a seven-day profile, that separation is 3.7 times larger than for a single reading.

\textbf{Report alert volume alongside every performance figure.} This thesis states, for its chosen operating point, how many alerts per clinician per week the system would generate for a cohort of a given size. A sensitivity figure without an alert volume is not an evaluation of a triage system, and this requirement extends the specification of \cref{sec:spec23}.

\textbf{Let the clinician set the threshold.} The operating point is a clinical resource decision, not a technical one. Exposing it as a configurable parameter---the CDSS-setting module of \textcite{yazdani2018} is the structural place for it---allows a service with one physician and a service with five to run the same system at different sensitivities.

%==========================================================================
\subsection{Feedback loops}\label{sec:feedback}

\subsubsection{To the physician: a proposal requiring approval}

The mechanism of \textcite{yazdani2018}---the system displays a proposed prescription, the physician approves or overrides, the approved prescription reaches the patient, every event is recorded---is the correct pattern, and three independent findings support it.

\textbf{It addresses therapeutic inertia}, which \cref{sec:inertia} identified on WHO authority as the mechanism behind the weakest link of the care cascade, 56\% of treated patients not at target. The WHO states that a simple standardised protocol, by clarifying the treatment steps for patient and clinician, contributes to reducing it. A system that presents the next protocol step at the moment the data show the target is not met is that protocol made operative.

\textbf{It is where the clinical value was measured.} \Textcite{bediang2021} found the remote expert modified the local provider's diagnosis in 46.2\% of cases against 17.9\% in controls. The value of the system lay in altering a clinical decision, not in transmitting data.

\textbf{It matches the mechanism of the only randomised trial on the target population.} \Textcite{mihevc2025} attributed their effect to improved communication, medication adherence, early detection of exacerbations and therapy adjustments based on accurate telemonitoring data, and explicitly not to behaviour change, which did not occur (\cref{sec:mechanism}).

The precedent reported by \textcite{yazdani2018} encodes the safety boundary precisely: automatic recommendations are sent to patients from a rule-based database, but recommendations about the amount of insulin to inject are sent first to the physician. Lifestyle advice may be automated; dosing may not.

\subsubsection{To the patient: what the evidence supports}

\Textcite{zhou2024} provide the only quantitative evidence here, from twenty randomised trials and 4\,168 participants. \Cref{tab:bct-implementation} maps the techniques most frequently employed in effective interventions onto their implementation in the proposed system.

\begin{table}[htbp]
\centering\small
\caption{Behaviour change techniques of \textcite{zhou2024} and their implementation here.}\label{tab:bct-implementation}
\begin{tabularx}{\textwidth}{@{}>{\raggedright\arraybackslash}p{5.4cm} c X@{}}
\toprule
\textbf{Technique} & \textbf{Frequency} & \textbf{Implementation} \\
\midrule
Self-monitoring of outcomes of behaviour & 17/20 & the measurement protocol itself \\
Feedback on outcomes of behaviour & 15/20 & returning the reading with its classification \\
Instruction on how to perform the behaviour & 15/20 & measurement guidance in the application \\
Pharmacological support & 15/20 & adherence tracking \\
Biofeedback & 14/20 & trend display \\
Prompts and cues & 13/20 & measurement reminder \\
Credible source & 13/20 & attribution to the clinical team \\
Action planning & 11/20 & scheduling of the next measurement \\
\bottomrule
\end{tabularx}
\end{table}

Their quantitative findings, carried from \cref{tab:subgroups}: interventions with at least eleven techniques achieved $-9.68$~mmHg against $-2.88$~mmHg for those with fewer; multifactorial interventions achieved $-6.50$~mmHg against a non-significant $-1.54$~mmHg; and only 30\% of trials declared a theoretical foundation. \Textcite{alessa2019} established that reminders appear in only 21\% of 186 commercial applications and feedback in 28\% (\cref{tab:app-features}), so the techniques with the strongest evidence are among the least implemented in the market.

\subsubsection{The asymmetry that must be stated}

The pooled effect of \textcite{zhou2024} on systolic pressure is $-5.78$~mmHg, at low GRADE certainty, with 75\% of included trials at high risk of bias. The clinician-coupled system of \textcite{mihevc2025} achieved $-6.9$~mmHg and found no behaviour change at all. \Textcite{mao2020} found mobile health alone at $-2.16$~mmHg, not significant, against $-6.17$~mmHg when combined with professional management (\cref{sec:oversight}).

The patient-facing loop is worth building and is not where the effect comes from. The design conclusion, stated once and carried into the implementation: behaviour change techniques are implemented because they are cheap---reminders, display, scheduling---and because \textcite{zhou2024} show their quantity correlates with effect size. The clinical decision loop is implemented because that is where the measured effect originates.

\subsubsection{The closed loop}\label{sec:closed-loop}

\Textcite{liu2025review} describe the full circuit: \bq{establishing effective feedback channels allows these responses to be used to further optimize algorithms and enhance services. A closed-loop feedback mechanism enables the system to update and optimize treatment plans in real-time based on the patient's actual response and health changes.}

This has an implication for the model that is easily overlooked. Every physician decision---approve, override, request consultation---is a label. An override on a red alert is a false positive confirmed by a clinician; an approval is a true positive. \Textcite{yazdani2018} store all of these as patient record. A system that captures them systematically accumulates, over months of operation, a validation set that is clinically adjudicated rather than threshold-derived---which is the escape from the circularity diagnosed in \cref{sec:task-not-algo}.

This is designed in from the start, because it costs nothing at build time and cannot be reconstructed retrospectively. It is also the most credible route to the kind of contextually valid training data that \textcite{segun2025} identified as absent from sub-Saharan Africa (\cref{sec:regional-landscape}).

\subsubsection{Safety constraints on automated advice}\label{sec:safety}

\begin{table}[htbp]
\centering\small
\caption{The boundary of automation implied collectively by the sources.}\label{tab:automation-boundary}
\begin{tabularx}{\textwidth}{@{}X X@{}}
\toprule
\textbf{May be automated} & \textbf{Must pass a clinician} \\
\midrule
Measurement reminders & any dose change \\
Return of the reading and its classification & any medication change \\
Trend display & any instruction to stop a treatment \\
Generic lifestyle guidance & any urgency-graded instruction \\
Scheduling the next measurement & any diagnosis \\
Non-transmission follow-up & any referral \\
\bottomrule
\end{tabularx}
\end{table}

Two reasons, both established. The best model of \textcite{mendoza2025} misses 49\% of uncontrolled cases, so autonomous action would act on an unreliable signal. And the WHO establishes that the effect of lifestyle modification is approximately 1.3~mmHg systolic (\cref{sec:inertia}), so automated lifestyle advice is both the safest thing to automate and the least consequential.

One caution on personalisation. \Textcite{liu2025review} report a self-care recommendation system for diabetes integrating clinical guidelines with social, economic, cultural and geographical conditions such as food intake and physical exercise. This is the right principle for a Cameroonian deployment---dietary advice derived from European guidelines is unusable where the foods it names are unavailable---but it is a principle, not a validated method, and adapting guideline content to local context requires clinical authorship, not an algorithm.

%==========================================================================
\subsection{Interoperability and the knowledge base}\label{sec:interop}

\Textcite{liu2025review} note that adherence to healthcare industry regulations is essential for platform compliance, naming the Health Insurance Portability and Accountability Act, the Health Information Technology for Economic and Clinical Health Act and the 21st Century Cures Act, and that the adoption of standardised medical terminology and coding systems such as ICD and SNOMED CT is crucial for supporting accurate medical records and analysis.

The regulatory instruments named are United States law and do not apply here. The terminology standards do. Coding diagnoses in ICD and clinical findings in SNOMED CT costs little at design time and is the difference between a record that can later be reconciled with a hospital information system and one that cannot. The nurses of \textcite{mucchi2021} raised precisely this, asking that interfacing with the management information systems in use in the healthcare sector be considered.

\Cref{sec:security} noted that the Cameroonian legal instrument governing health data processing must be identified from national sources, none of the documents in this corpus providing it. That requirement stands and now extends to the knowledge base: the version and date of every clinical rule must be recorded with it, since the framing document anchors the labels in guideline editions that will be superseded. A rule base without version metadata cannot be audited after a guideline revision.

%==========================================================================
\subsection{Specification}\label{sec:spec24}

\begin{table}[htbp]
\centering\footnotesize
\caption{Decision-support requirements established in \cref{sec:cdss}.}\label{tab:spec24}
\begin{tabularx}{\textwidth}{@{}X >{\raggedright\arraybackslash}p{5.9cm}@{}}
\toprule
\textbf{Requirement} & \textbf{Basis} \\
\midrule
Hybrid design: rule engine as floor, model as refinement & \textcite{liu2025review}; mandatory baseline, \cref{sec:spec23} \\
Knowledge base separable from platform, one rule set per pathology & \textcite{yazdani2018}; ontology precedent in \textcite{liu2025review} \\
Three-tier classification, not binary & \textcite{mihevc2025}; only validated protocol \\
Yellow tier triggers data collection, not clinician escalation & \textcite{mihevc2025}; volume regulator, \cref{sec:volume-response} \\
Alerts raised on window aggregates, not single readings & \cref{sec:frequency-decides}; variance reduction raises specificity \\
Red tier may escalate on a single extreme value & \textcite{mihevc2025}; immediate treatment at $\geq$ 160/100, \cref{tab:protocol} \\
Non-transmission is an alarmed event & \textcite{yazdani2018}; HEARTS retention indicator \\
Escalation produces a prioritised ordering, not a partition & \textcite{yazdani2018}; binary sort insufficient at volume \\
Model outputs a continuous score to support the ordering & \cref{sec:alert-volume} \\
System proposes, physician disposes; no autonomous prescribing & \textcite{yazdani2018}; $\mathrm{Se} = 0.51$ in \textcite{mendoza2025} \\
Physician may approve, override, or request consultation & \textcite{yazdani2018} \\
Every physician decision stored as an adjudicated label & \cref{sec:closed-loop} \\
Alert volume reported with every performance figure & \cref{sec:volume-arithmetic}; stated in no source \\
Operating threshold exposed as a clinical configuration parameter & \cref{sec:volume-response} \\
Time-series diagnostics before model selection & \cref{sec:ts-diagnostics} \\
Irregular sampling handled explicitly, not by naive resampling & \cref{sec:ts-diagnostics} \\
Per-patient robust normalisation by median and MAD & \textcite{liu2025review}; personalised threshold gap, \cref{sec:spec13} \\
Offline training path separated from online detection path & \textcite{liu2025review}; edge--cloud split in \textcite{mucchi2021} \\
At least eleven behaviour change techniques, including the eight most frequent & \textcite{zhou2024}: $-9.68$ vs $-2.88$~mmHg \\
Automated advice limited to the permitted column of \cref{tab:automation-boundary} & \textcite{mendoza2025}; WHO lifestyle effect 1.3~mmHg \\
ICD and SNOMED CT coding & \textcite{liu2025review} \\
Every rule versioned and dated & \cref{sec:interop} \\
All events recorded as patient record & \textcite{yazdani2018} \\
\bottomrule
\end{tabularx}
\end{table}

%==========================================================================
\subsection{What this section establishes}\label{sec:cdss-conclusion}

\textbf{Established.} The decision flow of a telemonitoring decision support system is specified completely in one source and validated clinically in another, and the two agree apart from the tier count---where the validated one has three. The rule engine coupled to a learned model is a described technique with named components. The physician-in-the-loop constraint is both designed and independently justified by the measured sensitivity of the best available model. The patient-facing feedback loop has quantified effect sizes and a quantified dose-response in the number of techniques employed.

\textbf{Established here and not in the sources.} The alert volume arithmetic of \cref{sec:alert-volume}, and its two consequences: that at 2\% prevalence an 80\% precision target requires 99.5\% specificity, and that per-reading alerting at sensor sampling frequency is arithmetically unusable regardless of model quality. These follow from the definitions and from the measured values of \textcite{mendoza2025}; they are stated in no source in this corpus.

\textbf{Not established.} No source measures alert fatigue, alert volume, or the time cost of an alert to a clinician. No source evaluates a decision support system of this kind in sub-Saharan Africa. The architecture of \textcite{yazdani2018}, which is the most complete specification available, has never been implemented. The closed feedback loop of \cref{sec:closed-loop} is described by \textcite{liu2025review} but demonstrated by no one. The decision support layer is the least evidenced layer of the entire system, which is notable given that \cref{sec:oversight} established it is the layer where the clinical effect originates.

%##########################################################################
% PART II
%##########################################################################
\part{Methodology, Modelling and System Architecture Design}

\chapter{Data Collection, Preprocessing and Predictive Risk Modelling}\label{ch:data}

\section{Multi-Source Data Collection Strategy and Dataset Construction}\label{sec:collection}

\subsection*{Purpose and status of the sources}

This section specifies the dataset the present work will construct. It differs in kind from \cref{ch:background,ch:techsota}: it is not a review of what others have done but a statement of what will be done, and it is therefore governed by the constraints those chapters established rather than by new evidence.

Two sources are used. \Textcite{cadrage2026} is the research framing document submitted to the supervisory committee; it states the intended contributions, hypotheses, limitations, dataset constitution and ethical framework, and it is the authority for everything this section commits to. \Textcite{efendi2025} supplies a worked example of a vital-signs record schema and of the consent and participant-characterisation apparatus around it; its limitations, established in \cref{sec:efendi-appraisal,sec:leakage}, mean it is used as a template for structure and not as a model for method.

\Cref{tab:carried-constraints} lists the constraints carried from the preceding chapters that govern what follows.

\begin{table}[htbp]
\centering\footnotesize
\caption{Constraints on the dataset established in \cref{ch:background,ch:techsota}.}\label{tab:carried-constraints}
\begin{tabularx}{\textwidth}{@{}X >{\raggedright\arraybackslash}p{4.4cm}@{}}
\toprule
\textbf{Constraint} & \textbf{Established in} \\
\midrule
Blood pressure only from a validated automated upper-arm cuff & \cref{sec:metrology,sec:contradiction} \\
Blood glucose only from a glucometer; no optical substitute & \cref{sec:device-landscape} \\
Camera photoplethysmography restricted to heart rate & \cref{tab:param-sources} \\
Measurement timestamp distinct from entry and arrival timestamp & \cref{sec:pipelines,sec:offline} \\
The three pipelines have distinct missingness mechanisms & \cref{sec:pipelines} \\
Reference value is the seven-day mean of second readings & \cref{sec:protocols} \\
The label must lie outside the input window & \cref{sec:principle} \\
Split by patient and forward in time & \cref{sec:forecasting} \\
Headline metric is recall on the critical class; report alert volume & \cref{sec:spec23,sec:volume-arithmetic} \\
Physician decisions captured as adjudicated labels & \cref{sec:closed-loop} \\
Data residency, consent and national legal instrument stated & \cref{sec:security,sec:interop} \\
\bottomrule
\end{tabularx}
\end{table}

%==========================================================================
\subsection{The three collection streams}\label{sec:streams}

\subsubsection{Stream A: clinical consultation records}

\textbf{Content.} The framing document specifies real, anonymised clinical data supplied by a partner physician in a hospital setting, structured as a time series per patient. From \cref{sec:pipelines}, this stream carries systolic and diastolic pressure, blood glucose, heart rate, laboratory results where available, and the background and context variables: age, sex, body mass index, history, treatments.

\textbf{Role.} This stream is the reference. Measurements are taken by trained staff on hospital equipment under the standardised conditions the WHO specifies and which \cref{sec:dx-requires} set out. It supplies the ground truth against which home measurements are assessed, and it supplies the clinical context---diagnosis, treatment regimen, comorbidity---that no home device can provide.

\textbf{Statistical character.} Low frequency, high reliability, missing not at random. A record exists because a consultation occurred; a consultation occurred because the patient attended; and \cref{sec:followup-requires} established that attendance in this setting falls far below the seven to eight annual contacts the protocol requires, with 33.3\% three-month attendance measured in Cameroonian control centres by \textcite{bediang2021}. A patient who deteriorates may consult more, or may disappear entirely. The sampling process is not independent of the outcome, and any model trained on this stream alone is partly modelling attendance behaviour.

\subsubsection{Stream B: home self-measurement}

\textbf{Content.} Systolic pressure, diastolic pressure, blood glucose and heart rate, with measurement timestamp, acquired by the patient at home on a validated cuff and a glucometer.

\textbf{Role.} This is the principal longitudinal series and the reason the system exists. It is the stream that converts the episodic record of Stream~A into something on which a trend can be computed.

\textbf{Protocol.} Following \textcite{mihevc2025}, the only clinically validated protocol available (\cref{sec:protocols}): blood pressure twice weekly, two readings morning and evening with a one-minute interval enforced by the device and the second retained; blood glucose weekly, fasting and ninety minutes after a meal; verification profiles on amber classification, being a seven-day blood pressure profile or a one-day six-point glycaemic profile.

\textbf{Statistical character.} Medium frequency, variable reliability, missing not at random in a second and distinct way. \Cref{sec:pipelines} established three specific failure modes: transcription error undetectable downstream; entry timestamp substituted for measurement timestamp; and selective omission, in which a patient omits a reading they find discouraging, which removes precisely the extreme values the model exists to detect. The first two are eliminated by Bluetooth transmission; the third is eliminated by no technical means and must be handled analytically.

\subsubsection{Stream C: instrumented sensors}

\textbf{Content.} Heart rate, oxygen saturation, body temperature, and, if present, activity and posture.

\textbf{Role.} Covariates only. \Cref{sec:device-landscape} established that no optical sensor in the reviewed literature delivers clinically validated blood pressure or blood glucose, and \cref{sec:frequency-decides} established a second and independent objection: at one-hertz sampling, a specificity of 0.98 produces more than eleven thousand false alerts per patient per week. Stream~C may characterise the context of a reading---whether the patient was moving, whether the reading was nocturnal---and may enrich the feature space. It may not carry a clinical decision.

\textbf{Statistical character.} High frequency, instrument-limited, missing completely at random through technical failure rather than behaviour.

\begin{figure}[htbp]
\centering
\begin{tikzpicture}[font=\footnotesize]
  \def\W{10.2}
  % Stream A
  \draw[thick,-{Stealth[length=2mm]}] (0,2.4) -- (\W,2.4);
  \node[anchor=east,font=\scriptsize] at (-0.1,2.4) {Stream A};
  \foreach \x in {0.6,3.4,6.9} \node[circle,fill=cRef,inner sep=2pt] at (\x,2.4) {};
  \node[anchor=west,font=\scriptsize,text=cGrey] at (\W+0.1,2.4) {$\tau \approx 90$~d};
  % Stream B
  \draw[thick,-{Stealth[length=2mm]}] (0,1.3) -- (\W,1.3);
  \node[anchor=east,font=\scriptsize] at (-0.1,1.3) {Stream B};
  \foreach \x in {0.4,0.9,1.5,2.1,2.6,3.2,3.8,4.3,4.9,5.5,6.0,6.6,7.2,7.7,8.3,8.9,9.4} 
    \node[circle,fill=cAfrica,inner sep=1.3pt] at (\x,1.3) {};
  \node[anchor=west,font=\scriptsize,text=cGrey] at (\W+0.1,1.3) {$\tau \approx 3.5$~d};
  % Stream C
  \draw[thick,-{Stealth[length=2mm]}] (0,0.3) -- (\W,0.3);
  \node[anchor=east,font=\scriptsize] at (-0.1,0.3) {Stream C};
  \draw[line width=3pt,cGrey!60] (0.3,0.3) -- (9.7,0.3);
  \node[anchor=west,font=\scriptsize,text=cGrey] at (\W+0.1,0.3) {$\tau = 1$~s};
  % horizon window
  \fill[cAfrica!12] (7.2,-0.1) rectangle (7.6,2.8);
  \node[anchor=south,font=\scriptsize,text=cAfrica] at (7.4,2.85) {48~h};
\end{tikzpicture}
\caption{The three streams on a common time axis over one quarter, with a 48-hour horizon marked. Stream~A contributes no observation inside the horizon; Stream~B contributes at most one; Stream~C contributes 172\,800 but cannot measure the parameters of interest.}\label{fig:streams-timeline}
\end{figure}

\subsubsection{What is actually available}\label{sec:scenarios}

The framing document states that the data are real anonymised clinical data from a partner physician in a hospital. This describes Stream~A. It does not describe Stream~B or Stream~C, and the readings differ in everything that follows (\cref{tab:scenarios}).

\begin{table}[htbp]
\centering\footnotesize
\caption{Collection scenarios and their consequences. The thesis must declare which obtains.}\label{tab:scenarios}
\begin{tabularx}{\textwidth}{@{}>{\raggedright\arraybackslash}p{2.3cm} >{\raggedright\arraybackslash}X >{\raggedright\arraybackslash}X >{\raggedright\arraybackslash}X >{\raggedright\arraybackslash}X@{}}
\toprule
 & \textbf{A: retrospective outpatient} & \textbf{B: prospective} & \textbf{C: mixed} & \textbf{D: retrospective inpatient} \\
\midrule
Source & existing consultation records & newly recruited cohort, home devices & records plus prospective sub-cohort & existing in-hospital observation charts \\
Streams & A only & A + B (+ C) & A + B on a subset & A only, at ward density \\
Interval $\tau$ & $\approx$ 90 days & 3--4 days & both & $\approx$ 1 day \\
Ethics burden & secondary use & prospective recruitment, individual consent & both & secondary use \\
Device cost & none & one cuff and one glucometer per patient & partial & none \\
Feasible horizon & next consultation & $\approx$ 7 days & depends on subset & 24--48 h, marginal \\
24--48 h horizon & impossible & marginal & marginal & marginal \\
Contribution C4 has an object & no & yes & yes & no \\
\bottomrule
\end{tabularx}
\end{table}

Scenario~D is not a variant of Scenario~A and must be distinguished from it. An inpatient observation chart records vital signs once or more per day for the duration of a stay, so a cohort of inpatients yields a series two orders of magnitude denser than the same number of outpatients, over a much shorter span. The target of the framing document---at least thirty patients and at least three hundred measurements, which is about ten measurements per patient---is consistent with a stay of about ten days at daily observation, and is not consistent with ninety-day outpatient intervals unless the records span several years per patient. If the available data are inpatient series, the horizon arithmetic of \cref{sec:temporal} changes materially and favourably, and the thesis should say so rather than describing the data generically as hospital records.

The point to register in all four cases is that the low-cost acquisition contribution presupposes Scenario~B or~C. A retrospective record review acquires nothing, and that contribution would have no object. Conversely, Scenario~B requires recruitment, individual consent, device provision and a collection period of several months, none of which the framing document budgets.

%==========================================================================
\subsection{The record schema}\label{sec:record-schema}

\subsubsection{Design principles}

Four principles, each traceable to a constraint established earlier.

\textbf{One row per measurement event, not one row per patient visit.} \Cref{sec:schema} established that a document store keyed on the current value cannot answer a window query, and that the rolling aggregates the framing document specifies require exactly that. The atomic unit is therefore the measurement.

\textbf{Provenance travels with every value}---which stream, which device, who performed it. \Cref{sec:pipelines} established that the three streams have different error structures; a model cannot account for that if the schema discards the distinction.

\textbf{Two timestamps, always.} Measurement time and record time. \Cref{sec:offline} established that offline-first synchronisation decouples them by up to days, and that a model trained on arrival order learns the connectivity pattern rather than the physiology.

\textbf{No derived field stored alongside its inputs in the same row.} \Cref{sec:task-not-algo} established label leakage as the dominant failure mode in this literature. Features are computed at training time from lagged windows; they are not persisted next to the label.

\begin{figure}[htbp]
\centering
\begin{tikzpicture}[node distance=1.2cm,
  t/.style={draw,rounded corners,align=left,font=\scriptsize,inner sep=4pt,minimum width=3.4cm},
  ar/.style={-{Stealth[length=2mm]},thick}]
  \node[t,fill=cRef!12] (pat) {\textbf{patient}\\ patient\_pseudonym (PK)\\ age\_band, sex, bmi\\ dx\_*\_year, diabetes\_type\\ comorbidities (ICD)\\ treatment\_regimen\\ regimen\_version\\ enrolment\_date};
  \node[t,fill=cGrey!14,right=2.4cm of pat] (meas) {\textbf{measurement}\\ measurement\_id (PK)\\ patient\_pseudonym (FK)\\ ts\_measurement, ts\_recorded\\ stream, parameter, value, unit\\ reading\_index, is\_retained\\ context, device\_id\\ device\_validated, operator\\ entry\_mode, protocol\_profile\\ quality\_flag};
  \node[t,fill=cAfrica!12,below=1.0cm of meas] (ev) {\textbf{event (adjudication)}\\ event\_id (PK)\\ patient\_pseudonym (FK)\\ ts\_triggered, trigger\_source\\ tier, score\\ ts\_reviewed\\ clinician\_verdict\\ action\_taken, comment};
  \draw[ar] (pat.east) -- node[above,font=\tiny]{1:N} (meas.west);
  \draw[ar] (pat.south) |- node[below left,font=\tiny,pos=0.75]{1:N} (ev.west);
\end{tikzpicture}
\caption{The three-table data model. The event table is the only source of labels not derived from the thresholds they are meant to predict (\cref{sec:closed-loop}).}\label{fig:schema-model}
\end{figure}

\subsubsection{Measurement record}

\begin{table}[htbp]
\centering\footnotesize
\caption{Fields of the measurement record.}\label{tab:measurement-record}
\begin{tabularx}{\textwidth}{@{}>{\raggedright\arraybackslash}p{3.1cm} >{\raggedright\arraybackslash}p{2.6cm} X@{}}
\toprule
\textbf{Field} & \textbf{Type} & \textbf{Notes} \\
\midrule
\texttt{measurement\_id} & UUID & primary key \\
\texttt{patient\_pseudonym} & string & pseudonymous; not derived from identity \\
\texttt{ts\_measurement} & ISO 8601 with offset & when the reading was taken \\
\texttt{ts\_recorded} & ISO 8601 with offset & when it entered the system \\
\texttt{stream} & enum & consultation $\mid$ home $\mid$ sensor \\
\texttt{parameter} & enum & sbp $\mid$ dbp $\mid$ glucose $\mid$ hr $\mid$ spo2 $\mid$ temp \\
\texttt{value} & float & --- \\
\texttt{unit} & enum & see \cref{tab:units} \\
\texttt{reading\_index} & int & 1 or 2 within a measurement pair \\
\texttt{is\_retained} & bool & true for the second reading, per protocol \\
\texttt{context} & enum & fasting $\mid$ postprandial\_90 $\mid$ morning $\mid$ evening $\mid$ unspecified \\
\texttt{device\_id} & string & links to the device register \\
\texttt{device\_validated} & bool & must be true for sbp, dbp, glucose \\
\texttt{operator} & enum & patient $\mid$ nurse $\mid$ physician \\
\texttt{entry\_mode} & enum & bluetooth $\mid$ manual $\mid$ transcribed \\
\texttt{protocol\_profile} & enum & routine $\mid$ bp\_7day $\mid$ glucose\_6point \\
\texttt{quality\_flag} & enum & ok $\mid$ out\_of\_range $\mid$ duplicate $\mid$ carried\_forward $\mid$ unit\_suspect \\
\bottomrule
\end{tabularx}
\end{table}

Two fields deserve comment. \texttt{device\_validated} enforces at the schema level the restriction that \cref{sec:metrology,sec:contradiction} derived from the WHO: a blood pressure value from an unvalidated device must not be insertable. \texttt{quality\_flag} encodes the three artefacts that \textcite{liu2025review} identify in electronic medical records---duplicates, unit errors, carried-forward values---as a first-class field rather than something discovered at cleaning time. \Cref{sec:records} identified the carried-forward value as the most dangerous of the three for a trend model, since it simulates a stability that did not occur.

\subsubsection{Patient record}

\begin{table}[htbp]
\centering\footnotesize
\caption{Fields of the patient record.}\label{tab:patient-record}
\begin{tabularx}{\textwidth}{@{}>{\raggedright\arraybackslash}p{3.4cm} >{\raggedright\arraybackslash}p{2.0cm} X@{}}
\toprule
\textbf{Field} & \textbf{Type} & \textbf{Notes} \\
\midrule
\texttt{patient\_pseudonym} & string & --- \\
\texttt{age\_band} & enum & five-year bands, not exact age (\cref{sec:reident}) \\
\texttt{sex} & enum & --- \\
\texttt{bmi} & float & as the framing document specifies \\
\texttt{dx\_hypertension\_year} & int & year only \\
\texttt{dx\_diabetes\_year} & int & year only \\
\texttt{diabetes\_type} & enum & t1 $\mid$ t2 $\mid$ unspecified \\
\texttt{comorbidities} & array & ICD-coded (\cref{sec:interop}) \\
\texttt{treatment\_regimen} & array & drug class, dose, start date \\
\texttt{regimen\_version} & int & increments on every change \\
\texttt{enrolment\_date} & date & offset per patient (\cref{sec:reident}) \\
\texttt{site} & string & excluded from the analysis file \\
\bottomrule
\end{tabularx}
\end{table}

\texttt{dx\_hypertension\_year} is included because \textcite{mendoza2025} found years since diagnosis to be the single strongest predictor of uncontrolled hypertension by permutation importance, at 11.9\% increase in mean squared error against 9.4 for age (\cref{sec:interpretability}). It is cheap to record and it is the best-evidenced feature available.

\texttt{regimen\_version} exists because \cref{sec:inertia} established therapeutic inertia as the mechanism behind the weakest link of the care cascade. Without a versioned regimen it is impossible to distinguish a patient whose pressure is drifting under an unchanged treatment from one whose treatment was intensified last week---and that distinction is the clinical content of the entire system.

\subsubsection{Units and conventions}

\begin{table}[htbp]
\centering\small
\caption{Units, plausible ranges and conventions.}\label{tab:units}
\begin{tabularx}{\textwidth}{@{}>{\raggedright\arraybackslash}p{2.4cm} >{\raggedright\arraybackslash}p{2.0cm} >{\raggedright\arraybackslash}p{2.6cm} X@{}}
\toprule
\textbf{Parameter} & \textbf{Unit} & \textbf{Plausible range} & \textbf{Note} \\
\midrule
SBP, DBP & mmHg & 60--260 / 30--160 & disjunctive rules, no coverage gaps (\cref{sec:chatrati}) \\
Blood glucose & mmol/L stored & 1.0--35.0 & Cameroonian practice uses g/L; convert at the boundary \\
Heart rate & bpm & 30--220 & --- \\
SpO$_2$ & \% & 50--100 & covariate only \\
Temperature & \textdegree C & 30.0--45.0 & covariate only \\
\bottomrule
\end{tabularx}
\end{table}

Blood glucose is the field most likely to produce silent unit errors, since three conventions are in use:
\begin{equation}
  1\ \text{mmol/L} = 18.0182\ \text{mg/dL} = 0.180\ \text{g/L}.
\end{equation}
The thresholds of \textcite{mihevc2025} in local units: fasting target below 7~mmol/L, about 1.26~g/L; postprandial target below 10~mmol/L, about 1.80~g/L; hypoglycaemia alert below 4~mmol/L, about 0.72~g/L; diagnostic threshold 11.1~mmol/L, about 2.00~g/L. Values are stored in one unit, converted at the boundary, and \texttt{quality\_flag} is set to \texttt{unit\_suspect} on any value whose magnitude is consistent with the wrong convention---a stored value of 1.26 where mmol/L is expected is almost certainly g/L.

\subsubsection{Event record: the adjudication table}\label{sec:adjudication}

\Cref{sec:closed-loop} established that every clinician decision is a clinically adjudicated label, that it is the only route in this corpus to training data not derived from the thresholds it is meant to predict, and that it cannot be reconstructed retrospectively.

\begin{table}[htbp]
\centering\footnotesize
\caption{Fields of the event record.}\label{tab:event-record}
\begin{tabularx}{\textwidth}{@{}>{\raggedright\arraybackslash}p{3.3cm} >{\raggedright\arraybackslash}p{2.0cm} X@{}}
\toprule
\textbf{Field} & \textbf{Type} & \textbf{Notes} \\
\midrule
\texttt{event\_id} & UUID & --- \\
\texttt{patient\_pseudonym} & string & --- \\
\texttt{ts\_triggered} & timestamp & --- \\
\texttt{trigger\_source} & enum & rule $\mid$ model $\mid$ non\_transmission $\mid$ manual \\
\texttt{tier} & enum & green $\mid$ amber $\mid$ red \\
\texttt{score} & float & model output, where applicable \\
\texttt{ts\_reviewed} & timestamp & null until reviewed \\
\texttt{clinician\_verdict} & enum & confirmed $\mid$ not\_clinically\_significant $\mid$ uncertain \\
\texttt{action\_taken} & enum & none $\mid$ advice $\mid$ dose\_change $\mid$ drug\_change $\mid$ consultation $\mid$ referral \\
\texttt{comment} & text & free text \\
\bottomrule
\end{tabularx}
\end{table}

\texttt{clinician\_verdict} is the label that breaks the circularity of \cref{sec:task-not-algo}. The difference $\texttt{ts\_reviewed} - \texttt{ts\_triggered}$ measures the response latency that \textcite{bediang2021} found to be the real bottleneck: 10.3 days against a 72-hour protocol, and not significantly faster than no system at all (\cref{sec:specialist-access}).

One scope note. This table is populated only by a system in operation. Under Scenario~A or~D of \cref{tab:scenarios} it is empty, and the secondary endpoint of \cref{sec:label-rules} is unavailable. The table is nonetheless specified here, because the system is designed to be deployed and the cost of adding it later is the loss of every event that occurred before it existed.

%==========================================================================
\subsection{Temporal design: the compatibility condition}\label{sec:temporal}

\subsubsection{The condition}

A forecast horizon $\Delta$ is meaningful only if the sampling interval $\tau$ satisfies $\tau \ll \Delta$. If $\tau \geq \Delta$, the forecast target falls between observations and the model is asked to predict a value against which it will never be scored. A reasonable floor is two or three observations within the horizon.

\Cref{sec:horizon} established this in its general form and \cref{fig:sampling-horizon} plots it. \Cref{tab:horizon-scenarios} applies it to the scenarios of \cref{tab:scenarios}.

\begin{table}[htbp]
\centering\small
\caption{Observations falling inside the forecast horizon, by sampling regime.}\label{tab:horizon-scenarios}
\begin{tabularx}{\textwidth}{@{}X >{\raggedright\arraybackslash}p{2.0cm} c c >{\raggedright\arraybackslash}p{2.0cm}@{}}
\toprule
\textbf{Regime} & \textbf{$\tau$} & \textbf{In 24~h} & \textbf{In 48~h} & \textbf{Compatible} \\
\midrule
Consultation records, quarterly & $\approx$ 90 d & 0 & 0 & no \\
Inpatient chart, daily & 1 d & 1 & 2 & marginal \\
Home routine protocol & $\approx$ 3.5 d & 0 & 0 or 1 & no \\
Home protocol, daily & 1 d & 1 & 2 & marginal \\
Verification profile, twice daily & 0.5 d & 2 & 4 & yes \\
Sensor stream, 1~Hz & 1 s & 86\,400 & 172\,800 & yes, but see \cref{sec:frequency-decides} \\
\bottomrule
\end{tabularx}
\end{table}

Under the routine home protocol this thesis has adopted, a 24-hour horizon contains zero observations and a 48-hour horizon at most one. The contribution as originally stated is not realisable on routine-frequency data.

\subsubsection{The resolution, restated}

\Cref{sec:horizon} resolved this and the resolution is restated here because it governs the dataset. The forecast horizon is not single-valued: it is matched to the regime the patient is in.

In the \textbf{routine regime}, the target is the classification of the next scheduled measurement, or whether the seven-day mean at the next window crosses threshold. This gives a horizon of three to seven days, which is clinically meaningful---the WHO confirmation visit is one to four weeks (\cref{tab:protocol}) and the verification profile of \textcite{mihevc2025} is seven days---and is compatible with the data.

In the \textbf{verification regime}, where the protocol prescribes two readings per day, a 24 to 48 hour horizon contains two to four observations and is estimable. This is also the regime in which a clinical decision is pending, so it is the regime in which a short-horizon forecast would change an action rather than merely be recorded.

The anticipatory character of the contribution is preserved by both; what is abandoned is the claim that a single 24 to 48 hour horizon applies uniformly. Should the available data prove to be inpatient series at daily observation (Scenario~D), the 24-hour horizon becomes marginally estimable in the routine regime as well, and the restriction relaxes accordingly.

\subsubsection{Sample size against model capacity}\label{sec:capacity}

The framing document targets at least 30 patients and at least 300 measurements, which is approximately ten observations per patient.

For feature engineering with a seven-day rolling window at two observations per week, the window contains two points and the first usable sample arrives after four observations, leaving roughly six usable samples per patient, or 180 in total.

For a recurrent model the position is worse. An LSTM layer has
\begin{equation}
  P = 4\bigl[u(f+u) + u\bigr]
\end{equation}
parameters, where $u$ is the unit count and $f$ the input dimension. For a deliberately small configuration---a single layer, $u = 16$, $f = 10$---this gives $4(16 \times 26 + 16) = 1\,728$ parameters, plus a small output head. Against 180 training sequences that is approximately ten parameters per sample, inverted from the conventional rule of thumb by two orders of magnitude. A two-layer network at $u = 32$ reaches about 13\,800 parameters, or seventy-five per sample.

Hypothesis H3---that the gain of the deep model may remain modest---is therefore not a prudent hedge. On the sample size as planned, a negative result for the recurrent model is close to arithmetically guaranteed, and the comparison will measure sample size rather than architecture. That can still be reported honestly, and \cref{sec:ml-conclusion} establishes that an honest negative result is worth more than an inflated positive one. But it must be framed as such in advance, not discovered at the defence.

A second risk compounds this. With group $k$-fold at five folds and 30 patients, each fold holds six patients. If critical events are rare---the framing document states they are---and clustered within a few patients, some folds will contain zero positive cases and the fold-level recall will be undefined. The split must therefore be stratified by patient-level event presence, not merely grouped.

\subsubsection{Four responses to the sample-size constraint}

\textbf{Increase the sampling frequency for a sub-cohort.} A subset measuring twice daily for a defined period yields the density the shorter horizon requires. This costs patient burden and must be weighed against the 10.9\% equipment-related dropout that \textcite{mihevc2025} observed at a far lighter protocol (\cref{sec:feasibility}).

\textbf{Use a public dataset for the comparative study and the local dataset for external validation.} \Textcite{liu2025review} catalogue public chronic disease datasets, several containing both blood pressure and glucose. The architecture comparison would be conducted where sample size permits, and the local dataset would answer the question \textcite{segun2025} raise---whether an externally trained model transfers to an African population (\cref{sec:regional-landscape}). This converts a sample-size weakness into a contribution, and it addresses a gap that \cref{sec:positioning} established is unfilled: no published African validation study exists for the tools in circulation.

\textbf{Reduce the scope of the architecture comparison.} Compare regularised logistic regression, random forest and XGBoost, drop the recurrent model, and state why. This is defensible---\textcite{mendoza2025} compared exactly these three---but it abandons a stated contribution.

\textbf{Report the recurrent model as a capacity-limited negative result}, with the parameter-per-sample ratio stated. This is the least satisfying option and the most honest one if none of the others is available.

%==========================================================================
\subsection{Labelling}\label{sec:labelling}

\subsubsection{Threshold sources and the convention conflict}

The framing document anchors the risk label in the ESH and ADA guideline editions, validated by the partner physician. \Cref{sec:interop} adds the requirement that every rule carry its version and date, since both will be superseded.

One convention conflict must be resolved explicitly. The WHO surveillance definition is 140/90~mmHg with a disjunctive rule (\cref{sec:hta-burden}); the ESH classification uses 130/80 as a threshold within its grading, which is also the definition \textcite{mendoza2025} adopt; and the home target of \textcite{mihevc2025} is 135/85 (\cref{sec:protocols}). Three different numbers, for three different purposes. \Cref{tab:threshold-allocation} fixes the allocation, and no section of this thesis departs from it.

\begin{table}[htbp]
\centering\small
\caption{Allocation of blood pressure thresholds by function.}\label{tab:threshold-allocation}
\begin{tabularx}{\textwidth}{@{}>{\raggedright\arraybackslash}p{3.6cm} >{\raggedright\arraybackslash}p{2.6cm} X@{}}
\toprule
\textbf{Function} & \textbf{Threshold} & \textbf{Source} \\
\midrule
Risk label for modelling & 130/80~mmHg & ESH classification, as specified in \textcite{cadrage2026} \\
Alerting tiers & 135/85~mmHg home target & \textcite{mihevc2025}, \cref{tab:thresholds} \\
Epidemiological comparison & 140/90~mmHg & \textcite{who2023hypertension}, \cref{ch:background} \\
\bottomrule
\end{tabularx}
\end{table}

\subsubsection{The circularity constraint, restated as a rule}\label{sec:label-rules}

If $y_t = h(x_t)$ and the features include $x_t$, the task is label recovery and not prediction. \Cref{sec:task-not-algo} demonstrated this on \textcite{efendi2025}, where three features were discretised using the very ranges that define the class and an accuracy of 0.973 measured nothing. The operational rules that follow:

\begin{enumerate}[nosep]
  \item The label is computed on $x_{t+\Delta}$ and the features on $x_{t-k}, \dots, x_t$, with no overlap.
  \item No feature expressing deviation from a threshold may be computed at the label instant.
  \item The clinical threshold rule applied to the last observed value is a mandatory baseline. The model must beat it, and the margin by which it does is the actual contribution.
  \item Where the adjudication table of \cref{sec:adjudication} has accumulated verdicts, performance is reported against clinician-adjudicated labels as a secondary endpoint. These are not derived from thresholds and are therefore immune to the circularity entirely.
\end{enumerate}

\subsubsection{Adjudication and inter-rater agreement}\label{sec:agreement}

The framing document commits to validation of the labelling by the partner physician. Two additions are needed for that commitment to be checkable.

A \textbf{written labelling protocol}, stating the rule, the inputs it sees and the handling of ambiguous cases, so that the labelling is reproducible by someone other than its author.

A \textbf{second reader on a sample}. \Textcite{zhou2024} report inter-rater agreement for their technique coding; \textcite{alessa2019} report a Cohen's kappa of 0.85; \textcite{dodoo2022} report 86\% agreement between two experts. A thesis whose central construct is a risk label, assigned by one person, with no agreement statistic, is weaker than any of its own sources on exactly this point. A second clinician reading 10 to 20\% of the records, with kappa reported, closes that gap at low cost.

%==========================================================================
\subsection{Ethical framework}\label{sec:ethics}

\subsubsection{What the framing document commits to}

Three commitments are recorded: the data are real and anonymised; direct identifiers are removed before any processing; collection proceeds with the agreement of the partner physician and the institution, for research purposes only. These are necessary and they are not sufficient.

\subsubsection{What is missing}

\begin{table}[htbp]
\centering\footnotesize
\caption{Elements of the ethical framework not yet documented.}\label{tab:ethics-gaps}
\begin{tabularx}{\textwidth}{@{}>{\raggedright\arraybackslash}p{4.6cm} >{\raggedright\arraybackslash}p{2.2cm} X@{}}
\toprule
\textbf{Required} & \textbf{Status} & \textbf{Note} \\
\midrule
Ethics committee reference number & absent & \textcite{mihevc2025} cite a trial registration; \textcite{bediang2021} cite a national ethics reference \\
Cameroonian legal instrument governing health data & absent & \textcite{mucchi2021} cite their national decree; flagged in \cref{sec:security,sec:interop} \\
Written institutional authorisation & referenced, not evidenced & must be appended \\
Individual consent, if prospective & absent & \textcite{efendi2025} obtained written consent; mandatory under Scenario~B or~C \\
Retention period and destruction date & absent & --- \\
Data residency region & absent & the endpoint of \textcite{efendi2025} was in south-east Asia \\
Re-identification risk assessment & absent & see \cref{sec:reident} \\
Secondary-use status, if retrospective & absent & see below \\
\bottomrule
\end{tabularx}
\end{table}

The last line matters. \Textcite{mendoza2025} could dispense with approval because their source survey is a public, already de-identified dataset. Hospital records are neither, and secondary use of them requires whatever national and institutional process applies---which must be identified rather than assumed.%
\verify{Obtain and cite the ethics committee reference, the institutional authorisation, and the Cameroonian legal instrument governing processing of personal health data. None of the documents in this corpus supplies the last of these, and it has been outstanding since \cref{sec:security}.}

\subsubsection{Re-identification risk in a small cohort}\label{sec:reident}

Removing direct identifiers is not anonymisation. The residual risk comes from quasi-identifiers: attributes that are not identifying alone but are jointly unique. In a cohort of thirty patients at a single facility, the combination of exact age, sex, exact consultation dates and site is almost certainly unique for every individual. A clinician at that facility, or anyone with access to its appointment book, could re-identify the entire dataset. The commitment to removing direct identifiers does not address this.

Four measures, each cheap:

\textbf{Age in five-year bands}, not exact years. This costs the model very little---\textcite{mendoza2025} used three age groups and found age among the strongest predictors nonetheless.

\textbf{Date shifting.} A random offset $\delta_i$, drawn once per patient and uniformly over $\pm 180$ days, is applied to every timestamp for that patient. Intervals within a patient are preserved exactly, which is all a time-series model needs, while absolute dates are destroyed. The offset table is held separately from the analysis file.

\textbf{Exclusion of the site identifier} from the analysis file. With one or two sites it carries no information and considerable risk.

\textbf{Suppression or banding of rare values.} A single patient with an exceptional diagnosis year or comorbidity is identifiable by that fact alone.

One consequence of date shifting must be recorded, because it conflicts with a requirement established earlier. Per-patient offsets destroy alignment \emph{between} patients, so no calendar-level event---a medicine stock-out, a strike, a seasonal effect---can be identified in the shifted data. \Cref{sec:scope-problem} required that medicine availability over the observation period be documented as a control variable, following the Indian example in which supply disruption discouraged return visits. That requirement is met by holding a separate calendar-event file, unshifted and containing no patient data, and by expressing each event in the analysis file as an elapsed-time flag relative to the patient's own enrolment.

A final consideration: this thesis will be published, and any figure plotting an individual patient's trajectory against real dates reproduces the risk the measures above remove. Figures plot shifted dates or elapsed time from enrolment.

\subsubsection{The analysis file}

The output of this section is one file, and its contents are stated exactly. The analysis file contains: the patient pseudonym; age band; sex; body mass index; years since each diagnosis; diabetes type; ICD-coded comorbidities; the treatment regimen with its version; and, per measurement, the shifted measurement timestamp, the stream, the parameter, the value and unit, the reading index and retention flag, the context, the operator, the entry mode, the protocol profile and the quality flag.

It contains no name, no record number, no telephone number, no address, no exact date of birth, no exact calendar date, no site identifier and no device serial number. The pseudonym is randomly generated and the mapping to identity is held by the partner physician and is not transferred.

%==========================================================================
\subsection{Quality controls at collection}\label{sec:quality}

Controls applied at write time are worth more than cleaning applied afterwards, because the information needed to correct an error is available at the point of entry and lost by the time analysis begins. \Cref{sec:schema} established that server-side schema validation is available and should be used.

\begin{table}[htbp]
\centering\footnotesize
\caption{Quality controls applied at write time.}\label{tab:quality-controls}
\begin{tabularx}{\textwidth}{@{}>{\raggedright\arraybackslash}p{2.9cm} X >{\raggedright\arraybackslash}p{2.9cm}@{}}
\toprule
\textbf{Control} & \textbf{Rule} & \textbf{Flag} \\
\midrule
Plausibility & value within the range of \cref{tab:units} & \texttt{out\_of\_range} \\
Internal coherence & SBP $>$ DBP, and SBP $-$ DBP $\geq$ 10~mmHg & \texttt{out\_of\_range} \\
Duplicate & same patient, parameter and timestamp to the minute & \texttt{duplicate} \\
Carried forward & three or more consecutive identical values of a continuous parameter & \texttt{carried\_forward} \\
Unit suspicion & magnitude consistent with another convention & \texttt{unit\_suspect} \\
Device validation & \texttt{device\_validated} false on sbp, dbp, glucose & reject at write \\
Temporal coherence & \texttt{ts\_measurement} $>$ \texttt{ts\_recorded} & \texttt{out\_of\_range} \\
Protocol conformity & reading pair present, interval 45--180~s & informational \\
\bottomrule
\end{tabularx}
\end{table}

Two notes. Flagged records are retained, not deleted---a carried-forward sequence is itself informative about data-collection behaviour, and deleting it hides a pattern. And the last control measures whether the protocol was actually followed; \textcite{mihevc2025} enforced the one-minute interval in the device, and the degree to which that enforcement held is a reportable quantity.

%==========================================================================
\subsection{Specification}\label{sec:spec31}

\begin{table}[htbp]
\centering\footnotesize
\caption{Dataset requirements established in \cref{sec:collection}.}\label{tab:spec31}
\begin{tabularx}{\textwidth}{@{}X >{\raggedright\arraybackslash}p{5.5cm}@{}}
\toprule
\textbf{Requirement} & \textbf{Basis} \\
\midrule
One row per measurement event & \cref{sec:schema} window queries \\
Both \texttt{ts\_measurement} and \texttt{ts\_recorded} stored & \cref{sec:pipelines,sec:offline} \\
\texttt{stream} recorded on every row & \cref{sec:pipelines} distinct error structures \\
\texttt{device\_validated} enforced at write for BP and glucose & \cref{sec:metrology} \\
\texttt{quality\_flag} as a first-class field & EMR artefacts in \textcite{liu2025review} \\
Glucose stored in one unit, converted at the boundary & \cref{tab:units} \\
\texttt{regimen\_version} incremented on every therapy change & therapeutic inertia, \cref{sec:inertia} \\
\texttt{dx\_hypertension\_year} recorded & \textcite{mendoza2025}: strongest predictor \\
Adjudication table in place from first deployment & \cref{sec:closed-loop}; cannot be reconstructed \\
Horizon matched to the regime, not single-valued & \cref{sec:horizon,sec:temporal} \\
Label on $x_{t+\Delta}$, features on $x_{t-k..t}$, no overlap & \cref{sec:principle} \\
Clinical threshold rule as mandatory baseline & \cref{sec:spec23} \\
Threshold allocation of \cref{tab:threshold-allocation} observed throughout & \cref{sec:labelling} \\
Guideline version and date recorded with every rule & \cref{sec:interop} \\
Written labelling protocol & \cref{sec:agreement} \\
Second reader on 10--20\% of records, kappa reported & $\kappa = 0.85$ in \textcite{alessa2019}; 86\% in \textcite{dodoo2022} \\
Stratified group $k$-fold, forward in time & \cref{sec:forecasting,sec:capacity} \\
Parameter-per-sample ratio stated for every model & \cref{sec:capacity} \\
Age in five-year bands & \cref{sec:reident} \\
Per-patient date shift, intervals preserved & \cref{sec:reident} \\
Separate unshifted calendar-event file for supply disruptions & \cref{sec:reident,sec:scope-problem} \\
Site identifier excluded from the analysis file & \cref{sec:reident} \\
Ethics reference, legal instrument, retention and residency stated & \cref{tab:ethics-gaps} \\
Figures plot shifted dates or elapsed time & \cref{sec:reident} \\
\bottomrule
\end{tabularx}
\end{table}

%==========================================================================
\subsection{Decisions and their consequences}\label{sec:decisions}

Three decisions govern everything downstream, and two of them remain open at the time of writing.

\textbf{The collection scenario} determines the streams available, the ethics pathway, the device budget, the collection period, and whether the low-cost acquisition contribution has an object at all. \Cref{tab:scenarios} sets out the four cases. This thesis declares its scenario in \cref{sec:preprocessing}, where the preprocessing pipeline is specified, because imputation strategy, normalisation scope and window construction all depend on which streams exist and at what interval.

\textbf{The forecast horizon} is settled. \Cref{sec:horizon} established the incompatibility between a uniform 24 to 48 hour horizon and routine-frequency sampling, and \cref{sec:temporal} adopts the regime-dependent resolution: next scheduled measurement or next seven-day window in the routine regime, 24 to 48 hours in the verification regime. The framing document is amended accordingly, and the anticipatory contribution is preserved in both.

\textbf{The sample-size strategy} remains open. At thirty patients and three hundred measurements, the architecture comparison will measure sample size rather than architecture, for the reasons computed in \cref{sec:capacity}. The strongest available response is to run the comparison on a public dataset and use the local cohort for external validation, which converts the limitation into a contribution and addresses the absence of African validation studies documented in \cref{sec:positioning}. Whichever response is adopted, the parameter-per-sample ratio is reported alongside every model.

%##########################################################################
% SECTION 3.2
%##########################################################################
\section{Ingestion, Multi-Channel Temporal Aggregation and Missing Data Handling}\label{sec:preprocessing}

\subsection*{Purpose, sources and a terminological correction}

\Cref{sec:collection} specified what is collected and in what form. This section specifies what is done to it between the database and the model. It is the stage at which most of the failures diagnosed in \cref{ch:techsota} are either prevented or built in permanently: label leakage, fabricated temporal structure, and the systematic bias that imputation introduces when data are missing for clinical reasons.

Two sources are used. \Textcite{cadrage2026} commits this work to masking, imputation, per-patient normalisation and robustness to missingness as responses to the irregular and missing measurements arising from intermittent connectivity, and to rolling aggregates, trends, variability and threshold deviations as the feature construction. \Textcite{liu2025review} describe the data-processing layer of a chronic-disease platform and name specific techniques for cleaning, outlier handling and time-series characterisation; the reservation on their method stands, and the techniques they name are used because each is independently verifiable, not on their authority.

A terminological correction is required before proceeding. There is no spatial dimension in this dataset. It is a single-site, single-cohort study: no patient coordinates, no geographic stratification and no multi-site structure is specified. The operation performed is temporal aggregation across several channels---blood pressure, blood glucose, heart rate---each sampled on its own schedule. This work therefore uses \emph{multi-channel temporal aggregation} throughout and does not describe the operation as spatiotemporal.

%==========================================================================
\subsection{The pipeline and its ordering constraint}\label{sec:pipeline}

\Textcite{liu2025review} describe the data-processing layer as comprising data cleansing, data conversion, feature engineering, data extraction and data integration, and state that it \bq{utilizes various preprocessing techniques, including data cleaning, format conversion, and outlier handling, to integrate and conduct preliminary analysis on the collected raw data}.

The list is correct but unordered, and the order is the part that determines validity. \Cref{sec:forecasting} established that every preprocessing step which learns a parameter from data must be fitted on the training fold alone. The pipeline therefore partitions into two phases separated by the split (\cref{fig:pipeline-phases}).

\begin{figure}[htbp]
\centering
\begin{tikzpicture}[node distance=0.5cm,
  ph/.style={draw,rounded corners,align=left,font=\scriptsize,inner sep=5pt,minimum width=5.0cm},
  sp/.style={draw,rounded corners,align=center,font=\scriptsize\bfseries,fill=cAfrica!18,
             minimum width=5.0cm,minimum height=0.85cm},
  ar/.style={-{Stealth[length=2.5mm]},thick}]
  \node[ph,fill=cGrey!12] (p1) {\textbf{Phase 1 --- deterministic, all data}\\[2pt]
    ingestion and schema validation\\ unit conversion to canonical unit\\
    range and coherence checks\\ deduplication\\ carried-forward detection\\
    protocol-conformity flagging};
  \node[sp,below=of p1] (split) {SPLIT: by patient, forward in time,\\ stratified on event presence};
  \node[ph,fill=cRef!12,below=of split] (p2) {\textbf{Phase 2 --- learned, training fold only}\\[2pt]
    per-patient normalisation statistics\\ pooled dispersion estimate\\
    outlier thresholds\\ any imputation model\\ window construction and features\\
    class-balance correction};
  \draw[ar] (p1) -- (split);
  \draw[ar] (split) -- (p2);
  \node[anchor=west,font=\scriptsize,text=cAfrica,align=left] at ([xshift=0.4cm]p2.east)
    {a step moved above\\ the split invalidates\\ the evaluation, and\\ the error is invisible\\ in the results};
\end{tikzpicture}
\caption{The two-phase preprocessing pipeline. Phase~1 operations learn nothing from the distribution and cannot leak; Phase~2 operations estimate a parameter and must therefore be fitted on the training fold alone.}\label{fig:pipeline-phases}
\end{figure}

A step placed on the wrong side of the split invalidates the evaluation, and the error is invisible in the results: it manifests only as an optimistic estimate that fails on deployment. Stating this partition explicitly, and implementing it as a pipeline object rather than as a sequence of notebook cells, is the cheapest guarantee of validity available to this work.

%==========================================================================
\subsection{Ingestion and deterministic validation}\label{sec:ingestion}

\Cref{sec:quality} specified the controls; this subsection specifies how they behave.

The governing principle is to flag and not to delete. \Textcite{liu2025review} report a finding that one cleaning tool was \bq{the only method capable of cleaning adult data and identifying pervasive carry-over and duplicate data in EHRs}, the errors concerned being \bq{unit errors, incorrect figures, duplicates, and carry-over values}. Three of those four are the artefacts specified in \cref{tab:quality-controls}, and the fourth---incorrect figures---is the transcription error that Bluetooth acquisition eliminates (\cref{sec:pipelines}).

The consequence of flagging rather than deleting is that the flags themselves become data. A patient whose series contains three carried-forward sequences is telling the system something about how their measurements are being produced, and a model that never sees that fact cannot account for it.

\begin{table}[htbp]
\centering\footnotesize
\caption{Ingestion operations, their phase and their behaviour on violation.}\label{tab:ingestion-ops}
\begin{tabularx}{\textwidth}{@{}X c c >{\raggedright\arraybackslash}p{4.0cm}@{}}
\toprule
\textbf{Operation} & \textbf{Learns from data?} & \textbf{Phase} & \textbf{On violation} \\
\midrule
Schema validation & no & 1 & reject at write \\
Unit conversion & no & 1 & convert; flag if magnitude suspect \\
Range check & no & 1 & flag \texttt{out\_of\_range}; retain \\
SBP $>$ DBP and SBP $-$ DBP $\geq$ 10~mmHg & no & 1 & flag; retain \\
Exact duplicate & no & 1 & retain one, flag the rest \\
Three or more identical consecutive values & no & 1 & flag \texttt{carried\_forward}; retain \\
Reading-pair interval 45--180~s & no & 1 & informational flag \\
Statistical outlier detection & yes & 2 & flag for review (\cref{sec:outliers}) \\
\bottomrule
\end{tabularx}
\end{table}

%==========================================================================
\subsection{Outliers: artefact or event?}\label{sec:outliers}

This subsection addresses a danger that none of the sources in this corpus raises, and that is specific to risk-detection systems.

\subsubsection{The problem}

\Textcite{liu2025review} report three approaches---box plots, absolute median differences and extreme value theory---and a use of robust linear regression for outlier detection \bq{allowing for the automatic cleaning of anomalies in longitudinal electronic health records}.

The phrase \emph{automatic cleaning of anomalies} is exactly the danger. In this system the anomaly is the event of interest. A systolic pressure of 185~mmHg is a statistical outlier in a patient whose median is 138. It is also a red-tier alert under the protocol of \cref{tab:thresholds}. An outlier-removal step applied before labelling deletes the positive class, and the deletion is silent: the model trains on a cleaned series in which hypertensive crises do not occur, achieves excellent apparent performance, and detects nothing in deployment.

The danger is amplified by the class imbalance the framing document anticipates. If critical states are rare, they are also by construction statistically extreme, and any automatic outlier rule targets them preferentially (\cref{sec:imbalance}).

\subsubsection{The discrimination rule}

\begin{table}[htbp]
\centering\footnotesize
\caption{Three categories of extreme value and their handling.}\label{tab:outlier-categories}
\begin{tabularx}{\textwidth}{@{}>{\raggedright\arraybackslash}p{3.4cm} X >{\raggedright\arraybackslash}p{4.2cm}@{}}
\toprule
\textbf{Category} & \textbf{Criterion} & \textbf{Action} \\
\midrule
Physiologically impossible & outside the range of \cref{tab:units} & flag \texttt{out\_of\_range}; exclude from features and from labels \\
Physiologically possible, extreme & within range, beyond the statistical threshold & retain; this may be the event \\
Internally contradicted & extreme value contradicted by a reading within hours, with no plausible trajectory & flag for clinician review; never auto-delete \\
\bottomrule
\end{tabularx}
\end{table}

A systolic reading of 400~mmHg is impossible and is an artefact. A reading of 195~mmHg is possible, and \cref{tab:protocol} records that the WHO prescribes same-day treatment initiation above 160/100---so it is both plausible and clinically urgent. The dividing line is physiological plausibility, not statistical rarity.

\subsubsection{The statistical threshold, and a constraint on it}\label{sec:mad-shrinkage}

Where a statistical threshold is used, for flagging and not for deletion, the median absolute deviation is preferred over the standard deviation because it is not itself inflated by the outliers it is meant to detect:
\begin{equation}
  \mathrm{MAD}_i = \operatorname{median}_t \bigl\lvert x_{i,t} - \operatorname{median}_t(x_{i,t}) \bigr\rvert,
  \qquad
  z_{i,t} = \frac{x_{i,t} - \operatorname{median}_t(x_{i,t})}{1.4826\,\mathrm{MAD}_i}.
\end{equation}
The constant 1.4826 makes the denominator consistent with the standard deviation under normality, so $z$ is interpretable on the usual scale.

The constraint is sample size. With approximately ten observations per patient (\cref{sec:capacity}), a per-patient median and MAD are estimated from ten points and are extremely unstable---a single extreme value among ten shifts the MAD substantially, which is the opposite of the robustness intended.

The response is shrinkage toward the pooled estimate:
\begin{equation}
  \widehat{\mathrm{MAD}}_i = \lambda_i\,\mathrm{MAD}_i + (1-\lambda_i)\,\mathrm{MAD}_{\text{pooled}},
  \qquad
  \lambda_i = \frac{n_i}{n_i + n_0},
\end{equation}
where $n_i$ is the observation count for patient $i$ and $n_0$ a tuning constant representing the sample size at which the individual estimate is weighted equally with the pooled one. With $n_0 = 20$, a patient with ten observations draws two-thirds of their dispersion estimate from the cohort and one-third from themselves. As the series lengthens the estimate becomes progressively personal---which is precisely the personalisation that \cref{sec:spec13} identified as the gap in the fixed-threshold protocol of \textcite{mihevc2025}, implemented in a way that degrades gracefully when data are scarce. The quantity $\mathrm{MAD}_{\text{pooled}}$ is computed on the training fold only.

%==========================================================================
\subsection{Irregular sampling}\label{sec:irregular}

\subsubsection{Three strategies, and why the obvious one is wrong}

\begin{table}[htbp]
\centering\small
\caption{Strategies for handling irregular sampling.}\label{tab:irregular-strategies}
\begin{tabularx}{\textwidth}{@{}>{\raggedright\arraybackslash}p{4.2cm} X >{\raggedright\arraybackslash}p{4.2cm}@{}}
\toprule
\textbf{Strategy} & \textbf{Mechanism} & \textbf{Cost} \\
\midrule
Resample to a regular grid & interpolate onto fixed intervals & fabricates observations \\
Preserve irregularity, encode the interval & keep raw timestamps; supply $\Delta t$ as a feature & model must handle variable spacing \\
Time-aware model & architecture consumes $\Delta t$ directly & more complex, more parameters \\
\bottomrule
\end{tabularx}
\end{table}

\Cref{sec:ts-diagnostics} established the objection to the first: at the routine protocol frequency, resampling onto a daily grid manufactures five interpolated points for every two real ones, and any trend feature computed on that grid is five-sevenths fiction. The fabricated points are also smoother than reality, since interpolation is a low-pass operation---so variability features computed on a resampled grid are systematically understated, and variability is one of the four feature families the framing document specifies. Resampling would thus bias downward precisely the quantity $\sigma_{\varepsilon,i}$ that \cref{sec:variability} identified as an independent risk factor.

The second strategy is the default for this work, with the third available for the recurrent model if the sample size permits, which \cref{sec:capacity} suggests it will not.

\subsubsection{What must be encoded}\label{sec:interval-features}

Preserving irregularity means the interval structure must enter the feature vector explicitly. Per channel $c$ and per sample at time $t$:

\begin{table}[htbp]
\centering\small
\caption{Interval-structure features, computed per channel.}\label{tab:interval-features}
\begin{tabularx}{\textwidth}{@{}>{\raggedright\arraybackslash}p{3.4cm} X@{}}
\toprule
\textbf{Quantity} & \textbf{Definition} \\
\midrule
$\Delta t^{(c)}_{\text{last}}$ & time since the most recent observation of channel $c$ \\
$n^{(c)}_W$ & observation count of channel $c$ within the window \\
$\mathrm{span}^{(c)}_W$ & $t^{(c)}_{\max} - t^{(c)}_{\min}$ actually covered within the window \\
$\overline{\Delta t}^{(c)}_W$ & mean inter-observation interval within the window \\
$\mathrm{CV}(\Delta t)^{(c)}_W$ & coefficient of variation of intervals: regularity of adherence \\
\bottomrule
\end{tabularx}
\end{table}

The per-channel subscript is essential. Under the protocol of \cref{sec:streams}, blood pressure arrives twice weekly and blood glucose weekly, so the channels are never synchronous. A representation that assumes a shared time index silently aligns measurements taken days apart.

\subsubsection{Time-series characterisation before modelling}

\Textcite{liu2025review} prescribe a diagnostic step identifying periodicity, drift and stationarity. This belongs in the experimental chapter, before any model is fitted, with one qualification repeated from \cref{sec:ts-diagnostics}: autocorrelation and unit-root tests assume regular sampling. Applying them here requires either resampling, which the preceding subsection rejects, or methods for unevenly spaced series such as the Lomb--Scargle periodogram. This thesis states which was used and does not report an autocorrelation function computed on a resampled series as though it described the raw one.

\begin{table}[htbp]
\centering\footnotesize
\caption{What the pre-modelling diagnostic should establish.}\label{tab:ts-diagnostic}
\begin{tabularx}{\textwidth}{@{}>{\raggedright\arraybackslash}p{4.0cm} >{\raggedright\arraybackslash}p{4.2cm} X@{}}
\toprule
\textbf{Question} & \textbf{If yes} & \textbf{Implication} \\
\midrule
Is there circadian periodicity in blood pressure? & morning and evening readings differ systematically & \texttt{context} must be a feature; do not pool morning and evening \\
Is there drift under stable therapy? & the median moves across windows & trend features are informative; this is the forecast target \\
Is the series stationary? & no & differencing required before trend features are interpretable \\
Does variability change with level? & yes & consider a variance-stabilising transform \\
\bottomrule
\end{tabularx}
\end{table}

%==========================================================================
\subsection{Window construction}\label{sec:windows}

\subsubsection{The sample}

A training sample is a pair: an input window and a label computed outside it,
\begin{equation}
  \mathcal{S}_{i,t} = \Bigl(
  \underbrace{\{x_{i,s} : t - L \leq s \leq t\}}_{\text{input window of length } L},
  \;
  \underbrace{g\bigl(x_{i,t'}\bigr)}_{\text{label}}
  \Bigr),
  \qquad t' \in [\,t + \varepsilon,\; t + \Delta\,],
\end{equation}
where $\varepsilon > 0$ is an embargo between the end of the window and the earliest admissible label observation, so that no observation contributes to both. Because sampling is irregular, the label is not located at a fixed instant $t+\Delta$ but is the first qualifying observation inside the label search interval; if no observation falls in that interval the sample has no label and is discarded.

The embargo is the operational form of the rule derived in \cref{sec:principle}: the label must lie outside the input window by construction, so that the model cannot recover it (\cref{fig:window-embargo}).

\begin{figure}[htbp]
\centering
\begin{tikzpicture}[font=\footnotesize]
  \draw[thick,-{Stealth[length=2mm]}] (0,0) -- (11.4,0) node[right,font=\scriptsize]{time};
  % input window
  \fill[cRef!18] (0.6,-0.3) rectangle (5.4,0.3);
  \node[anchor=south,font=\scriptsize,text=cRef!60!black] at (3.0,0.35) {input window, length $L$};
  \foreach \x in {0.9,2.2,3.4,4.0,5.1} \node[circle,fill=cRef,inner sep=1.6pt] at (\x,0) {};
  % embargo
  \fill[cGrey!30] (5.4,-0.3) rectangle (6.4,0.3);
  \node[anchor=north,font=\scriptsize,text=cGrey!70!black] at (5.9,-0.35) {embargo $\varepsilon$};
  % label search interval
  \fill[cAfrica!16] (6.4,-0.3) rectangle (9.6,0.3);
  \node[anchor=south,font=\scriptsize,text=cAfrica] at (8.0,0.35) {label search interval, to $t+\Delta$};
  \node[circle,fill=cAfrica,inner sep=2.2pt] at (7.5,0) {};
  \node[anchor=north,font=\scriptsize,text=cAfrica] at (7.5,-0.35) {label};
  % markers
  \draw[dotted] (5.4,-0.55) -- (5.4,0.55); \node[font=\scriptsize,anchor=north] at (5.4,-0.6) {$t$};
  \draw[dotted] (9.6,-0.55) -- (9.6,0.55); \node[font=\scriptsize,anchor=north] at (9.6,-0.6) {$t+\Delta$};
  \node[anchor=west,font=\scriptsize,align=left] at (0.0,-1.25)
    {retained only if $n^{(c)}_W \geq n^{(c)}_{\min}$ for the required channels,\\
     and if at least one observation falls in the label search interval};
\end{tikzpicture}
\caption{Construction of one training sample. Irregular spacing means the label is the first qualifying observation inside the search interval, not a value at a fixed instant.}\label{fig:window-embargo}
\end{figure}

\subsubsection{Time-based, not count-based}

\begin{table}[htbp]
\centering\small
\caption{Window definitions compared.}\label{tab:window-type}
\begin{tabularx}{\textwidth}{@{}>{\raggedright\arraybackslash}p{3.4cm} X X@{}}
\toprule
 & \textbf{Count-based: last $k$ observations} & \textbf{Time-based: last $L$ days} \\
\midrule
Information content & constant & variable \\
Time span & variable & constant \\
Can be empty & no & yes \\
Clinical meaning & weak: ``the last five readings'' may span a year & strong: ``the seven-day mean'' is the protocol unit \\
\bottomrule
\end{tabularx}
\end{table}

Time-based windows are correct here, because every clinical reference in this work is expressed in time: the seven-day profile of \textcite{mihevc2025}, the one-to-four-week confirmation visit of \cref{tab:protocol}, the monthly follow-up until target of \cref{sec:followup-requires}. A count-based window of five readings could span four days in an adherent patient and seven months in a lapsed one, and the resulting feature would mean different things for different patients.

The cost is that a time-based window may contain too few observations to compute a feature. The response is a minimum-observation condition, retained if and only if $n^{(c)}_W \geq n^{(c)}_{\min}$ for the required channels, with $n_{\min} \geq 2$ for any level feature and $n_{\min} \geq 3$ for any slope. Samples failing the condition are discarded, and the number discarded is reported: it is a direct measure of adherence and of how much the dataset shrinks under the protocol.

\subsubsection{Window length, and what each length can support}\label{sec:window-length}

The choice of $L$ is constrained by the minimum-observation condition just stated. At the routine frequency of two retained blood pressure values per week, \cref{tab:window-length} gives what each candidate length can actually support.

\begin{table}[htbp]
\centering\small
\caption{Window length against feature feasibility, at two retained readings per week.}\label{tab:window-length}
\begin{tabularx}{\textwidth}{@{}c c X@{}}
\toprule
\textbf{$L$ (days)} & \textbf{Expected $n_W$} & \textbf{Features supported under $n_{\min}$} \\
\midrule
7 & 2 & level only; no slope ($n_{\min} = 3$ not met), no variability \\
14 & 4 & level and slope; variability marginal \\
28 & 8 & level, slope and variability \\
\bottomrule
\end{tabularx}
\end{table}

This corrects a natural assumption. Seven days is the obvious candidate because it is the verification profile of \textcite{mihevc2025} and the averaging period that \cref{sec:protocols} specified for the reference value---but at routine frequency a seven-day window holds two observations and cannot support the slope feature on which the entire anticipatory contribution rests. The working choice is therefore $L = 14$ days for the routine regime, with $L = 7$ retained only inside the verification regime, where the protocol prescribes two readings per day and a seven-day window holds fourteen. Results are reported for more than one $L$, as a cheap sensitivity analysis that pre-empts the question.%
\verify{The routine frequency is read here as two retained blood pressure values per week, consistent with the rate used in \cref{sec:horizon}. The source is ambiguous between two measurement sessions and two sessions each yielding a morning and an evening value, which would double the count and make $L=7$ viable. This is the same ambiguity flagged in \cref{sec:protocols} and it must be resolved before $L$ is fixed.}

\subsubsection{Overlapping windows and the independent sample count}\label{sec:overlap}

With approximately ten observations per patient, the temptation to use a small stride and generate overlapping windows is strong, because it multiplies the apparent sample count. Overlapping windows are near-duplicates: two windows sharing three of four observations are not independent samples. Three consequences follow.

Overlapping windows must never straddle the train--test boundary; a window in the test set sharing observations with a training window leaks. The effective sample size for any confidence interval is the number of independent windows, not the number generated. And where overlap is used to augment training data, it is used in training only, with validation and test built from non-overlapping windows.

The arithmetic deserves to be stated, because it sharpens the capacity estimate of \cref{sec:capacity}. A patient contributing ten observations at two per week spans about thirty-five days. The number of \emph{non-overlapping} windows of length $L$ is approximately $\mathrm{span}/L$, which gives two to three per patient at $L = 14$, or sixty to ninety across thirty patients. The figure of 180 usable samples quoted in \cref{sec:capacity} is the overlapping count, obtained by sliding the window forward one observation at a time; the independent count is between a third and a half of it. Every confidence interval in the experimental chapter is computed on the independent count, and both figures are reported.

%==========================================================================
\subsection{Missingness: mechanism, and why imputation is the wrong default}\label{sec:missingness}

\subsubsection{The three mechanisms}

\begin{table}[htbp]
\centering\small
\caption{Missingness mechanisms and the validity of imputation.}\label{tab:missing-mechanisms}
\begin{tabularx}{\textwidth}{@{}c X c@{}}
\toprule
\textbf{Mechanism} & \textbf{Definition} & \textbf{Imputation valid?} \\
\midrule
MCAR & missingness independent of all values & yes \\
MAR & missingness depends on observed values & yes, under the right model \\
MNAR & missingness depends on the unobserved value itself & no \\
\bottomrule
\end{tabularx}
\end{table}

\subsubsection{Which mechanism governs each stream}

\Cref{sec:pipelines,sec:streams} established the answers, and they are unfavourable.

\textbf{Stream A is MNAR.} A record exists because the patient attended. \Cref{sec:followup-requires} established attendance far below the required frequency, and \textcite{bediang2021} measured 33.3\% three-month attendance in Cameroonian controls. Attendance depends on the patient's state: deterioration may drive attendance up, or may drive the patient out of the system entirely. The probability that a value is observed depends on the value.

\textbf{Stream B is MNAR by a different route.} \Cref{sec:pipelines} identified selective omission: a patient may decline to record a reading they find discouraging. The omitted readings are disproportionately the abnormal ones.

\textbf{Stream C is approximately MCAR.} Battery exhaustion, connectivity loss and device failure are, to a first approximation, independent of the physiological value.

\subsubsection{The direction of the bias}\label{sec:bias-direction}

Suppose observation occurs with probability depending on the value, and that extreme values are less likely to be recorded. Then
\begin{equation}
  \mathbb{E}\bigl[x \mid \text{observed}\bigr] \neq \mathbb{E}[x],
\end{equation}
and specifically, if omission is concentrated on high values, $\mathbb{E}[x \mid \text{observed}] < \mathbb{E}[x]$.

Mean imputation fills every gap with $\bar{x}_{\text{observed}}$, which is biased toward normality. Forward-fill carries the last observed value, which---since the patient was willing to record it---is also disproportionately a normal one. Either way, the model learns that absence predicts stability. The truth is the reverse: absence predicts deterioration.

This is not a quantitative inaccuracy that larger samples would fix. It is a sign error in the learned relationship, and it is unrecoverable by any amount of data under the same mechanism.

It is also, independently, a clinical error, and here the two analyses converge. \Cref{sec:ltfu} established loss to follow-up at 50\% or more as the central mechanism of care-cascade failure. \Cref{sec:decision-flow} established that non-transmission must itself be an alarmed event. A preprocessing step that imputes silence as stability contradicts the system's own alerting logic.

\subsubsection{Why the reported pipelines do not transfer}

\Textcite{liu2025review} describe a pipeline that \bq{addressed missing data by imputing features, standardized the range of dataset values through scaling, selected features to eliminate redundancy with minimal contribution to the prediction outcome, and augmented features to correct class imbalance}. That work uses the Pima Indians Diabetes Database---a cross-sectional dataset in which missingness is a recording artefact, not a behavioural signal. Imputation is defensible there and is not defensible here. It is also the same dataset on which \textcite{chatrati2022} trained, and whose transferability \cref{sec:chatrati} found undiscussed.

They also report a differentially private synthesis using autoencoders and generative adversarial networks, \bq{capable of capturing temporal information and feature correlations inherent in the original patient data}. Synthetic generation is a plausible response to small sample size, but it inherits whatever bias the source data carry: a generator trained on MNAR-missing data reproduces the missingness mechanism rather than correcting it. It is not a solution to this problem and is not presented as one.

%==========================================================================
\subsection{Masking: representing absence as information}\label{sec:masking}

The framing document commits to masking, and the correct implementation inverts the problem: absence is not a hole to be filled but a channel to be read.

\subsubsection{The representation}

For each channel $c$ and each sample, the feature vector carries, alongside whatever values are present:

\begin{table}[htbp]
\centering\small
\caption{Missingness features, per channel.}\label{tab:mask-features}
\begin{tabularx}{\textwidth}{@{}>{\raggedright\arraybackslash}p{4.2cm} X@{}}
\toprule
\textbf{Feature} & \textbf{Meaning} \\
\midrule
$m^{(c)} \in \{0,1\}$ & was channel $c$ observed in this window at all \\
$n^{(c)}_W$ & how many times \\
$\Delta t^{(c)}_{\text{last}}$ & how long since the last observation \\
$r^{(c)} = n^{(c)}_W / n^{(c)}_{\text{expected}}$ & adherence ratio against the protocol \\
\bottomrule
\end{tabularx}
\end{table}

The fourth deserves emphasis. The protocol specifies a known number of measurements per week, so the expected count in a window of length $L$ is known, and the ratio of observed to expected is a direct, interpretable measure of protocol adherence available at every time point without additional instrumentation. Given that \cref{sec:cameroon-joint} established retention as the one robust finding across both Cameroonian studies, a feature that measures it continuously is likely to be among the most informative in the set---and it exists only because missingness was treated as signal.

\subsubsection{Model-specific handling}\label{sec:mask-models}

\textbf{Tree ensembles handle this natively, and that is an argument in their favour.} XGBoost learns a default branch direction at each split for instances with a missing value, so missingness is used as a splitting signal without any imputation at all. Combined with the explicit mask features above, a gradient-boosted model receives the missingness information twice: implicitly through the default direction and explicitly through the mask channel.

This is a substantive point for the hypotheses of \cref{sec:forecasting}. \Cref{sec:capacity} established that the sample size makes a negative recurrent-model result near-certain; this subsection adds a second and independent reason to expect tree ensembles to do well here, namely that the dominant irregularity of the data is one that trees absorb and that recurrent networks must be specially modified to accommodate.

\textbf{Recurrent models require explicit construction.} A sequence model consumes a fixed-width vector at each step and has no native representation of ``not measured''. Three elements are needed: the mask as an additional input channel; a decay on the carried value, so that the influence of the last observed reading diminishes as the gap widens rather than persisting unchanged; and loss masking, so that no gradient is computed on positions imputed for shape.

A decay formulation implementing the second, proposed here as a design choice and not taken from any source in this corpus:
\begin{equation}
  \tilde{x}^{(c)}_t = m^{(c)}_t x^{(c)}_t
  + \bigl(1 - m^{(c)}_t\bigr)\Bigl[\gamma^{(c)}_t\, x^{(c)}_{t'}
  + \bigl(1-\gamma^{(c)}_t\bigr)\,\mu^{(c)}_i\Bigr],
  \qquad
  \gamma^{(c)}_t = \exp\bigl(-\kappa_c\,\Delta t^{(c)}_{\text{last}}\bigr),
\end{equation}
where $x^{(c)}_{t'}$ is the last observed value of the channel, $\mu^{(c)}_i$ the patient's running mean and $\kappa_c > 0$ a learned or fixed decay rate. For a short gap $\gamma \to 1$ and the last value carries; for a long gap $\gamma \to 0$ and the representation falls back to the patient's own baseline rather than asserting that the last reading still holds.

The decay rate should differ by channel, because the channels have different physiological persistence: a blood pressure reading three days old is more informative about today than a fasting glucose reading three days old. Setting $\kappa_c$ per channel, and reporting the fitted values, is itself a result.

One caution on this construction. The fallback target $\mu^{(c)}_i$ is the patient's own baseline, so a long gap returns the representation to that patient's normal---which is, in a weaker form, the same assertion of stability that \cref{sec:bias-direction} rejects. The mask channel and the gap feature are what prevent the model from reading that fallback as evidence; they must be present, and the decay must never be used without them.

\subsubsection{The one case where imputation is acceptable}

Sensor-stream gaps, which are approximately MCAR, and only for covariate channels carrying no clinical decision. Even there the mask is retained and the imputation is fitted on training data only.

Blood pressure and blood glucose are never imputed. They carry the decision; absence is reported as absence.

%==========================================================================
\subsection{Per-patient normalisation without leakage}\label{sec:normalisation}

The framing document commits to per-patient normalisation. It is the right instinct---\cref{sec:variability} established that the clinically meaningful quantities are each patient's own habitual level $\mu_i$ and variability $\sigma_{\varepsilon,i}$, not a population constant---but the naive implementation leaks.

\subsubsection{The leak}

Computing a patient's mean and dispersion over their whole series, then normalising and splitting, gives every training sample access to statistics computed partly from test-period observations. The leak is small in magnitude and systematic in direction, and it inflates the result.

\subsubsection{Expanding-window normalisation}\label{sec:expanding}

The correct construction is causal: at time $t$, normalise using only observations strictly before $t$,
\begin{gather}
  \hat\mu^{(c)}_{i,t} = \operatorname{median}\bigl\{x^{(c)}_{i,s} : s < t\bigr\},
  \qquad
  \hat\sigma^{(c)}_{i,t} = 1.4826\cdot\operatorname{MAD}\bigl\{x^{(c)}_{i,s} : s < t\bigr\},\\
  z^{(c)}_{i,t} = \frac{x^{(c)}_{i,t} - \hat\mu^{(c)}_{i,t}}
                       {\max\bigl(\hat\sigma^{(c)}_{i,t},\ \sigma_{\text{floor}}\bigr)}.
\end{gather}

This has four properties that make it the right choice independently of the leakage argument. It is leak-free by construction, since no future observation enters any statistic. It works for patients never seen in training, which the naive approach does not: under group $k$-fold a test-fold patient has no training statistics, and expanding-window normalisation simply builds them from that patient's own early observations. It matches deployment, since in production the system has only the past. And it degrades gracefully when combined with the shrinkage of \cref{sec:mad-shrinkage}: early in a patient's series $\hat\sigma$ is dominated by the pooled estimate, and as observations accumulate it becomes progressively personal.

Two practical requirements follow. A burn-in period is needed---the first $n_{\text{burn}}$ observations of each patient cannot be normalised against their own history and must either be dropped or normalised against the pooled statistics; the choice is stated and the number dropped reported. And $\sigma_{\text{floor}}$ must be set to prevent division by a near-zero dispersion, which arises exactly in the carried-forward sequences that \cref{sec:ingestion} flags.

%==========================================================================
\subsection{The feature set}\label{sec:features}

The framing document specifies four families---rolling aggregates, trends, variability, threshold deviations---to which this section adds a fifth, missingness, and a sixth, static context.

\begin{table}[htbp]
\centering\footnotesize
\caption{The feature set, per channel $c$ and window $W$ of length $L$.}\label{tab:feature-set}
\begin{tabularx}{\textwidth}{@{}>{\raggedright\arraybackslash}p{1.9cm} >{\raggedright\arraybackslash}p{4.0cm} X@{}}
\toprule
\textbf{Family} & \textbf{Feature} & \textbf{Definition} \\
\midrule
Level & last value & $x^{(c)}_{t}$ \\
 & window median & $\operatorname{median}\{x^{(c)}_s : s \in W\}$ \\
 & reference mean & mean of retained second readings over seven days (\cref{sec:protocols}) \\
\addlinespace
Trend & slope & OLS coefficient of $x$ on $t$ over $W$, in units per day \\
 & net change & $\bigl(x^{(c)}_{t} - x^{(c)}_{t-L}\bigr)/L$ \\
\addlinespace
Variability & robust dispersion & $1.4826\cdot\operatorname{MAD}$ over $W$ \\
 & mean absolute successive difference & $\frac{1}{n-1}\sum_{k=2}^{n}\bigl\lvert x^{(c)}_k - x^{(c)}_{k-1}\bigr\rvert$ \\
\addlinespace
Deviation & personal $z$ & $z^{(c)}_{i,t}$ per \cref{sec:expanding}, expanding window \\
 & guideline margin & $x^{(c)}_{t} - \theta_c$, lagged only \\
\addlinespace
Missingness & mask, count, gap, adherence ratio & \cref{tab:mask-features} \\
\addlinespace
Context & time of day, fasting state, operator, entry mode & from the record schema \\
\addlinespace
Static & age band, sex, BMI, years since diagnosis, regimen version, drug count & \cref{tab:patient-record} \\
\bottomrule
\end{tabularx}
\end{table}

Three notes.

The \textbf{guideline-margin} feature is the one most likely to leak, since $\theta_c$ is the same quantity used to construct the label (\cref{tab:threshold-allocation}). It is admissible only when computed on a value strictly inside the input window, never at the label instant. If its inclusion proves difficult to police, the safest course is to omit it: the personal $z$ carries most of the same information without the exposure.

\textbf{Regimen version and drug count} are among the most informative static features available and are frequently omitted. \Cref{sec:inertia} established therapeutic inertia as the mechanism behind the weakest link of the cascade; a feature encoding whether therapy has recently changed distinguishes a drifting patient under an unchanged regimen from one whose treatment was just intensified.

\textbf{Years since diagnosis} is the best-evidenced feature in the entire set. \Textcite{mendoza2025} found it the strongest predictor of uncontrolled hypertension by permutation importance, and it costs one integer to record.

%==========================================================================
\subsection{What must be reported}\label{sec:reporting}

A preprocessing chapter that does not quantify what it did is not checkable.

\begin{table}[htbp]
\centering\footnotesize
\caption{Quantities reported for the preprocessing stage.}\label{tab:preprocessing-report}
\begin{tabularx}{\textwidth}{@{}X >{\raggedright\arraybackslash}p{5.2cm}@{}}
\toprule
\textbf{Quantity} & \textbf{Why} \\
\midrule
Patients enrolled, patients retained, reasons for loss & 10.9\% equipment-related dropout in \textcite{mihevc2025} \\
Observations per patient: median, interquartile range, minimum, maximum & the sample-size claim of \cref{sec:capacity} \\
Inter-observation interval per channel: distribution & justifies the window length \\
Observed-to-expected ratio per channel & protocol adherence \\
Missingness per channel, overall and per patient & \cref{sec:missingness} \\
Flag counts by type & \cref{sec:ingestion} \\
Windows generated, windows discarded for $n_{\min}$, with reason & \cref{sec:windows} \\
Independent against overlapping window counts & \cref{sec:overlap} \\
Burn-in observations dropped & \cref{sec:expanding} \\
Class distribution before and after any balancing & \cref{sec:imbalance} \\
\bottomrule
\end{tabularx}
\end{table}

The observed-to-expected ratio is the row to foreground. It is simultaneously a data-quality statistic, a clinical finding about adherence in this cohort, and a model feature. It is the kind of number a jury will ask for and that cannot be reconstructed after the fact.

%==========================================================================
\subsection{Specification}\label{sec:spec32}

\begin{table}[htbp]
\centering\footnotesize
\caption{Preprocessing requirements established in \cref{sec:preprocessing}.}\label{tab:spec32}
\begin{tabularx}{\textwidth}{@{}X >{\raggedright\arraybackslash}p{5.2cm}@{}}
\toprule
\textbf{Requirement} & \textbf{Basis} \\
\midrule
Pipeline partitioned into deterministic and learned phases & \cref{sec:forecasting,sec:pipeline} \\
Pipeline implemented as an object, not as notebook cell order & \cref{sec:pipeline} \\
Flags retained, never silently deleted & \cref{sec:ingestion} \\
Impossible values excluded; extreme values retained & \cref{tab:outlier-categories} \\
Statistical outlier detection flags for review, never auto-deletes & \cref{sec:outliers} \\
MAD shrunk toward a pooled estimate fitted on training only & \cref{sec:mad-shrinkage} \\
No resampling onto a regular grid & \cref{sec:ts-diagnostics,sec:irregular} \\
Inter-observation intervals encoded per channel & \cref{tab:interval-features} \\
Periodicity tested with methods valid for uneven sampling & \cref{sec:irregular} \\
Time-based windows with a minimum-observation condition & \cref{sec:windows} \\
$L = 14$ days in the routine regime, $L = 7$ in the verification regime & \cref{tab:window-length} \\
Embargo between window end and label search interval & \cref{sec:principle,sec:windows} \\
Overlapping windows confined to training; independent count reported & \cref{sec:overlap} \\
Blood pressure and blood glucose never imputed & \cref{sec:masking} \\
Missingness encoded as mask, count, gap and adherence ratio & \cref{tab:mask-features} \\
Per-channel decay for recurrent models, used only with mask and gap; $\kappa_c$ reported & \cref{sec:mask-models} \\
Imputation permitted only for MCAR covariate channels, mask retained & \cref{sec:masking} \\
Expanding-window per-patient normalisation & \cref{sec:expanding} \\
Burn-in stated and count reported & \cref{sec:expanding} \\
Dispersion floor to guard carried-forward sequences & \cref{sec:expanding} \\
Guideline-margin feature lagged only, or omitted & \cref{sec:features} \\
Full reporting table of \cref{tab:preprocessing-report} included & \cref{sec:reporting} \\
\bottomrule
\end{tabularx}
\end{table}

%==========================================================================
\subsection{Dependency on the collection scenario}\label{sec:scenario-dependency}

This section is written parametrically in $L$, $\Delta$, $\varepsilon$ and $n_{\min}$, and those parameters are fixed by the scenario that \cref{tab:scenarios} sets out. \Cref{tab:scenario-params} states the consequence of each.

\begin{table}[htbp]
\centering\footnotesize
\caption{Preprocessing parameters by collection scenario.}\label{tab:scenario-params}
\begin{tabularx}{\textwidth}{@{}>{\raggedright\arraybackslash}p{2.6cm} >{\raggedright\arraybackslash}p{2.0cm} >{\raggedright\arraybackslash}p{2.4cm} X@{}}
\toprule
\textbf{Scenario} & \textbf{$L$} & \textbf{$\Delta$} & \textbf{Consequence} \\
\midrule
A: retrospective outpatient & not applicable & next consultation & a 14-day window is empty for almost every sample; the design reduces to visit-to-visit features and the recurrent model is not applicable \\
B: prospective home & 14 days & next measurement or next 7-day window & workable; the discard rate under $n_{\min}$ becomes a headline statistic \\
C: mixed & two definitions & per stream & the streams require different window definitions and must be joined at the feature level, not the observation level; this is the only configuration in which the three-pipeline design is actually exercised \\
D: retrospective inpatient & 3--7 days & 24--48 h & daily observation makes short windows viable; the series is short, so few windows per patient \\
\bottomrule
\end{tabularx}
\end{table}

The scenario is declared at the head of the experimental chapter, once the data are in hand, and the parameters of \cref{tab:spec32} are fixed accordingly at that point. Nothing in the present section changes with the scenario except the numerical values of those four parameters; the ordering constraint, the treatment of missingness, the normalisation and the reporting requirements hold in all four cases.

%##########################################################################
% SECTION 3.3
%##########################################################################
\section{Feature Engineering for Physiological Parameters}\label{sec:feature-eng}

\subsection*{Purpose, sources and a notational convention}

\Cref{sec:preprocessing} specified how raw measurements become windows. This section specifies what is computed from each window and how those quantities are assembled into the vector the model consumes. It is the last stage at which a design error remains invisible: a badly constructed feature set produces a model that trains, scores and reports, and that is wrong in ways no metric reveals.

Two sources are used. \Textcite{cadrage2026} commits this work to rolling aggregates, trends, variability and threshold deviations, and anchors the labelling in the ESH and ADA guidelines. \Textcite{liu2025review} describe feature engineering within the data-processing layer and survey feature-selection methods used in the chronic-disease literature; the standing reservation on their method applies, and the methods they name are used because each is independently verifiable.

Two findings from earlier chapters govern everything here. \Cref{sec:principle} established that a feature computed at the label instant converts prediction into label recovery. \Cref{sec:overlap} established that the number of \emph{independent} windows available is between sixty and ninety, and that figure is the binding constraint of \cref{sec:budget}.

A notational convention is fixed here and used for the remainder of the thesis. The symbol $\lambda$ denotes the penalty parameter of a regularisation path, as in the LASSO procedure of \textcite{mendoza2025}. The current value of a channel is written $x_{\text{last}}$, and the shrinkage weight introduced in \cref{sec:mad-shrinkage} is written $w_i$.

%==========================================================================
\subsection{The information ceiling of a window}\label{sec:ceiling}

This subsection derives a constraint that determines the window length, and it is the reason the rest of the section is shaped as it is.

A window of channel $c$ contains $n$ numbers. Every descriptor computed from it---mean, standard deviation, slope, last value, median, dispersion, successive difference---is a function of those same $n$ numbers. No set of descriptors can carry more than $n$ independent quantities. Computing seven descriptors from a two-point window produces seven columns, not seven pieces of information.

\subsubsection{The degeneracy at two observations}

The degeneracy is exact and can be shown. For $n = 2$, with observations $x_1, x_2$ at times $t_1, t_2$, the mean is $\bar{x} = (x_1+x_2)/2$ and $x_{\text{last}} = x_2$. The sample standard deviation is
\begin{equation}
  s = \sqrt{\frac{(x_1-\bar x)^2 + (x_2-\bar x)^2}{n-1}}
    = \sqrt{\frac{(x_2-x_1)^2}{2}}
    = \frac{\lvert x_2 - x_1\rvert}{\sqrt{2}},
\end{equation}
while the mean absolute successive difference is $\mathrm{MASD} = \lvert x_2 - x_1 \rvert$, so that
\begin{equation}
  \boxed{\;s = \frac{\mathrm{MASD}}{\sqrt{2}}\;}
\end{equation}
exactly. The two are proportional with correlation one, by construction and not by coincidence. The ordinary least squares slope is
\begin{equation}
  \beta = \frac{x_2 - x_1}{t_2 - t_1} = \frac{\pm\,\mathrm{MASD}}{\Delta t},
\end{equation}
which adds only the sign and a rescaling by the interval. And the median of two points is their mean.

With two observations, the four descriptors the framing document specifies---level, dispersion, slope, last value---are not four quantities. They are two: a level and a signed difference. Supplying all four provides two pieces of information in four columns, which inflates the dimension without inflating the content, and \cref{sec:budget} shows that dimension is the scarce resource here (\cref{fig:degeneracy}).

\begin{figure}[htbp]
\centering
\begin{tikzpicture}
\begin{axis}[width=0.62\textwidth,height=5.2cm,
  xlabel={$\mathrm{MASD}$ (mmHg)},ylabel={$s$ (mmHg)},
  xmin=0,xmax=30,ymin=0,ymax=22,grid=major,grid style={gray!20},
  legend pos=north west,legend style={font=\scriptsize}]
\addplot[thick,cAfrica,domain=0:30,samples=2]{x/sqrt(2)};
\addplot[only marks,mark=*,mark size=2pt,cRef]
  coordinates {(4,2.83)(9,6.36)(14,9.90)(21,14.85)(27,19.09)};
\legend{$s=\mathrm{MASD}/\sqrt{2}$, two-point windows}
\end{axis}
\end{tikzpicture}
\caption{At $n=2$ every two-point window lies exactly on the line $s = \mathrm{MASD}/\sqrt{2}$. The two descriptors are one quantity in two columns.}\label{fig:degeneracy}
\end{figure}

\subsubsection{The consequence for window length}

The rule that follows resolves a parameter that would otherwise be set arbitrarily: $L$ must satisfy $n_W^{(c)}(L) \geq k_c$, where $k_c$ is the number of descriptors intended for channel $c$. Under the routine protocol of two blood pressure values and one glucose value per week, \cref{tab:ceiling} gives the ceiling.

\begin{table}[htbp]
\centering\small
\caption{Observations per window and the ceiling on independent descriptors.}\label{tab:ceiling}
\begin{tabularx}{\textwidth}{@{}c c c X@{}}
\toprule
\textbf{$L$} & \textbf{$n_W$, blood pressure} & \textbf{$n_W$, glucose} & \textbf{Independent descriptors available, blood pressure} \\
\midrule
7 days & 2 & 1 & 2 \\
14 days & 4 & 2 & 4 \\
28 days & 8 & 4 & up to 8 \\
\bottomrule
\end{tabularx}
\end{table}

To support the four descriptors the framing document specifies, the window must be at least fourteen days---and glucose reaches four observations only at twenty-eight. This confirms from an independent direction the choice of $L = 14$ made in \cref{tab:window-length}, which was derived there from the minimum-observation condition rather than from the information ceiling. The seven-day window inherited from the verification profile of \textcite{mihevc2025} is the right unit for \emph{alerting}, where it averages readings to reduce variance, and is too short for \emph{feature construction}, where it supports two descriptors. These are different operations and use different windows; stating that explicitly is better than silently using one length for both.

%==========================================================================
\subsection{Statistical descriptors}\label{sec:descriptors}

\subsubsection{Definitions}

For channel $c$, patient $i$, window $W$ of length $L$ ending at $t$ and containing $n$ observations $\{(t_k, x_k)\}$:

\textbf{Level.}
\begin{equation}
  \mu^{(c)} = \frac{1}{n}\sum_{k=1}^{n} x_k,
  \qquad
  \tilde{x}^{(c)} = \operatorname{median}_k x_k,
  \qquad
  x^{(c)}_{\text{last}} = x_n.
\end{equation}

\textbf{Trend.} The ordinary least squares slope on time, which is correct for irregular spacing and requires no resampling:
\begin{equation}
  \beta^{(c)} = \frac{\sum_{k}(t_k - \bar t)(x_k - \bar x)}{\sum_{k}(t_k - \bar t)^2},
\end{equation}
in units of the channel per day. This estimator is undefined when all $t_k$ coincide and unstable when they nearly do, which occurs when a patient records several readings in one sitting and none in the rest of the window. Its sampling variance is proportional to $\sigma^2 / \sum_k (t_k - \bar t)^2$, so it explodes precisely when the window is temporally clustered. A guard is therefore required: $\beta$ is computed only when $\mathrm{span}_W^{(c)} \geq L/2$, and the mask is set otherwise.

\textbf{Variability.}
\begin{equation}
  s^{(c)} = \sqrt{\frac{1}{n-1}\sum_k (x_k - \bar x)^2},
  \qquad
  \mathrm{MAD}^{(c)} = \operatorname{median}_k\bigl\lvert x_k - \tilde x\bigr\rvert,
  \qquad
  \mathrm{MASD}^{(c)} = \frac{1}{n-1}\sum_{k=2}^{n}\bigl\lvert x_k - x_{k-1}\bigr\rvert.
\end{equation}
The third measures variation between consecutive readings rather than around the mean, and it is the only one of the three that uses the ordering. For $n > 2$ it is not a rescaling of $s$: a series that oscillates has high MASD and moderate $s$, while one that drifts monotonically has low MASD relative to $s$. At $n = 2$ the distinction vanishes, which is a further reason to prefer $L \geq 14$ days.

\textbf{Personal deviation.} The expanding-window standardised deviation $z^{(c)}_{i,t}$ of \cref{sec:expanding}.

\subsubsection{Small-sample behaviour}\label{sec:small-sample}

Each estimator degrades differently as $n$ falls, and the differences matter at $n = 2$ to $4$.

\begin{table}[htbp]
\centering\small
\caption{Behaviour of each descriptor at small $n$.}\label{tab:small-sample}
\begin{tabularx}{\textwidth}{@{}>{\raggedright\arraybackslash}p{3.0cm} c X@{}}
\toprule
\textbf{Descriptor} & \textbf{Minimum $n$} & \textbf{Behaviour at small $n$} \\
\midrule
$\mu$, $x_{\text{last}}$ & 1 & stable \\
$\tilde{x}$ & 1 & equals $\mu$ at $n=2$; no robustness gained \\
$s$ & 2 & very high variance, $\operatorname{Var}(s^2) \propto 1/(n-1)$ \\
MAD & 3 & degenerate at $n=2$, where it equals half the range \\
$\beta$ & 2 & variance $\propto \sigma^2/\sum_k(t_k-\bar t)^2$; explodes on a temporally clustered window \\
MASD & 2 & equals $\sqrt{2}\,s$ at $n=2$ \\
\bottomrule
\end{tabularx}
\end{table}

The robust descriptors---median and MAD---are chosen for their resistance to outliers, but that resistance requires observations to be resistant with. At $n = 2$ the median is the mean and the MAD is half the range, so nothing is gained and a column is spent. Robust statistics belong in this design at the level of the patient's history, where $n$ is larger (\cref{sec:expanding}), and not at the level of a short window.

\subsubsection{Why learned representations are not used}

The account that \textcite{liu2025review} give of physiological feature extraction is largely about deep unsupervised representation learning. They report a convolutional model applied \bq{using an unsupervised feature learning method to learn meaningful feature representations from unlabeled physiological signals}, with the aim of eliminating \bq{the limitations of current feature extraction and feature selection in physiological signal anomaly detection}, the representations then feeding an anomaly detector.

The architectural goal is the same as this work's. The method does not transfer. Learned representations require sample sizes orders of magnitude beyond those available here, and convolutional feature learning is designed for densely sampled signals---electrocardiograms, photoplethysmographic waveforms---not for two blood pressure readings per week. That direction is recorded here as where the field is going, and hand-constructed descriptors are adopted as the correct choice at this sample size. The reasoned rejection is itself part of the methodological contribution, since it is a decision taken rather than an option overlooked.

%==========================================================================
\subsection{Clinical threshold features}\label{sec:threshold-features}

\subsubsection{The family}

For channel $c$ with guideline threshold $\theta_c$, computed on window observations only:

\begin{table}[htbp]
\centering\footnotesize
\caption{Threshold-derived features.}\label{tab:threshold-features}
\begin{tabularx}{\textwidth}{@{}>{\raggedright\arraybackslash}p{3.0cm} >{\raggedright\arraybackslash}p{5.0cm} X@{}}
\toprule
\textbf{Feature} & \textbf{Definition} & \textbf{Interpretation} \\
\midrule
Exceedance indicator & $\mathbb{1}\bigl[x_{\text{last}} > \theta_c\bigr]$ & binary, current \\
Signed margin & $x_{\text{last}} - \theta_c$ & how far, in clinical units \\
Normalised margin & $(x_{\text{last}} - \theta_c)/\theta_c$ & comparable across channels \\
Time-above fraction & $\frac{1}{n}\sum_k \mathbb{1}[x_k > \theta_c]$ & proportion of the window out of range \\
Exceedance burden & $\sum_{k=2}^{n} \max\bigl(0,\ \bar{x}_{k-1,k} - \theta_c\bigr)(t_k - t_{k-1})$ & magnitude $\times$ duration, in mmHg$\cdot$days \\
Grade band & ordinal index of the guideline band containing $x_{\text{last}}$ & preserves the graded structure \\
\bottomrule
\end{tabularx}
\end{table}

The fifth deserves comment. Exceedance burden integrates how far above threshold the patient was and for how long, approximating the integral of $\max(0, x-\theta)$ by evaluating the excess at the midpoint of each inter-observation interval. It is the natural quantity when the clinical concern is cumulative exposure rather than instantaneous level---and \cref{sec:transition} established exactly that framing, since the WHO burden figures are expressed in disability-adjusted life years, which are themselves a magnitude integrated over time. Two qualifications. The midpoint approximation differs from the true integral on any segment that crosses $\theta$, and the error is small only when the inter-observation interval is short, which at routine frequency it is not. And no source in this corpus defines this quantity: it is presented as an original construction to be evaluated, not as an established index.

The sixth matters because guidelines are graded, not binary. Collapsing a graded classification to a single exceedance discards the distinction between grade~1 and grade~3, which is precisely the distinction that determines urgency---and \textcite{kingue2013}, reported in \cref{sec:telemedcam}, found the effect of their intervention confined to stage~III hypertension and absent for stages~I and~II. A feature set that cannot represent the grade cannot represent where the effect lives.

\subsubsection{The leakage boundary, precisely}\label{sec:leak-boundary}

\Cref{sec:principle} established the rule. Applied here it requires a distinction that is easily blurred.

\textbf{Admissible:} the exceedance of the threshold computed on observations inside the input window. If the label is a threshold crossing at $t+\Delta$ and the feature is the exceedance at $t$, the two are correlated through the autocorrelation of the physiology, which is a real relationship and not a leak.

\textbf{Inadmissible:} the exceedance computed at or after the label instant, or on any observation inside the embargo of \cref{sec:windows}.

The distinction is enforced by construction: features are computed from the window, labels from a disjoint interval, and no code path reads the label interval during feature computation. This is enforced by the pipeline object of \cref{sec:pipeline} and not by discipline.

\subsubsection{The threshold table as configuration}\label{sec:threshold-table}

\Cref{sec:interop} required every rule to carry its version and date. The same applies here: thresholds are a configuration artefact, not constants embedded in feature code. The table carries, per row, the channel, the context (home or office), the band label, the lower and upper bounds, the unit, the source, the source version and the date from which it is effective.

Three reasons. The guidelines will be superseded and the features must be recomputable under the new table. Three conventions are in play---the WHO surveillance threshold, the ESH grading and the home target of \textcite{mihevc2025}---and \cref{tab:threshold-allocation} fixed which is used for which function, so the table makes that allocation machine-readable. And \cref{sec:feature-limits} records that two of the three sources are not yet in hand, so the table will be populated later and the code must not assume otherwise.

\Cref{tab:verified-thresholds} records what the corpus does support, usable as verified placeholders until the guidelines are obtained.

\begin{table}[htbp]
\centering\footnotesize
\caption{Threshold values verifiable from the present corpus.}\label{tab:verified-thresholds}
\begin{tabularx}{\textwidth}{@{}>{\raggedright\arraybackslash}p{2.9cm} >{\raggedright\arraybackslash}p{2.0cm} X@{}}
\toprule
\textbf{Source} & \textbf{Channel} & \textbf{Values} \\
\midrule
\textcite{who2023hypertension} & office BP & surveillance $\geq$ 140/90; same-day treatment $\geq$ 160/100; confirm 140--159/90--99 at one to four weeks \\
\textcite{mihevc2025} & home BP & green 110--134/60--84; amber 135--174/85--104 and $<$110/$<$60; red $\geq$ 175/$\geq$ 105 \\
\textcite{mihevc2025} & glucose & green 4--7 pre-meal, 7--10 post-meal; amber 7.1--10 and 10.1--14; red $<$4, $>$10 pre-meal, $>$14 post-meal (mmol/L) \\
\bottomrule
\end{tabularx}
\end{table}

%==========================================================================
\subsection{Cross-channel and derived features}\label{sec:cross-channel}

The multi-parametric fusion the framing document proposes is partly this: quantities computed across channels rather than within one.

\begin{table}[htbp]
\centering\footnotesize
\caption{Cross-channel features and their justification.}\label{tab:cross-features}
\begin{tabularx}{\textwidth}{@{}>{\raggedright\arraybackslash}p{3.3cm} >{\raggedright\arraybackslash}p{3.3cm} X@{}}
\toprule
\textbf{Feature} & \textbf{Definition} & \textbf{Justification} \\
\midrule
Pulse pressure & $\mathrm{SBP} - \mathrm{DBP}$ & marker of arterial stiffness; \cref{sec:comorbidity} established stiffness as the physiological link between diabetes and hypertension \\
Mean arterial pressure & $\mathrm{DBP} + \tfrac{1}{3}(\mathrm{SBP}-\mathrm{DBP})$ & the perfusion-relevant summary \\
Glycaemic excursion & post-prandial minus fasting, same day & distinguishes an elevated baseline from poor post-prandial control \\
Rate--pressure product & $\mathrm{HR} \times \mathrm{SBP}$ & combines the two channels the system reliably acquires \\
\bottomrule
\end{tabularx}
\end{table}

A technical justification is required for including exact functions of existing features. A decision tree partitions by axis-aligned splits on single variables. To represent the condition $\mathrm{SBP} - \mathrm{DBP} > 60$ it must approximate a diagonal boundary with a staircase of splits, which costs depth and sample efficiency---and \cref{tab:inductive-bias} recorded that axis-aligned splits are precisely the inductive bias of the tree families. Supplying pulse pressure explicitly converts that staircase into a single split. For a linear model the feature adds nothing, being a linear combination of existing columns; for a recurrent model it may ease optimisation without adding information. The right answer therefore depends on the model family, and this thesis states so rather than presenting the derived features as universally beneficial.

One asynchrony caution. The glycaemic excursion requires a fasting and a post-prandial reading on the same day. Under the routine protocol both occur weekly, so the feature is computable once per week at best and its mask will be set frequently. Pulse pressure requires systolic and diastolic values, which always arrive together from the same cuff reading, so it is available whenever blood pressure is. Cross-channel features inherit the worst availability of their inputs, and their masks are computed accordingly.

%==========================================================================
\subsection{Assembling the vector}\label{sec:vector}

The feature vector for sample $(i,t)$ concatenates six blocks,
\begin{equation}
  \mathbf{v}_{i,t} = \Bigl[\;
  \underbrace{\mathbf{d}^{(c)}}_{\text{descriptors}} \;\big\|\;
  \underbrace{\boldsymbol{\theta}^{(c)}}_{\text{thresholds}} \;\big\|\;
  \underbrace{\mathbf{x}^{\text{cross}}}_{\text{cross-channel}} \;\big\|\;
  \underbrace{\mathbf{m}^{(c)}}_{\text{missingness}} \;\big\|\;
  \underbrace{\mathbf{k}}_{\text{context}} \;\big\|\;
  \underbrace{\mathbf{s}_i}_{\text{static}} \;\Bigr],
\end{equation}
with every block carrying its own availability mask per \cref{tab:mask-features}. Two assembly rules.

\textbf{Channel-wise, not timestamp-wise.} Each channel is windowed on its own schedule and its descriptors computed independently (\cref{sec:interval-features}). Concatenation happens at the feature level, after each channel has been reduced to a fixed-width summary, and never at the observation level, which would require aligning measurements taken days apart.

\textbf{Every feature carries its mask.} A missing descriptor is represented as a sentinel plus an explicit indicator, never as zero. Zero is a valid blood pressure slope.

%==========================================================================
\subsection{The dimensionality budget}\label{sec:budget}

This subsection contains the operative constraint of the section, and no source in this corpus states it.

\subsubsection{The count}

A feature set constructed without restraint reaches the count of \cref{tab:dimension-count}.

\begin{table}[htbp]
\centering\small
\caption{Feature count of an unrestrained set.}\label{tab:dimension-count}
\begin{tabularx}{\textwidth}{@{}X c@{}}
\toprule
\textbf{Block} & \textbf{Count} \\
\midrule
4 clinical channels $\times$ 7 descriptors & 28 \\
4 channels $\times$ 4 threshold features & 16 \\
Cross-channel & 4 \\
4 channels $\times$ 4 missingness features & 16 \\
\midrule
Window-dependent subtotal & 64 \\
Context & 4 \\
Static & 9 \\
\midrule
One window length & 77 \\
Three window lengths (7, 14, 28 days): $64\times3 + 13$ & 205 \\
\bottomrule
\end{tabularx}
\end{table}

Only the window-dependent blocks multiply with the number of window lengths; context and static features do not, which is why three lengths give 205 rather than $3 \times 77$.

\subsubsection{Why $p \approx n$ is fatal, and what $n$ actually is}\label{sec:np}

When the feature count approaches the sample count, any sufficiently flexible model interpolates the training data exactly. Training error goes to zero, cross-validation variance becomes enormous, and the selected model is determined by which noise directions happened to align with the labels in a particular fold. The reported performance then measures the random seed.

The working constraint for tabular problems with modest signal is $n/p \geq 10$, preferably $n/p \geq 20$.

The value of $n$ must be the \emph{independent} window count, not the generated one. \Cref{sec:overlap} established that sliding a window forward one observation at a time yields about 180 windows from the planned cohort, but that non-overlapping windows number only sixty to ninety. Overlapping windows inflate the row count without inflating the information, so they cannot relax a budget whose purpose is to control overfitting. The budget is therefore computed at $n = 60$ to $90$, and \cref{tab:np-budget} gives the admissible dimension.

\begin{table}[htbp]
\centering\small
\caption{Admissible feature count under the two sample-size conventions.}\label{tab:np-budget}
\begin{tabularx}{\textwidth}{@{}X c c c@{}}
\toprule
\textbf{Convention} & \textbf{$n$} & \textbf{$p$ at $n/p = 10$} & \textbf{$p$ at $n/p = 20$} \\
\midrule
Generated windows (overlapping) & 180 & 18 & 9 \\
Independent windows, upper & 90 & 9 & 4 \\
Independent windows, lower & 60 & 6 & 3 \\
\bottomrule
\end{tabularx}
\end{table}

\begin{figure}[htbp]
\centering
\begin{tikzpicture}
\begin{axis}[width=0.78\textwidth,height=5.6cm,
  xlabel={Number of features $p$},ylabel={$n/p$},
  xmin=2,xmax=40,ymin=0,ymax=30,grid=major,grid style={gray!20},
  legend pos=north east,legend style={font=\scriptsize}]
\addplot[thick,cGrey,domain=2:40,samples=100]{180/x};
\addplot[thick,cRef,domain=2:40,samples=100]{90/x};
\addplot[thick,cAfrica,domain=2:40,samples=100]{60/x};
\addplot[dashed,black,domain=2:40,samples=2]{10};
\addplot[dotted,black,domain=2:40,samples=2]{20};
\draw[thick,cGrey!70] (axis cs:8,0) -- (axis cs:8,30);
\draw[thick,cGrey!70] (axis cs:16,0) -- (axis cs:16,30);
\node[font=\scriptsize,anchor=south west] at (axis cs:8.2,25) {$p=8$};
\node[font=\scriptsize,anchor=south west] at (axis cs:16.2,25) {$p=16$};
\legend{$n=180$ generated,$n=90$ independent,$n=60$ independent}
\end{axis}
\end{tikzpicture}
\caption{The dimensionality budget. Dashed and dotted lines mark $n/p = 10$ and $20$. A sixteen-feature set clears the threshold only under the generated-window count; against independent windows it does not.}\label{fig:np-budget}
\end{figure}

The unrestrained set of 205 features exceeds even the most permissive bound by a factor of more than eleven. The comparison with the corpus is instructive: \textcite{mendoza2025} had $n = 1\,308$ and retained about eleven predictors after LASSO, giving $n/p \approx 119$---and their best model still reached only an area under the curve of 0.75 and a sensitivity of 0.51 (\cref{tab:mendoza-models}). A study operating near $n/p \approx 5$ cannot expect to do better.

\subsubsection{Two candidate sets}\label{sec:candidate-sets}

\Cref{tab:feature-sets} gives two sets, with $L = 14$ days so that four descriptors are supported.

\begin{table}[htbp]
\centering\footnotesize
\caption{Primary and extended feature sets.}\label{tab:feature-sets}
\begin{tabularx}{\textwidth}{@{}>{\raggedright\arraybackslash}p{2.1cm} X c X c@{}}
\toprule
\textbf{Block} & \textbf{Primary set} & \textbf{$p$} & \textbf{Extended set} & \textbf{$p$} \\
\midrule
SBP & $x_{\text{last}}$, $\beta$, signed margin & 3 & $x_{\text{last}}$, $\beta$, $s$, signed margin & 4 \\
Glucose & $x_{\text{last}}$, signed margin & 2 & $x_{\text{last}}$, $\beta$, $s$, signed margin & 4 \\
DBP & --- & 0 & $x_{\text{last}}$, $\beta$ & 2 \\
Cross-channel & --- & 0 & pulse pressure & 1 \\
Missingness & adherence ratio & 1 & adherence ratio, maximum gap & 2 \\
Static & years since diagnosis, regimen version & 2 & age band, years since diagnosis, regimen version & 3 \\
\midrule
\textbf{Total} & & \textbf{8} & & \textbf{16} \\
\textbf{$n/p$ at $n=90$} & & \textbf{11.3} & & \textbf{5.6} \\
\textbf{$n/p$ at $n=60$} & & \textbf{7.5} & & \textbf{3.8} \\
\bottomrule
\end{tabularx}
\end{table}

The primary set of eight features is the one reported as the main result. It clears $n/p = 10$ at the upper independent count and falls marginally short at the lower, which is stated rather than concealed. The extended set of sixteen is reported as a secondary analysis, with its $n/p$ ratio quoted alongside every figure, so that a reader can discount it appropriately.

Four structural economies make the primary budget work, and each is a decision rather than an omission.

\textbf{One window length, not three.} A three-length sweep triples the window-dependent dimension for information that is largely redundant, since a seven-day and a fourteen-day mean on the same series correlate strongly. Sensitivity to $L$ is reported by re-running the pipeline, not by concatenating lengths into one vector.

\textbf{Asymmetric treatment of channels.} Systolic pressure and glucose carry the clinical decision and receive the most columns; diastolic pressure enters only the extended set; oxygen saturation and temperature are excluded entirely, consistent with \cref{tab:pipelines} establishing them as context rather than decision variables.

\textbf{Dispersion dropped from the primary set.} At $L = 14$ and four observations, $s$ is estimated from four points and \cref{tab:small-sample} shows its variance is then very high. It is retained in the extended set so that its contribution can be measured, not assumed.

\textbf{Missingness aggregated.} Per-channel masks are retained in the data but enter the model as one or two summary features rather than sixteen.

\subsubsection{Reporting}

This thesis reports $n$, $p$ and $n/p$ for the final feature set alongside every performance figure, and states which sample-size convention $n$ refers to. A model result without a dimensionality statement is not interpretable at this sample size, and stating it pre-empts the first methodological question a jury will ask.

%==========================================================================
\subsection{Feature selection without leakage}\label{sec:selection}

\subsubsection{What the corpus reports}

\Textcite{liu2025review} survey the methods in use: correlation feature selection reported as superior in accuracy and execution time for chronic disease prediction, best-first search among wrapper methods, a hybrid of enhanced $k$-means clustering with filter and wrapper stages, recursive feature elimination with gradient boosting and with boosted decision trees, a genetic algorithm, and one coronary heart disease model retaining sixty-five variables.

\Textcite{mendoza2025} provide the best-documented procedure: LASSO regression with $\alpha = 1$ and a binomial family, the optimal penalty chosen by ten-fold cross-validation at $\lambda = 0.0056528$, and variables with zero coefficients excluded. Of twelve candidates, engagement in physical activity was eliminated.

\subsubsection{The leak these methods invite}\label{sec:nested}

Feature selection is a learned operation and must sit inside the cross-validation loop. Selecting on the full dataset and then cross-validating the selected set gives every fold's test data a say in which features exist, and the resulting estimate is optimistic. The bias grows as $p/n$ grows, which is exactly the regime of \cref{sec:budget}.

The correct construction is nested cross-validation: the outer loop estimates performance; the inner loop, on the outer training fold only, performs selection and tuning; the selected set is applied to the outer test fold without modification. The selected set may differ between outer folds, and reporting how often each feature is selected across folds is more informative than reporting one set---it measures stability, which at this $n/p$ is the quantity in doubt.

A note on \textcite{mendoza2025}. Their methods describe LASSO with ten-fold cross-validation, and separately an 80/20 train--test split for the tree models. The paper does not state whether selection preceded or followed the split. If it preceded, the selected set has seen the test partition. At $n = 1\,308$ with twelve candidates the resulting bias is small; the point is the reporting gap and not a demonstrated error, and this thesis states its own ordering explicitly rather than leaving the same ambiguity.

\subsubsection{The procedure adopted}

The feature set is fixed a priori on clinical and structural grounds, as the primary set of \cref{tab:feature-sets}. Performance on that fixed set is the primary result. Nested selection is run as a secondary analysis, to test whether a data-driven set improves on it, and the selection frequency of each feature across outer folds is reported.

Fixing the set in advance is the stronger design at this sample size. It removes a source of variance, it makes the result reproducible, and it means the primary claim does not depend on a selection procedure whose stability cannot be established on sixty to ninety independent samples.

%==========================================================================
\subsection{Scaling and model-specific preparation}\label{sec:scaling}

\begin{table}[htbp]
\centering\small
\caption{Preparation required by model family.}\label{tab:model-prep}
\begin{tabularx}{\textwidth}{@{}X c c X@{}}
\toprule
\textbf{Model} & \textbf{Scaling required} & \textbf{Monotone transform tolerated} & \textbf{Missing handled natively} \\
\midrule
Logistic regression & yes & no, it affects coefficients & no \\
Random forest & no & yes, splits are rank-based & partially \\
XGBoost & no & yes & yes, learned default direction \\
LSTM & yes & no & no, requires \cref{sec:mask-models} \\
\bottomrule
\end{tabularx}
\end{table}

Scaling is fitted on the training fold only (\cref{sec:pipeline}). For the recurrent model, standardisation to zero mean and unit variance per feature is conventional; for the penalised linear baseline it is required for the regularisation path to be meaningful, since the LASSO penalises coefficients on the scale of their features.

One asymmetry is recorded as a finding rather than a detail. XGBoost requires no scaling, tolerates monotone transforms and handles missing values natively through a learned default split direction. The recurrent model requires all three to be solved by hand. Given that \cref{sec:missingness} established missingness as the dominant irregularity of this dataset, and \cref{sec:capacity} established that the sample size disfavours the recurrent model, this is a third independent reason to expect the tree ensemble to prevail in the architecture comparison. The three reasons---sample size, native missingness handling, and the absence of preparation requirements---are stated together here rather than discovered in the results.

%==========================================================================
\subsection{The baselines this feature set must beat}\label{sec:baselines}

\Cref{sec:spec23} required the clinical threshold rule as a mandatory baseline. Feature engineering introduces a second, and it is the harder one.

\textbf{Persistence.} Predict that the classification at $t+\Delta$ equals the classification at $t$. No features, no model, no training.

Physiological series are strongly autocorrelated over the horizons considered here, so persistence is a genuinely strong baseline---and because the label is a threshold function of the same channel that supplies $x_{\text{last}}$, the elaborate feature set is competing against simply carrying the current value forward. Any trend, variability or threshold feature earns its place only by the margin it adds over that.

\begin{table}[htbp]
\centering\small
\caption{The baseline ladder, in order.}\label{tab:baselines}
\begin{tabularx}{\textwidth}{@{}c X@{}}
\toprule
\textbf{Rung} & \textbf{Description} \\
\midrule
0 & majority class \\
1 & persistence: the classification at $t$ carried to $t+\Delta$ \\
2 & clinical rule: guideline thresholds applied to $x_{\text{last}}$ \\
3 & penalised logistic regression on level features only \\
4 & penalised logistic regression on the primary set \\
5 & random forest and XGBoost on the primary set \\
6 & LSTM on the raw masked sequence \\
\bottomrule
\end{tabularx}
\end{table}

The contribution of this chapter is the gap between rungs~2 and~5. If that gap is small, the honest conclusion is that at this sampling density and sample size engineered features add little over the guideline rule---which is a result, is publishable, and is considerably more useful to a clinician than an inflated accuracy figure. \Cref{sec:ml-conclusion} established that the literature is full of the latter.

%==========================================================================
\subsection{Specification}\label{sec:spec33}

\begin{table}[htbp]
\centering\footnotesize
\caption{Feature-engineering requirements established in \cref{sec:feature-eng}.}\label{tab:spec33}
\begin{tabularx}{\textwidth}{@{}X >{\raggedright\arraybackslash}p{5.2cm}@{}}
\toprule
\textbf{Requirement} & \textbf{Basis} \\
\midrule
$L \geq 14$ days for feature construction; 7 days retained for alerting & \cref{sec:ceiling} \\
Descriptors per channel not exceeding observations per window & \cref{sec:ceiling} \\
$s$ and MASD never both included when $n = 2$ & $s = \mathrm{MASD}/\sqrt{2}$ exactly \\
Median and MAD at patient level, not window level & \cref{tab:small-sample} \\
OLS slope guarded by $\mathrm{span}_W \geq L/2$ & \cref{sec:descriptors} \\
Threshold features on window observations only; embargo enforced in code & \cref{sec:leak-boundary} \\
Threshold table as versioned configuration, not constants & \cref{sec:threshold-table} \\
Grade band retained as ordinal, not collapsed to binary & effect confined to stage III in \textcite{kingue2013} \\
Pulse pressure supplied explicitly for tree models & \cref{sec:cross-channel} \\
Cross-channel features inherit the worst input mask & \cref{sec:cross-channel} \\
Budget computed on independent windows, not generated ones & \cref{sec:overlap,sec:np} \\
Primary set $p = 8$; extended set $p = 16$ reported separately & \cref{tab:feature-sets} \\
One window length per run; sensitivity by re-running & \cref{sec:candidate-sets} \\
$n$, $p$, $n/p$ and the convention for $n$ reported with every result & \cref{sec:budget} \\
Feature set fixed a priori as the primary analysis & \cref{sec:selection} \\
Any selection nested inside cross-validation & \cref{sec:nested} \\
Selection frequency per feature reported across folds & \cref{sec:nested} \\
Scaling fitted on the training fold only & \cref{sec:pipeline} \\
Persistence baseline reported & \cref{sec:baselines} \\
Contribution stated as the margin over the clinical rule & \cref{tab:baselines} \\
\bottomrule
\end{tabularx}
\end{table}

%==========================================================================
\subsection{Limits of this specification}\label{sec:feature-limits}

Three limitations are recorded here rather than left to be discovered.

\textbf{The guideline sources are not yet in hand.} The framing document names the ESH and ADA guidelines as the labelling sources, and neither is in the corpus assembled for this thesis. The threshold table of \cref{sec:threshold-table} cannot be populated from them, and the values of \cref{tab:verified-thresholds} are operational thresholds from \textcite{mihevc2025}, adapted for an older population, standing as placeholders.%
\verify{Obtain the ESH guideline edition and the ADA Standards of Care edition named in \textcite{cadrage2026}, populate the threshold table from them with version and date, and recompute every threshold-derived feature. Until then no grade-band feature can be constructed, since the bands are not known.}

\textbf{No descriptor in this set is validated against a clinical outcome by any source in this corpus.} The mean, slope and variability features are standard constructions, but no document assembled for this thesis demonstrates that any of them predicts anything in hypertension or diabetes. The only feature in the entire set with published predictive evidence is years since diagnosis, from \textcite{mendoza2025}. Blood pressure variability as a risk marker has a literature in cardiology; it is not represented here, and until it is, the variability features are presented as candidates proposed on physiological reasoning and not as established markers.%
\verify{Obtain at least one primary source on blood pressure variability as an independent cardiovascular risk marker, or state in the discussion that the variability features rest on physiological reasoning alone.}

\textbf{The exceedance burden is an original construction.} It is defined in \cref{tab:threshold-features} by analogy with the magnitude-integrated-over-time logic of the disability-adjusted life year, and it is evaluated here for the first time. It is reported as a candidate feature whose contribution is measured, not as an index with prior standing.

%##########################################################################
% SECTION 3.4
%##########################################################################
\section{Learning Problem Formulation, Model Training and Risk Score Derivation}\label{sec:learning}

\subsection*{Purpose, sources and working assumptions}

\Cref{sec:collection,sec:preprocessing,sec:feature-eng} produced a feature matrix. This section specifies what is predicted from it, how the models are fitted and compared, and how a model output becomes the score that drives the three-tier alert of \cref{sec:decision-flow}. It closes the chain from measurement to clinical action.

Three sources are used. \Textcite{cadrage2026} states the contributions and hypotheses and commits to class weighting, oversampling, a recall-centred metric and group $k$-fold by patient. \Textcite{mendoza2025} supply the only fully documented gradient-boosting implementation in the corpus with reported discrimination and confidence intervals. \Textcite{efendi2025} supply a second implementation whose configuration and results were shown in \cref{sec:task-not-algo} to be unusable, and which is retained here as a worked example of what to avoid.

This section is written under the assumptions set out in \cref{sec:assumptions}, which also states what would have to be revised if any of them is overturned. Where an assumption bears directly on a quantity, it is marked at that point.

%==========================================================================
\subsection{Formalising the learning problem}\label{sec:formulation}

\subsubsection{Three candidate formulations}

For patient $i$ with window ending at $t$ and feature vector $\mathbf{v}_{i,t}$ constructed per \cref{sec:feature-eng}:

\textbf{Formulation 1, forecast the parameter.} Regression on the continuous value,
\begin{equation}
  \hat{x}^{(c)}_{i,t+\Delta} = f\bigl(\mathbf{v}_{i,t}\bigr),
  \qquad \mathcal{L} = \bigl(\hat{x} - x\bigr)^2,
\end{equation}
with the risk class obtained by applying the clinical rule to the forecast. The target is continuous, so every sample is informative and the class-imbalance problem of \cref{sec:imbalance} does not arise at training time. Against this, squared error weights a 5~mmHg miss at 120 the same as at 175, where it crosses a threshold.

\textbf{Formulation 2, classify the future state.} Direct classification onto the four risk levels. The target is the clinical decision itself, but the imbalance is severe and all information about how far a patient lies from a boundary is discarded.

\textbf{Formulation 3, forecast the margin.} Regression on the signed distance to a reference threshold,
\begin{equation}
  \hat{m}^{(c)}_{i,t+\Delta} = \hat{x}^{(c)}_{i,t+\Delta} - \theta_c^{\text{ref}},
\end{equation}
after which the tier is obtained by applying the full band table of \cref{sec:threshold-table} to the reconstructed value $\hat{x} = \hat{m} + \theta_c^{\text{ref}}$. This is Formulation~1 recentred: it inherits the sample efficiency of the continuous target while placing the principal clinical boundary at zero. The reference threshold is a presentational choice---it makes the score interpretable---and the multi-tier decision still comes from the band table, not from the single margin.

\subsubsection{Disambiguating the dual-task formulation}\label{sec:dual-task}

The phrase used in the framing document admits two readings, and this thesis states which it means, since the ambiguity is visible to a reader and costs nothing to remove.

\textbf{Reading A, two heads on one model.} A shared representation with a regression head forecasting the parameter and a classification head predicting the risk class, trained jointly under
\begin{equation}
  \mathcal{L} = \alpha\,\mathcal{L}_{\text{reg}} + (1-\alpha)\,\mathcal{L}_{\text{cls}}.
\end{equation}
The regression head supplies a dense training signal that regularises the classification head, which is valuable precisely when positives are rare: the regression loss is informative on every sample, while the classification loss is informative mainly on the few positives.

\textbf{Reading B, two pathologies.} Separate targets for blood pressure and for blood glucose, predicted jointly from a shared feature vector. This is justified by \cref{sec:comorbidity}, which established the comorbidity at 58.1\% of diabetic patients, and by the WHO target coupling the two (\cref{sec:comorbidity}).

The two are not exclusive, and this thesis adopts both: Reading~B as the outer structure, one target per pathology, with Reading~A inside each where the model family permits a shared representation.

\subsubsection{The horizon}\label{sec:horizon-adopted}

\Cref{sec:horizon,sec:temporal} established that a uniform 24 to 48 hour horizon contains zero or one observations under the routine protocol, and adopted a regime-dependent resolution. In the routine regime, which is the regime of the training data,
\begin{equation}
  \Delta \;\longrightarrow\; \text{the next scheduled measurement occasion},
\end{equation}
so that the task is: given the fourteen-day window ending at $t$, predict the classification of the next measurement. At the routine frequency this is a horizon of three to four days. It is anticipatory, it is clinically meaningful---the confirmation interval of \cref{tab:protocol} is one to four weeks and the verification profile of \textcite{mihevc2025} is seven days---and it is estimable. The 24 to 48 hour horizon is retained only for patients inside the verification regime, where two readings per day make it estimable; that case is reported separately and on far fewer samples.

This is presented as a deliberate adaptation to the sampling density available, with the arithmetic of \cref{sec:temporal} as its justification. That is a stronger position than retaining an unrealisable horizon.

\subsubsection{The formulation adopted}

\begin{equation}
  \boxed{\;\hat{m}^{(c)}_{i,\,t+\Delta} = f\bigl(\mathbf{v}_{i,t}\bigr),
  \qquad c \in \{\mathrm{SBP},\ \text{glucose}\},
  \qquad \Delta = \text{next occasion}\;}
\end{equation}

with the tier derived by applying the versioned band table of \cref{sec:threshold-table} to the reconstructed forecast, and a classification head trained jointly where the model family permits it.

Three reasons. It maximises the information extracted from a small sample, since every window contributes a continuous residual rather than a label that is almost always negative. It places the principal clinical boundary at zero, so the score is directly interpretable in clinical units. And it decouples the model from the thresholds, so a guideline revision changes the decision layer without retraining---which is the separation that \cref{sec:dsts-def} required of the knowledge base.

%==========================================================================
\subsection{The label}\label{sec:label-design}

\subsubsection{Ordinal, not nominal}

The framing document specifies four levels, and they are ordered:
\begin{equation}
  \text{stable} \prec \text{moderate} \prec \text{high} \prec \text{critical}.
\end{equation}

The consequence is in the loss. Under multi-class cross-entropy, predicting \emph{critical} when the truth is \emph{stable} costs exactly what predicting \emph{moderate} costs. Clinically it does not. Either an ordinal formulation is used---a cumulative-link model, or a set of nested binary classifiers $\mathbb{1}[y \succeq k]$---or Formulation~3 is adopted, in which the ordering is carried by the continuous forecast and no choice of ordinal loss is required. This is a second argument for the margin formulation.

\subsubsection{Two label sources, reported separately}\label{sec:two-labels}

\Cref{sec:label-rules} established two label definitions, and both are used.

\begin{table}[htbp]
\centering\small
\caption{The two label sources and their status.}\label{tab:label-sources}
\begin{tabularx}{\textwidth}{@{}>{\raggedright\arraybackslash}p{3.6cm} X >{\raggedright\arraybackslash}p{4.2cm}@{}}
\toprule
\textbf{Label} & \textbf{Source} & \textbf{Status} \\
\midrule
Rule-derived & band table applied to $x_{t+\Delta}$ & primary; available for every sample \\
Clinician-adjudicated & \texttt{clinician\_verdict} of \cref{tab:event-record} & secondary; available only for triggered events \\
\bottomrule
\end{tabularx}
\end{table}

Performance against the second is the more credible result, because it is immune to the circularity of \cref{sec:task-not-algo}: the label was not computed from the same thresholds the features see. It will be available on far fewer samples, and it is reported as a secondary endpoint with its own sample size rather than merged with the primary. Under a retrospective collection scenario it is unavailable entirely (\cref{sec:adjudication}).

\subsubsection{Prevalence}

Class prevalence determines everything downstream and is measured, not assumed. It is reported per class, overall and per patient, before any balancing. The 30.3\% of \textcite{mendoza2025} is the only comparable figure in the corpus and is for a cross-sectional outcome; the prevalence of a transition to the critical tier over a three-day horizon will be far lower, and the regime to plan for is the $\pi = 0.02$ of \cref{tab:ppv}.

%==========================================================================
\subsection{Training protocol}\label{sec:training}

\subsubsection{The cross-validation structure}\label{sec:cv-structure}

The constraints established earlier compose into one structure: split by patient, forward in time, stratified on patient-level event presence, with preprocessing, feature selection and hyperparameter tuning all fitted inside the training fold (\cref{sec:forecasting,sec:capacity,sec:pipeline,sec:nested}).

The resulting design is nested cross-validation with $K_{\text{out}} = 5$ outer folds estimating performance and $K_{\text{in}} = 4$ inner folds, on the outer training set only, performing selection and tuning. The outer test fold is touched once, after the inner loop has chosen everything (\cref{fig:nested-cv}).

\begin{figure}[htbp]
\centering
\begin{tikzpicture}[font=\scriptsize]
  % outer
  \node[anchor=east] at (-0.2,2.2) {outer};
  \foreach \i in {0,...,4} {
    \pgfmathsetmacro{\xa}{\i*2.1}
    \ifnum\i=2
      \fill[cAfrica!25] (\xa,1.9) rectangle (\xa+1.9,2.5);
      \node at (\xa+0.95,2.2) {test};
    \else
      \fill[cRef!18] (\xa,1.9) rectangle (\xa+1.9,2.5);
      \node at (\xa+0.95,2.2) {train};
    \fi
    \draw (\xa,1.9) rectangle (\xa+1.9,2.5);
  }
  \node[anchor=west,font=\scriptsize] at (10.7,2.2) {6 patients each};
  % arrow
  \draw[-{Stealth[length=2mm]},thick] (5.0,1.75) -- (5.0,1.15);
  \node[anchor=west,font=\scriptsize] at (5.15,1.45) {the 24 training patients are split again};
  % inner
  \node[anchor=east] at (-0.2,0.6) {inner};
  \foreach \i in {0,...,3} {
    \pgfmathsetmacro{\xb}{\i*2.1}
    \ifnum\i=1
      \fill[cGrey!35] (\xb,0.3) rectangle (\xb+1.9,0.9);
      \node at (\xb+0.95,0.6) {validation};
    \else
      \fill[cRef!12] (\xb,0.3) rectangle (\xb+1.9,0.9);
      \node at (\xb+0.95,0.6) {train};
    \fi
    \draw (\xb,0.3) rectangle (\xb+1.9,0.9);
  }
  \node[anchor=west,font=\scriptsize,align=left] at (8.6,0.6)
    {selection, tuning,\\ calibration, thresholds};
  \node[anchor=west,font=\scriptsize,text=cAfrica] at (0,-0.35)
    {the outer test fold is scored once, after everything has been chosen; the whole procedure repeats $R=10$ times};
\end{tikzpicture}
\caption{Nested cross-validation. Every learned quantity---preprocessing statistics, selected features, hyperparameters, the calibrator and the decision thresholds---is estimated inside the inner loop.}\label{fig:nested-cv}
\end{figure}

With thirty patients each outer fold holds six. Stratification on event presence is mandatory, because an unstratified fold may contain no positives at all and the fold-level recall is then undefined (\cref{sec:capacity}).

\subsubsection{Repeated cross-validation, because variance is the problem}\label{sec:repeats}

At the independent sample size established in \cref{sec:np}, with six patients per fold, the fold-to-fold variance of any metric is large, and a single five-fold estimate is one draw from a wide distribution. Reporting that draw as a point estimate overstates what is known.

The entire nested procedure is therefore repeated $R = 10$ times with different fold assignments, and the distribution is reported: median and interquartile range across repetitions, not a single number. This is nearly free at this scale and it converts an unfounded point estimate into an honest interval. It also supplies the confidence intervals that \cref{sec:spec23} required and that \textcite{efendi2025} omitted entirely.

\subsubsection{Gradient-boosting hyperparameters}\label{sec:xgb-config}

\Textcite{mendoza2025} report their configuration precisely, and it is the right template---but their values are tuned for $n = 1\,308$ and must not be copied to a sample two orders of magnitude smaller.

\begin{table}[htbp]
\centering\footnotesize
\caption{Gradient-boosting configuration, against the reference implementation.}\label{tab:xgb-config}
\begin{tabularx}{\textwidth}{@{}>{\raggedright\arraybackslash}p{3.0cm} >{\raggedright\arraybackslash}p{2.2cm} >{\raggedright\arraybackslash}p{2.6cm} X@{}}
\toprule
\textbf{Parameter} & \textbf{\textcite{mendoza2025}} & \textbf{Adopted here} & \textbf{Reason} \\
\midrule
\texttt{max\_depth} & 6 & 2--3 & depth 6 permits 64 leaves; at the sample size available that is one to three samples per leaf \\
\texttt{eta} & 0.3 & 0.03--0.05 & small steps with more rounds generalise better at small $n$ \\
rounds & 100 & early stopping on inner validation & fixed round counts overfit or underfit unpredictably \\
\texttt{min\_child\_weight} & not stated & 5--10 & forces leaves to be supported by several samples \\
\texttt{subsample} & not stated & 0.7--0.8 & variance reduction \\
\texttt{colsample\_bytree} & not stated & 0.7--0.8 & decorrelates trees; matters at the $p/n$ of \cref{sec:np} \\
$\lambda$ (L2) & not stated & 1--10 & explicit regularisation, the term $\Omega(f)$ of \cref{sec:families} \\
objective & logistic & logistic or squared error & per the formulation adopted \\
seed & 123, fixed & fixed and stated & reproducibility \\
\bottomrule
\end{tabularx}
\end{table}

The depth reduction is the important one. A depth-six tree on sixty to ninety independent samples is a memorisation device, and the discipline imposed in \cref{sec:budget} on feature count is undone if the model is then permitted that much capacity.

The practice of fixing and reporting the seed is copied exactly from \textcite{mendoza2025}. \Textcite{efendi2025} report none.

\subsubsection{Recurrent-model configuration}\label{sec:lstm-config}

\Cref{sec:capacity} computed that a single layer of sixteen units on ten input channels has about 1\,800 parameters. Against the independent window count of sixty to ninety, that is twenty to thirty parameters per sample, not the ten per sample obtained from the generated count.

\begin{table}[htbp]
\centering\small
\caption{Recurrent-model configuration.}\label{tab:lstm-config}
\begin{tabularx}{\textwidth}{@{}>{\raggedright\arraybackslash}p{3.4cm} >{\raggedright\arraybackslash}p{3.4cm} X@{}}
\toprule
\textbf{Parameter} & \textbf{Adopted} & \textbf{Reason} \\
\midrule
Layers & 1 & two layers triple the parameter count \\
Units & 8--16 & the smallest that can represent the dynamics \\
Dropout & 0.3--0.5 & heavy, given the ratio \\
Recurrent dropout & 0.2--0.3 & --- \\
Sequence length & 4--8 observations & all a fourteen-day window contains \\
Masking & mandatory & no native missing handling (\cref{sec:mask-models}) \\
Early stopping & on inner-validation recall & not on loss \\
Optimiser & Adam, learning rate $10^{-3}$ & conventional \\
Seed & fixed and reported & --- \\
\bottomrule
\end{tabularx}
\end{table}

This thesis states in advance that this configuration is expected to underperform. Four independent reasons have now accumulated: the parameter-to-sample ratio just computed; the absence of native missing-value handling where missingness is the dominant irregularity (\cref{sec:mask-models}); the need for manual scaling and decay construction (\cref{tab:model-prep}); and the window length of fourteen days containing at most four observations, which is a very short sequence for an architecture whose advantage is long-range dependence.

Pre-stating this is not pessimism. A negative finding predicted in advance from sample-size arithmetic is evidence about the data regime; the same finding reported as a surprise is evidence about nothing.

\subsubsection{Why the second implementation transfers nothing}

The configuration of \textcite{efendi2025}---three hundred epochs, a learning rate of 0.025, a sigmoid activation and the Adam optimiser---was shown in \cref{sec:task-not-algo} to be inapplicable to three of the five models they fitted, which do not train by gradient descent. Their reported performance was shown in the same place to measure label recovery on a task whose features were discretised by the ranges that define the label.

That paper is cited in this thesis for its system architecture and its device validation, never for its model configuration or its results. Its one transferable element is the deployment pattern of a trained model applied to incoming readings---which \cref{sec:edge-triage} relocated from the application layer to the analysis layer.

\subsubsection{Compute is not the constraint}

\Textcite{efendi2025} report their hardware as though it bounded the work. At this sample size and dimension it does not. One gradient-boosting fit is a matter of milliseconds, and the full nested design---five outer folds, four inner, ten repetitions, a grid of fifty configurations, of the order of $10^4$ fits---completes in minutes on an ordinary laptop.

Compute is therefore not presented as a constraint in this work. The constraint is sample size, and reporting hardware specifications in its place invites the question of why the sample is not larger.

%==========================================================================
\subsection{Class imbalance at training time}\label{sec:imbalance-training}

\subsubsection{Three mechanisms}

\begin{table}[htbp]
\centering\small
\caption{Mechanisms for handling imbalance and their effect on calibration.}\label{tab:imbalance-mechanisms}
\begin{tabularx}{\textwidth}{@{}>{\raggedright\arraybackslash}p{3.6cm} X c@{}}
\toprule
\textbf{Mechanism} & \textbf{Implementation} & \textbf{Effect on calibration} \\
\midrule
Class weighting & a positive-class weight $w$ in the loss & distorts \\
Oversampling the minority & duplicate or synthesise positives & distorts \\
Threshold adjustment & leave training untouched; move the decision point & none \\
\bottomrule
\end{tabularx}
\end{table}

The framing document commits to the first two. The third is absent from it and is in several respects the best: it leaves the probability model intact and moves only the operating point, which \cref{sec:volume-response} already established is a clinical configuration parameter rather than a modelling choice.

\subsubsection{Inside the fold}

Any resampling is fitted on the training fold only. Oversampling before the split places duplicates of the same patient's window in both training and test, and the resulting estimate is meaningless. This is among the most common errors in small-sample clinical machine learning and it is invisible in the output.

Under group $k$-fold the error is partly prevented by the patient-level split, but only partly: synthetic samples generated from a training patient remain in training, while duplicated windows from an overlapping stride (\cref{sec:overlap}) can still straddle the boundary if the stride and the split are applied in the wrong order.

\subsubsection{The calibration cost}

Weighting and oversampling both change the prior the model sees, so the output is a probability under a prevalence that does not exist. A score of 0.6 from a model trained with a positive-class weight of twenty does not mean a sixty per cent chance of the event.

If the score is only ever compared against a threshold this does not matter, since any monotone distortion preserves the ranking. If the score is displayed to a clinician as a risk it matters entirely, because a number presented as a probability that is not one is a safety problem and not a technical imprecision. \Cref{sec:cdss} established that the clinician decides; a clinician deciding on a miscalibrated probability is deciding on a misleading number. \Cref{sec:risk-score} resolves this.

%==========================================================================
\subsection{Deriving the risk score}\label{sec:risk-score}

\subsubsection{From margin to probability}

A gradient-boosted model with a logistic objective outputs a margin $\eta = \sum_k f_k(\mathbf{v})$, mapped to a score by the logistic function $s(\mathbf{v}) = \sigma(\eta)$. This is a score in $[0,1]$. It is not a calibrated probability, for three reasons: the boosting objective optimises ranking and loss reduction rather than calibration; regularisation shrinks margins toward zero, compressing scores toward one half; and any class weighting shifts the implied prior.

Under the margin formulation adopted in \cref{sec:formulation}, the model outputs a forecast in clinical units, which is interpretable but is not a risk. Conversion requires a second step mapping the forecast and its uncertainty to the probability of crossing:
\begin{equation}
  P\bigl(x_{t+\Delta} > \theta\bigr) = \Phi\!\left(\frac{\hat{m}_{t+\Delta}}{\hat{\sigma}_{\text{resid}}}\right),
\end{equation}
where $\hat{\sigma}_{\text{resid}}$ is the residual standard deviation estimated on the inner validation fold and the residuals are assumed approximately normal, which is checked. This is a strength of the margin formulation: the uncertainty is estimated explicitly rather than implied by an uncalibrated classifier output.

\subsubsection{Prior correction after weighting}\label{sec:prior-correction}

Where class weighting is used the distortion has a closed form and is undone. With a positive-class weight $w$, the fitted model estimates odds inflated by $w$, so that $\mathrm{odds}'(\mathbf{v}) = w\,\mathrm{odds}(\mathbf{v})$. Inverting, with $s'$ the raw output and $s$ the corrected probability:
\begin{equation}
  \boxed{\;s(\mathbf{v}) = \frac{s'(\mathbf{v})}{s'(\mathbf{v}) + w\bigl(1 - s'(\mathbf{v})\bigr)}\;}
\end{equation}

Two worked checks. With $w = 20$ and a raw output of $s' = 0.5$, the corrected probability is $0.5/(0.5 + 20 \times 0.5) = 1/21 \approx 0.048$: a raw score suggesting an even chance corresponds to a true probability of about five per cent, so a clinician shown $0.5$ would be shown a number about ten and a half times too large. The second check is sharper. The conventional setting is $w = (1-\pi)/\pi$, which at the planning prevalence $\pi = 0.02$ gives $w = 49$; then $s' = 0.5$ maps to $0.5/(0.5 + 49 \times 0.5) = 1/50 = 0.02$, exactly the base rate---which is what a raw score of one half ought to mean, and is nowhere near what it appears to mean (\cref{fig:prior-correction}).

\begin{figure}[htbp]
\centering
\begin{tikzpicture}
\begin{axis}[width=0.66\textwidth,height=5.6cm,
  xlabel={Raw model output $s'$},ylabel={Corrected probability $s$},
  xmin=0,xmax=1,ymin=0,ymax=1,grid=major,grid style={gray!20},
  legend pos=north west,legend style={font=\scriptsize}]
\addplot[dashed,cGrey,domain=0:1,samples=2]{x};
\addplot[thick,cRef,domain=0:1,samples=200]{x/(x + 5*(1-x))};
\addplot[thick,cAfrica,domain=0:1,samples=200]{x/(x + 20*(1-x))};
\addplot[thick,black,domain=0:1,samples=200]{x/(x + 49*(1-x))};
\legend{no weighting,$w=5$,$w=20$,$w=49$}
\end{axis}
\end{tikzpicture}
\caption{The prior correction. At $w=49$, the setting conventional at a two per cent prevalence, a raw output of $0.5$ corresponds to a true probability of $0.02$.}\label{fig:prior-correction}
\end{figure}

The same correction applies to oversampling by a factor $r$, with $w$ replaced by $r$.

\subsubsection{Calibration}\label{sec:calibration}

Prior correction fixes the location; it does not fix the shape. A second step is required, fitted on held-out data.

\begin{table}[htbp]
\centering\small
\caption{Calibration methods and their suitability.}\label{tab:calibration-methods}
\begin{tabularx}{\textwidth}{@{}>{\raggedright\arraybackslash}p{3.4cm} >{\raggedright\arraybackslash}p{4.2cm} X@{}}
\toprule
\textbf{Method} & \textbf{Form} & \textbf{Suited to} \\
\midrule
Platt scaling & $\hat{p} = \sigma(a\eta + b)$, two parameters & small samples; assumes a sigmoidal distortion \\
Isotonic regression & monotone step function & larger samples; flexible but overfits at small $n$ \\
Beta calibration & three parameters on $[0,1]$ & intermediate flexibility \\
\bottomrule
\end{tabularx}
\end{table}

At the sample size available, Platt scaling is the correct choice. Isotonic regression fits a free monotone function and will overfit a calibration set of a few dozen samples.

The calibrator is fitted on the inner validation fold and applied to the outer test fold. Fitting it on the test fold is a leak; fitting it on the training fold yields a calibrator fitted to data the model has already memorised, which is useless.

\subsubsection{Assessing calibration}\label{sec:brier}

Discrimination and calibration are distinct, and a model may have a good area under the curve with poor calibration, since that measure depends only on ranking.

The Brier score is the mean squared error of the probability, $\mathrm{BS} = \frac{1}{n}\sum_i (\hat{p}_i - y_i)^2$, and it decomposes as
\begin{equation}
  \mathrm{BS} = \underbrace{\mathrm{REL}}_{\text{miscalibration, lower better}}
  - \underbrace{\mathrm{RES}}_{\text{discrimination, higher better}}
  + \underbrace{\mathrm{UNC}}_{\pi(1-\pi),\ \text{fixed}}.
\end{equation}
Reporting the decomposition rather than the raw score separates the two failure modes, which matters here because they have different remedies: miscalibration is fixed by Platt scaling, poor resolution is not fixable at all.

The calibration curve bins the predictions and plots observed frequency against mean predicted probability, with the diagonal as reference. At this sample size quantile bins of at least twenty samples are used, giving at most four or five bins against the independent count, and a confidence interval is plotted on each---a calibration curve without intervals at this sample size is uninterpretable.

A caution on the uncertainty term. Since $\mathrm{UNC} = \pi(1-\pi)$ depends only on prevalence, at $\pi = 0.02$ it equals $0.0196$, so a model predicting the base rate for everyone achieves a Brier score of $0.0196$. A low Brier score at low prevalence is not evidence of anything, and it is always reported alongside the base-rate value.

%==========================================================================
\subsection{Threshold selection and the three tiers}\label{sec:tiers}

\subsubsection{Two cut-points}

The three-tier protocol of \cref{tab:triage-levels} requires two cut-points on the calibrated score,
\begin{equation}
  \text{tier}(\mathbf{v}) =
  \begin{cases}
    \text{green} & s(\mathbf{v}) < \tau_{\text{amber}},\\
    \text{amber} & \tau_{\text{amber}} \leq s(\mathbf{v}) < \tau_{\text{red}},\\
    \text{red} & s(\mathbf{v}) \geq \tau_{\text{red}},
  \end{cases}
\end{equation}
with distinct selection criteria, because the two tiers have different costs. Amber costs the patient a verification profile and costs the clinician nothing, so it is set for high sensitivity. Red costs clinician time, so it is set for acceptable precision.

\subsubsection{Selection by target sensitivity}\label{sec:threshold-selection}

The default of one half is a convention with no clinical content. \Cref{sec:imbalance} established that the threshold is set from the cost of error. Operationally: a target sensitivity on the critical class is fixed in advance, for instance $\mathrm{Se}^\star = 0.90$ for amber; on the inner validation fold the smallest $\tau$ achieving it is found; that $\tau$ is applied unchanged to the outer test fold; and the sensitivity actually obtained, the resulting precision and the alert volume are reported.

The third step is the one usually violated. Choosing the threshold on the test fold inflates every metric and is a second form of the leak that the nested design of \cref{sec:cv-structure} exists to prevent.

\subsubsection{Alert volume at the chosen operating point}

\Cref{sec:volume-arithmetic} established this as a mandatory reporting requirement. At the selected $\tau_{\text{red}}$ the weekly alert count is
\begin{equation}
  A_{\text{week}} = N m \bigl[\pi\,\mathrm{Se} + (1-\pi)(1-\mathrm{Sp})\bigr]
\end{equation}
for a stated cohort size $N$ and measurement rate $m$. For a pilot cohort of thirty patients at three readings per week, with $\pi = 0.02$, $\mathrm{Se} = 0.90$ and $\mathrm{Sp} = 0.98$:
\begin{equation}
  A_{\text{week}} = 30 \times 3 \times \bigl[0.02(0.90) + 0.98(0.02)\bigr]
  = 90 \times 0.0376 \approx 3.4 \text{ alerts per week},
\end{equation}
of which approximately 1.6 are genuine. A sensitivity figure without this calculation is not an evaluation of a triage system, and it appears next to every reported operating point.

\subsubsection{The rule engine as a floor}\label{sec:rule-floor}

\Cref{sec:dsts-def} established the hybrid design. At the decision layer this means the model may raise a tier but may not lower one set by the clinical rule:
\begin{equation}
  \text{tier}_{\text{final}} = \max\Bigl(\text{tier}_{\text{rule}}(x_{\text{last}}),\ \text{tier}_{\text{model}}\bigl(s(\mathbf{v})\bigr)\Bigr)
\end{equation}
under the ordering green $\prec$ amber $\prec$ red.

Three justifications, all established earlier. A patient at 185/110~mmHg is red under the same-day-treatment criterion of \cref{tab:protocol} regardless of any model. The best model of \textcite{mendoza2025} misses 49\% of positives, so no model in this corpus is reliable enough to override a guideline. And the composition makes the system's worst case equal to guideline conformity: if the model fails entirely, the system degrades to the standard of care rather than to nothing. That property is what makes the system defensible to a clinician, and it is stated explicitly rather than left implicit.

%==========================================================================
\subsection{Reporting and pre-registration}\label{sec:reporting-prereg}

\subsubsection{What is reported}

\begin{table}[htbp]
\centering\footnotesize
\caption{Reporting requirements for the modelling stage.}\label{tab:model-reporting}
\begin{tabularx}{\textwidth}{@{}X >{\raggedright\arraybackslash}p{5.2cm}@{}}
\toprule
\textbf{Item} & \textbf{Why} \\
\midrule
$n$, $p$, $n/p$, the convention for $n$, class prevalence per class & \cref{sec:budget} \\
The full baseline ladder, rungs 0 to 6 & \cref{tab:baselines} \\
The persistence baseline explicitly & \cref{sec:baselines}; the hardest to beat \\
Per-class sensitivity, precision and $F_1$, never aggregates alone & \cref{sec:sens-spec} \\
Positive class stated for every metric & \cref{sec:positive-class} \\
AUC with confidence interval, chance level marked & \cref{sec:metrics} \\
Distribution across the $R$ repetitions, not a point estimate & \cref{sec:repeats} \\
Brier score with its base-rate value and decomposition & \cref{sec:brier} \\
Calibration curve with intervals & \cref{sec:brier} \\
$\tau_{\text{amber}}$, $\tau_{\text{red}}$ and how they were chosen & \cref{sec:threshold-selection} \\
Alert volume at the operating point & \cref{sec:volume-arithmetic} \\
Seeds, library versions, the full hyperparameter grid & \textcite{mendoza2025} \\
Secondary endpoint against clinician-adjudicated labels & \cref{sec:two-labels} \\
\bottomrule
\end{tabularx}
\end{table}

\subsubsection{Pre-registration}\label{sec:prereg}

The failures diagnosed throughout \cref{ch:techsota}---selection on the full dataset, thresholds chosen on test data, metrics reported for whichever class flatters, post-hoc sensitivity analyses presented as primary---share one remedy: declaring the analysis before seeing the test data.

Four items are fixed in advance and stated here, in the methodology, rather than in the results:

\begin{enumerate}[nosep]
  \item \textbf{The primary endpoint:} recall on the critical class at the pre-specified $\tau_{\text{amber}}$.
  \item \textbf{The primary comparison:} gradient boosting on the primary feature set against the clinical rule baseline. One comparison, not a matrix.
  \item \textbf{The threshold rule:} target sensitivity 0.90, selected on inner validation.
  \item \textbf{The feature set:} the eight-feature primary set fixed a priori per \cref{tab:feature-sets}.
\end{enumerate}

Everything else is secondary and labelled as such. This costs nothing, it is verifiable by a reader comparing this chapter with the results, and it is the strongest available answer to the question of whether the result was found or constructed.

%==========================================================================
\subsection{Specification}\label{sec:spec34}

\begin{table}[htbp]
\centering\footnotesize
\caption{Modelling requirements established in \cref{sec:learning}.}\label{tab:spec34}
\begin{tabularx}{\textwidth}{@{}X >{\raggedright\arraybackslash}p{5.2cm}@{}}
\toprule
\textbf{Requirement} & \textbf{Basis} \\
\midrule
Target is the forecast margin to a reference threshold, not the class & \cref{sec:formulation} \\
Tier derived from the band table, not from a single margin & \cref{sec:threshold-table} \\
Horizon is the next measurement occasion in the routine regime & \cref{sec:horizon-adopted} \\
The dual-task formulation disambiguated explicitly & \cref{sec:dual-task} \\
Ordinal structure preserved in the loss or carried by the margin & \cref{sec:label-design} \\
Nested cross-validation, 5 outer by 4 inner, grouped, stratified, forward in time & \cref{sec:cv-structure} \\
Repeated $R = 10$ times; the distribution reported & \cref{sec:repeats} \\
Tree depth 2--3, not 6 & \cref{tab:xgb-config} \\
Early stopping on inner validation, not fixed rounds & \cref{tab:xgb-config} \\
Recurrent model at most one layer and sixteen units; underperformance pre-stated & \cref{sec:lstm-config} \\
Seeds fixed and reported & \textcite{mendoza2025} \\
Resampling strictly inside the training fold & \cref{sec:imbalance-training} \\
Prior correction applied after any weighting & \cref{sec:prior-correction} \\
Platt scaling, fitted on inner validation & \cref{sec:calibration} \\
Brier score reported with base rate and decomposition & \cref{sec:brier} \\
Two thresholds, selected on inner validation by target sensitivity & \cref{sec:threshold-selection} \\
Alert volume reported at every operating point & \cref{sec:volume-arithmetic} \\
$\text{tier}_{\text{final}} = \max(\text{rule}, \text{model})$ & \cref{sec:rule-floor} \\
The four pre-registration items stated in the methodology & \cref{sec:prereg} \\
\textcite{efendi2025} cited for architecture only & \cref{sec:task-not-algo} \\
\bottomrule
\end{tabularx}
\end{table}

%==========================================================================
\subsection{Assumptions on which this section rests}\label{sec:assumptions}

\Cref{sec:collection,sec:preprocessing,sec:feature-eng} were written parametrically in quantities that the collection scenario fixes. This section cannot be, since a training protocol must name its sample size. \Cref{tab:assumptions} states the values assumed, and what would have to be revised if each is overturned.

\begin{table}[htbp]
\centering\footnotesize
\caption{Working assumptions and their consequences.}\label{tab:assumptions}
\begin{tabularx}{\textwidth}{@{}>{\raggedright\arraybackslash}p{3.2cm} >{\raggedright\arraybackslash}p{3.4cm} X@{}}
\toprule
\textbf{Quantity} & \textbf{Assumed} & \textbf{What changes if overturned} \\
\midrule
Collection scenario & B: prospective cohort on the protocol of \cref{sec:streams} & under a retrospective scenario the adjudicated secondary endpoint of \cref{sec:two-labels} is unavailable; under the inpatient variant the horizon shortens and the sequence lengthens \\
Horizon $\Delta$ & next measurement occasion, three to four days & a denser protocol restores the 24--48 hour horizon of the original contribution \\
Feature window $L$ & 14 days & \cref{tab:ceiling} fixes the minimum; a denser protocol permits 7 days \\
Independent samples $n$ & 60 to 90 (\cref{sec:overlap}) & every $n/p$ figure, the depth limit and the bin count for calibration \\
Feature count $p$ & 8 primary, 16 extended & \cref{tab:feature-sets} \\
Planning prevalence $\pi$ & 0.02 & the weight $w$, the alert volume, and the Brier base rate \\
\bottomrule
\end{tabularx}
\end{table}

Two items are not assumptions but missing documents, carried forward from \cref{sec:feature-limits}: the guideline editions that define the labels. The band table of \cref{sec:threshold-table} cannot be populated without them, and the labelling on which this entire chapter rests is defined by them.%
\verify{Until the guideline editions are obtained, every threshold-derived quantity in \cref{sec:learning}---the reference threshold $\theta_c^{\text{ref}}$, the band boundaries, the rule baseline and the rule floor---is computed from the operational values of \cref{tab:verified-thresholds} and must be recomputed once the guidelines are in hand.}

%##########################################################################
% CHAPTER 4
%##########################################################################
\chapter{Overall System Architecture and Software Solution Design}\label{ch:architecture}

\section{Multi-Tier System Architecture}\label{sec:architecture}

\subsection*{Purpose, sources and a correction on applicable law}

\Cref{ch:data} specified the data and the model. This chapter specifies the system that carries them. \Cref{sec:architecture} fixes the layer decomposition, the data flow between layers, the latency budget, the failure modes, and the security and legal framework.

Four sources are used, all examined in \cref{ch:techsota}. \Textcite{mucchi2021} supply the four-layer reference architecture, derived from the ETSI SmartBAN high-level model, and the security principle adopted in \cref{sec:security-design}. \Textcite{efendi2025} supply the only implemented stack in the corpus with a working cloud backend; the reservations of \cref{sec:task-not-algo} bear on their modelling, not on their architecture, which is why they are cited here. \Textcite{liu2025review} survey security techniques and interoperability standards, under the standing reservation of \cref{sec:cloud-gaps}. \Textcite{mihevc2025} supply the only clinically validated system in the corpus and the only statement of a regulatory framework with which the authors actually complied.

\textbf{The correction.} \Textcite{mihevc2025} state that their data \bq{were encrypted and complied with the General Data Protection Regulation}. Their trial was conducted in Slovenia, within the European Union, where that regulation is binding. A system deployed in Cameroon, processing the data of Cameroonian patients, operated by a Cameroonian institution, falls outside its territorial scope. Asserting compliance with an instrument that imposes no obligation is an error of category, and worse, it displaces the work of identifying the instrument that does apply. \Cref{sec:legal} treats this properly rather than reproducing the claim.

%==========================================================================
\subsection{The Four Layers}\label{sec:layers}

\subsubsection{The model and its provenance}

The decomposition is that of \textcite{mucchi2021}, with one modification. Their architecture comprises BODY, ENVIRONMENT, EDGE and CLOUD layers. This system has no environmental sensors, so that layer is omitted; and the application tier, which \textcite{mucchi2021} fold into the cloud layer as \bq{E-Health Services}, is promoted to a layer of its own, following \textcite{efendi2025} and \textcite{liu2025review}.

\begin{table}[htbp]
\centering\small
\caption{Layer correspondence across the four reference architectures.}\label{tab:layer-map}
\begin{tabularx}{\textwidth}{@{}c >{\raggedright\arraybackslash}p{2.4cm} >{\raggedright\arraybackslash}p{2.2cm} >{\raggedright\arraybackslash}p{3.4cm} X@{}}
\toprule
\textbf{\#} & \textbf{\textcite{mucchi2021}} & \textbf{\textcite{efendi2025}} & \textbf{\textcite{liu2025review}} & \textbf{This system} \\
\midrule
1 & BODY & Physical & Access & Physical \\
--- & ENVIRONMENT & --- & Access & omitted \\
2 & EDGE & (absent) & Data transfer & Edge \\
3 & CLOUD & Network & Storage, processing, analysis & Cloud \\
4 & (within CLOUD) & Application & Application service, decision support & Application \\
\bottomrule
\end{tabularx}
\end{table}

\Cref{sec:multitier} established that the layer on which the four sources disagree is the Edge, and that its presence or absence is the architectural decision that matters. This system has one, for three reasons developed below: it is where triage happens (\cref{sec:edge-triage}); it is where transmission priority is decided (\cref{sec:transmission-priority}); and it is what allows the system to function during the connectivity interruptions that \cref{sec:digital-barriers} quantified.

\begin{figure}[htbp]
\centering
\begin{tikzpicture}[font=\footnotesize,
  lay/.style={draw,rounded corners=2pt,minimum width=9.2cm,minimum height=1.15cm,align=center},
  ann/.style={font=\scriptsize,text=cGrey,anchor=west,align=left}]

  \node[lay,fill=cRef!10] (app) at (0,0)
    {\textbf{APPLICATION}\\\scriptsize patient view \textbar{} nurse view \textbar{} physician queue};
  \node[lay,fill=cRef!18] (cloud) at (0,-1.6)
    {\textbf{CLOUD}\\\scriptsize repository \textbar{} control \textbar{} features \textbar{} scoring \textbar{} alerts \textbar{} registry};
  \node[lay,fill=cAfrica!14] (edge) at (0,-3.2)
    {\textbf{EDGE} \;\scriptsize(patient smartphone: HUB and Gateway)\\\scriptsize collection \textbar{} validation \textbar{} rule tiering \textbar{} local store \textbar{} transmission};
  \node[lay,fill=cGrey!15] (phy) at (0,-4.8)
    {\textbf{PHYSICAL}\\\scriptsize validated upper-arm cuff \textbar{} glucometer \textbar{} (optional heart-rate sensor)};
  \node[draw,dashed,text=cGrey,minimum width=9.2cm,minimum height=0.6cm] (env) at (0,-5.9)
    {\scriptsize ENVIRONMENT --- omitted, no environmental sensors in scope};

  % upstream
  \draw[-{Stealth[length=2.2mm]},thick,cAfrica] (phy.east)  -- ++(0.7,0) |- (edge.east)  node[ann,pos=0.25,xshift=4mm]{BLE};
  \draw[-{Stealth[length=2.2mm]},thick,cAfrica] (edge.east)  -- ++(1.4,0) |- (cloud.east) node[ann,pos=0.25,xshift=11mm]{MQTT / TLS};
  \draw[-{Stealth[length=2.2mm]},thick,cAfrica] (cloud.east) -- ++(0.7,0) |- (app.east)   node[ann,pos=0.25,xshift=4mm]{REST};
  % downstream
  \draw[-{Stealth[length=2.2mm]},thick,cRef] (app.west)   -- ++(-0.7,0) |- (cloud.west) node[ann,pos=0.25,xshift=-26mm]{decision};
  \draw[-{Stealth[length=2.2mm]},thick,cRef] (cloud.west) -- ++(-1.4,0) |- (edge.west)  node[ann,pos=0.25,xshift=-30mm]{WebSocket push};
  \draw[-{Stealth[length=2.2mm]},thick,cRef] (edge.west)  -- ++(-0.7,0) |- (phy.west)   node[ann,pos=0.25,xshift=-26mm]{prompt};
\end{tikzpicture}
\caption{Layer decomposition and data flow. Red denotes the upstream measurement path, blue the downstream decision path. The Edge is the layer on which the four reference architectures of \cref{tab:layer-map} disagree.}\label{fig:layers}
\end{figure}

\subsubsection{Physical layer}

\begin{table}[htbp]
\centering\small
\caption{Physical-layer specification.}\label{tab:physical-layer}
\begin{tabularx}{\textwidth}{@{}>{\raggedright\arraybackslash}p{3.6cm} X >{\raggedright\arraybackslash}p{2.9cm}@{}}
\toprule
\textbf{Component} & \textbf{Specification} & \textbf{Basis} \\
\midrule
Blood pressure monitor & automated upper-arm cuff, independently validated, Bluetooth Low Energy & \cref{sec:device-landscape}; \textcite{who2023hypertension} \\
Glucometer & strip-based, Bluetooth Low Energy & \cref{sec:device-landscape} \\
Measurement enforcement & two readings, interval $\geq$ 45~s enforced by the device, second retained & \cref{sec:mihevc} \\
Optional covariate sensor & heart rate only; never blood pressure, never SpO$_2$ & \cref{sec:cuffless} \\
Device register & make, model, serial, validation reference, calibration date & \cref{sec:record-schema} \\
\bottomrule
\end{tabularx}
\end{table}

Excluded by design: cuffless wrist devices, smartphone-camera blood pressure, and any inferential blood-pressure method. \Cref{sec:device-landscape} established the WHO position, that no alternative to the validated upper-arm cuff existed at the time of the 2023 report, and \cref{sec:cuffless} established that three sources converge on it. This is Position~A of \cref{sec:acq-decision} expressed as a hardware specification.

\subsubsection{Edge layer}\label{sec:edge-layer}

\Textcite{mucchi2021} decompose the Edge into a Gateway containing a reception and transmission unit, and a HUB containing data collection, data processing, local storage and transmission units. In this system all of these run on the patient's smartphone, which is simultaneously HUB and Gateway; their \emph{Intra-Com-3} link is therefore absent, being the link that exists only when HUB and Gateway are separate devices.

\begin{table}[htbp]
\centering\small
\caption{Edge-layer specification.}\label{tab:edge-layer}
\begin{tabularx}{\textwidth}{@{}>{\raggedright\arraybackslash}p{2.8cm} X >{\raggedright\arraybackslash}p{3.1cm}@{}}
\toprule
\textbf{Unit} & \textbf{Function} & \textbf{Basis} \\
\midrule
Data collection & BLE ingestion; capture of \texttt{ts\_measurement} at source & \cref{sec:record-schema} \\
Data processing & deterministic validation; rule-engine tier evaluation; aggregation of any high-frequency covariate stream & \cref{sec:ingestion,sec:rule-model,sec:alert-volume} \\
Local storage & offline-first queue, encrypted at rest, surviving application restart & \cref{sec:offline} \\
Transmission & MQTT upstream; tier-dependent priority; SMS fallback on red & \cref{sec:protocols-cloud,sec:transmission-priority} \\
Patient interface & measurement prompts, reading display, trend view, therapy view & \cref{sec:feedback} \\
\bottomrule
\end{tabularx}
\end{table}

Three functions sit at the Edge that most designs place in the Cloud, and each has a reason.

The \textbf{rule-engine tier evaluation} runs here because the tier determines transmission urgency (\cref{sec:transmission-priority}). A rule evaluated only in the cloud cannot tell the phone that a reading is urgent, which is circular: the phone would have to transmit in order to learn whether to transmit quickly.

\textbf{Aggregation of high-frequency streams} runs here because \cref{sec:volume-arithmetic} established that per-reading alerting on a 1~Hz stream generates of the order of ten thousand false alerts per patient per week. The Edge reduces such a stream to a daily summary before anything is transmitted or scored.

\textbf{Local storage} is what makes the system usable under intermittent connectivity, which \cref{sec:digital-barriers} established as a structural feature of the deployment context rather than an occasional fault.

\subsubsection{Cloud layer}\label{sec:cloud-layer}

\begin{table}[htbp]
\centering\small
\caption{Cloud-layer specification.}\label{tab:cloud-layer}
\begin{tabularx}{\textwidth}{@{}>{\raggedright\arraybackslash}p{2.8cm} X >{\raggedright\arraybackslash}p{3.1cm}@{}}
\toprule
\textbf{Unit} & \textbf{Function} & \textbf{Basis} \\
\midrule
Repository & longitudinal record keyed for window queries; both \texttt{ts\_measurement} and \texttt{ts\_recorded} stored & \cref{sec:schema,sec:record-schema} \\
Control & identity, roles, access policy, audit log & \cref{sec:security} \\
Data processing & feature construction; quality flags; expanding-window normalisation & \cref{sec:preprocessing,sec:feature-eng} \\
Analysis & model scoring; prior correction; Platt calibration; tier mapping & \cref{sec:risk-score,sec:tiers} \\
Alert management & queue construction, priority ordering, escalation, non-transmission detection & \cref{sec:decision-flow,sec:alert-volume} \\
Offline training & periodic retraining from the warehouse; model registry with versions & \cref{sec:feedback} \\
Services & REST for clients; WebSocket for push & \cref{sec:protocols-cloud} \\
\bottomrule
\end{tabularx}
\end{table}

The separation of online scoring from offline training is taken from \textcite{liu2025review}; \cref{sec:feedback} established that the separation transfers while their tooling does not, being dimensioned for a throughput this cohort will never approach.

The \textbf{model registry} appears in no source and is required by \cref{sec:leak-boundary,sec:interop}: every score must be attributable to a model version and a threshold-table version, or no result can be reproduced once either changes. Since \cref{sec:assumptions} records that the guideline editions defining the threshold table are not yet in hand, the registry is also the mechanism by which scores computed under provisional thresholds remain identifiable after the definitive table is installed.

\subsubsection{Application layer}\label{sec:application-layer}

Two profiles follow the SELF and PRO distinction of \textcite{mucchi2021}, with a third added for the task-sharing arrangement that \cref{sec:hearts-spec} established as the regional norm.

\begin{table}[htbp]
\centering\small
\caption{Application-layer profiles and capabilities.}\label{tab:profiles}
\begin{tabularx}{\textwidth}{@{}>{\raggedright\arraybackslash}p{2.1cm} X@{}}
\toprule
\textbf{Profile} & \textbf{Capabilities} \\
\midrule
Patient & enter or confirm readings; view own history and trend; receive prompts; view prescribed therapy; receive the outcome of a clinician decision \\
Nurse & the above for assigned patients; perform and attribute measurements; view flagged readings; \emph{cannot} modify therapy \\
Physician & priority-ordered queue; patient record with trend; model score with its explanation; proposed protocol step; approve, override or request consultation; every action written to the adjudication table \\
\bottomrule
\end{tabularx}
\end{table}

The priority-ordered queue is from \textcite{yazdani2018}, and \cref{sec:alert-volume} established that an ordering rather than a partition is what the alert arithmetic requires. The adjudication write-back is from \cref{sec:safety} and is the mechanism that accumulates the clinician-adjudicated labels of \cref{sec:two-labels}.

\subsubsection{Traceability}\label{sec:traceability}

\begin{table}[htbp]
\centering\footnotesize
\caption{Provenance of each design element.}\label{tab:traceability}
\begin{tabularx}{\textwidth}{@{}X >{\raggedright\arraybackslash}p{5.4cm}@{}}
\toprule
\textbf{Design element} & \textbf{Origin} \\
\midrule
Four-layer decomposition & \textcite{mucchi2021}, ETSI SmartBAN \\
Edge with local storage and processing & \textcite{mucchi2021} \\
MQTT upstream, HTTP downstream, WebSocket push & \textcite{efendi2025} \\
Offline-first with deferred merge & \textcite{efendi2025} \\
Three-tier alerting with verification profile & \textcite{mihevc2025} \\
Priority-ordered clinician queue & \textcite{yazdani2018} \\
Configurable checkup period and rule selection & \textcite{yazdani2018} \\
Offline training, online scoring split & \textcite{liu2025review} \\
ICD and SNOMED CT coding & \textcite{liu2025review} \\
Validated-cuff restriction & \textcite{who2023hypertension}; \cref{sec:device-landscape} \\
Tier-dependent transmission priority & this work, \cref{sec:transmission-priority} \\
Alert-volume-driven aggregation at the Edge & this work, \cref{sec:volume-arithmetic} \\
Rule-engine floor $\max(\text{rule},\text{model})$ & this work, \cref{sec:rule-floor} \\
Model and threshold-table version registry & this work, \cref{sec:cloud-layer} \\
Latency decomposition and the dominance of $t_5$ & this work, \cref{sec:latency} \\
Failure-mode and degradation table & this work, \cref{sec:failure} \\
\bottomrule
\end{tabularx}
\end{table}

%==========================================================================
\subsection{Inter-Layer Data Flow}\label{sec:dataflow}

\subsubsection{The nominal path}

\begin{equation}
\underbrace{\text{cuff, glucometer}}_{\text{Physical}}
\xrightarrow{\;\text{BLE}\;}
\underbrace{\text{phone}}_{\text{Edge}}
\xrightarrow{\;\text{MQTT/TLS}\;}
\underbrace{\text{broker} \to \text{repository} \to \text{features} \to \text{score} \to \text{tier}}_{\text{Cloud}}
\to \underbrace{\text{queue} \to \text{clinician}}_{\text{Application}}
\end{equation}

\subsubsection{Where each decision sits}

\begin{table}[htbp]
\centering\small
\caption{Decision placement by layer and latency class.}\label{tab:decision-placement}
\begin{tabularx}{\textwidth}{@{}X >{\raggedright\arraybackslash}p{2.8cm} >{\raggedright\arraybackslash}p{2.4cm} >{\raggedright\arraybackslash}p{2.1cm}@{}}
\toprule
\textbf{Decision} & \textbf{Layer} & \textbf{Latency class} & \textbf{Basis} \\
\midrule
Is the reading well-formed? & Edge & milliseconds & \cref{sec:ingestion} \\
Which tier under the clinical rule? & Edge & milliseconds & \cref{sec:transmission-priority} \\
How urgently to transmit? & Edge & milliseconds & this work \\
Which tier under the model? & Cloud & seconds & \cref{sec:tiers} \\
Final tier $=\max(\text{rule},\text{model})$ & Cloud & seconds & \cref{sec:rule-floor} \\
Where in the queue? & Cloud & seconds & \cref{sec:alert-volume} \\
Confirm, override or escalate? & Application, human & hours to days & \cref{sec:decision-flow} \\
\bottomrule
\end{tabularx}
\end{table}

The last row dominates the total, by a margin that \cref{sec:latency} quantifies.

\subsubsection{Tiered transmission priority}\label{sec:transmission-priority}

This is the design decision that justifies running the rule engine at the Edge, and it appears in no source in the corpus.

\Cref{sec:offline} established that offline-first storage with deferred merge removes one cause of missing data, since readings taken offline are delayed rather than lost. It also noted the cost: arrival time decouples from measurement time, sometimes by days, which is why \cref{sec:record-schema} stores both.

For routine readings that trade is favourable. For a red-tier reading it is not. A hypertensive crisis queued on a handset until the next period of connectivity is a clinical failure that the architecture itself caused, and the fact that the reading eventually arrives is no mitigation.

The resolution is to make transmission priority a function of the tier, which requires the tier to be known at the Edge.

\begin{table}[htbp]
\centering\small
\caption{Transmission behaviour by tier, evaluated at the Edge.}\label{tab:transmission-priority}
\begin{tabularx}{\textwidth}{@{}>{\raggedright\arraybackslash}p{2.0cm} X@{}}
\toprule
\textbf{Tier} & \textbf{Transmission behaviour} \\
\midrule
Green & queue locally; batch synchronisation on next connectivity; no retry urgency \\
Amber & attempt immediate send; exponential backoff; synchronisation within hours \\
Red & immediate send; aggressive retry; SMS fallback after $n$ failures; local on-device warning to the patient regardless of connectivity \\
\bottomrule
\end{tabularx}
\end{table}

The SMS fallback is available because \textcite{dodoo2022} report 816 million mobile connections in sub-Saharan Africa against a projection of 1.05 billion, and because \textcite{yazdani2018} already use SMS in their reminder loop. SMS reaches where data does not, and a red alert is short enough to fit one message.

Two cautions attach to that justification and are stated rather than elided. A connection count is not a count of people, since multiple subscriptions per user are common in the region, so the figure overstates population coverage; and a national connection total establishes nothing about signal availability at a particular patient's residence, which is the quantity that actually governs whether the fallback works.%
\verify{The \textcite{dodoo2022} projection is stated to a horizon now past. Confirm whether the realised figure is available, and replace the projection with it; a projection to a date already elapsed invites the obvious question.}

The local on-device warning matters independently of transmission. A patient whose reading is critical should be told so by the handset, immediately, even when nothing can be sent. That instruction is generic and therefore safe under \cref{sec:safety}: advising a patient to seek care is not a dose change, and the prohibition established there is on autonomous therapeutic action, not on autonomous information.

\subsubsection{The reverse path}

\begin{equation}
\text{clinician decision} \to \text{adjudication table} \to \text{Cloud}
\xrightarrow{\;\text{push}\;} \text{phone} \to \text{patient}
\end{equation}

Two requirements. Any therapy change passes a clinician, and nothing on this path is generated autonomously (\cref{sec:safety}). And every decision is written to the adjudication table of \cref{sec:adjudication} \emph{before} it is delivered, so that the label is captured even when delivery fails --- which is the common case precisely for the patients whose connectivity is worst, and therefore the case in which losing the label would bias the adjudicated sample.

%==========================================================================
\subsection{The Latency Budget}\label{sec:latency}

No source in this corpus decomposes end-to-end latency, and \textcite{bediang2021} supply the one field measurement that makes the decomposition worth performing.

\subsubsection{Decomposition}

\begin{table}[htbp]
\centering\small
\caption{Latency segments.}\label{tab:latency}
\begin{tabularx}{\textwidth}{@{}X c >{\raggedright\arraybackslash}p{2.6cm} >{\raggedright\arraybackslash}p{3.4cm}@{}}
\toprule
\textbf{Segment} & \textbf{Symbol} & \textbf{Expected} & \textbf{Governed by} \\
\midrule
Measurement to phone & $t_1$ & seconds & BLE \\
Phone to cloud, online & $t_2$ & seconds & network \\
Phone to cloud, offline & $t_2'$ & hours to days & connectivity; \cref{sec:transmission-priority} \\
Ingestion, features, scoring & $t_3$ & $<1$~s & \cref{sec:assumptions}; \cref{sec:training} \\
Queue to clinician view & $t_4$ & seconds & push \\
Clinician review & $t_5$ & hours to days & staffing \\
Decision to patient & $t_6$ & seconds to hours & reverse path \\
\bottomrule
\end{tabularx}
\end{table}

The bound on $t_3$ follows from \cref{sec:training}: at the independent sample size and the primary feature count fixed in \cref{tab:assumptions}, a single forward pass through a gradient-boosted ensemble of shallow trees is a sub-millisecond operation, and the second is consumed by ingestion and feature construction, not by the model.

\subsubsection{Where the time actually goes}

\Textcite{bediang2021} measured both ends of this pathway in a Cameroonian district-hospital setting.

\begin{table}[htbp]
\centering\small
\caption{Measured delays, from \textcite{bediang2021}.}\label{tab:bediang-latency}
\begin{tabularx}{\textwidth}{@{}X c c c@{}}
\toprule
\textbf{Measured quantity} & \textbf{Intervention} & \textbf{Control} & \textbf{$p$} \\
\midrule
Time to perform the examination & $1.3 \pm 4.1$ days & $7.0 \pm 14.5$ days & $<0.01$ \\
Time to receive expert opinion & $10.3 \pm 11.6$ days & $11.2 \pm 27.9$ days & $0.82$ \\
\bottomrule
\end{tabularx}
\end{table}

The protocol under which these were collected required expert opinion within 72 hours. The measured mean in the intervention arm was 10.3 days, which exceeds that target by a factor of approximately 3.4.

The technology accelerated acquisition by a factor of about five, from 7.0 days to 1.3. It did not accelerate the clinical decision at all: the expert-opinion delay was not distinguishable from the control pathway, in which patients travelled to a referral centre.

One caution on these statistics. In three of the four cells the standard deviation exceeds the mean, by a factor of three in the intervention examination time and by a factor of two and a half in the control expert-opinion time. A non-negative quantity with that dispersion cannot be approximately normal; the distributions are heavily right-skewed, the means are therefore poor central estimates, and a $p$-value computed on means and standard deviations under a normality assumption is not reliable.%
\verify{Confirm whether \textcite{bediang2021} report medians, interquartile ranges or a non-parametric test alongside these means, and whether the fifty-day maximum attributed to them in \cref{sec:failure} appears in the paper or is inferred. The qualitative conclusion --- acquisition accelerated, decision not --- survives either way, since a factor of five and a null result are both robust to the distributional shape; the point estimates are what should not be quoted precisely.}

This is the quantified form of the argument of \cref{sec:alert-volume}: $t_1$ through $t_4$ are engineering problems, solved and small, while $t_5$ is a resource problem and dominates the total by orders of magnitude (\cref{fig:latency}).

\begin{figure}[htbp]
\centering
\begin{tikzpicture}
\begin{axis}[width=0.92\textwidth,height=5.4cm,
  xbar,
  xmode=log,
  xmin=0.1,xmax=5e6,
  xlabel={Elapsed time (seconds, logarithmic)},
  symbolic y coords={$t_3$ scoring,$t_1$ BLE,$t_2$ network,$t_4$ push,$t_6$ to patient,$t_2'$ offline queue,$t_5$ clinician},
  ytick=data,
  y=0.52cm,
  nodes near coords={},
  enlarge y limits=0.12,
  xmajorgrids,grid style={gray!20},
  tick label style={font=\scriptsize},label style={font=\scriptsize}]
\addplot[fill=cRef!40,draw=cRef] coordinates {
  (0.5,$t_3$ scoring) (5,$t_1$ BLE) (2,$t_2$ network) (2,$t_4$ push) (3600,$t_6$ to patient)};
\addplot[fill=cGrey!40,draw=cGrey] coordinates {(86400,$t_2'$ offline queue)};
\addplot[fill=cAfrica!45,draw=cAfrica] coordinates {(889920,$t_5$ clinician)};
\end{axis}
\end{tikzpicture}
\caption{Latency segments on a logarithmic scale. The engineering segments $t_1$ to $t_4$ occupy under ten seconds in total; $t_5$, taken at the 10.3-day mean of \textcite{bediang2021}, exceeds them by about five orders of magnitude.}\label{fig:latency}
\end{figure}

\subsubsection{What follows for the architecture}

Optimising $t_1$ through $t_4$ has almost no effect on the outcome. A design change reducing scoring latency from 800~ms to 200~ms saves 0.6~s against a total of 10.3 days, which is $8.9 \times 10^5$~s: a saving of about seven parts in ten million, or $7 \times 10^{-5}$ per cent of the total. No plausible engineering improvement to those segments is clinically detectable.

Reducing $t_5$ is the only material improvement available, and the architecture can act on it in exactly two ways.

\textbf{Reduce queue length, by triage.} This is the arithmetic of \cref{sec:alert-volume}, delivered by the tiering of \cref{sec:tiers}: fewer items in the queue means less waiting per item, for a fixed clinician capacity.

\textbf{Reduce per-item review time, by what is presented.} A clinician who must open a record, scroll a table and mentally reconstruct a trajectory spends minutes per item. One shown the fourteen-day trajectory, the threshold line, the calibrated score with its two or three contributing features, and the proposed protocol step, spends seconds. The dashboard design of the next section is therefore a latency intervention and is justified on those grounds, not as a presentational refinement. This is the one place in the system where interface design and clinical outcome are connected by an arithmetic argument rather than by assertion.

$t_2'$ is bounded by the tiering of \cref{sec:transmission-priority}, which is the reason that mechanism exists: it caps the offline delay for the readings where delay is harmful, while leaving it unbounded where it is not.

%==========================================================================
\subsection{Failure Modes and Degradation}\label{sec:failure}

No source in this corpus analyses what happens when a layer fails. Given that \cref{sec:regional-landscape} documented a regional telemedicine service ceasing operations after reaching 2.5 million users, and that \textcite{dodoo2022} identify sustainability as the characteristic failure mode of the region's projects, graceful degradation is not an optional refinement.

\begin{table}[htbp]
\centering\footnotesize
\caption{Failure modes, detection and degraded behaviour.}\label{tab:failure-modes}
\begin{tabularx}{\textwidth}{@{}>{\raggedright\arraybackslash}p{2.9cm} >{\raggedright\arraybackslash}p{2.3cm} X >{\raggedright\arraybackslash}p{3.0cm}@{}}
\toprule
\textbf{Failure} & \textbf{Detected by} & \textbf{Degradation} & \textbf{Residual risk} \\
\midrule
Device battery exhausted & absence of readings & non-transmission alarm & patient unmonitored until noticed \\
Device lost or broken & as above & as above & as above \\
Phone unavailable & as above & none at the Edge & no local alerting at all \\
Connectivity lost & Edge & local queue; red escalates by SMS & amber delayed \\
Cloud unavailable & Edge retry & local queue; Edge rule still tiers; red by SMS & no model score, no queue \\
Model service fails & Cloud & rule floor: $\text{tier}=\text{tier}_{\text{rule}}$ & loss of personalisation only \\
Calibration drifts & monitoring & thresholds revert to rule & over- or under-alerting \\
Clinician unavailable & queue age & escalation after a defined interval & unbounded $t_5$ \\
Guideline superseded & manual & versioned table swapped; registry preserves provenance & historical scores not comparable \\
\bottomrule
\end{tabularx}
\end{table}

Three rows carry the design argument.

\textbf{Model failure degrades to guideline conformity, not to nothing.} This is the composition $\text{tier}_{\text{final}} = \max(\text{tier}_{\text{rule}}, \text{tier}_{\text{model}})$ of \cref{sec:rule-floor}, and it is the property that makes the system defensible to a clinician: its worst case equals the current standard of care. Given that the best model of \textcite{mendoza2025} attains a sensitivity of 0.51, a design whose failure mode is to behave like the guideline is the correct one, and the composition is what converts that observation into a structural guarantee rather than a hope.

\textbf{Phone unavailability is a single point of failure.} All five Edge functions of \cref{tab:edge-layer} run on one device. No source addresses this, and the honest statement is that a patient without a working handset is outside the system. The only available mitigation is the non-transmission alarm, which detects the condition without remedying it --- and which depends on the Cloud noticing an absence, this being the one detection that cannot by construction be performed at the Edge. A device that has failed cannot report its own failure.

\textbf{Clinician unavailability requires an escalation rule}, which no source specifies. The HEARTS indicators of \cref{sec:hearts-spec} include patients with no visit in three and in twelve months; the analogous system quantity is queue age. A red item unreviewed after a defined interval must escalate to a named alternate. Without such a rule $t_5$ is unbounded, and the extreme tail of the \textcite{bediang2021} distribution becomes the expected behaviour rather than an outlier --- which is precisely what the skew noted in \cref{sec:latency} implies, since in a right-skewed delay distribution the long waits are not rare accidents but a stable fraction of the caseload.

%==========================================================================
\subsection{Security}\label{sec:security-design}

\subsubsection{The organising principle}

\Textcite{mucchi2021} state it, and it is the correct organising idea:

\begin{quote}\small
\bq{Health data can flow from patient to cloud and from cloud to others by using different channels, which not always can be assumed secure. A trusted and secure domain for the health data should be assured, so that there is no segment which has a lower security than others.}
\end{quote}

Security is determined by the weakest segment. The design task is therefore to enumerate the segments and specify each, rather than to accumulate techniques.

\subsubsection{The seven segments}

\begin{table}[htbp]
\centering\footnotesize
\caption{Security segments and their protection.}\label{tab:security-segments}
\begin{tabularx}{\textwidth}{@{}c >{\raggedright\arraybackslash}p{2.6cm} X >{\raggedright\arraybackslash}p{2.6cm}@{}}
\toprule
\textbf{\#} & \textbf{Segment} & \textbf{Protection} & \textbf{Basis} \\
\midrule
1 & Device to phone & BLE with LE Secure Connections; pairing bound to the patient account; legacy pairing refused & --- \\
2 & Phone at rest & operating-system full-disk encryption required; application-sandboxed store; no identifiers in logs or crash reports & \cref{sec:reident} \\
3 & Phone to cloud & TLS 1.2 or later; certificate pinning; MQTT over TLS & \textcite{efendi2025} \\
4 & Cloud at rest & encryption at rest; server-side access rules enforcing role separation & \textcite{efendi2025,liu2025review} \\
5 & Cloud to clinician & TLS; authenticated session; short expiry; re-authentication before a therapy change & \cref{sec:application-layer} \\
6 & Clinician device & outside system control & --- \\
7 & Export and backup & the analysis file of \cref{sec:reident} & \cref{sec:reident} \\
\bottomrule
\end{tabularx}
\end{table}

\subsubsection{The two weakest segments, which no source addresses}

\textbf{Segment 6.} The clinician reviews on a personal handset or a shared hospital workstation. The pilot of \textcite{mucchi2021} installed their application on the nurses' own smartphones. The system cannot enforce device encryption, screen locking or the absence of malicious software there, and in practice this is the weakest link in almost every deployment of this kind.

What the architecture can do: enforce short session expiry; require re-authentication before any therapeutic action; avoid caching patient data on the client; and log access, so that a compromise is detectable even where it is not preventable. What it cannot do is secure the device, and this thesis states that rather than implying otherwise.

\textbf{Segment 7.} The research export is the one copy of the data that leaves every protection of segments 1 to 5, being a file on a laptop. \Cref{sec:reident} specified its contents and the measures that reduce re-identification risk. Those measures are the security control for this segment and are presented as such, not solely as an ethical matter.

\subsubsection{What the survey offers, and what is proportionate}

The survey of \textcite{liu2025review} is extensive: differentially private synthesis using convolutional autoencoders and generative adversarial networks; $k$-anonymity through the Mondrian algorithm; RFID-based authentication for remote patient monitoring providing \bq{mutual authentication, user untraceability, prevention of replay attacks, forward and backward secrecy, and data integrity}; lightweight encryption with access tokens; permissioned blockchain with elliptic-curve digital signatures; blockchain combined with a distributed file system; private set intersection with $k$-anonymity; and steganographic concealment of clinical text within medical images.

Almost none of this is proportionate to a thirty-patient pilot, and adopting it would be security theatre: a catalogue of mechanisms whose presence in a dissertation signals diligence while protecting nothing that is actually at risk here. The proportionate set is transport encryption, encryption at rest, role-separated authentication, access logging, and the de-identification measures of \cref{sec:reident}.

Two items from the list nonetheless deserve mention, for different reasons. \textbf{Access logging for non-repudiation}, which both \textcite{mucchi2021} and one of the surveyed works reach for blockchain to achieve, is obtainable from an append-only audit table at negligible cost, and \cref{sec:security} already required it. And \textbf{$k$-anonymity} is directly relevant, because \cref{sec:reident} established that thirty patients with exact ages and dates are uniquely identifiable; the banding and date-shifting specified there are the operational form of that idea, and naming the correspondence is more useful than citing the algorithm.

\subsubsection{Implemented versus recommended}

\Cref{sec:security-gap} identified the characteristic pattern of this literature: a security section cataloguing techniques the authors did not use. \Textcite{mucchi2021} are candid about it, writing that \bq{we only relied on the application layer security protocols, and on the fact that the data where stored in a private server without public access} --- having just recommended physical-layer security, cryptographic key extraction from biosignals, and blockchain audit trails.

This thesis therefore presents two separate statements: what the system implements, and what a production deployment would require beyond it. Conflating them is the failure the corpus demonstrates, and the partition of \cref{sec:built-specified} is where that separation is made concrete.

One empirical note is worth including. \Textcite{efendi2025} surveyed their users across four dimensions, and security scored lowest, at 4.06 out of 5, or 81.2\%, against 4.56 for effectiveness. Users notice. A security section that is candid about its limits is also the one most likely to be believed.%
\verify{A mean of 4.06 is not attainable from ten respondents on a single five-point item, since any such mean is a multiple of 0.1. It is attainable as the mean of a five-item subscale over ten respondents, which yields multiples of 0.02. Confirm the instrument's item count before quoting the figure, and report $n$ and the item count alongside it.}

%==========================================================================
\subsection{Legal Framework}\label{sec:legal}

\subsubsection{The General Data Protection Regulation does not apply}\label{sec:gdpr}

\Textcite{mihevc2025} complied with the GDPR because they processed the data of European Union residents within the Union. This system does neither. To write that it is \bq{GDPR compliant} would assert compliance with an instrument imposing no obligation here, and --- the more serious cost --- would stand in place of identifying the instrument that does.

The position adopted is instead the following. The design adopts the data-protection principles of the GDPR --- lawfulness, purpose limitation, data minimisation, storage limitation, integrity and confidentiality --- as a voluntary design standard, while recording that the applicable legal framework is Cameroonian. That framing is accurate, it preserves the substantive protections, and it is not open to the objection that the claim invites.

\subsubsection{The applicable instrument}\label{sec:applicable-law}

No document in the corpus states the Cameroonian instrument governing the processing of health data, and this thesis does not assert one: an incorrect statutory reference is a more serious defect than an acknowledged gap, because it is a positive misstatement of law rather than an absence.

What the corpus does establish around the gap is informative. \Textcite{mucchi2021} name their national instrument explicitly, which is the practice to follow. \Textcite{liu2025review} name HIPAA, the HITECH Act and the 21st Century Cures Act, all United States law and none applicable here --- which is the same error as the GDPR claim, committed by a survey rather than a trial. And \textcite{dodoo2022} identify the legal and regulatory dimension as one of six barriers to telemedicine in sub-Saharan Africa, covering \bq{policies, legislation, and standards governing confidentiality and privacy}. The gap is therefore a documented feature of the regional context and not an oversight peculiar to this work, which is the correct way to present it.

Four instruments must be identified and cited before deployment: the Cameroonian personal-data and health-data instrument; the institutional policy of the partner facility; the national or regional ethics committee reference for this study; and any African Union instrument on data protection that Cameroon has ratified.%
\verify{These four references are outstanding and are the author's to supply. The third has a documented pathway: \textcite{bediang2021} cite ethical clearance CE00398/N\textdegree CRERSHC/2016, so a Cameroonian study with an identifiable Cameroonian ethics reference exists within this corpus, and the procedure is therefore known. Until the first is supplied, \cref{sec:gdpr} is the operative statement of the legal position and \cref{sec:residency} cannot be resolved.}

\subsubsection{Data residency is a legal choice}\label{sec:residency}

The published endpoint of \textcite{efendi2025} was hosted in a south-east Asian region. Where data physically reside determines which jurisdiction's law governs their storage, so selecting a region is a legal decision taken by means of a configuration flag --- which is exactly the kind of decision that is made by default and never revisited.

Three consequences. Selecting a European region may bring European law into scope unnecessarily, for a system that otherwise has no European nexus, and would make true the claim that \cref{sec:gdpr} declines to make --- by accident rather than by design. Selecting a region distant from Cameroon adds latency to $t_2$ and $t_4$, which \cref{sec:latency} showed to be negligible in the total, so latency is not a legitimate ground for the choice in either direction. And the chosen region must be stated in the thesis, because an unstated region is an unstated jurisdiction.

\subsubsection{Consent}\label{sec:consent}

Under the working assumption of Scenario~B recorded in \cref{tab:assumptions}, individual consent is required and is not optional. \Textcite{efendi2025} record that \bq{each participant assign the letter of consent that willing to be involved in the experiment}, following an interview establishing health history; \textcite{mihevc2025} obtained consent under a registered protocol.

The consent document must cover eight points: what is collected; who sees it; where it is stored and under which jurisdiction; the retention period and destruction date; that participation does not replace clinical care; that withdrawal is possible; what happens to already-collected data on withdrawal; and that the pseudonym mapping is held by the partner physician and is not transferred (\cref{sec:reident}).

The third and fourth of these depend on \cref{sec:residency,sec:applicable-law} respectively, so the consent document cannot be finalised before those are resolved. This dependency is stated because it determines the order of the remaining work: the legal instrument is not a formality to be appended at the end but a precondition for the consent form, which is itself a precondition for collection.

%==========================================================================
\subsection{Specification}\label{sec:spec41}

\begin{table}[htbp]
\centering\footnotesize
\caption{Architectural requirements established in \cref{sec:architecture}.}\label{tab:spec41}
\begin{tabularx}{\textwidth}{@{}X >{\raggedright\arraybackslash}p{5.0cm}@{}}
\toprule
\textbf{Requirement} & \textbf{Basis} \\
\midrule
Four layers: Physical, Edge, Cloud, Application & \textcite{mucchi2021} \\
Environment layer omitted, and the omission stated & \cref{sec:layers} \\
Physical layer restricted to validated upper-arm cuff and glucometer & \textcite{who2023hypertension}; \cref{sec:acq-decision} \\
Edge performs validation, rule tiering, aggregation, local storage, transmission & \cref{sec:edge-layer} \\
Transmission priority determined by tier, at the Edge & \cref{sec:transmission-priority} \\
SMS fallback on red tier, with its coverage caveat stated & \cref{sec:transmission-priority} \\
Local on-device warning on red, independent of connectivity & \cref{sec:transmission-priority,sec:safety} \\
Cloud: repository, control, processing, analysis, alert management, offline training & \cref{sec:cloud-layer} \\
Model and threshold-table version registry & \cref{sec:leak-boundary,sec:interop} \\
Three application profiles; physician sees a priority-ordered queue & \textcite{yazdani2018} \\
Clinician decisions written to the adjudication table before delivery & \cref{sec:safety} \\
Latency budget reported; $t_5$ identified as dominant & \textcite{bediang2021}; \cref{sec:latency} \\
Point estimates from \textcite{bediang2021} quoted with their skew caveat & \cref{sec:latency} \\
Dashboard design justified as a $t_5$ intervention & \cref{sec:latency} \\
Failure table published; model failure degrades to the rule & \cref{sec:failure} \\
Queue-age escalation rule for unreviewed red items & \cref{sec:failure} \\
Handset single point of failure stated explicitly & \cref{sec:failure} \\
Seven security segments enumerated, each specified & \cref{tab:security-segments} \\
Clinician device acknowledged as outside system control & \cref{sec:security-design} \\
Append-only access log & \cref{sec:security} \\
Implemented security stated separately from recommended & \cref{sec:security-gap,sec:built-specified} \\
GDPR not claimed; its principles adopted voluntarily, and said so & \cref{sec:gdpr} \\
Cameroonian instrument, institutional policy and ethics reference cited & \cref{sec:applicable-law} \\
Data residency region stated & \cref{sec:residency} \\
Consent covering the eight points, with its dependencies stated & \cref{sec:consent} \\
\bottomrule
\end{tabularx}
\end{table}

%==========================================================================
\subsection{What Is Built and What Is Specified}\label{sec:built-specified}

\Cref{sec:arch-validation} established that of the four reference architectures in this corpus, two have no implementation at all, one was validated on six healthy volunteers over ten minutes, and the fourth was validated by a small panel on simulated patients with no real patients involved. None has been evaluated in routine care.

This thesis therefore states plainly which parts of the architecture above are built, which are specified, and which are deferred. That distinction is what separates a design chapter from a list of aspirations, and the corpus demonstrates that it is frequently not drawn.

\begin{table}[htbp]
\centering\footnotesize
\caption{Implementation status of the architecture.}\label{tab:build-status}
\begin{tabularx}{\textwidth}{@{}>{\raggedright\arraybackslash}p{3.0cm} X@{}}
\toprule
\textbf{Status} & \textbf{Components} \\
\midrule
Built and exercised & Physical acquisition; Edge ingestion, validation, local storage, rule tiering, transmission; Cloud repository, feature pipeline, model scoring, tier mapping; Application patient and physician views \\
Specified, not built & SMS fallback; queue-age escalation; offline retraining loop; model registry; ICD and SNOMED CT coding \\
Deferred, with reasons & Environment layer; nurse profile; multi-site deployment; production security beyond the proportionate set of \cref{sec:security-design} \\
\bottomrule
\end{tabularx}
\end{table}

Nothing in the second or third column weakens the work. An unbuilt component that is specified, justified and scoped is a contribution; one that is silently absent is a gap. The difference is visible to a reader only if the partition is published, which is why it appears here rather than being left to inference.

Two document dependencies now block \cref{ch:architecture} as they blocked \cref{ch:data}: the guideline editions required by \cref{tab:assumptions} for the threshold table, and the Cameroonian legal instrument required by \cref{sec:applicable-law}. Neither can be substituted by reasoning from the corpus, and both are recorded as outstanding rather than provisionally filled.

%##########################################################################
% SECTION 4.2
%##########################################################################
\section{IoT Hardware Subsystem and Communication Protocols}\label{sec:iot}

\subsection*{Purpose, sources and scope}\label{sec:iot-scope}

\Cref{sec:architecture} placed four layers. This section specifies the optional instrumented branch of the Physical layer and the protocol that carries its output to the Edge.

Two sources are used. \Textcite{efendi2025} supply the component inventory, pin lists and electrical characteristics, and the MQTT/HTTP protocol split; the reservations established in \cref{sec:efendi-appraisal,sec:task-not-algo} mean that their components are taken while their configuration and results are not. \Textcite{cadrage2026} state contribution~C4, a deliberately spare acquisition device using manual entry and smartphone-camera photoplethysmography, and identify measurement noise and intermittent connectivity as the technological limitations of the approach.

A scoping statement opens this section, because without it the subsystem appears to contradict \cref{ch:techsota}.

\begin{quote}
The instrumented subsystem specified here acquires heart rate, oxygen saturation and body temperature. It does \emph{not} acquire blood pressure or capillary glycaemia, which \cref{sec:device-landscape} established may come only from independently validated upper-arm cuffs and glucometers. The subsystem is therefore a covariate channel, optional to the system's clinical function; its purpose is to enrich the feature space and to demonstrate contribution~C4. The clinical decision does not depend on it.
\end{quote}

Stating this at the head of the section forecloses the reading in which the microcontroller appears to be the system's primary acquisition device, which would contradict \cref{sec:acquisition} and is the first point a reader would challenge.

%==========================================================================
\subsection{What the Subsystem Measures, and Why That Is a Limited Claim}\label{sec:measures}

\begin{table}[htbp]
\centering\small
\caption{Parameter provenance and clinical status.}\label{tab:parameter-provenance}
\begin{tabularx}{\textwidth}{@{}>{\raggedright\arraybackslash}p{4.0cm} X >{\raggedright\arraybackslash}p{3.4cm}@{}}
\toprule
\textbf{Parameter} & \textbf{Acquired by} & \textbf{Enters the clinical decision?} \\
\midrule
Systolic, diastolic pressure & validated BLE cuff & yes, primary \\
Capillary glycaemia & BLE glucometer & yes, primary \\
Heart rate & cuff pulse reading; or camera PPG; or MAX30100 & no, covariate \\
Oxygen saturation & MAX30100 & no, covariate \\
Body temperature & MAX30205 & no, covariate \\
\bottomrule
\end{tabularx}
\end{table}

Three consequences follow and are stated rather than left implicit.

\textbf{The subsystem is removable.} If it is not built, or fails in deployment, the system loses feature columns and loses nothing clinically. The failure table of \cref{tab:failure-modes} already reflects this.

\textbf{Its validation burden is correspondingly lighter.} \Cref{sec:device-validation} required Bland--Altman comparison against a validated reference, under movement, for any device carrying a clinical decision. A covariate channel requires that the measurement be reproducible and plausible, not that it be clinically accurate --- a weaker and attainable standard.

\textbf{It is where contribution C4 lives.} The low-cost acquisition claim of \textcite{cadrage2026} cannot rest on the cuff and the glucometer, which are commercial validated instruments costing what they cost. It rests on this subsystem and on the camera alternative of \cref{sec:camera-ppg}.

A fourth consequence is less comfortable and is stated here because a reader will reach it unaided. An automated oscillometric cuff reports pulse rate alongside blood pressure, so heart rate is already available at every clinical measurement occasion without any additional instrument. Neither the camera nor the MAX30100 therefore adds a parameter the system lacks; what they add is temporal density, heart rate at moments when no cuff reading is taken. That is a real but modest gain, and C4 is defensible on that basis and on the device-cost argument of \cref{sec:options}, not on the claim that these channels are otherwise unobtainable.

%==========================================================================
\subsection{Component Specification}\label{sec:components}

Taken from \textcite{efendi2025}, who give electrical characteristics and pin lists but no wiring diagram, no GPIO assignment and no register configuration.

\begin{table}[htbp]
\centering\small
\caption{Microcontroller characteristics, after \textcite{efendi2025}.}\label{tab:esp32}
\begin{tabularx}{\textwidth}{@{}>{\raggedright\arraybackslash}p{4.4cm} X@{}}
\toprule
\textbf{Characteristic} & \textbf{Value} \\
\midrule
Processor & Tensilica Xtensa \\
ROM & 448~kB \\
SRAM & 520~kB \\
Sleep-mode memory & RTC fast SRAM 8~kB; RTC slow SRAM 8~kB \\
Pin count & 36, comprising EN, GND, VIN, 3.3~V and GPIO \\
Connectivity & integrated Wi-Fi and Bluetooth \\
Wireless security & IEEE 802.11, supporting WAPI, WFA, WPA and WPA2 \\
\bottomrule
\end{tabularx}
\end{table}

The RTC sleep memory is the characteristic \textcite{efendi2025} report and do not use, and \cref{sec:power} shows it to be the one that determines whether the device is deployable at all.

\begin{table}[htbp]
\centering\small
\caption{Pulse-oximetry sensor characteristics, after \textcite{efendi2025}.}\label{tab:max30100}
\begin{tabularx}{\textwidth}{@{}>{\raggedright\arraybackslash}p{4.4cm} X@{}}
\toprule
\textbf{Characteristic} & \textbf{Value} \\
\midrule
Function & heart rate and oxygen saturation \\
Principle & reflective dual-wavelength photoplethysmography \\
Wavelengths & 660~nm (red), 880~nm (infrared) \\
Supply & 3.3--5~V optimal \\
Interface & I\textsuperscript{2}C \\
Pins (7) & VIN, SCL, SDA, INT, IRD, RD, GND \\
INT & interrupt; pulse detection \\
IRD & infrared LED indicator \\
\bottomrule
\end{tabularx}
\end{table}

\begin{table}[htbp]
\centering\small
\caption{Temperature sensor characteristics, after \textcite{efendi2025}.}\label{tab:max30205}
\begin{tabularx}{\textwidth}{@{}>{\raggedright\arraybackslash}p{4.4cm} X@{}}
\toprule
\textbf{Characteristic} & \textbf{Value} \\
\midrule
Function & body temperature \\
Range & up to 50~\textdegree C \\
Accuracy & 0.1~\textdegree C \\
Supply & 2.7--3.3~V \\
Operating current & 600~\textmu A optimal \\
Converter & high-resolution sigma-delta ADC \\
Interface & I\textsuperscript{2}C-compatible two-wire \\
Pins (8) & VCC, GND, SDA, SCL, OS, A0, A1, A2 \\
A0--A2 & address selection \\
Alarm functions & over-temperature alarm, interrupt, shutdown \\
\bottomrule
\end{tabularx}
\end{table}

One electrical constraint is derivable from \cref{tab:max30100,tab:max30205} and is not stated by \textcite{efendi2025}. The temperature sensor's supply range tops out at 3.3~V; the pulse-oximetry sensor's optimal range begins at 3.3~V. The two intervals intersect at exactly one point. The bus and both sensors must therefore run at 3.3~V, which is also the microcontroller's logic level --- convenient, but it places one sensor at the bottom of its stated range and the other at the top of its. A single regulated 3.3~V rail is the only configuration satisfying both, and supplying the oximeter from 5~V while the temperature sensor runs at 3.3~V would impose a level-translation requirement on the shared bus that neither device provides.

%==========================================================================
\subsection{Interconnection}\label{sec:interconnect}

\subsubsection{The bus and the address constraint}

Both sensors are I\textsuperscript{2}C and both expose SDA and SCL. They therefore share one two-wire bus, which requires distinct slave addresses.

The pin lists resolve this, and the resolution is visible in an asymmetry \textcite{efendi2025} report without remarking on it: the temperature sensor has A0, A1 and A2 address-selection pins, and the oximeter has none. The oximeter's address is fixed in silicon; the temperature sensor's is configurable.

The address-selection pins therefore exist to avoid exactly this collision, and must be strapped deliberately rather than left floating. A floating address pin on an I\textsuperscript{2}C device produces an indeterminate address and intermittent bus failure --- which is the class of fault that manifests as missing readings and would be indistinguishable, in the collected data, from a patient who did not measure. That indistinguishability matters beyond the circuit: \cref{sec:missingness} classified missingness by mechanism, and a hardware fault masquerading as patient non-adherence corrupts that classification at source.

\subsubsection{Pin assignment}\label{sec:pinout}

\Textcite{efendi2025} give pin names and no GPIO assignment. The following uses the conventional default I\textsuperscript{2}C mapping for this microcontroller. It is a design choice, not a specification taken from the source, and is presented as such.

\begin{table}[htbp]
\centering\small
\caption{Proposed GPIO assignment. This table is design, not source material.}\label{tab:pinout}
\begin{tabularx}{\textwidth}{@{}>{\raggedright\arraybackslash}p{2.8cm} >{\raggedright\arraybackslash}p{2.0cm} >{\raggedright\arraybackslash}p{2.2cm} X@{}}
\toprule
\textbf{Microcontroller} & \textbf{MAX30100} & \textbf{MAX30205} & \textbf{Function} \\
\midrule
3.3~V & VIN & VCC & supply \\
GND & GND & GND & reference \\
GPIO21 & SDA & SDA & I\textsuperscript{2}C data \\
GPIO22 & SCL & SCL & I\textsuperscript{2}C clock \\
GPIO19 (input) & INT & --- & pulse, data-ready interrupt \\
GPIO18 (input) & --- & OS & over-temperature alarm \\
--- & --- & A0, A1, A2 & strapped to GND or 3.3~V; must not float \\
\bottomrule
\end{tabularx}
\end{table}

\begin{figure}[htbp]
\centering
\begin{tikzpicture}[font=\scriptsize,
  chip/.style={draw,rounded corners=2pt,minimum height=1.5cm,align=center,fill=cGrey!8},
  wire/.style={thick},
  lab/.style={font=\tiny,text=cGrey,inner sep=1pt}]

  \node[chip,minimum width=2.6cm] (mcu) at (0,0) {\textbf{ESP32}\\\tiny MCU + radio};
  \node[chip,minimum width=2.4cm] (ox)  at (5.6,0.95) {\textbf{MAX30100}\\\tiny HR, SpO$_2$};
  \node[chip,minimum width=2.4cm] (tmp) at (5.6,-0.95) {\textbf{MAX30205}\\\tiny temperature};

  % shared I2C bus
  \draw[wire,cAfrica] (mcu.east) ++(0,0.35) -- ++(1.3,0) coordinate (bus);
  \draw[wire,cAfrica] (bus) -- ++(0,1.0) -| (ox.west);
  \draw[wire,cAfrica] (bus) -- ++(0,-1.4) -| (tmp.west);
  \node[lab,text=cAfrica,anchor=south] at (1.3,0.45) {SDA / SCL, 3.3~V bus};

  % interrupts
  \draw[wire,cRef,-{Stealth[length=1.6mm]}] (ox.south west) ++(0.2,0) |- ($(mcu.east)+(0,0.0)$);
  \node[lab,text=cRef,anchor=west] at (3.0,0.1) {INT $\to$ GPIO19};
  \draw[wire,cRef,-{Stealth[length=1.6mm]}] (tmp.north west) ++(0.2,0) |- ($(mcu.east)+(0,-0.35)$);
  \node[lab,text=cRef,anchor=west] at (3.0,-0.3) {OS $\to$ GPIO18};

  % straps
  \draw[wire,dashed,cGrey] (tmp.south) -- ++(0,-0.55) node[lab,anchor=north]{A0--A2 strapped, never floating};

  % rail
  \node[draw,fill=cRef!12,minimum width=1.6cm,minimum height=0.6cm] (reg) at (0,-2.2) {3.3~V reg.};
  \draw[wire] (reg.north) -- (mcu.south);
  \draw[wire] (reg.east) -| ($(tmp.south)+(-0.8,0)$);
  \node[lab,anchor=west] at (1.0,-2.2) {single rail: the only voltage both sensors admit (\cref{sec:components})};
\end{tikzpicture}
\caption{Interconnection. One shared I\textsuperscript{2}C bus on a single 3.3~V rail, with both sensor interrupt lines wired so that the processor need not poll.}\label{fig:interconnect}
\end{figure}

Two notes on the interrupt lines. Using INT rather than polling the oximeter is what makes the duty-cycled operation of \cref{sec:power} possible: the processor sleeps until the sensor signals data ready, instead of waking on a timer and polling. Using OS as a hardware over-temperature alarm yields a threshold crossing that does not require the processor to be awake at all, which is the hardware analogue of the red-tier escalation of \cref{sec:transmission-priority}.

\subsubsection{Datasheet dependencies}

The source does not supply the following, which must be obtained from manufacturer documentation.

\begin{table}[htbp]
\centering\small
\caption{Values required from datasheets, not available in the corpus.}\label{tab:datasheet-deps}
\begin{tabularx}{\textwidth}{@{}>{\raggedright\arraybackslash}p{5.6cm} X@{}}
\toprule
\textbf{Missing} & \textbf{Why it matters} \\
\midrule
I\textsuperscript{2}C base addresses of both devices & address collision; A0--A2 strapping \\
Pull-up resistor values for SDA and SCL & bus rise time; reliability at the chosen bus speed \\
Oximeter register configuration & LED current, pulse width, sample rate, averaging \\
Oximeter supply and LED currents & the power budget of \cref{sec:power} \\
Microcontroller active and deep-sleep currents & the power budget of \cref{sec:power} \\
Whether the breakout board carries its own regulator and pull-ups & may conflict with the single-rail decision of \cref{sec:components} \\
\bottomrule
\end{tabularx}
\end{table}

These are stated as datasheet dependencies rather than filled with invented values. A fabricated resistor value in a thesis is worse than an acknowledged gap, because it will be tested by anyone who builds the circuit.%
\verify{The datasheets for all three components are publicly available from the manufacturer and are not in the present corpus. Obtain them and replace the illustrative currents of \cref{tab:power-regimes} with measured or specified values before the power budget is quoted as a result rather than as an order-of-magnitude argument.}

%==========================================================================
\subsection{Power and Duty Cycle}\label{sec:power}

This subsection contains the derivation that determines whether the subsystem is deployable, and it follows from a limitation \textcite{efendi2025} state but do not act upon.

\subsubsection{The limitation they identify}

They write that \bq{one of the main limitations is the dependence on battery power, which can limit the duration of continuous monitoring}. Their own protocol explains why: data collected every second for ten minutes, producing 600 records per participant, with the radio and both sensors continuously active.

\subsubsection{The arithmetic}\label{sec:power-arithmetic}

With duty cycle $d$, active current $I_{\text{act}}$, sleep current $I_{\text{slp}}$ and battery capacity $C$:
\begin{equation}
  d = \frac{t_{\text{act}}}{T}, \qquad
  I_{\text{avg}} = d\,I_{\text{act}} + (1-d)\,I_{\text{slp}}, \qquad
  \text{life} = \frac{C}{I_{\text{avg}}}.
\end{equation}

Under the protocol adopted in \cref{sec:streams}, two measurement occasions per week, with a 60-second active window covering sensor settling, acquisition and transmission:
\begin{equation}
  T = 3.5~\text{days} = 302\,400~\text{s},
  \qquad d = \frac{60}{302\,400} = 1.98 \times 10^{-4},
\end{equation}
a duty cycle of 0.020\%.

\begin{table}[htbp]
\centering\small
\caption{Power regimes. The currents are illustrative, of the right order for this class of device, and are to be replaced by datasheet values.}\label{tab:power-regimes}
\begin{tabularx}{\textwidth}{@{}X c c@{}}
\toprule
\textbf{Regime} & \textbf{$I_{\text{avg}}$} & \textbf{Life on 2\,000~mAh} \\
\midrule
Continuous, as \textcite{efendi2025} operate & 120~mA & 17~h \\
Duty-cycled at $d = 1.98\times10^{-4}$ & 0.0338~mA & $\approx 6.7$~years \\
\bottomrule
\end{tabularx}
\end{table}

With $I_{\text{act}} = 120$~mA and $I_{\text{slp}} = 10$~\textmu A, the duty-cycled average is
\begin{equation}
  I_{\text{avg}} = \underbrace{1.98\times10^{-4} \times 120}_{0.0238~\text{mA, active}}
  + \underbrace{0.0100}_{\text{mA, sleep}} = 0.0338~\text{mA}.
\end{equation}

The improvement is a factor of $120/0.0338 = 3.5 \times 10^3$: between three and four orders of magnitude, obtained not by better hardware but by matching the sampling regime to the clinical protocol. It is worth noting the ceiling: even with a sleep current of exactly zero the factor would be $1/d = 5.0\times10^3$, so no amount of sleep-current optimisation reaches four orders at this duty cycle.

\subsubsection{Which term dominates, and the design rule that follows}\label{sec:power-crossover}

The decomposition above shows that the active term contributes 70\% of the average current and the sleep term 30\%. The average is therefore governed principally by the active window, not by the sleep floor --- which inverts the intuition that a duty cycle of two parts in ten thousand must make the active phase irrelevant.

The crossover is derivable. The two terms are equal when $d\,I_{\text{act}} = I_{\text{slp}}$, that is at an active window of
\begin{equation}
  t^\star = \frac{I_{\text{slp}}\,T}{I_{\text{act}}}
  = \frac{0.010 \times 302\,400}{120} \approx 25~\text{s}.
\end{equation}

The design rule follows directly and is more useful than the lifetime figure. Above a 25-second active window, shortening the window is the effective lever and reducing sleep current buys little; below it, the sleep floor dominates and further shortening buys nothing (\cref{fig:power}). Since 60 seconds is above the crossover, the first optimisation to attempt is a shorter acquisition --- faster sensor settling, a pre-established network association, a smaller payload --- and not a lower-leakage sleep mode.

\begin{figure}[htbp]
\centering
\begin{tikzpicture}
\begin{axis}[width=0.72\textwidth,height=5.4cm,
  xlabel={Active window $t_{\text{act}}$ (s)},
  ylabel={$I_{\text{avg}}$ (mA)},
  xmin=0,xmax=180,ymin=0,ymax=0.08,
  grid=major,grid style={gray!20},
  legend pos=north west,legend style={font=\scriptsize},
  tick label style={font=\scriptsize},label style={font=\scriptsize}]
\addplot[thick,cAfrica,domain=0:180,samples=100]{0.010 + x*120/302400};
\addplot[dashed,cGrey,domain=0:180,samples=2]{0.010};
\addplot[only marks,mark=*,cRef] coordinates {(25.2,0.020)};
\addplot[only marks,mark=square*,black] coordinates {(60,0.0338)};
\legend{$I_{\text{avg}}(t_{\text{act}})$,sleep floor 10~\textmu A,crossover $t^\star\approx25$~s,adopted window 60~s}
\end{axis}
\end{tikzpicture}
\caption{Average current against active window. The adopted 60-second window lies above the crossover, so the active term dominates and shortening the window --- not lowering sleep current --- is the effective optimisation.}\label{fig:power}
\end{figure}

\subsubsection{The same finding as the alert-volume argument}

\Cref{sec:volume-arithmetic} established that one-hertz sampling makes alerting arithmetically unusable: at a specificity of 0.98, continuous streaming generates of the order of twelve thousand false alerts per patient per week, against approximately 0.06 under a three-readings-per-week protocol.

The battery limitation and the alert-volume impossibility are the same error viewed from two directions. Sampling continuously at a frequency the clinical protocol does not require exhausts the battery and floods the alerting layer. Duty-cycling to the protocol resolves both at once, and \cref{sec:mqtt-payload} shows it resolves a third problem as well. This convergence is stated explicitly because it is the kind of argument that demonstrates the architecture was reasoned rather than assembled.

\subsubsection{Implementation}

The RTC slow SRAM, 8~kB retained through deep sleep, holds the state needed across cycles: last transmission timestamp, retry counter, queued readings not yet acknowledged, and the tier of the last evaluation. \Textcite{efendi2025} report that this memory exists; using it is what makes deep sleep practical, since without retained state every wake is a cold start.

The wake schedule should be event-driven where possible and timed where not: wake on the cuff's BLE advertisement, so that the covariate channel samples when the clinical measurement happens rather than on an independent timer firing at unrelated moments. This also guarantees temporal alignment between the covariate and the parameter it contextualises, which \cref{sec:irregular} identified as a problem when channels are sampled on independent schedules.

%==========================================================================
\subsection{MQTT Client Design}\label{sec:mqtt-design}

\subsubsection{What the source specifies}

\Textcite{efendi2025} establish the split: MQTT upstream from device to gateway, HTTP downstream to clients, with the gateway translating into cloud functions or the realtime database. Their justification is that MQTT \bq{enables resource-constrained devices to send sensor data to a gateway}. \Cref{sec:protocols-cloud} accepted the choice and recorded that no source in this corpus quantifies the advantage: no bandwidth, latency, power or reliability comparison against HTTP is reported anywhere in the corpus.

The elements below are design specification, not source material, and are presented as such.

\subsubsection{Topic structure}

\begin{center}
\texttt{telemon/v1/<site>/<patient\_pseudonym>/<channel>/<measurement\_id>}
\end{center}

Four properties. The version prefix permits schema migration without breaking deployed devices. The pseudonym, rather than a name or record number, satisfies \cref{sec:reident} at the transport layer, which matters because topic names appear in broker logs. The channel segment allows subscription filtering, so that the scoring service subscribes to clinical channels and ignores covariates. And the hierarchy supports wildcard subscription for monitoring.

\subsubsection{Quality of service, and why exactly-once is the wrong choice for red}\label{sec:qos}

MQTT defines three delivery guarantees, and \cref{sec:transmission-priority} established that transmission urgency should follow the alert tier. The natural mapping and the one adopted differ, and the difference is worth the paragraph.

\begin{table}[htbp]
\centering\small
\caption{Quality of service by tier.}\label{tab:qos}
\begin{tabularx}{\textwidth}{@{}>{\raggedright\arraybackslash}p{3.2cm} c >{\raggedright\arraybackslash}p{2.6cm} X@{}}
\toprule
\textbf{Tier} & \textbf{QoS} & \textbf{Guarantee} & \textbf{Rationale} \\
\midrule
Green, routine covariate & 0 & at most once & loss of one heart-rate sample is immaterial \\
Amber, clinical reading & 1 & at least once & must arrive; duplicates removed by \texttt{measurement\_id} \\
Red, critical & 1 & at least once & see below; the real guarantee is the SMS fallback \\
\bottomrule
\end{tabularx}
\end{table}

Assigning QoS~2 to the red tier is the obvious choice and is the wrong one. Exactly-once delivery is obtained by a four-step handshake, two round trips more than at-least-once, and the links on which red-tier delay is actually harmful are precisely the degraded links where those extra round trips are slowest and most likely to fail. Meanwhile the property QoS~2 provides, suppression of duplicates, is already provided at the application layer: \cref{sec:record-schema} makes \texttt{measurement\_id} a UUID, the repository deduplicates on it, and \cref{sec:ingestion} flags duplicates rather than silently discarding them. QoS~2 therefore purchases, at the cost of latency, a guarantee the system already holds.

The resolution is QoS~1 with application-layer idempotency for both amber and red, with urgency expressed through retry policy and the SMS fallback of \cref{sec:transmission-priority} rather than through the delivery mode. This is a correction to the intuitive mapping and is stated as such.

\subsubsection{Last Will and Testament, and the non-transmission alarm}

MQTT permits a client to register, at connection time, a message the broker will publish on its behalf should the connection drop without a clean disconnect.

This is the transport-layer implementation of the requirement that silence be alarmed. \Cref{sec:decision-flow} established, from the repeated SMS reminder of \textcite{yazdani2018} and the HEARTS retention indicator, that non-transmission must itself trigger an action; \cref{tab:failure-modes} placed battery exhaustion, device loss and handset unavailability in the failure table with the non-transmission alarm as the detection mechanism. The Last Will supplies that signal directly: the broker announces the disconnection, and the Cloud escalates according to how long the device has been absent relative to its configured checkup period, which is the configurable parameter of \textcite{yazdani2018}.

The distinction between a graceful disconnect, the device sleeping as designed, and an ungraceful one is exactly the distinction between normal duty cycling and failure, and the protocol supplies it at no cost. It is also the one detection that \cref{sec:failure} identified as impossible at the Edge --- a failed device cannot report its own failure --- so locating it in the broker rather than in either endpoint is not merely convenient but necessary.

\subsubsection{Payload and the data-cost derivation}\label{sec:mqtt-payload}

The payload is JSON, per \cref{sec:schema}, carrying the fields of \cref{sec:record-schema}. A record is of the order of 300 bytes. A measurement occasion stores three records: the retained blood-pressure reading, the glycaemia reading, and the covariate bundle --- approximately 900 bytes per occasion.

\begin{table}[htbp]
\centering\small
\caption{Transmitted volume by regime, at 300 bytes per record.}\label{tab:data-volume}
\begin{tabularx}{\textwidth}{@{}X c c@{}}
\toprule
\textbf{Regime} & \textbf{Per patient, per month} & \textbf{30 patients, per month} \\
\midrule
Continuous 1~Hz, as \textcite{efendi2025} operate & 778~MB & 23~GB \\
Duty-cycled, 2 occasions per week & 7.7~kB & 231~kB \\
\bottomrule
\end{tabularx}
\end{table}

The derivation, stated so that the two columns can be checked against each other. Continuous: $86\,400 \times 300~\text{B} = 25.9$~MB per patient per day, $\times 30 = 778$~MB per month, $\times 30$ patients $= 23.3$~GB. Duty-cycled: $2 \times 900~\text{B} = 1.8$~kB per week, $\times 4.29 = 7.7$~kB per patient per month, $\times 30 = 231$~kB. The ratio is $1.0 \times 10^5$: five orders of magnitude.

On Cameroonian mobile data tariffs this is the difference between a recurring cost borne by the patient and a cost that rounds to zero. This is the third convergence on the same design decision, and the one tied most directly to \cref{ch:background}: \cref{sec:cloud-gaps} recorded that no source in the corpus states what any of these architectures costs to operate, and \cref{sec:regional-landscape} documented a service ceasing operations, for financial reasons, after reaching approximately 2.5 million registered users. A system consuming 23~GB per month for thirty patients does not scale; one consuming 231~kB does. Duty-cycling is a sustainability decision as much as a power decision.

\subsubsection{Security}

Per segment~3 of \cref{tab:security-segments}: MQTT over TLS, client authentication by per-device credential, and broker-side access control restricting each device to its own topic subtree. \Textcite{efendi2025} report that the microcontroller supports IEEE 802.11 with WPA and WPA2, which protects the wireless link but not the application-layer connection; the two are distinct and both are required. Reporting the former as though it addressed the latter is a smaller instance of the pattern identified in \cref{sec:security-gap}.

%==========================================================================
\subsection{The Camera-PPG Alternative}\label{sec:camera-ppg}

\subsubsection{What it may and may not measure}

\Cref{sec:cuffless} established the position across four sources, and it is not reopened here.

\begin{table}[htbp]
\centering\small
\caption{Admissibility of camera photoplethysmography by parameter.}\label{tab:camera-admissibility}
\begin{tabularx}{\textwidth}{@{}>{\raggedright\arraybackslash}p{3.4cm} >{\raggedright\arraybackslash}p{2.6cm} X@{}}
\toprule
\textbf{Parameter} & \textbf{Camera PPG} & \textbf{Basis} \\
\midrule
Heart rate & admissible & periodicity requires no absolute calibration (\cref{sec:ppg}) \\
Oxygen saturation & not admissible & requires controlled dual-wavelength illumination and empirical calibration (\cref{sec:ppg}) \\
Blood pressure & not admissible & \textcite{who2023hypertension}; \textcite{alessa2019}: \bq{highly inaccurate and unfeasible} \\
\bottomrule
\end{tabularx}
\end{table}

The camera alternative is therefore a substitute for one function of the instrumented subsystem, heart rate, and for neither clinical parameter. Contribution~C4 is defensible on that basis and on no wider one.

\subsubsection{Signal chain}\label{sec:ppg-chain}

Proposed design; no source in this corpus validates it.

\begin{enumerate}[nosep,label=\arabic*.]
  \item \textbf{Acquisition.} Fingertip over the rear camera, LED flash as illumination, 30~fps, fixed exposure and white balance. Automatic exposure compensation is the dominant artefact source and must be disabled, since the control loop responds to the very intensity variation that constitutes the signal.
  \item \textbf{Reduction.} Mean intensity of the red channel over a central region of interest, one value per frame, giving a 30~Hz series. Nyquist is 15~Hz, or 900~beats per minute, far above the physiological band.
  \item \textbf{Detrend.} Remove the slow baseline, the DC term of \cref{sec:ppg}, by moving-average subtraction or first differencing.
  \item \textbf{Bandpass.} 0.7--3.5~Hz, corresponding to 42--210~beats per minute, which covers the physiological range with margin at both ends.
  \item \textbf{Rate estimation.} Dominant frequency by autocorrelation or spectral peak, cross-checked against inter-peak intervals in the time domain. Agreement between the two estimators is itself a quality signal.
  \item \textbf{Aggregation.} Median over a 30-second acquisition, with the interquartile range retained as an uncertainty measure.
\end{enumerate}

The passband is chosen on physiological grounds rather than from the reference ranges of \textcite{efendi2025}, because \cref{sec:efendi-appraisal} established that those ranges contradict one another between two of their tables and cannot be used to set a filter boundary.

\subsubsection{Quality gating}

A camera reading that is not gated is worse than no reading, because it enters the series indistinguishable from a good one. Three gates, each cheap:

\begin{table}[htbp]
\centering\small
\caption{Quality gates for camera acquisition.}\label{tab:ppg-gates}
\begin{tabularx}{\textwidth}{@{}>{\raggedright\arraybackslash}p{2.4cm} X >{\raggedright\arraybackslash}p{3.4cm}@{}}
\toprule
\textbf{Gate} & \textbf{Test} & \textbf{On failure} \\
\midrule
Contact & DC level within an expected band & reject; prompt re-acquisition \\
Motion & frame-to-frame variance above threshold & reject \\
Periodicity & spectral peak prominence below threshold & reject \\
\bottomrule
\end{tabularx}
\end{table}

A rejected acquisition sets the mask of \cref{sec:masking} rather than producing a value, which is the behaviour \cref{sec:missingness} requires: absence is reported as absence. The rejection rate is itself reported, per \cref{sec:ppg-validation}.

\subsubsection{The validation study that comes free}\label{sec:ppg-validation}

\Cref{sec:device-validation} required Bland--Altman comparison against a validated reference. For heart rate that reference already exists within the system at no additional cost: an automated oscillometric cuff reports pulse rate alongside blood pressure, and under the protocol of \cref{sec:streams} a cuff reading is taken twice weekly.

Pairing a camera acquisition with every cuff reading therefore yields, over thirty patients across a collection span of approximately five weeks, of the order of $30 \times 2 \times 5 = 300$ paired observations.

One methodological caution attaches to that count and must be respected in the analysis. Those 300 observations come from thirty patients, approximately ten per patient, so they are not independent: a systematic per-patient offset, from skin tone, perfusion or technique, inflates agreement if the repetitions are treated as independent samples. The limits of agreement must therefore be computed by the variant of the Bland--Altman method for repeated measurements, which partitions within-subject from between-subject variance, and the effective sample size for the between-subject component is thirty, not 300.

\begin{table}[htbp]
\centering\small
\caption{What the method-comparison sub-study reports.}\label{tab:ppg-reporting}
\begin{tabularx}{\textwidth}{@{}>{\raggedright\arraybackslash}p{6.0cm} X@{}}
\toprule
\textbf{Statistic} & \textbf{Why} \\
\midrule
Bland--Altman bias and limits of agreement, repeated-measures variant & separates systematic offset from random disagreement, without treating repetitions as independent \\
Proportion within $\pm5$ and $\pm10$~bpm & clinically interpretable \\
Rejection rate by gate & measures usability, not only accuracy \\
Between- and within-subject variance components & identifies whether disagreement is per-patient or per-reading \\
Not MAPE alone & \cref{sec:efendi-appraisal}: averages away the structure of the disagreement, and its denominator is ambiguous \\
\bottomrule
\end{tabularx}
\end{table}

This converts the weakest element of contribution~C4, an unvalidated acquisition method, into a self-contained method-comparison sub-study with a cost-free reference. It is the strongest available answer to the objection that camera photoplethysmography is unproven, and it must be planned into the protocol rather than attempted afterwards --- because the pairing requires the two acquisitions to be simultaneous, which cannot be reconstructed retrospectively.

%==========================================================================
\subsection{The Three Options, Compared}\label{sec:options}

\Cref{sec:acq-decision} adopted Position~A and left the hardware question formally open. The analysis above closes it on cost and risk rather than on preference.

\begin{table}[htbp]
\centering\footnotesize
\caption{Acquisition options compared.}\label{tab:options}
\begin{tabularx}{\textwidth}{@{}>{\raggedright\arraybackslash}p{3.0cm} >{\raggedright\arraybackslash}p{3.0cm} >{\raggedright\arraybackslash}p{3.2cm} X@{}}
\toprule
 & \textbf{A: spare} & \textbf{B: instrumented} & \textbf{C: staged hybrid} \\
\midrule
BP and glucose & validated BLE devices & validated BLE devices & validated BLE devices \\
Heart rate & cuff pulse, camera PPG & MAX30100 & cuff pulse, camera PPG \\
SpO$_2$, temperature & --- & MAX30100, MAX30205 & bench only \\
Hardware build & none & full, 30 units & prototype, not deployed \\
Device cost per patient & zero above the cuff & MCU, sensors, battery, enclosure & zero deployed \\
Validation burden & camera against cuff, free & Bland--Altman, both sensors, under movement & camera only \\
Power design required & no & yes, \cref{sec:power} & bench only \\
C4 demonstrated by & camera PPG, manual entry & the hardware & both, staged \\
Clinical risk on failure & none, covariates only & none, covariates only & none \\
\bottomrule
\end{tabularx}
\end{table}

Option~C is adopted. Its relationship to the earlier decision should be stated plainly, because it looks like a reversal and is not: the \emph{deployed} configuration under C is identical to Position~A, and C differs only in adding a bench-evaluated prototype that no patient receives. The decision of \cref{sec:acq-decision} is therefore unchanged for the cohort; what C adds is an artefact and a set of measurements, not a change to what patients use.

The reasoning, in the form it should take in the defence:

\begin{quote}
The instrumented subsystem is specified and prototyped in order to demonstrate the hardware design and the duty-cycled MQTT transport, and is evaluated on the bench. The deployed cohort uses spare acquisition --- validated Bluetooth cuff and glucometer, with camera photoplethysmography for heart rate --- because the covariate channels the hardware would add do not enter the clinical decision, because the per-patient device cost would be incurred for no clinical gain, and because the camera alternative admits a cost-free method-comparison study against the cuff's own pulse reading. The hardware path is documented as the extension route for a deployment in which continuous covariate monitoring is clinically indicated.
\end{quote}

That position preserves contribution~C4, preserves the hardware work as a contribution, and does not require thirty devices to be built and maintained by one researcher.

%==========================================================================
\subsection{Specification}\label{sec:spec42}

\begin{table}[htbp]
\centering\footnotesize
\caption{Requirements established in \cref{sec:iot}.}\label{tab:spec42}
\begin{tabularx}{\textwidth}{@{}X >{\raggedright\arraybackslash}p{5.0cm}@{}}
\toprule
\textbf{Requirement} & \textbf{Basis} \\
\midrule
Subsystem acquires covariates only; scoping statement opens the section & \cref{sec:device-landscape,sec:iot-scope} \\
Redundancy of camera heart rate with cuff pulse acknowledged & \cref{sec:measures} \\
Single regulated 3.3~V rail for both sensors and the bus & \cref{sec:components}, derived from the two supply ranges \\
Shared I\textsuperscript{2}C bus; A0--A2 strapped, never floating & \cref{sec:interconnect} \\
INT and OS used as interrupts, not polled & \cref{sec:pinout} \\
GPIO assignment presented as design, not as sourced & \cref{sec:pinout} \\
Datasheet dependencies listed; values not invented & \cref{tab:datasheet-deps} \\
Duty-cycled operation matched to the clinical protocol & \cref{sec:power-arithmetic} \\
Active window identified as the dominant term; $t^\star \approx 25$~s & \cref{sec:power-crossover} \\
RTC slow SRAM used for cross-cycle state & \textcite{efendi2025}; \cref{sec:power} \\
Wake on the cuff's BLE advertisement, not an independent timer & \cref{sec:power} \\
Power budget recomputed with datasheet currents before quoting & \cref{tab:power-regimes} \\
MQTT topic versioned; pseudonym only, never an identifier & \cref{sec:reident,sec:mqtt-design} \\
QoS 0 for green, QoS 1 for amber and red; QoS 2 rejected with reasons & \cref{sec:qos} \\
Idempotency on \texttt{measurement\_id} & \cref{sec:record-schema} \\
Last Will registered; drives the non-transmission alarm & \cref{sec:decision-flow,sec:failure} \\
JSON payload; data volume computed, with the derivation shown & \cref{sec:mqtt-payload} \\
MQTT over TLS; per-device credential; topic-scoped broker ACL & \cref{tab:security-segments} \\
Camera PPG restricted to heart rate only & \cref{sec:cuffless} \\
Automatic exposure disabled during acquisition & \cref{sec:ppg-chain} \\
Passband set physiologically, not from \textcite{efendi2025} ranges & \cref{sec:ppg-chain,sec:efendi-appraisal} \\
Three quality gates; rejection sets the mask, not a value & \cref{sec:masking,sec:camera-ppg} \\
Camera PPG validated against cuff pulse, repeated-measures Bland--Altman & \cref{sec:ppg-validation} \\
Rejection rate and variance components reported alongside agreement & \cref{sec:ppg-validation} \\
Option~C adopted; its identity with Position~A when deployed stated & \cref{sec:options,sec:acq-decision} \\
\bottomrule
\end{tabularx}
\end{table}

%==========================================================================
\subsection{Limits of What This Section Establishes}\label{sec:iot-limits}

Three elements of this section are design proposals rather than established results, and are labelled as such wherever they are used.

\textbf{The circuit is incomplete.} \Textcite{efendi2025} give pin names and electrical characteristics; they give no wiring diagram, no GPIO assignment, no pull-up values, no decoupling, no power-supply design and no register configuration. \Cref{tab:pinout} and \cref{fig:interconnect} are design proposals built from their pin lists together with the conventional default I\textsuperscript{2}C mapping, and the remaining values come from datasheets outside this corpus (\cref{tab:datasheet-deps}).

\textbf{Camera photoplethysmography is unvalidated within this corpus.} No source provides accuracy, precision or failure rate for camera-based heart rate. The signal chain of \cref{sec:ppg-chain} is a proposal; \cref{sec:ppg-validation} is how the evidence is to be obtained; and until that study is run the claim is stated as untested rather than as supported.

\textbf{The protocol choice is unquantified.} \Cref{sec:protocols-cloud} recorded that \textcite{efendi2025} assert the advantage of MQTT over HTTP and measure nothing, and that position is unchanged. No efficiency benefit is claimed in this work on the strength of the corpus.%
\verify{A measurement of the MQTT/HTTP comparison is within reach and would convert an inherited assertion into an original result: both protocols are available on the same hardware, the comparison is bytes on the wire, round trips and energy per transmitted record, and the measurement is a few hours of bench work. If it is run, it belongs in \cref{ch:results} and this paragraph should be rewritten to cite it.}

Two document dependencies remain outstanding from \cref{ch:data} and are unchanged: the guideline editions required by \cref{tab:assumptions}, and the Cameroonian legal instrument required by \cref{sec:applicable-law}.

%##########################################################################
% SECTION 4.3
%##########################################################################
\section{Cloud Backend Infrastructure and Real-Time Management}\label{sec:backend}

\subsection*{Purpose and sources}

\Cref{sec:cloud-layer} specified the Cloud layer's units; \cref{sec:iot} specified what arrives at it. This section specifies the backend itself: how the record is stored, who may read and write it, what fires when data arrive, and how changes reach the clinician.

Two sources are used. \Textcite{efendi2025} supply the only implemented backend in the corpus, including a published live schema, the protocol split and the offline synchronisation behaviour; the reservations of \cref{sec:efendi-appraisal} stand, and their schema in particular is used as a counter-example rather than as a template. \Textcite{yazdani2018} supply the functional decomposition --- five modules over a three-layer structure --- which \cref{sec:modules} shows maps cleanly onto a serverless realisation despite having been designed for a client--server architecture.

%==========================================================================
\subsection{The Property That Governs Everything}\label{sec:tree-property}

\Textcite{efendi2025} describe their platform as \bq{a cloud-hosted database provided by Google} supporting \bq{various protocols including HTTP and MQTT}, in which \bq{data transmitted by IoT devices will be stored as JSON which is then accessed by users in real-time}, and enumerate four advantages: integrated features, zero maintenance, connection to the wider cloud platform, and code privacy.

One architectural property underlies all of it and determines the design: the realtime database is a single JSON tree, and a read of any node returns that node together with everything beneath it.

Three consequences follow, and the schema of \cref{sec:schema-design} exists to satisfy them.

\textbf{Depth and fan-out are performance decisions, not organisational ones.} A node that accumulates children indefinitely becomes unreadable, because every read of its parent downloads the whole accumulation. In a longitudinal system this is not a corner case: history grows monotonically by design.

\textbf{Query capability is deliberately narrow.} The tree supports ordering by key, by child value or by value, combined with range and limit operators --- but one ordering per query and no compound conditions. Anything the system must query by, the key must encode.

\textbf{Denormalisation is mandatory.} Data are duplicated into whatever shapes the readers require, and consistency is maintained at write time rather than by joining at read time.

%==========================================================================
\subsection{Why the Published Schema Cannot Be Used}\label{sec:schema-counterexample}

\Textcite{efendi2025} publish their live database, and it is worth reproducing because it states the problem precisely:

\begin{center}\small
\begin{tabular}{@{}l@{}}
\texttt{<endpoint>.asia-southeast1.firebasedatabase.app/} \\
\quad\texttt{|-- heart-rate: 91} \\
\quad\texttt{|-- oxygen-saturation: 97} \\
\quad\texttt{`-- temperature: 37} \\
\end{tabular}
\end{center}

Three keys at the root, each holding a scalar.

This structure stores the present and discards the past. \Cref{sec:schema} established the objection and it is not reopened: \textcite{cadrage2026} specify rolling aggregates, trends, variability and threshold deviations, and \cref{sec:feature-eng} made the fourteen-day window the unit of feature construction. A structure holding one value per channel cannot answer a window query, because the window does not exist within it.

Two further defects follow from the same shape. There is no patient dimension: the tree is single-tenant, and admitting a second patient would require restructuring every path. And there is no distinction between measurement time and record time, which \cref{sec:metrology,sec:offline} established as mandatory once offline synchronisation decouples them.

Their own system exhibits the consequence. They report that the heart-rate graph \bq{will update every 5 min}, and their application displays a current value with a short live trace. That is all the schema supports. It is adequate for a live monitor and inadequate for a longitudinal record, which is what this system is. The distinction is worth stating in these terms rather than as a criticism: their schema is correct for their stated purpose, and the purposes differ.

%==========================================================================
\subsection{The Schema}\label{sec:schema-design}

\subsubsection{Design rules}

\begin{table}[htbp]
\centering\small
\caption{Schema design rules and their justification.}\label{tab:schema-rules}
\begin{tabularx}{\textwidth}{@{}>{\raggedright\arraybackslash}p{6.2cm} X@{}}
\toprule
\textbf{Rule} & \textbf{Reason} \\
\midrule
A separate top-level tree per access pattern & a read returns the whole subtree (\cref{sec:tree-property}) \\
The key encodes what is queried by & one ordering per query, no compound conditions \\
Time series keyed by measurement timestamp, at millisecond resolution & range query by key; window construction (\cref{sec:feature-eng}) \\
No unbounded fan-out under a node that is read whole & read cost grows with history \\
Duplicate freely; maintain consistency at write time & denormalisation is the model \\
Composite keys zero-padded to fixed width & keys sort lexicographically (\cref{sec:key-design}) \\
Nothing identifying anywhere in the tree & \cref{sec:reident} \\
\bottomrule
\end{tabularx}
\end{table}

\subsubsection{The seven trees}

\begin{figure}[htbp]
\centering
\begin{tikzpicture}[font=\scriptsize,
  t/.style={draw,rounded corners=2pt,fill=cRef!12,minimum height=0.62cm,align=left,inner sep=3pt},
  n/.style={draw,rounded corners=2pt,fill=cAfrica!12,minimum height=0.62cm,align=left,inner sep=3pt},
  lab/.style={font=\tiny,text=cGrey,anchor=west}]
  \node[font=\footnotesize\bfseries] (root) at (0,0.9) {\texttt{/} (single JSON tree)};
  \node[t,text width=6.4cm] (m)  at (0,0.15) {\texttt{/measurements/\{pid\}/\{channel\}/\{ts\_ms\}}};
  \node[t,text width=6.4cm] (p)  at (0,-0.62) {\texttt{/patients/\{pid\}}};
  \node[n,text width=6.4cm] (q)  at (0,-1.39) {\texttt{/queue/\{clinician\}/\{pri\}\_\{ts\}\_\{eid\}}};
  \node[n,text width=6.4cm] (e)  at (0,-2.16) {\texttt{/events/\{event\_id\}}};
  \node[t,text width=6.4cm] (d)  at (0,-2.93) {\texttt{/devices/\{device\_id\}}};
  \node[n,text width=6.4cm] (ct) at (0,-3.70) {\texttt{/config/thresholds/\{version\}}};
  \node[n,text width=6.4cm] (cm) at (0,-4.47) {\texttt{/config/models/\{version\}}};
  \node[lab] at (3.4,0.15)  {range query by key $\to$ the 14-day window};
  \node[lab] at (3.4,-0.62) {per-patient parameters, one read};
  \node[lab,text=cAfrica] at (3.4,-1.39) {ordered read, no sort, no index};
  \node[lab,text=cAfrica] at (3.4,-2.16) {the adjudication table};
  \node[lab] at (3.4,-2.93) {heartbeat, validation status};
  \node[lab,text=cAfrica] at (3.4,-3.70) {versioned rule table};
  \node[lab,text=cAfrica] at (3.4,-4.47) {model registry};
  \node[lab,text=cAfrica,anchor=west] at (-3.1,-5.25) {red: trees absent from \textcite{efendi2025}, required by \cref{ch:data,sec:architecture}};
\end{tikzpicture}
\caption{The seven top-level trees. Each exists to serve one access pattern, because a read returns an entire subtree.}\label{fig:trees}
\end{figure}

\begin{table}[htbp]
\centering\footnotesize
\caption{Tree contents.}\label{tab:tree-contents}
\begin{tabularx}{\textwidth}{@{}>{\raggedright\arraybackslash\ttfamily\scriptsize}p{4.0cm} X@{}}
\toprule
\normalfont\textbf{Path} & \textbf{Fields} \\
\midrule
/measurements/\{pid\}/ \{channel\}/\{ts\_ms\} &
\texttt{value, unit, ts\_recorded, reading\_index, is\_retained, context, device\_id, operator, entry\_mode, protocol\_profile, quality\_flag, measurement\_id, supersedes} \\
\addlinespace
/patients/\{pid\} &
\texttt{age\_band, sex, bmi, dx\_htn\_year, dx\_dm\_year, dm\_type, comorbidities[], regimen[], regimen\_version, enrolment\_date\_shifted, checkup\_period\_days, assigned\_clinician, threshold\_version, last\_seen\_ts} \\
\addlinespace
/queue/\{clinician\_id\}/ \{pri\}\_\{ts\}\_\{eid\} &
\texttt{pid, tier, score, trigger\_source, summary} \\
\addlinespace
/events/\{event\_id\} &
\texttt{pid, ts\_triggered, trigger\_source, tier, score, model\_version, threshold\_version, ts\_reviewed, clinician\_verdict, action\_taken, comment} \\
\addlinespace
/devices/\{device\_id\} &
\texttt{make, model, serial\_hash, validated, calibration\_date, assigned\_pid, last\_heartbeat\_ts} \\
\addlinespace
/config/thresholds/ \{version\} &
\texttt{channel, context, band\_label, lower, upper, unit, source, source\_version, effective\_from} \\
\addlinespace
/config/models/\{version\} &
\texttt{algorithm, trained\_at, feature\_set\_hash, calibration\_params, tau\_amber, tau\_red, metrics} \\
\bottomrule
\end{tabularx}
\end{table}

Four of these trees do not appear in \textcite{efendi2025} and are required by earlier chapters. \texttt{/events} is the adjudication table of \cref{sec:adjudication}, the only route to the clinician-adjudicated labels of \cref{sec:two-labels}. \texttt{/config/thresholds} is the versioned rule table of \cref{sec:interop,sec:leak-boundary}. \texttt{/config/models} is the registry of \cref{sec:cloud-layer}, without which a score cannot be attributed to the model that produced it. And \texttt{/queue} is the priority ordering of \textcite{yazdani2018}, which \cref{sec:key-design} shows the key design delivers directly.

\Textcite{yazdani2018} specify that the database holds \bq{physiological data\ldots along with the physician's prescriptions, diagnosis and recommendations of decision support system}, and that \bq{all events in this telehealth monitoring architecture, including physician and CDSS prescriptions, are stored as patient record}. The \texttt{/events} tree is that requirement, with the addition that each entry carries the model and threshold versions that produced it --- an addition \cref{sec:assumptions} makes necessary, since scores computed under provisional thresholds must remain identifiable once the definitive table is installed.

\subsubsection{The key design, and a correction to the composite key}\label{sec:key-design}

\textbf{The window query.} With the measurement timestamp as the key, a fourteen-day window is a key range:

\begin{center}\small
\texttt{/measurements/\{pid\}/sbp .orderByKey() .startAt(t\_minus\_14d) .endAt(t)}
\end{center}

This returns exactly the window of \cref{sec:feature-eng} and nothing else. Under the schema of \cref{sec:schema-counterexample} it is not expressible at all; under a schema keyed by an opaque push identifier it would require ordering by a child value, which is slower and requires an index.

\textbf{The queue query.} The intention is that the key be composed so that lexicographic ordering yields a list sorted by priority and then by age within priority, making the read

\begin{center}\small
\texttt{/queue/\{clinician\_id\} .orderByKey() .limitToFirst(20)}
\end{center}

This works, but only under a constraint the draft design omits, and omitting it produces a queue that is silently mis-ordered rather than one that fails visibly.

Keys sort as strings. An unpadded numeric priority therefore sorts in dictionary order, not numeric order, and a timestamp of varying digit width does the same within a priority band.

\begin{table}[htbp]
\centering\small
\caption{Lexicographic ordering of composite keys, unpadded and padded.}\label{tab:key-sort}
\begin{tabularx}{\textwidth}{@{}>{\ttfamily\footnotesize}p{4.2cm} >{\ttfamily\footnotesize}p{4.2cm} X@{}}
\toprule
\normalfont\textbf{Unpadded, as sorted} & \normalfont\textbf{Padded, as sorted} & \normalfont\textbf{Correct?} \\
\midrule
1\_1760\_a & p1\_0000001760\_a & red, older \\
10\_1759\_d & p1\_0000009817\_b & red, newer \\
2\_1759\_b & p2\_0000001759\_c & amber \\
2\_9817\_c & p3\_0000001759\_d & green \\
\bottomrule
\end{tabularx}
\end{table}

In the unpadded column a priority-10 item sorts above every priority-2 item, and within priority~2 the ordering by timestamp happens to be correct only because both values have the same digit count. The defect is invisible in testing with a small, homogeneous set of keys and appears in deployment as a queue whose head is not the most urgent item --- which defeats the entire purpose of \cref{sec:alert-volume}, where the ordering was established as the response to alert volume.

The corrected specification: priority is a fixed-width alphanumeric token with the urgent tier sorting first (\texttt{p1} red, \texttt{p2} amber, \texttt{p3} green), the timestamp is epoch milliseconds zero-padded to fixed width, and the event identifier breaks remaining ties. The non-numeric prefix is also required for a second reason: a key consisting only of digits may be treated as an array index rather than an object key, which changes the ordering semantics.

Under that correction, the priority ordering required by \cref{sec:alert-volume} is obtained for free from the key, with no sort, no index and no secondary query. This is the kind of result that falls out only if the access pattern is designed before the schema, which is why \cref{tab:schema-rules} places that rule first.

\subsubsection{Why timestamp keys also solve the offline merge}

\Textcite{efendi2025} describe the offline behaviour: data are \bq{stored locally and real-time events continue to occur even when the device is offline}, and are \bq{synchronized between local data change remotely when the device is offline and auto merging when the device is going to be online}.

Automatic merge resolves conflicts leaf by leaf, last write winning. For mutable state that is a hazard. For an append-only series keyed by measurement timestamp it is harmless, because every key is written exactly once by exactly one writer, so no two writers ever contend for the same leaf.

That guarantee depends on the key being collision-free, which is why \cref{tab:schema-rules} specifies millisecond resolution. At second resolution two legitimate records can share a key --- a nurse entering a reading twice in the same second, or a device batching --- and the append-only rule of \cref{sec:auth} would then \emph{reject the second record as a duplicate}. The interaction between two otherwise sound rules would thus discard real data silently, and the resolution is resolution: a key fine enough that collision does not occur, with \texttt{measurement\_id} available as the disambiguator if it ever does.

The timestamp key is therefore correct for three independent reasons --- the range query, the queue ordering, and conflict-free offline merge --- which is a reasonable indication that it is the right structure rather than merely a convenient one.

%==========================================================================
\subsection{Authentication and Authorisation}\label{sec:auth}

\subsubsection{Two distinct mechanisms}

\Textcite{efendi2025} report that the platform offers \bq{built-in authentication and authorization mechanisms, ensuring secure access control for IoT devices and users}, and that enabling \bq{secure access directly to the database from client-side} supports collaborative applications.

These are two mechanisms, and conflating them is a common error. Authentication establishes identity and issues a token. Authorisation is enforced by server-side rules evaluated on every read and write. Since clients address the database directly, the rules constitute the entire access-control layer: there is no application server between client and data in which to place policy. \Cref{tab:security-segments} placed this at segment~4.

\subsubsection{Roles belong in the token, not in the database}\label{sec:roles}

The role --- patient, nurse, physician, administrator --- must be carried as a custom claim on the authentication token, not as a field in the user's database record.

Two reasons. A database field must itself be protected by a rule, and a rule that reads a role from the database in order to decide whether that role may be read is circular. And a role stored as data is a role the client may attempt to write; a role minted server-side into a signed token is not.

One consequence must be recorded. Token claims are cached until the token refreshes, so a role revocation does not take effect immediately but at the next refresh. This bounds how quickly access can be withdrawn and is the concrete reason \cref{sec:security-design} required short session expiry: the expiry interval \emph{is} the revocation latency, and choosing it long for convenience lengthens the window in which a withdrawn clinician retains access.

\Textcite{yazdani2018} identify four user types --- physician, healthcare team, patient, system administrator --- which maps onto four claims. This system adds the constraint of \cref{tab:profiles}: a nurse may measure and view but may not modify therapy.

\subsubsection{Rule structure, and the cascade}\label{sec:rules}

\begin{table}[htbp]
\centering\footnotesize
\caption{Access matrix. Rules grant; they cannot revoke (see below).}\label{tab:access-matrix}
\begin{tabularx}{\textwidth}{@{}>{\raggedright\arraybackslash\ttfamily\scriptsize}p{3.4cm} >{\raggedright\arraybackslash}p{2.4cm} >{\raggedright\arraybackslash}p{2.2cm} >{\raggedright\arraybackslash}p{2.2cm} X@{}}
\toprule
\normalfont\textbf{Path} & \textbf{Patient} & \textbf{Nurse} & \textbf{Physician} & \textbf{Device} \\
\midrule
/measurements/\{own\} & read; append-only write & read, write if assigned & read if assigned & write own channel \\
/measurements/\{other\} & none & if assigned & if assigned & none \\
/patients/\{own\} & read, no write & read, limited write & read, write regimen & none \\
/queue/\{clinician\} & none & none & own queue only & none \\
/events & none & read if assigned & read, write verdict & none \\
/config/* & none & read & read & read thresholds \\
/devices & own device & read assigned & read & own heartbeat \\
\bottomrule
\end{tabularx}
\end{table}

One property of the rule language governs how this table must be implemented and is a frequent source of error. Rules cascade downward and are monotone: a read granted at an ancestor path cannot be narrowed at a descendant. A permissive rule high in the tree therefore silently overrides every restriction beneath it, and \cref{tab:access-matrix} is implementable only because no row grants access at a path above another row that needs less. This must be verified when the rules are written, not assumed from the table, since the table reads as though each row were independent and the rule language does not treat them so.

Three constraints deserve emphasis.

\textbf{Measurement writes are append-only.} A rule denying writes where the key already exists makes the clinical record immutable, which is what a clinical record must be.

\textbf{Immutability requires an amendment path.} An immutable record with no way to correct an entry error is unusable --- a transposed reading entered as 1200 instead of 120 would be permanent --- so correction proceeds as it does in clinical practice: a new record superseding the old by a \texttt{supersedes} pointer, with both retained and the original flagged rather than erased. The feature pipeline of \cref{sec:preprocessing} reads only non-superseded records; the audit trail retains both. This is why \texttt{supersedes} appears in \cref{tab:tree-contents} and did not appear in the draft schema.

\textbf{Only a physician may write \texttt{clinician\_verdict}}, because \cref{sec:safety} establishes that the system proposes and the physician disposes, and the adjudication table is the record of that division. And no client writes \texttt{/config}: threshold and model versions are deployment artefacts, not user data.

\subsubsection{Write-time validation}\label{sec:write-validation}

\Cref{sec:quality} required a set of controls applied at write time. The rule language supports validation expressions alongside access conditions, and this is where that requirement becomes executable rather than aspirational.

\begin{table}[htbp]
\centering\small
\caption{Write-time validation predicates.}\label{tab:write-validation}
\begin{tabularx}{\textwidth}{@{}>{\raggedright\arraybackslash}p{6.4cm} X@{}}
\toprule
\textbf{Validation} & \textbf{Enforces} \\
\midrule
Required fields present & schema conformity \\
\texttt{value} within the plausible range of \cref{sec:record-schema} & range check \\
\texttt{unit} matches the canonical unit for the channel & unit-error prevention \\
Key parses as a timestamp, not in the future & temporal coherence \\
Key does not already exist & append-only \\
For \texttt{sbp}, \texttt{dbp}, \texttt{glucose}: referenced device has \texttt{validated = true} & the WHO restriction, enforced in the database \\
\bottomrule
\end{tabularx}
\end{table}

The last row is the one worth drawing attention to. The restriction to validated upper-arm cuffs, established in \cref{sec:device-landscape} and carried through four chapters, reduces to a single server-side predicate that makes a non-conforming write impossible. A constraint enforced by the database cannot be forgotten by a developer, which is a materially stronger guarantee than a constraint enforced by convention --- and it is the only point in the system where a clinical guideline becomes a mechanical impossibility rather than a documented intention.

%==========================================================================
\subsection{Cloud Functions: The Event-Driven Pipeline}\label{sec:functions}

\Textcite{efendi2025} mention cloud functions as one of the facilities their gateway writes into, alongside the realtime database, but do not use them. This system does, and the decomposition follows the module structure of \textcite{yazdani2018}.

\subsubsection{Triggers}

\begin{table}[htbp]
\centering\footnotesize
\caption{Event-driven functions.}\label{tab:triggers}
\begin{tabularx}{\textwidth}{@{}>{\raggedright\arraybackslash\ttfamily\scriptsize}p{3.2cm} >{\raggedright\arraybackslash}p{3.0cm} X@{}}
\toprule
\normalfont\textbf{Trigger} & \textbf{Fires on} & \textbf{Action} \\
\midrule
onMeasurementWrite & create under \texttt{/measurements} & re-validate; set \texttt{quality\_flag}; recompute the rule tier server-side; record any discrepancy with the Edge tier \\
onScoringDue & a measurement completing a window & build features (\cref{sec:feature-eng}); score; prior-correct and calibrate (\cref{sec:risk-score}); map to tier; apply $\max(\text{rule},\text{model})$ (\cref{sec:rule-floor}) \\
onTierEscalation & tier amber or red & write \texttt{/events}; fan out to \texttt{/queue/\{clinician\}} with the composed key of \cref{sec:key-design} \\
onRedAlert & tier red & push notification; SMS fallback after $n$ failed deliveries (\cref{sec:transmission-priority}) \\
onVerdictWrite & \texttt{clinician\_verdict} set & close the event; append to the adjudicated label store (\cref{sec:two-labels}); notify the patient of the outcome \\
onDeviceWillMessage & MQTT Last Will (\cref{sec:mqtt-design}) & mark the device offline; start the non-transmission clock \\
\bottomrule
\end{tabularx}
\end{table}

The first deserves a note. The rule tier is computed twice: at the Edge, to set transmission priority (\cref{sec:transmission-priority}), and again in the Cloud. The second computation is not redundant, because it uses the authoritative threshold version, and a discrepancy between the two indicates a device running a stale configuration. Recording the discrepancy makes configuration drift detectable, which is otherwise a silent failure --- and a silent failure of exactly the kind \cref{sec:failure} identified as the dangerous class, since a device tiering by superseded thresholds produces plausible output that is wrong.

Under the fourteen-day window and the twice-weekly protocol, every measurement completes a rolling window, so \texttt{onScoringDue} fires on essentially every arrival. The separation from \texttt{onMeasurementWrite} is nonetheless kept, because the two have different failure consequences: a failed validation must block ingestion, whereas a failed scoring must not, the rule floor of \cref{sec:rule-floor} being sufficient on its own.

\subsubsection{Two scheduled sweeps}\label{sec:sweeps}

Neither is event-driven, and both implement requirements from earlier chapters that no source implements.

\textbf{Non-transmission sweep.} For each patient, compare the time since the last measurement against \texttt{checkup\_period\_days} --- the configurable parameter of \textcite{yazdani2018}, which they define as the intervals \bq{in which patient data should be collected and sent are determined by medical staff}. On exceedance, raise an event with \texttt{trigger\_source = non\_transmission}. This is the requirement of \cref{sec:decision-flow} that silence be alarmed, and the repeated SMS reminder of \textcite{yazdani2018} is its patient-facing half.

\textbf{Queue-age sweep.} For each unreviewed event, compare its age against a tier-dependent limit and escalate to a named alternate on exceedance. \Cref{sec:failure} identified clinician unavailability as the failure with unbounded consequence, and \cref{sec:latency} established both the measured expert-opinion delay and the right-skew of its distribution. The skew is the argument for the sweep: in a right-skewed delay distribution the long waits are a stable fraction of the caseload rather than rare accidents, so without this sweep the tail is the system's expected behaviour for some patients every month.

\subsubsection{Idempotency}

Serverless triggers may fire more than once for a single event, and \cref{sec:qos} assigns at-least-once delivery to clinical readings, which permits duplicates by construction. Every function must therefore be safe to re-run: keyed writes rather than appends, \texttt{measurement\_id} as the deduplication key, and event creation conditional on absence.

Without this, a retried trigger produces a duplicate alert --- which, given the arithmetic of \cref{sec:volume-arithmetic}, is the one failure mode this system can least afford, since the entire alerting design rests on keeping the weekly count near three.

%==========================================================================
\subsection{Synchronisation and Notification}\label{sec:sync}

\subsubsection{Who should hold a persistent connection}\label{sec:connections}

\Textcite{efendi2025} describe the mechanism: \bq{Firebase provides Realtime Database and designed for real-time updates. Whenever current sensor data arrives, Firebase sends it to linked devices via persistent WebSocket connections. This enables clients to examine the health vitals virtually immediately.}

The capability is real and \cref{sec:protocols-cloud} accepted it. The design question is which clients should hold one, and the answer is not all of them.

\begin{table}[htbp]
\centering\small
\caption{Connection policy by client.}\label{tab:connection-policy}
\begin{tabularx}{\textwidth}{@{}>{\raggedright\arraybackslash}p{3.6cm} >{\raggedright\arraybackslash}p{2.4cm} X@{}}
\toprule
\textbf{Client} & \textbf{Persistent?} & \textbf{Reason} \\
\midrule
Clinician dashboard & yes & the queue must update live; this is a $t_5$ intervention (\cref{sec:latency}) \\
Patient application & no & connect, synchronise, disconnect \\
Device or gateway & no & MQTT carries the upstream; no subscription required \\
\bottomrule
\end{tabularx}
\end{table}

Under the protocol of \cref{sec:streams} the patient application transmits twice weekly. Holding an always-open socket for a client with something to say for sixty seconds a week consumes a connection continuously to serve a duty cycle of $2\times10^{-4}$, which is the argument of \cref{sec:power} applied to the network rather than to the battery.

\subsubsection{Push, and the tier mapping}

Notification is the delivery mechanism, and it follows the tier exactly as transmission priority does.

\begin{table}[htbp]
\centering\small
\caption{Notification by tier and recipient.}\label{tab:notification}
\begin{tabularx}{\textwidth}{@{}>{\raggedright\arraybackslash}p{1.6cm} X >{\raggedright\arraybackslash}p{4.4cm}@{}}
\toprule
\textbf{Tier} & \textbf{Patient} & \textbf{Clinician} \\
\midrule
Green & silent; visible in the application & not notified; appears in the queue \\
Amber & prompt for the verification profile & queued at normal priority \\
Red & immediate high-priority push; local on-device warning regardless of connectivity (\cref{sec:transmission-priority}) & immediate push; SMS fallback \\
\bottomrule
\end{tabularx}
\end{table}

The Notification module of \textcite{yazdani2018} has exactly two functions, patient notification and physician notification, and \cref{tab:notification} is that module parameterised by tier.

Two safety constraints from \cref{sec:safety} apply to the content. A patient notification may say that a reading is high and that the patient should contact their health centre; it may not convey a dose change, a medication change, or any instruction to stop a treatment. And a clinician notification carries the proposal, never an executed action.

\subsubsection{Concurrent connections are the dominant cost term}\label{sec:connection-cost}

\Cref{sec:cloud-gaps} recorded that no source in this corpus states operating cost, and \cref{sec:regional-landscape} documented a service ceasing operations for financial reasons. The realtime database bills on three axes --- stored volume, downloaded volume, and concurrent connections --- and the third is the one a careless design inflates.

The saving must be computed on the right quantity. The billed figure is \emph{concurrent} connections, a peak, not a time average. Thirty clinician-style persistent sockets is a peak of thirty held continuously. Thirty patient clients each connecting twice weekly for sixty seconds gives a time-average of $30 \times 1.98\times10^{-4} \approx 0.006$ connections, but the peak depends entirely on whether those connections coincide: randomly distributed, the peak is one or two, and the saving is between one and two orders of magnitude.

This yields a design requirement that is invisible unless the peak is computed rather than the average. If the application prompts every patient to measure at the same hour --- which is the natural implementation of a measurement reminder, and the one \textcite{yazdani2018} describe --- then all thirty synchronise together, the peak returns to thirty, and the entire saving is lost. Measurement prompts must therefore be jittered across a window, which couples the reminder schedule to the cost model. Matching system behaviour to the clinical protocol produces the saving; matching it carelessly gives it back.

This is the third place in this chapter where that matching produces a large reduction, after battery (\cref{sec:power}) and transmitted volume (\cref{sec:mqtt-payload}).

%==========================================================================
\subsection{The Five Modules, Realised}\label{sec:modules}

The architecture of \textcite{yazdani2018} is a client--server design in XML. Its functional decomposition survives translation to a serverless stack essentially unchanged, which is worth stating: the module boundaries were drawn on functional grounds and are therefore portable, while the technology was not.

\begin{table}[htbp]
\centering\footnotesize
\caption{The modules of \textcite{yazdani2018} and their realisation here.}\label tab:modules}
\begin{tabularx}{\textwidth}{@{}>{\raggedright\arraybackslash}p{3.0cm} >{\raggedright\arraybackslash}p{3.6cm} X@{}}
\toprule
\textbf{Module} & \textbf{Components specified} & \textbf{Realisation here} \\
\midrule
Membership management & registration; checkup period; rule selection & \texttt{/patients} with \texttt{checkup\_period\_days} and \texttt{threshold\_version} \\
Communication transmission & deliver data; ensure device connection & MQTT ingest (\cref{sec:mqtt-design}); \texttt{/devices} heartbeat; Last Will \\
Supervisor & XML parser; decision-support detection & \texttt{onMeasurementWrite}, \texttt{onScoringDue}; routing by channel and assigned rule version \\
Decision-support setting & add system; update rules & \texttt{/config/thresholds/\{version\}}, administrator-write only \\
Notification & patient; physician & \texttt{onTierEscalation}, \texttt{onRedAlert}, \texttt{onVerdictWrite} \\
\bottomrule
\end{tabularx}
\end{table}

Two of their elements are more valuable than their age suggests.

\textbf{Checkup period as a per-patient parameter.} They define it as the interval at which data must be collected and sent, set by medical staff. \Cref{sec:followup-requires} established the WHO follow-up requirement: monthly after initiation or change until target, three- to six-monthly once controlled. A per-patient configurable interval makes that requirement a setting rather than a constant, and it is the input to the non-transmission sweep of \cref{sec:sweeps}.

\textbf{Rule selection per patient.} They attach a decision-support system and its rules to a patient according to disease and symptoms. For a system covering two pathologies with two guideline sources, this is what allows a hypertensive patient, a diabetic patient and a comorbid patient to be managed on one platform under different rule sets --- the scalability argument of \cref{sec:rule-model}, implemented as a reference from \texttt{/patients} into \texttt{/config/thresholds}. Given that \cref{sec:comorbidity} put the comorbidity at 58.1\% of diabetic patients, the comorbid case is not an edge case and the per-patient reference is not an optional refinement.

%==========================================================================
\subsection{Cost, Portability and Residency}\label{sec:cost}

\subsubsection{Cost structure}

The billing axes are named above; the figures must come from current provider documentation and are not asserted here. What the design controls:

\begin{table}[htbp]
\centering\small
\caption{Cost drivers and the design decisions that control them.}\label{tab:cost-drivers}
\begin{tabularx}{\textwidth}{@{}>{\raggedright\arraybackslash}p{3.4cm} >{\raggedright\arraybackslash}p{3.8cm} X@{}}
\toprule
\textbf{Axis} & \textbf{Driver} & \textbf{Controlled by} \\
\midrule
Stored volume & measurement count $\times$ record size & duty cycling (\cref{sec:power}); 231~kB per month for 30 patients (\cref{tab:data-volume}) \\
Downloaded volume & window reads per dashboard session & range queries, never subtree reads (\cref{sec:key-design}) \\
Concurrent connections & clients holding sockets, at peak & \cref{sec:connections}; clinicians only, with jittered patient prompts (\cref{sec:connection-cost}) \\
\bottomrule
\end{tabularx}
\end{table}

An estimated monthly cost per patient should be reported at the pilot scale and at a hypothetical thousand-patient scale. \Cref{sec:cloud-gaps} recorded that no source in the corpus does this; doing it is a contribution in its own right, and the case documented in \cref{sec:regional-landscape} establishes why it matters more here than the absence of such figures from the literature would suggest.%
\verify{Current pricing, quotas and free-tier limits are not in the corpus and change over time. Obtain them from provider documentation and cite the access date. The structure of \cref{tab:cost-drivers} is stable; only the unit prices are not.}

\subsubsection{Portability}

\Textcite{efendi2025} list zero maintenance among the platform's advantages. The corresponding cost is that security rules, function bindings and the tree structure are provider-specific and do not port.

The mitigation is a boundary. The domain logic --- the rule engine, the feature pipeline of \cref{sec:feature-eng}, the scoring and calibration of \cref{sec:risk-score} --- is kept in plain, provider-agnostic code, and the cloud functions do nothing but receive an event, call that code, and write the result. Under that arrangement the irreplaceable part of the system is portable and only the plumbing is not. This also has a consequence for \cref{ch:results}: the model can be evaluated offline, outside the platform entirely, which is how the nested cross-validation of \cref{sec:cv-structure} will in fact be run.

A scheduled export to a neutral format serves three purposes simultaneously: portability insurance, the research analysis file of \cref{sec:reident}, and a backup. It is also segment~7 of \cref{tab:security-segments}, the segment no source addresses, so its protection is specified there rather than left to the backup's defaults.

\subsubsection{Residency is a one-time irreversible decision}\label{sec:residency-irreversible}

The published endpoint of \textcite{efendi2025} was hosted in a south-east Asian region, which is visible in their own figure. \Cref{sec:residency} established that the storage region determines the governing jurisdiction, making it a legal decision taken by a configuration setting.

A database instance's region is fixed at creation and cannot subsequently be changed. Moving the data later requires creating a new instance and migrating, with a new endpoint and reconfiguration of every client. The decision must therefore be taken deliberately, before the first write, in conjunction with the legal determination required by \cref{sec:applicable-law} --- which makes that outstanding document not merely a citation gap but a precondition for creating the database. Combined with the consent dependency identified in \cref{sec:consent}, the applicable instrument now gates two separate steps that must precede any collection.

%==========================================================================
\subsection{Specification}\label{sec:spec43}

\begin{table}[htbp]
\centering\footnotesize
\caption{Backend requirements established in \cref{sec:backend}.}\label{tab:spec43}
\begin{tabularx}{\textwidth}{@{}X >{\raggedright\arraybackslash}p{4.8cm}@{}}
\toprule
\textbf{Requirement} & \textbf{Basis} \\
\midrule
Seven top-level trees, one per access pattern & \cref{sec:schema-design} \\
Time series keyed by \texttt{ts\_measurement} at millisecond resolution & \cref{sec:offline,sec:record-schema} \\
Window retrieved by key range, never a subtree read & \cref{sec:key-design} \\
Queue key composed with fixed-width padded fields and a non-numeric priority prefix & \cref{sec:key-design} \\
\texttt{/events} adjudication tree present from first deployment & \cref{sec:safety,sec:adjudication} \\
\texttt{/config/thresholds} and \texttt{/config/models} versioned & \cref{sec:interop,sec:cloud-layer} \\
No identifying field anywhere in the tree & \cref{sec:reident} \\
Roles as custom token claims, never database fields & \cref{sec:roles} \\
Session expiry chosen as the revocation latency it implies & \cref{sec:roles,sec:security-design} \\
Rule cascade verified: no grant above a path needing a narrower grant & \cref{sec:rules} \\
Measurements append-only; existing key rejected & \cref{sec:rules} \\
Corrections by superseding record, never by edit & \cref{sec:rules} \\
Only physicians write \texttt{clinician\_verdict}; no client writes \texttt{/config} & \cref{sec:safety,sec:rules} \\
Write-time validation rejects unvalidated devices for BP and glucose & \cref{sec:device-landscape,sec:quality,sec:write-validation} \\
Rule tier recomputed server-side; discrepancy recorded as drift & \cref{sec:functions} \\
$\max(\text{rule},\text{model})$ applied in the scoring function & \cref{sec:rule-floor} \\
Non-transmission sweep against \texttt{checkup\_period\_days} & \cref{sec:decision-flow}; \textcite{yazdani2018} \\
Queue-age escalation sweep, justified by the delay skew & \cref{sec:failure,sec:latency} \\
All functions idempotent; \texttt{measurement\_id} as deduplication key & \cref{sec:qos,sec:functions} \\
Persistent connection for clinicians only & \cref{sec:connections} \\
Patient measurement prompts jittered, to protect the peak & \cref{sec:connection-cost} \\
Push content constrained by the safety rule & \cref{sec:safety} \\
Domain logic provider-agnostic; functions thin & \cref{sec:cost} \\
Scheduled export to a neutral format, protected as segment 7 & \cref{tab:security-segments,sec:reident} \\
Estimated cost per patient per month reported at two scales & \cref{sec:cloud-gaps} \\
Database region chosen deliberately before first write; irreversible & \cref{sec:residency,sec:residency-irreversible} \\
\bottomrule
\end{tabularx}
\end{table}

%==========================================================================
\subsection{Limits of What This Section Establishes}\label{sec:backend-limits}

Three limits apply and are stated wherever the corresponding material is used.

\textbf{Pricing and quotas are outside the corpus.} \Cref{tab:cost-drivers} gives the structure of the cost model; the unit figures must come from current provider documentation with the access date cited.

\textbf{The choice between the realtime database and a document store is not settled here.} \Textcite{efendi2025} used the realtime database and it suits this workload well, a timestamp-keyed range query being exactly what a single JSON tree does cheaply. A document store offers richer querying under different billing, and the comparison must be made against current documentation rather than from this corpus. The schema of \cref{sec:schema-design} would require restructuring into collections and documents under that alternative; the access patterns it is designed around would not change, and those are the part worth preserving. This is the practical value of designing the access pattern first: it survives the storage decision.

\textbf{Nothing here has been validated in routine care.} \Cref{sec:arch-validation} established that of the four architectures in the corpus, two have no implementation, one was validated on six healthy volunteers over ten minutes, and the fourth on a small panel with simulated patients. This backend design has no precedent in the corpus that has been tested in service, and that is stated rather than implied.

The three outstanding documents are unchanged --- the two guideline editions required by \cref{tab:assumptions} and the Cameroonian instrument required by \cref{sec:applicable-law} --- with the instrument now determining the irreversible residency decision of \cref{sec:residency-irreversible} in addition to the consent document of \cref{sec:consent}.

%##########################################################################
% SECTION 4.4
%##########################################################################
\section{Application Interfaces Design}\label{sec:interfaces}

\subsection*{Purpose and sources}

\Cref{sec:architecture,sec:iot,sec:backend} specified the architecture, the hardware and the backend. This section specifies the two interfaces through which the system is actually used, and it is the point at which the system either works in practice or is abandoned.

\Textcite{efendi2025} supply the only implemented mobile interface in the corpus: registration, login and monitoring screens, a prediction display, a live chart, and a minimum device specification. \Textcite{yazdani2018} supply the clinician side --- the priority-ordered physician view, the role structure, and the configurable checkup period. \Textcite{cadrage2026} supply the resource-constrained deployment context that motivates local inference. \Textcite{zhou2024} supply the only quantified evidence in the corpus about which interface features change outcomes.

One source is excluded. The link compilation examined in \cref{sec:comparability} contains several incorrect author attributions and identifiers, and a document whose attributions cannot be relied upon cannot support a design claim. Where a paper it lists is used below, it is cited from the paper itself.%
\verify{The cardiologist-density figures attributed to \textcite{segun2025} in \cref{sec:queue-row} must be checked against that paper directly rather than against the link compilation, since that compilation is the document whose attributions are in question. If the figures are not in the paper, the argument of \cref{sec:queue-row} stands on clinician scarcity generally and the specific ratios should be removed.}

%==========================================================================
\subsection{What the Corpus Offers, and the Gap}\label{sec:corpus-gap}

The application of \textcite{efendi2025} has one user type. Registration collects username, password, full name, gender, address and date of birth with age derived, together with a condition dropdown offering optimal and less optimal. Login offers credentials, password recovery and a registration link. The monitoring screen shows heart rate, oxygen saturation, temperature, a prediction field reading \emph{Stable} or \emph{Not Stable}, a heart-rate graph updating every five minutes, and a logout control. Their flowchart runs registration and login, then sensing, microcontroller and radio, the cloud database, a REST interface, and the monitoring screen, concluding that the \bq{user receives real-time information}.

There is no clinician anywhere in it. The monitored person is the only user, and the prediction is displayed to that person directly.

\Textcite{yazdani2018} have the clinician but no implementation. Their architecture is given in UML: three layers, five modules, an activity diagram, XML exchange and SMS notification, with an interface that \bq{can be installed on a smart device (tablet, cell phone) or be presented as web-based}. Their single most valuable element for this section is the priority-ordered physician view, carried forward since \cref{sec:alert-volume}.

The honest position is therefore this: the clinician dashboard specified below has no implemented precedent in this corpus. One source implemented an interface without a clinician; the other designed a clinician interface without implementing it. \Cref{sec:interface-limits} restates this as a limitation, and the dashboard is not presented as following established practice.

%==========================================================================
\subsection{The Design Argument That Governs the Patient Application}\label{sec:app-argument}

Two findings in the corpus, placed together, determine where effort should go. Neither source makes the comparison.

\Textcite{alessa2019} reviewed 186 hypertension self-management applications. 58.1\% performed a single function. Of those supporting blood-pressure tracking, 73 of 84 required manual entry. Only 3.2\% reported any health-professional input in development, and 70\% shared data with third parties. Their behaviour-change-technique count has a median of 9; the reported mean of 18.4 cannot be reconciled with their stated range of 6--17, since a mean cannot exceed a maximum, and is not used here (\cref{sec:apps}).

\Textcite{zhou2024} meta-analysed 20 randomised trials with 4\,168 participants and report a systolic reduction of $-5.78$~mmHg (95\% CI $-7.97$ to $-3.59$) with 85.2\% retention. Their subgroup analyses are the operative result: interventions employing eleven or more behaviour-change techniques achieved $-9.68$~mmHg against $-2.88$ for ten or fewer, and multi-component interventions achieved $-6.50$ against $-1.54$ for single-component.

Set side by side, the median deployed application sits at nine techniques, which places it in the stratum that achieved $-2.88$ --- below the count at which a clinically meaningful effect appears. The market does not merely underperform; it clusters immediately beneath the point at which the evidence says benefit begins. That is the gap this design must close, and \cref{sec:bct-mapping} shows it closes it by construction rather than by content.

Two cautions attach and are stated here rather than appended. \Textcite{zhou2024} rate the systolic estimate as low certainty, 75\% of included trials were at high risk of bias, and only 30\% were theory-based; the diastolic estimate of $-3.28$ is moderate certainty. The dose--response in technique count is a subgroup finding from a low-certainty body of evidence, so it is adopted here as a design target and not as a predicted effect size.

A third finding constrains what the application should promise. \Textcite{mihevc2025} achieved a systolic reduction of $-6.9$~mmHg (95\% CI $-11$ to $-2.7$, $p = .001$) and an HbA1c reduction of $-0.5$\% ($-0.8$ to $-0.2$, $p = .002$) --- and report that no behaviour change occurred. The mechanism was medication titration by the clinician in response to transmitted readings. \Cref{sec:mechanism} drew the conclusion and it governs here: the patient application's primary function is reliable capture and prompt escalation, not behaviour modification. \Textcite{mihevc2025} also lost 10.9\% of participants to inability to use the equipment, which makes entry burden an attrition mechanism rather than a usability nicety.

%==========================================================================
\subsection{Patient Mobile Application}\label{sec:patient-app}

\subsubsection{Platform}\label{sec:platform}

\Textcite{efendi2025} built a native Android application. A cross-platform toolkit offers a single codebase and a single widget set.

The usual argument for cross-platform development, covering a second mobile platform, is weak in this deployment, where one platform's share is overwhelming and a build for the other would serve a small minority. The real reasons to prefer it are development speed for a single researcher and consistent rendering across heterogeneous low-end hardware. The clinician dashboard should not be built in the same framework: it is a dense tabular interface used in a desktop browser, and a plain web application serves it better than a cross-platform widget toolkit.

One check must be performed before committing, and it cannot be performed from the corpus: the minimum platform version of \textcite{efendi2025} may already fall below the floor supported by current toolchains.%
\verify{Obtain the current minimum API levels for the chosen cross-platform framework and for the on-device inference runtime, and state the API level this work actually targets. If the floor of \textcite{efendi2025} is unreachable, \cref{tab:device-spec} must be revised upward and the implication for device cost stated.}

\subsubsection{Manual entry: three safety requirements}\label{sec:manual-entry}

Manual entry is the dominant acquisition mode in the reviewed market --- 73 of 84 applications --- and \cref{sec:options} adopted Option~C, under which manual entry carries blood pressure and glycaemia while instrumentation carries covariates only. The entry form is therefore the system's principal data path and not a fallback, which is a reversal of how such forms are usually treated.

The registration form of \textcite{efendi2025} is a reasonable starting point, but three requirements appear neither in it nor in any source, and all three are safety-critical.

\textbf{The unit must be fixed per patient, displayed, and never selectable per entry.} Glucose in mg/dL and in mmol/L differ by a factor of 18.0, so 7~mmol/L and 126~mg/dL are the same quantity. A value of 7 entered under the wrong unit is not a noisy reading; it is a reading that will cross or fail to cross every threshold in the system, and it is perfectly plausible under either unit, so no range check can catch it. \Cref{sec:write-validation} places a unit-matching predicate in the database rules, which catches a mismatch server-side; the interface must make the error unreachable client-side, by binding the unit to the patient record at enrolment and rendering it as fixed text beside the input.

\textbf{The form must capture duplicate readings, not one.} \Cref{sec:device-landscape} established the guideline convention of duplicate measurement with averaging. A form offering a single field silently produces a different quantity from the one on which the thresholds are defined. Two fields with the average computed and displayed, together with \texttt{reading\_index} recorded per \cref{tab:tree-contents}, is the minimum conforming design.

\textbf{Context is mandatory, not optional.} \Cref{sec:threshold-features} and the record schema require measurement context --- at rest, post-exercise, post-meal, pre-meal --- because the threshold applied depends on it. A context field the patient may skip yields records that cannot be labelled, which removes them from the training set entirely rather than merely degrading them.

Two further behaviours follow from earlier sections. Plausibility feedback at entry uses the ranges of \cref{sec:record-schema} and \emph{queries rather than rejects}: a hard rejection discards the genuine extreme value, which is precisely the value the system exists to catch, so the correct response to an implausible entry is a confirmation prompt whose confirmation is itself recorded. And the timestamp defaults to the present and may be edited backwards only, since \cref{sec:write-validation} rejects future timestamps at write time --- a client that permits a future timestamp produces a write the server will refuse, which the patient experiences as data loss.

\subsubsection{Trends and logs}\label{sec:trends}

The heart-rate graph of \textcite{efendi2025} updates every five minutes. That cadence belongs to a live monitor and is the wrong unit for this system, where the protocol of \cref{sec:streams} produces approximately two blood-pressure occasions per week.

The chart displays the fourteen-day window of \cref{sec:feature-eng} with the applicable threshold band drawn behind the series. Drawing the band is what makes the chart answer a clinical question rather than merely display numbers: the patient sees position relative to target, which is the discrepancy-between-current-state-and-goal technique counted in \cref{tab:bct-mapping}. The window is already available locally through offline persistence (\cref{sec:schema-design}), so the chart renders without a network round trip --- which matters because a chart that blanks when connectivity drops teaches the patient that the application is unreliable.

The log is the same data in tabular form, showing both timestamps of \cref{sec:record-schema}, measurement time and record time, because a patient who entered a backdated reading should be able to see that the system recorded it as backdated.

\textbf{What the patient is not shown.} \Textcite{efendi2025} display their classifier's output to the monitored person directly. This system does not, and the departure is deliberate: \cref{sec:never-display} establishes that even a clinician should not be shown a raw score from a model fitted at this sample size, and the patient has no means at all of discounting it. The patient sees the measured values, their position relative to the band, and --- on amber or red --- the instruction permitted by \cref{sec:safety}. No probability, no score, and no label of their future state.

The threshold band cannot be drawn, and no target can be displayed, until the guideline editions recorded as outstanding in \cref{tab:assumptions} are supplied. This is the most concrete consequence yet of that gap: the central visual element of the patient interface is unbuildable without them.

\subsubsection{Self-management advice, and how the technique target is met}\label{sec:bct-mapping}

\Cref{sec:safety} is the binding constraint and is not negotiable: the application may tell the patient that a reading is high and that they should contact their health centre; it may not convey a dose change, a medication change, or an instruction to stop a treatment. The finding of \textcite{alessa2019} that 3.2\% of applications had any professional input, and that three claimed cuffless blood-pressure measurement without supporting evidence, is what unconstrained advice content produces.

The reconciliation between the eleven-technique target of \textcite{zhou2024} and the absence of behaviour change in \textcite{mihevc2025} is the design insight of this section. Many of the highest-value behaviour-change techniques are structural properties of a telemonitoring system rather than items of advice content. A system that measures, charts, reminds, and routes a clinician's response accumulates techniques by existing. The mapping below is therefore how the target is reached without a library of tips.

\begin{table}[htbp]
\centering\small
\caption{Behaviour-change techniques delivered structurally. Ten of the twelve are architecture rather than content.}\label{tab:bct-mapping}
\begin{tabularx}{\textwidth}{@{}>{\raggedright\arraybackslash}p{6.0cm} X@{}}
\toprule
\textbf{Technique} & \textbf{Delivered by} \\
\midrule
Self-monitoring of outcome & the measurement protocol itself (\cref{sec:streams}) \\
Feedback on outcome & the trend chart with threshold band (\cref{sec:trends}) \\
Prompts and cues & scheduled reminder from \texttt{checkup\_period\_days} (\cref{sec:sweeps}) \\
Goal setting, outcome & the target band displayed \\
Discrepancy between current state and goal & the reading plotted against that band \\
Review of outcome goals & the clinician verdict relayed back (\cref{sec:functions}) \\
Action planning & a fixed when-and-where measurement schedule \\
Habit formation & the same schedule, repeated \\
Social support, practical & nurse-assisted measurement (\cref{tab:profiles}) \\
Credible source & advice attributed to the clinician and the guideline version \\
Information about health consequences & threshold explanation at the band \\
Behavioural contract & the enrolment protocol agreement \\
\bottomrule
\end{tabularx}
\end{table}

This mapping is a construction of this work. It has not been coded against the taxonomy by a trained coder, which is the standard the intervention-reporting literature expects, and the taxonomy itself is not in the supplied corpus. The count is therefore reported as a design intention subject to formal coding, not as an achieved score.%
\verify{Obtain and cite the behaviour-change technique taxonomy, and have \cref{tab:bct-mapping} coded by a second reader against it. Until then no numeric technique count should appear in an abstract or a results chapter, because a self-assigned count against an uncited taxonomy is not a measurement.}

\subsubsection{Notifications}\label{sec:app-notifications}

The tier mapping of \cref{tab:notification} governs delivery. Green is silent and visible in the application; amber prompts the verification profile of \cref{sec:mihevc}; red raises an immediate high-priority notification together with a local on-device warning regardless of connectivity, per \cref{sec:transmission-priority} --- which is one of the two functions local inference exists to serve (\cref{sec:local-alert}).

\Textcite{yazdani2018} add the reminder in the other direction: \bq{if no data is sent within the time limits set by the physician, a SMS reminder will be sent}. That is the patient-facing half of the non-transmission sweep of \cref{sec:sweeps}, and SMS rather than push is the correct channel for it, since the condition being detected is precisely that the application is not communicating --- a push notification to a silent client is a message sent down the channel known to be broken.

One interaction with \cref{sec:connection-cost} must be respected. Measurement prompts must be jittered across a window rather than fired at a common hour, because a synchronised prompt synchronises the subsequent connections and restores the concurrent-connection peak the design exists to avoid.

\subsubsection{Minimum device specification}\label{sec:device-spec}

The specification of \textcite{efendi2025} is the only one in the corpus and is adopted as a floor, subject to the toolchain check of \cref{sec:platform}.

\begin{table}[htbp]
\centering\small
\caption{Minimum device specification, after \textcite{efendi2025}, with one requirement relaxed.}\label{tab:device-spec}
\begin{tabularx}{\textwidth}{@{}>{\raggedright\arraybackslash}p{4.0cm} X@{}}
\toprule
\textbf{Requirement} & \textbf{Value} \\
\midrule
Operating system & as \textcite{efendi2025} specify, subject to \cref{sec:platform} \\
Free storage & 1~GB \\
RAM & 2~GB \\
Processor & 1.2~GHz quad-core \\
Connectivity & \emph{intermittent, any generation} --- not 4G \\
\bottomrule
\end{tabularx}
\end{table}

The connectivity requirement is the one that must not be carried over. \Cref{sec:mqtt-payload} established a transmitted volume of 231~kB per month for thirty patients under duty cycling, or 7.7~kB per patient. At the effective throughput of a second-generation mobile connection, the entire cohort's monthly traffic is of the order of a minute and a half of transfer time, and a single patient's is a few seconds. Specifying a fourth-generation connection as a floor, as \textcite{efendi2025} do for a continuously streaming monitor, would exclude exactly the resource-constrained settings \textcite{cadrage2026} targets --- and would do so to serve a volume that a far older network carries without difficulty. The requirement here is intermittent connectivity of any generation, which is what \cref{sec:offline,sec:edge-inference} exist to accommodate.

%==========================================================================
\subsection{Clinician Dashboard}\label{sec:dashboard}

\subsubsection{A queue, not a patient list}\label{sec:queue-view}

The unit of clinical work is an unreviewed event, not a patient. \Cref{sec:schema-design} therefore stores the queue under a composed key, and \cref{sec:key-design} showed that reading that path with a key ordering returns the list already sorted by priority and then by age within priority. The dashboard performs no sorting: the ordering specified by \textcite{yazdani2018} is a property of the data, not a feature of the view.

The justification is the latency budget of \cref{sec:latency}. Decomposing end-to-end delay showed that $t_5$, clinician review, exceeds every other term by orders of magnitude. \Textcite{bediang2021} supply the measurement: their intervention accelerated access to examination from 33.3\% to 98.9\% (a relative risk of 2.97) and modified the diagnosis in 46.2\% of cases, yet the delay to expert opinion was 10.3 against 11.2 days ($p = 0.82$), statistically indistinguishable, against a protocol requiring 72 hours. Automating acquisition while leaving review untouched produced no improvement in the term that dominates.

The dashboard is therefore the system's principal latency intervention, and the queue-age escalation sweep of \cref{sec:sweeps} is what prevents the tail of that delay distribution from becoming this system's ordinary behaviour for some patients --- which, as \cref{sec:latency} established from the dispersion of those measurements, is what a right-skewed delay distribution produces if nothing bounds it.

\subsubsection{Three tiers, not two}\label{sec:three-tiers}

\Textcite{cadrage2026} specifies ranking by risk tier, normal against alert. This should be three tiers, for two reasons drawn from the corpus.

\Textcite{mihevc2025} --- the only trial in the corpus reporting both a protocol and a significant effect --- used green, amber and red, and \cref{sec:mihevc} built the two-stage verification profile on the amber tier. A binary scheme must collapse amber into one of the other two. Folded into normal, the verification profile never triggers and borderline readings go unexamined. Folded into alert, every borderline reading enters the queue and the arithmetic of \cref{sec:accuracy-trap} applies: at a 2\% event rate and 90\% sensitivity, a precision of 80\% requires a specificity of 99.5\%, which no model in this corpus approaches, so the queue fills with false alerts and ceases to be read.

The second reason is the experience of \textcite{efendi2025} themselves. Their binary classifier labelled a heart rate of 70.1 as stable, and their two stated reference bands for the stable class disagree with each other, so the label cannot be checked against a single stated criterion.%
\verify{\Textcite{efendi2025} state the stable heart-rate band differently in two tables, and the lower bound differs between those statements. Record both bands verbatim with their table numbers, and state whether 70.1 falls inside or outside each. The argument here requires only that the two bands disagree, which is established; the stronger claim that the label contradicts a stated band requires knowing which lower bound applies.}

The design point survives either reading. A two-state display offers no position in which to express uncertainty, so a classifier whose boundaries are incoherent emits a confident label rather than a flag for review --- and a confident wrong label is harder for a clinician to discount than an acknowledged borderline.

\subsubsection{What a queue row carries}\label{sec:queue-row}

Clinician attention is the scarcest resource in this deployment. \Textcite{segun2025} report that a substantial proportion of sub-Saharan African countries have no cardiologist at all, and specialist-to-population ratios of the order of one per several hundred thousand. Every element on the dashboard must therefore reduce the time to a triage decision, which means the row must be triageable without opening the patient record.

\begin{table}[htbp]
\centering\small
\caption{Fields carried by a queue row.}\label{tab:queue-row}
\begin{tabularx}{\textwidth}{@{}>{\raggedright\arraybackslash}p{6.0cm} X@{}}
\toprule
\textbf{Field} & \textbf{Why} \\
\midrule
Tier & the triage decision itself \\
Trigger source: rule, model, or non-transmission & \cref{sec:never-display} \\
The triggering value, with unit and context & the clinical fact \\
Days since last review & the escalation condition (\cref{sec:sweeps}) \\
Direction of the fourteen-day trend & distinguishes an isolated excursion from a deterioration \\
Patient pseudonym and age band & identification without identity (\cref{sec:reident}) \\
Model and threshold versions & \cref{tab:tree-contents} \\
\bottomrule
\end{tabularx}
\end{table}

\subsubsection{What must never be displayed}\label{sec:never-display}

Two prohibitions, both derived in earlier chapters, and both easy to violate by building the obvious interface.

\textbf{Never display a raw score from a class-weighted model.} \Cref{sec:prior-correction} derived the correction $s = s'/(s' + w(1-s'))$. At a weight of $w = 20$ a raw output of $s' = 0.5$ corresponds to a true probability of $1/21 \approx 0.048$, so a clinician shown $0.5$ would be reading a number about ten and a half times too large; at the conventional setting $w = (1-\pi)/\pi = 49$ the same raw output corresponds to exactly the base rate of $0.02$. Only the prior-corrected, calibrated value may be displayed. Given the independent sample size recorded in \cref{tab:assumptions}, even the corrected value carries false precision at two decimal places, and the tier together with the contributing reason communicates more honestly than any decimal.

\textbf{Always display which mechanism fired.} \Cref{sec:rule-floor} set $\text{tier}_{\text{final}} = \max(\text{tier}_{\text{rule}}, \text{tier}_{\text{model}})$. A rule trigger is a threshold crossing traceable to a published guideline and is defensible on its own terms. A model trigger, from the sample size of \cref{tab:assumptions}, is weak evidence. Rendering both in the same colour with no attribution tells the clinician that two very different grades of evidence are the same thing, which is the specific failure that makes a decision-support system distrusted after its first false alarm. The trigger source is therefore a primary field and not metadata.

\subsubsection{Verdict capture is the label pipeline}\label{sec:verdict-capture}

\Cref{sec:adjudication,sec:two-labels} established that the physician's adjudication is the only source of clinically validated labels in this design. The verdict control is therefore not administrative overhead; it is the training-data acquisition mechanism for the second contribution.

The consequence is a hard interface requirement: if recording a verdict is tedious, the label store stays empty and that contribution fails. The control must resolve an event in two or three interactions, with options mapping directly onto the label definition --- true alert, false alert, indeterminate --- plus the action taken, and an optional free-text comment that is never required. The requirement of \textcite{yazdani2018} that \bq{all events\ldots including physician and CDSS prescriptions, are stored as patient record} is satisfied by writing \texttt{clinician\_verdict} into the events tree, which \cref{sec:functions} makes the trigger for closing the event and appending to the label store.

That the two-or-three-interaction figure is an argument rather than a measurement is stated in \cref{sec:interface-limits}.

\subsubsection{Roles}\label{sec:dash-roles}

\Textcite{yazdani2018} identify physician, healthcare team, patient and system administrator. \Cref{tab:profiles,tab:access-matrix} constrain them: the physician holds a queue and may write a verdict and modify therapy; the nurse may measure and view assigned patients but holds no queue and may not modify therapy; the administrator manages enrolment and configuration but reads no clinical values.

The dashboard renders from the token claim of \cref{sec:roles}, and the nurse view is not a greyed-out physician view but a different screen --- a measurement and assistance worklist, matching what \cref{tab:profiles} assigns that role. A disabled-control version of the physician screen communicates that the nurse is a deficient physician, which is both wrong about the task-sharing arrangement of \cref{sec:hearts-spec} and an invitation to work around the restriction.

\subsubsection{It must replace a task, not add one}\label{sec:replace-task}

\Textcite{dodoo2022} identified the decisive adoption barrier across 66 studies: \bq{healthcare professionals view telemedicine as an additional task}. \Cref{sec:records,sec:time} documented what the dashboard is competing against --- paper registers, incomplete records, and readings recited from a card at a consultation.

The design conclusion is that the dashboard must become the place where home readings are reviewed, in place of the register, and not a second system consulted after it. Concretely: the queue row must carry enough that no paper is required to triage; the patient view must show the same fourteen-day window the patient sees, so that the consultation and the dashboard do not disagree in front of the patient; and the verdict must be the single act that both closes the alert and records the clinical note. A dashboard requiring the clinician to transcribe anything into another system has added a task, and the evidence of \textcite{dodoo2022} predicts it will be abandoned.

%==========================================================================
\subsection{Edge Inference}\label{sec:edge-inference}

No source in the corpus implements this. \Textcite{efendi2025} compute the prediction in the cloud and display it, and the decision support of \textcite{yazdani2018} is entirely server-side. What follows is a design proposal with no tested precedent here.

\subsubsection{What actually has to be exported}\label{sec:export-what}

\Textcite{cadrage2026} describes exporting the trained model to the mobile application. The model is the easy part and the smallest part. A score reproduced on the device requires the whole path from \cref{sec:feature-eng,sec:learning} to be reproduced identically.

\begin{table}[htbp]
\centering\small
\caption{What on-device scoring requires, by difficulty.}\label{tab:edge-requirements}
\begin{tabularx}{\textwidth}{@{}>{\raggedright\arraybackslash}p{6.4cm} X@{}}
\toprule
\textbf{Component} & \textbf{Difficulty on device} \\
\midrule
The fourteen-day window & already solved by offline persistence (\cref{sec:schema-design}) \\
Feature construction (\cref{sec:feature-eng}) & must match the cloud exactly \\
Per-patient normalisation constants & must be the training-time constants, shipped and versioned (\cref{sec:expanding}) \\
Masking and the decay $\gamma_t = \exp(-\kappa_c\,\Delta t)$ & deterministic, portable \\
The model & trivially small (\cref{sec:model-size}) \\
Calibration and prior correction (\cref{sec:risk-score}) & mandatory, or the local score is an order of magnitude wrong \\
Thresholds $\tau_{\text{amber}}$, $\tau_{\text{red}}$ & from the model registry (\cref{tab:tree-contents}) \\
\bottomrule
\end{tabularx}
\end{table}

The difficulty is fidelity, not capacity. A feature pipeline reimplemented for the device against one written in the training language will diverge --- in how a window boundary is treated, how a tie is broken, how a missing value is masked, which normalisation constants are used. A device producing silently different scores from the cloud is worse than a device producing none, because the discrepancy is invisible to everyone, including the clinician who sees one number and the researcher who computes another.

The only defensible arrangements are a single shared implementation compiled for both targets, or an exported pipeline artefact that both sides execute, in either case with conformance tests asserting tolerance-bounded agreement on a fixed test set as a release gate. The feature-set hash already recorded in \cref{tab:tree-contents} is what makes such a gate checkable at runtime rather than only at build time.

\subsubsection{The rule engine is the part worth exporting}\label{sec:rule-export}

This follows directly from the composition of \cref{sec:rule-floor}.

The rule engine is a short list of comparisons of a value against a threshold, selected by channel and context. It is deterministic, has no numerical drift, requires no normalisation constants, and is a few hundred bytes of configuration. The model requires everything in \cref{tab:edge-requirements}.

The rule engine therefore delivers nearly all of the offline value at a small fraction of the risk. In an outage, a device carrying only the rule engine degrades to guideline conformity --- which is precisely the floor \cref{sec:rule-floor} was constructed to guarantee --- whereas a device carrying a mis-ported model pipeline degrades to an unknown function, and an unknown function is not a degradation mode that can be reasoned about or defended.

The recommendation is to ship the rule engine in the first version, treat local model scoring as a second stage gated on the conformance tests of \cref{sec:export-what} passing, and present it in that order.

\subsubsection{Version stamping and reconciliation}\label{sec:version-stamp}

An offline device may run a stale configuration for days. Every locally produced tier must therefore be stamped with the threshold version and, where a model was used, the model version and feature-set hash.

On reconnection the cloud re-scores the same window under the authoritative versions and records any discrepancy --- exactly the mechanism \cref{sec:functions} applies to the rule tier computed at the Edge, applied here to the device. This converts configuration drift from a silent failure into a measurable one, and the discrepancy rate becomes a reportable validation result in its own right.

\subsubsection{What a local alert may actually do}\label{sec:local-alert}

Narrower than it first appears, and the narrowness is a safety property rather than a limitation.

The clinician is not on the device, and \cref{sec:safety} forbids the application from issuing therapeutic instructions. A locally computed red tier may therefore do exactly two things: warn the patient to seek care, and queue the event for transmission at the earliest opportunity, with the local warning raised regardless of connectivity per \cref{sec:transmission-priority}. It may not titrate, advise a dose, or stand in for review.

Under the transport design of \cref{sec:qos} the queued event transmits at least-once with application-layer deduplication when connectivity returns, and the Last Will mechanism of \cref{sec:mqtt-design} means the backend already knows the device went silent --- so a red tier computed during an outage reaches the clinician through two independent routes, the queued event and the absence alarm.

\subsubsection{Size is not the constraint}\label{sec:model-size}

Worth stating because it is routinely assumed to be. A gradient-boosted ensemble of the size contemplated in \cref{sec:xgb-config} --- depth two to three, with early stopping --- is tens of kilobytes. The recurrent model of \cref{sec:lstm-config} has, at sixteen units over ten input channels, exactly
\begin{equation}
  P = 4\bigl[u(f+u) + u\bigr] = 4\bigl[16(10+16) + 16\bigr] = 1\,728
\end{equation}
parameters, which is under 7~kB at single precision and under 2~kB quantised. Against the storage and memory floor of \cref{tab:device-spec} the model is roughly five orders of magnitude below the available budget.

The binding constraints are the fidelity requirement of \cref{sec:export-what} and the toolchain floor of \cref{sec:platform} --- not memory, storage or compute.

%==========================================================================
\subsection{Specification}\label{sec:spec44}

\begin{table}[htbp]
\centering\footnotesize
\caption{Interface requirements established in \cref{sec:interfaces}.}\label{tab:spec44}
\begin{tabularx}{\textwidth}{@{}X >{\raggedright\arraybackslash}p{4.8cm}@{}}
\toprule
\textbf{Requirement} & \textbf{Basis} \\
\midrule
Patient application cross-platform; dashboard a plain web application & \cref{sec:platform,sec:queue-view} \\
Unit fixed per patient at enrolment, displayed, never selected per entry & \cref{sec:manual-entry}; validated server-side in \cref{sec:write-validation} \\
Entry form captures duplicate readings with computed average and \texttt{reading\_index} & \cref{sec:device-landscape,tab:tree-contents} \\
Context field mandatory & \cref{sec:threshold-features} \\
Implausible values queried and confirmed, never rejected & \cref{sec:record-schema} \\
Timestamp defaults to now, editable backwards only & \cref{sec:write-validation} \\
Trend chart is the fourteen-day window with threshold band drawn & \cref{sec:feature-eng,sec:trends} \\
Log displays both timestamps & \cref{sec:record-schema} \\
Patient shown values, band position and permitted instruction --- never a score or label & \cref{sec:trends,sec:safety,sec:never-display} \\
Advice content may never convey a dose or medication change & \cref{sec:safety} \\
Technique count targeted at eleven or more, met structurally; formally coded before any claim & \textcite{zhou2024}; \cref{sec:bct-mapping} \\
Non-transmission reminder by SMS, not push & \cref{sec:sweeps}; \textcite{yazdani2018} \\
Measurement prompts jittered across a window & \cref{sec:connection-cost,sec:app-notifications} \\
Connectivity floor is intermittent, any generation --- not 4G & \cref{sec:mqtt-payload,sec:device-spec} \\
Dashboard is a queue of events, read pre-ordered, no client-side sort & \cref{sec:key-design}; \textcite{yazdani2018} \\
Three tiers, not two & \textcite{mihevc2025}; \cref{sec:mihevc,sec:accuracy-trap} \\
Row triageable without opening the patient record & \cref{sec:queue-row} \\
Never display a raw weighted-model score & \cref{sec:prior-correction,sec:never-display} \\
Trigger source displayed as a primary field & \cref{sec:rule-floor,sec:never-display} \\
Model and threshold versions shown on every event & \cref{tab:tree-contents} \\
Verdict resolvable in two or three interactions; comment never required & \cref{sec:adjudication,sec:verdict-capture} \\
Nurse view a distinct worklist, not a restricted physician view & \cref{tab:profiles,sec:dash-roles} \\
Dashboard must replace the register, not supplement it & \textcite{dodoo2022}; \cref{sec:replace-task} \\
Rule engine exported first; local model scoring gated on conformance tests & \cref{sec:rule-floor,sec:rule-export} \\
Shared or artefact-based feature pipeline, with agreement tests as a release gate & \cref{sec:export-what} \\
Local scores version-stamped; cloud re-scores and records the discrepancy & \cref{sec:functions,sec:version-stamp} \\
Local red tier may only warn and queue & \cref{sec:safety,sec:transmission-priority} \\
\bottomrule
\end{tabularx}
\end{table}

%==========================================================================
\subsection{Limits of What This Section Establishes}\label{sec:interface-limits}

Five limits apply.

\textbf{No usability evaluation of this design exists.} \Textcite{efendi2025} evaluated their own application on their own cohort with their own instrument, and the dimension that scored lowest was security, at 4.06 of 5 --- by their own measurement, the weakest aspect of the only implemented interface in the corpus. No instrument results exist for the design above, and the two-or-three-interaction requirement of \cref{sec:verdict-capture} is an argument, not a measurement. A usability study is the obvious validation activity and is proposed in \cref{ch:results}.

\textbf{There is no validated clinician interface to design against.} \Cref{sec:corpus-gap} established the gap. The dashboard is reasoned from the latency budget, the alert arithmetic and the adoption literature; it is not copied from a working system.

\textbf{Edge inference is unprecedented in this corpus}, which is why \cref{sec:rule-export} recommends shipping only the rule engine first.

\textbf{The technique count is unverified.} The mapping of \cref{tab:bct-mapping} is a construction of this work, uncoded, against a taxonomy not in the corpus.

\textbf{Current toolchain minimum platform versions are unknown}, which determines whether the floor of \textcite{efendi2025} is reachable (\cref{sec:platform}).

The outstanding documents now block a visible interface element rather than an internal convention. The guideline editions recorded in \cref{tab:assumptions} are required in order to draw the threshold band and to display any target: without them the central visual of the patient application cannot be built, which makes them a dependency of the implementation and not only of the labelling. The Cameroonian instrument required by \cref{sec:applicable-law} now governs two further items --- what the patient may be shown about their own record, and the consent screen at enrolment --- in addition to the residency decision of \cref{sec:residency-irreversible} and the consent document of \cref{sec:consent}. The finding of \textcite{alessa2019} that 70\% of reviewed applications shared data with third parties is the reason that instrument, and not a European regulation, must be identified before the consent screen is written.

%==========================================================================
% BIBLIOGRAPHY
%==========================================================================
\printbibliography[heading=bibintoc,title={References}]

\end{document}

